% !TEX TS-program = lualatex
% encoding : utf8
% Copyright 2011-2026  Alain Matthes
% SPDX-License-Identifier: LPPL-1.3c
% Maintainer: Alain Matthes

\PassOptionsToPackage{unicode}{hyperref}

\documentclass[DIV         = 14,
               fontsize    = 10,
               index       = totoc,
               twoside,
               cadre,
               headings    = small,
               ]{tkz-doc-euclide}
\gdef\tkznameofpack{tkz-euclide}
\gdef\tkzversionofpack{5.14c}
\gdef\tkzdateofpack{\today}
\gdef\tkznameofdoc{doc-tkz-euclide}
\gdef\tkzversionofdoc{5.14c}
\gdef\tkzdateofdoc{\today}
\gdef\tkzauthorofpack{Alain Matthes}
\gdef\tkzadressofauthor{}
\gdef\tkznamecollection{AlterMundus}
\gdef\tkzurlauthor{http://altermundus.fr}
\gdef\tkzengine{lualatex}
\gdef\tkzurlauthorcom{http://altermundus.fr}
\gdef\tkztranslatorofdoc{翻译：金永权}
\gdef\tkztranslateraddressofdoc{夏邑县第一高级中学}
\nameoffile{\tkznameofpack}
% -- Packages ---------------------------------------------------
\usepackage[dvipsnames,svgnames]{xcolor}
\usepackage{calc}
\usepackage{tkz-base}
\usepackage[lua]{tkz-euclide}
\usepackage{pgfornament}
\usetikzlibrary{backgrounds}
\usepackage[colorlinks,pdfencoding=auto, psdextra]{hyperref}
\hypersetup{
      linkcolor=Gray,
      citecolor=Green,
      filecolor=Mulberry,
      urlcolor=NavyBlue,
      menucolor=Gray,
      runcolor=Mulberry,
      linkbordercolor=Gray,
      citebordercolor=Green,
      filebordercolor=Mulberry,
      urlbordercolor=NavyBlue,
      menubordercolor=Gray,
      runbordercolor=Mulberry,
      pdfsubject={欧几里得几何},
      pdfauthor={\tkzauthorofpack},
      pdftitle={\tkznameofpack},
      pdfcreator={\tkzengine},
      bookmarksdepth=3
}
\usepackage{tkzexample}
\usepackage{fontspec}
\setmainfont{texgyrepagella}[
  Extension = .otf,
  UprightFont = *-regular ,
  ItalicFont  = *-italic  ,
  BoldFont    = *-bold    ,
  BoldItalicFont = *-bolditalic
]
\setsansfont{texgyreheros}[
  Extension = .otf,
  UprightFont = *-regular ,
  ItalicFont  = *-italic  ,
  BoldFont    = *-bold    ,
  BoldItalicFont = *-bolditalic ,
]

\setmonofont{lmmono10-regular.otf}[
  Numbers={Lining,SlashedZero},
  ItalicFont=lmmonoslant10-regular.otf,
  BoldFont=lmmonolt10-bold.otf,
  BoldItalicFont=lmmonolt10-boldoblique.otf,
]
\newfontfamily\ttcondensed{lmmonoltcond10-regular.otf}
%% (La)TeX font-related declarations:
\linespread{1.05}      % Pagella needs more space between lines
\usepackage[math-style=literal,bold-style=literal]{unicode-math}
\usepackage{fourier-otf}
\let\rmfamily\ttfamily
\usepackage{multicol,lscape}
\usepackage[UTF8]{ctex}
\AtBeginDocument{%
  \renewcommand{\contentsname}{目录}%
  \renewcommand{\indexname}{索引}%
}
\usepackage[normalem]{ulem}
\usepackage{multirow,multido,booktabs,cellspace}
\usepackage{shortvrb,fancyvrb,bookmark}
\usepackage{makeidx}
\makeindex

%<--------------------------------------------------------------------------->
% settings styles
\tkzSetUpColors[background=white,text=black]
\tkzSetUpCompass[color=orange, line width=.2pt,delta=10]
\tkzSetUpArc[color=gray,line width=.2pt,delta=10]
\tkzSetUpPoint[size=2,color=teal]
\tkzSetUpLine[line width=.2pt,color=teal]
\tkzSetUpStyle[color=orange,line width=.2pt]{new}
\tikzset{every picture/.style={line width=.2pt}}
\tikzset{label angle style/.append style={color=teal,font=\footnotesize}}
\tikzset{label style/.append style={below,color=teal,font=\scriptsize}}
\tikzset{new/.style={color=orange,line width=.2pt}}

\AtBeginDocument{\MakeShortVerb{\|}} % link to shortvrb
\setlength{\emergencystretch}{3em}% fewer overfull lines with long macro names
\def\code{\texttt}
\newcommand*{\ItkzPopt}[2]{\texttt{#2}\index{#1_3@\texttt{#1: options}!\texttt{#2}}}

\begin{document}

\parindent=0pt
\tkzTitleFrame{tkz-euclide v5.14c\\Euclidean Geometry\\中文版}
\clearpage

\defoffile{\lefthand\ \tkzname{\tkznameofpack} 是一组便捷的宏，用于在带有笛卡尔坐标系的平面中绘制图形。它处理欧几里得几何中最经典的情形。\tkzname{\tkznameofpack} 构建于 PGF 及其配套前端 \TIKZ\ 之上，是一个对 (La)TeX 友好的绘图宏包。它提供了一个高层接口，让你能够轻松地构建图形：你可以尽可能地按照手工作图的方式，一步一步地构建它们。\\
第 5 版的 \tkzname{\tkznameofpack} 可以使用 Lua 进行某些计算；请参阅本手册末尾的 \code{News} 和 \code{Lua} 两节。\\
请注意：英语不是我的母语，因此可能存在一些错误。
}

\presentation

\vspace*{1cm}
\lefthand\ 首先，我要感谢 \textbf{Till Tantau} 编写了如此出色的宏包 \href{https://github.com/pgf-tikz/pgf}{\TIKZ}。

\vspace*{12pt}
\lefthand\ 致谢：我收到了许多宝贵的建议、意见、更正和示例，来自

\tkzimp{Jean-Côme Charpentier}, \tkzimp{Josselin Noirel}, \tkzimp{Manuel Pégourié-Gonnard}, \tkzimp{Franck Pastor}, \tkzimp{David Arnold},

\tkzimp{Ulrike Fischer}, \tkzimp{Stefan Kottwitz}, \tkzimp{Christian Tellechea}, \tkzimp{Nicolas Kisselhoff}, \tkzimp{Wolfgang Büchel},

\tkzimp{John Kitzmiller}, \tkzimp{Dimitri Kapetas}, \tkzimp{Gaétan Marris}, \tkzimp{Mark Wibrow}, \tkzimp{Yves Combe}, \tkzimp{Paul Gaborit}, \tkzimp{Laurent Van Deik} 和 \tkzimp{Muzimuzhi Z}。

\vspace*{12pt}
\lefthand\ 我还要感谢 Eric Weisstein，MathWorld 的创建者：
\href{http://mathworld.wolfram.com/about/author.html}{MathWorld}。

\vspace*{12pt}
\lefthand\ 你可以在我的网站上找到一些示例：
\href{http://altermundus.fr}{altermundus.fr}。 \hspace{2cm} 正在建设中！

\vfill
如有关于本手册的错别字或其他意见，请发送至 \href{mailto:al.ma@mac.com}{\textcolor{blue}{Alain Matthes}}。

本文件可按照 \LaTeX{} 项目公共许可证的条款进行再分发和/或修改，该许可证从 \href{http://www.ctan.org/}{CTAN}\  档案库分发。

\clearpage
\setcounter{tocdepth}{2}% sections and subsections; subsubsections stay in the PDF bookmarks
\tableofcontents
\clearpage
\newpage

%\parindent=2em
\part{入门}
\section{介绍与概览}

\begin{tkzexample}[latex=5cm,small]
  \begin{tikzpicture}[scale=.25]
  \tkzDefPoints{0/0/A,12/0/B,6/12*sind(60)/C}
  \foreach \density in {20,30,...,240}{%
    \tkzDrawPolygon[fill=teal!\density](A,B,C)
    \pgfnodealias{X}{A}
    \tkzDefPointWith[linear,K=.15](A,B) \tkzGetPoint{A}
    \tkzDefPointWith[linear,K=.15](B,C) \tkzGetPoint{B}
    \tkzDefPointWith[linear,K=.15](C,X) \tkzGetPoint{C}}
  \end{tikzpicture}
\end{tkzexample}

\vspace*{12pt}

\subsection{为什么需要 \tkzname{\tkznameofpack}？}
我的最初目标是为其他数学教师和我自己提供一个工具，以便快速创建欧几里得几何图形，而无需花费太多精力学习一门新的编程语言。
\tkzname{\tkznameofpack} 面向使用 \LATEX 的数学教师：它让在 \LATEX 中轻松创建正确的绘图成为可能。

最简单的方法似乎是重现手工作图的方式。
要描述一个作图过程，你当然必须定义对象，还要定义你所执行的操作。在我看来，一种接近数学家及其学生语言的语法会更容易理解，而且这种语法也应该接近 \LaTeX 的语法。
对象当然是点、线段、直线、三角形、多边形和圆。至于操作，我认为五种就足够了：定义、创建、绘制、标记和标注。

语法也许略显冗长，但我相信它很容易学习。
因此，学生和教师都能轻松使用这个工具。

\subsection{\tkzname{\TIKZ} 与 \tkzname{\tkznameofpack}}

我喜欢用 \TIKZ 编程，如果没有 \TIKZ，我绝不会想到创建 \tkzname{\tkznameofpack}。请不要忘记，\TIKZ 在幕后工作，你随时可以插入 \TIKZ 代码：\tkzname{\tkznameofpack} 并不妨碍你使用 \TIKZ。
话虽如此，我不认为混合语法是一件好事。

没有必要比较 \TIKZ 和 \tkzname{\tkznameofpack}：它们面向的受众不同。前者能让你做很多事情，后者只绘制几何图形。前者可以做到后者所做的一切，但后者会更轻松地完成你想要的事情。

主要目的是定义点以创建几何图形。\tkzname{\tkznameofpack} 允许你从这些点出发绘制欧几里得几何的基本对象，但对于某些操作（如给面着色）可能不够用。在这种情况下，你随时可以使用 \TIKZ。

下面是基于 \PGF 手册中几何示例对 \tkzname{\TIKZ} 和 \tkzname{\tkznameofpack} 的一些比较。
古希腊人用来构造几何图形和角的两个主要工具是圆规和直尺。我的想法是让你尽可能自然地一步一步跟随用这两个工具手工作图的过程。

\subsubsection{第一卷，命题 I，\_几何原本\_ }

\begin{tikzpicture}
\node [mybox,title={第一卷，命题 I  \_几何原本\_}] (box){%
    \begin{minipage}{0.90\textwidth}
{\emph{在一条给定的有限直线上作一个等边三角形。}
} 
    \end{minipage}
};
\end{tikzpicture}% 


解释：

\emph{PgfManual} 的第四个教程是关于几何作图的。\emph{T. Tantau} 用他出色的工具 Ti\emph{k}Z 得到了这个图。这里我用 \emph{tkz-euclide} 给出同样的作图。Ti\emph{k}Z 代码以绿色显示，\emph{tkz-euclide} 代码以红色显示。

\medskip

\vbox{\color{green!50!black} |\usepackage{tikz}|\\
|\usetikzlibrary{calc,intersections,through,backgrounds}|}

\medskip
\vbox{\color{red} |\usepackage{tkz-euclide}|}

\medskip
我们如何得到直线 $AB$？我们使用两个固定点。\\

\medskip
\vbox{\color{green!50!black} 
|\coordinate [label=left:$A$] (A) at (0,0);|\\
|\coordinate [label=right:$B$] (B) at (1.25,0.25);|\\
|\draw (A) -- (B);|}

\medskip
\vbox{\color{red}
|\tkzDefPoint(0,0){A}|\\
|\tkzDefPoint(1.25,0.25){B}|\\
|\tkzDrawSegment(A,B)|\\
|\tkzLabelPoint[left](A){$A$}|\\
|\tkzLabelPoint[right](B){$B$}|}

我们想在每个点 $A$ 和 $B$ 周围画一个圆，其半径为线段 $AB$ 的长度。
\medskip

\vbox{\color{green!50!black}
|\draw let \p1 = ($ (B) - (A) $),|\\
|\n2 = {veclen(\x1,\y1)} in|\\
|          (A) circle (\n2)|\\
|          (B) circle (\n2);|}

\medskip
\vbox{\color{red} 
|\tkzDrawCircles(A,B B,A)|
}

圆 $\mathcal{D}$ 和 $\mathcal{E}$ 的交点：

\medskip

\vbox{\color{green!50!black} 
|draw [name path=A--B] (A) -- (B);|\\
|node (D) [name path=D,draw,circle through=(B),label=left:$D$] at (A) {}; |\\
|node (E) [name path=E,draw,circle through=(A),label=right:$E$] at (B) {};|\\
|path [name intersections={of=D and E, by={[label=above:$C$]C,[label=below:$C'$]C'}}]; |\\
|draw [name path=C--C',red] (C) -- (C');|\\
|path [name intersections={of=A--B and C--C',by=F}];|\\
|node [fill=red,inner sep=1pt,label=-45:$F$] at (F) {};|\\}

\medskip
\vbox{\color{red} |\tkzInterCC(A,B)(B,A) \tkzGetPoints{C}{X}|\\}


如何绘制点：

\medskip
\vbox{\color{green!50!black} |\foreach \point in {A,B,C}|\\
|\fill [black,opacity=.5] (\point) circle (2pt);|\\}

\medskip
\vbox{\color{red}| \tkzDrawPoints[fill=gray,opacity=.5](A,B,C)|\\}

\subsubsection{使用 \pkg{tkz-euclide} 的完整代码}

我们需要定义一些颜色：

|\colorlet{input}{red!80!black} |\\
|\colorlet{output}{red!70!black}|\\
|\colorlet{triangle}{green!50!black!40}  |

\begin{tkzexample}[vbox,small]
  \colorlet{input}{red!80!black} 
  \colorlet{output}{red!70!black}
  \colorlet{triangle}{green!50!black!40}
  \begin{tikzpicture}[scale=1.25,thick,help lines/.style={thin,draw=black!50}]
  \tkzDefPoint(0,0){A}     
  \tkzDefPoint(1.25+rand(),0.25+rand()){B}      
  \tkzInterCC(A,B)(B,A) \tkzGetPoints{C}{X}

  \tkzFillPolygon[triangle,opacity=.5](A,B,C)
  \tkzDrawSegment[input](A,B) 
  \tkzDrawSegments[red](A,C B,C)  
  \tkzDrawCircles[help lines](A,B B,A)
  \tkzDrawPoints[fill=gray,opacity=.5](A,B,C)
  
  \tkzLabelPoints(A,B)
  \tkzLabelCircle[below=12pt](A,B)(180){$\mathcal{D}$}
  \tkzLabelCircle[above=12pt](B,A)(180){$\mathcal{E}$}
  \tkzLabelPoint[above,red](C){$C$}
      
  \end{tikzpicture}
\end{tkzexample}

\subsubsection{第一卷，命题 II，\_几何原本\_}

\begin{tikzpicture}
\node [mybox,title={第一卷，命题 II  \_几何原本\_}] (box){%
\begin{minipage}{0.90\textwidth}
  {\emph{把一条等于给定直线的直线放置，使其一端位于给定点。}}
\end{minipage}
};
\end{tikzpicture}% 

解释

在第一部分中，我们需要线段 $AB$ 的中点。使用 \TIKZ 时，我们可以使用库 \tkzname{calc}：

\medskip
\vbox{\color{green!50!black} |\coordinate [label=left:$A$] (A) at (0,0);|\\
|\coordinate [label=right:$B$] (B) at (1.25,0.25);|\\
|\draw (A) -- (B);|\\
|\node [fill=red,inner sep=1pt,label=below:$X$] (X) at ($ (A)!.5!(B) $) {};|\\}

使用 \pkg{tkz-euclide}，我们有宏 \tkzcname{tkzDefMidPoint}；我们用 \tkzcname{tkzGetPoint} 得到点 $X$，但下一步不需要这个点。


\medskip
\vbox{\red |\tkzDefPoints{0/0/A,0.75/0.25/B,1/1.5/C}|\\  
|\tkzDefMidPoint(A,B) \tkzGetPoint{X}|}

\medskip
然后我们需要构造一个等边三角形。用 \pkg{tkz-euclide} 很容易。用 TikZ 则需要更多功夫，因为你需要中点 $X$ 来通过三角计算得到点 $D$。

\medskip
\vbox{\color{green!50!black}
|\node [fill=red,inner sep=1pt,label=below:$X$] (X) at ($ (A)!.5!(B) $) {}; | \\
|\node [fill=red,inner sep=1pt,label=above:$D$] (D) at                      |  \\
|($ (X) ! {sin(60)*2} ! 90:(B) $) {};                                       |  \\
|\draw (A) -- (D) -- (B);                                                   |  \\
}                                                                           

\medskip
\vbox{\color{red} |\tkzDefTriangle[equilateral](A,B) \tkzGetPoint{D}|}

我们可以在图的末尾用以下命令绘制三角形：

\medskip
\vbox{\color{red} |\tkzDrawPolygon{A,B,C}|}

\medskip
我们知道如何绘制以 $B$ 为圆心、经过 $C$ 的圆 $\mathcal{H}$，以及如何放置点 $E$ 和 $F$：

\medskip
\vbox{\color{green!50!black} 
|\node (H) [label=135:$H$,draw,circle through=(C)] at (B) {};|          \\
|\draw (D) -- ($ (D) ! 3.5 ! (B) $) coordinate [label=below:$F$] (F);|  \\
|\draw (D) -- ($ (D) ! 2.5 ! (A) $) coordinate [label=below:$E$] (E);|} 

\medskip

\vbox{\color{red} |\tkzDrawCircle(B,C)|\\
|\tkzDrawLines[add=0 and 2](D,A D,B)|}

\medskip
我们可以在图的末尾放置点 $E$ 和 $F$。现在不需要它们。

求直线与圆的交点：这里我们寻找以 $B$ 为圆心、经过 $C$ 的圆与直线 $DB$ 的交点。
无限直线 $DB$ 与圆相交，但使用 \TIKZ 时我们需要延长线段 $DB$，这可以通过 partway 计算来完成。我们得到点 $F$，而 $BF$ 或 $DF$ 与圆相交：

\medskip
\vbox{\color{green!50!black}| \node (H) [label=135:$H$,draw,circle through=(C)] at (B) {}; |  \\
|\path let \p1 = ($ (B) - (C) $) in|                                     \\
|  coordinate [label=left:$G$] (G) at ($ (B) ! veclen(\x1,\y1) ! (F) $); |  \\
|\fill[red,opacity=.5] (G) circle (2pt);|}                                

\medskip
至于两个圆的交点，用 \pkg{tkz-euclide} 求直线与圆的交点也很容易。我们不需要 $F$：

\medskip
\vbox{\color{red} | \tkzInterLC(B,D)(B,C)\tkzGetFirstPoint{G}|}

\medskip
没有更多困难了。下面是经过一些简化的最终代码。
我们绘制以 $D$ 为圆心、经过 $G$ 的圆 $\mathcal{K}$。它与直线 $AD$ 交于点 $L$。$AL = BC$。

\vbox{\color{red} | \tkzDrawCircle(D,G)|}
\vbox{\color{red} | \tkzInterLC(D,A)(D,G)\tkzGetSecondPoint{L}|}

\begin{tkzexample}[latex=7cm,small]
\begin{tikzpicture}[scale=1.5]
\tkzDefPoint(0,0){A}
\tkzDefPoint(0.75,0.25){B}  
\tkzDefPoint(1,1.5){C} 
\tkzDefTriangle[equilateral](A,B)\tkzGetPoint{D}
\tkzInterLC[near](D,B)(B,C)      \tkzGetSecondPoint{G}
\tkzInterLC[near](A,D)(D,G)      \tkzGetFirstPoint{L}
\tkzDrawCircles(B,C D,G)
\tkzDrawLines[add=0 and 2](D,A D,B)
\tkzDrawSegment(A,B) 
\tkzDrawSegments[red](A,L B,C) 
\tkzDrawPoints[red](D,L,G)
\tkzDrawPoints[fill=gray](A,B,C)
\tkzLabelPoints[left,red](A)
\tkzLabelPoints[below right,red](L)
\tkzLabelCircle[above](B,C)(20){$\mathcal{(H)}$}
\tkzLabelPoints[above left](D)
\tkzLabelPoints[above](G)
\tkzLabelPoints[above,red](C)
\tkzLabelPoints[right,red](B)
\tkzLabelCircle[below](D,G)(-90){$\mathcal{(K)}$}
\end{tikzpicture}
\end{tkzexample}

% 与第 3、4 版的比较现在位于"新特性与兼容性"一节中。

\subsection{如何使用 \tkzname{\tkznameofpack} 宏包？}
\subsubsection{让我们看一个经典示例}
为了展示正确的方式，我们将构建一个等边三角形。有几种方法可行；我们将遵循欧几里得的步骤。

\begin{itemize}
\item   首先，你需要一个文档类。要测试你的代码，最好的选择是使用 \tkzname{standalone}\index{standalone} 类创建一个单独的图。
\begin{verbatim}  
\documentclass{standalone}
\end{verbatim}
\item 然后加载 \tkzname{\tkznameofpack} 宏包：
\begin{verbatim}  
\usepackage{tkz-euclide} or \usepackage[lua]{tkz-euclide}
\end{verbatim}

 你不需要加载 \TIKZ：\tkzname{\tkznameofpack} 宏包在 TikZ 之上工作并会加载它。

 \item 开始文档并打开一个 TikZ picture 环境：
\begin{verbatim}
\begin{document}
\begin{tikzpicture}
\end{verbatim}

\item 现在我们定义两个固定点：
\begin{verbatim}
\tkzDefPoint(0,0){A}
\tkzDefPoint(5,2){B}
\end{verbatim}

\item 两个点定义两个圆；我们来使用它们：

 以 $A$ 为圆心经过 $B$ 的圆和以 $B$ 为圆心经过 $A$ 的圆。这两个圆有两个公共点。
\begin{verbatim}
\tkzInterCC(A,B)(B,A)
\end{verbatim}
我们用以下命令得到交点
\begin{verbatim}
\tkzGetPoints{C}{D}
\end{verbatim}

\item 所有必要的点现在都已定义；我们可以进入最后的步骤，包括绘图。
\begin{verbatim}
\tkzDrawCircles[gray,dashed](A,B B,A)
\tkzDrawPolygon(A,B,C)% The triangle
\end{verbatim}
\item 绘制所有点 $A$、$B$、$C$ 和 $D$：
\begin{verbatim}
\tkzDrawPoints(A,...,D)
\end{verbatim}

\item 在最后一步中，我们为点添加标签并使用选项来定位它们：\\
\begin{verbatim}
\tkzLabelSegments[swap](A,B){$c$}
\tkzLabelPoints(A,B,D)
\tkzLabelPoints[above](C)
\end{verbatim}
\item 最后，我们关闭两个环境：
\begin{verbatim}
\end{tikzpicture}
\end{document}
\end{verbatim}

\item 完整代码：

\begin{tkzexample}[latex=8cm,small]
 \begin{tikzpicture}[scale=.5]
   % fixed points
  \tkzDefPoint(0,0){A}
  \tkzDefPoint(5,2){B}
  % calculus
  \tkzInterCC(A,B)(B,A)
  \tkzGetPoints{C}{D}
  % drawings
  \tkzDrawCircles(A,B B,A)
  \tkzDrawPolygon(A,B,C)
  \tkzDrawPoints(A,...,D)
  % marking
  \tkzMarkSegments[mark=s||](A,B B,C C,A)
  % labelling
  \tkzLabelSegments[swap](A,B){$c$}
  \tkzLabelPoints(A,B,D)
  \tkzLabelPoints[above](C)
\end{tikzpicture}
\end{tkzexample}

 \end{itemize}

\subsubsection{第一部分：黄金三角形}
\begin{center}
\begin{tikzpicture}
  
\tkzDefPoint(0,0){C} % 可以用 \tkzDefPoint[label=below:$C$](0,0){C} 但不要这样做
\tkzDefPoint(2,6){B}
% 我们通过旋转得到 D 和 E
\tkzDefPointBy[rotation= center B angle 36](C) \tkzGetPoint{D} 
\tkzDefPointBy[rotation= center B angle 72](C) \tkzGetPoint{E} 
% 为了得到 A，我们使用直线的交点
\tkzInterLL(B,E)(C,D) \tkzGetPoint{A}
\tkzInterLL(C,E)(B,D) \tkzGetPoint{H}

% 角
\tkzMarkAngles[size=2](C,B,D E,A,D) %这是为了绘制弧
\tkzLabelAngles[pos=1.5](C,B,D E,A,D){$\alpha$}
\tkzMarkRightAngle(B,H,C)
\tkzDrawPoints(A,...,E)

% 绘图
\tkzDrawArc[delta=10](B,C)(E)
\tkzDrawPolygon(C,B,D)
\tkzDrawSegments(D,A B,A C,E)

% 现在才添加标签
\tkzLabelPoints[below left](C,A)
\tkzLabelPoints[below right](D)
\tkzLabelPoints[above](B,E)
\end{tikzpicture}
\end{center}

让我们分析这个图：
\begin{enumerate}
  \item $CBD$ 和 $DBE$ 是等腰三角形；
  
  \item $BC=BE$，且 $(BD)$ 是角 $CBE$ 的平分线；
  
  \item 由此我们推出角 $CBD$ 和 $DBE$ 相等；设 $\alpha$ 为它们的度量。
   \[\widehat{BAC} +\widehat{ABC} + \widehat{BCA}=180^\circ \ \text{在三角形中}\ BAC \]
   \[3\alpha + \widehat{BCA}=180^\circ\  \text{在三角形中}\ CBD\]
   于是
     \[\alpha + 2\widehat{BCA}=180^\circ \] 
   或者
     \[\widehat{BCA}=90^\circ -\alpha/2 \] 
    
    \item 最终   \[\widehat{CBD}=\alpha=36^\circ \]
     三角形 $CBD$ 是一个 \code{golden} 三角形。
\end{enumerate}

\vspace*{24pt}
我们如何构造一个黄金三角形，或者一个 $36^\circ$ 的角？

\begin{enumerate}
  \item 我们放置固定点 $C$ 和 $D$。|\tkzDefPoint(0,0){C}| 和 |\tkzDefPoint(4,0){D}|；
  \item 我们构造一个正方形 $CDef$，并构造 $[Cf]$ 的中点 $m$；
  
 我们可以用圆规和直尺完成所有这些；
  \item 然后我们画一个以 $m$ 为圆心、经过 $e$ 的弧。这个弧与直线 $(Cf)$ 交于 $n$；
  \item 现在以 $C$ 和 $D$ 为圆心、半径为 $Cn$ 的两个弧定义点 $B$。
\end{enumerate}

\begin{tkzexample}[latex=7cm,small]
\begin{tikzpicture}
  \tkzDefPoint(0,0){C}
  \tkzDefPoint(4,0){D}
  \tkzDefSquare(C,D)                     
  \tkzGetPoints{e}{f}
  \tkzDefMidPoint(C,f)                   
  \tkzGetPoint{m}
  \tkzInterLC(C,f)(m,e)                  
  \tkzGetSecondPoint{n}
  \tkzInterCC[with nodes](C,C,n)(D,C,n) 
  \tkzGetFirstPoint{B}
  \tkzDrawSegment[brown,dashed](f,n)
  \pgfinterruptboundingbox% from tikz
  \tkzDrawPolygon[brown,dashed](C,D,e,f)
  \tkzDrawArc[brown,dashed](m,e)(n)
  \tkzCompass[brown,dashed,delta=20](C,B)
  \tkzCompass[brown,dashed,delta=20](D,B)
  \endpgfinterruptboundingbox 
  \tkzDrawPolygon(B,...,D)
  \tkzDrawPoints(B,C,D,e,f,m,n)
  \tkzLabelPoints[above](B)
  \tkzLabelPoints[left](f,m,n)
  \tkzLabelPoints(C,D)
  \tkzLabelPoints[right](e)
\end{tikzpicture}
\end{tkzexample}


在构建黄金三角形 $BCD$ 之后，我们通过注意到 $BD=DA$ 来构建点 $A$。然后我们得到点 $E$，最后得到点 $F$，作为已定义对象（一条直线和一个圆）的交点。


\subsubsection{第二部分：使用黄金三角形和欧几里得三角形的另外两种方法}

\tkzname{\tkznameofpack} 知道如何定义 \code{golden} 或 \code{euclide} 三角形。我们可以将 $BCD$ 和 $BCA$ 定义为黄金三角形。


  \begin{center}
    \begin{tkzexample}[code only,small]
      \begin{tikzpicture}
        \tkzDefPoint(0,0){C}
        \tkzDefPoint(4,0){D}
        \tkzDefTriangle[golden](C,D)
        \tkzGetPoint{B}
        \tkzDefTriangle[golden](B,C)
        \tkzGetPoint{A}
        \tkzInterLC[near](A,B)(B,D) \tkzGetFirstPoint{E}
        \tkzInterLL(B,D)(C,E) \tkzGetPoint{F}
        \tkzDrawPoints(C,D,B)
        \tkzDrawPolygon(B,...,D)  
        \tkzDrawPolygon(B,C,D)
        \tkzDrawSegments(D,A A,B C,E)
        \tkzDrawArc[delta=10](B,C)(E)
        \tkzDrawPoints(A,...,F) 
        \tkzMarkRightAngle(B,F,C)  
        \tkzMarkAngles(C,B,D E,A,D)
        \tkzLabelAngles[pos=1.5](C,B,D E,A,D){$\alpha$} 
        \tkzLabelPoints[below](A,C,D,E)
        \tkzLabelPoints[above right](B,F)
      \end{tikzpicture} 
    \end{tkzexample}
  \end{center}

下面是使用旋转的最后一种方法：

\begin{center}
  \begin{tkzexample}[code only,small]
  \begin{tikzpicture} 
  \tkzDefPoint(0,0){C} % possible 
  % \tkzDefPoint[label=below:$C$](0,0){C} 
  % but don't do this
  \tkzDefPoint(2,6){B}
  % We get D and E with a rotation
  \tkzDefPointBy[rotation= center B angle 36](C) \tkzGetPoint{D} 
  \tkzDefPointBy[rotation= center B angle 72](C) \tkzGetPoint{E} 
  % To get A we use an intersection of lines
  \tkzInterLL(B,E)(C,D) \tkzGetPoint{A}
  \tkzInterLL(C,E)(B,D) \tkzGetPoint{H}
  % drawing
  \tkzDrawArc[delta=10](B,C)(E)
  \tkzDrawPolygon(C,B,D)
  \tkzDrawSegments(D,A B,A C,E)
  % angles 
  \tkzMarkAngles(C,B,D E,A,D) %this is to draw the arcs
  \tkzLabelAngles[pos=1.5](C,B,D E,A,D){$\alpha$}
  \tkzMarkRightAngle(B,H,C)
  \tkzDrawPoints(A,...,E)
  % Label only now
  \tkzLabelPoints[below left](C,A)
  \tkzLabelPoints[below right](D)
  \tkzLabelPoints[above](B,E)
  \end{tikzpicture}
  \end{tkzexample}
\end{center}


\subsubsection{完整但最小的示例}


选定一个长度单位后，本示例展示如何用直尺和圆规从长度为 $a$ 的线段得到长度为 $\sqrt{a}$ 的线段。

$IB=a$，$AI=1$

\vspace{12pt}
\hypertarget{firstex}{}
\begin{tkzexample}[vbox,small]
\begin{tikzpicture}[scale=1,ra/.style={fill=gray!20}]
   % fixed points
   \tkzDefPoint(0,0){A}
   \tkzDefPoint(1,0){I}
   % calculation
   \tkzDefPointBy[homothety=center A ratio  10 ](I) \tkzGetPoint{B}  
   \tkzDefMidPoint(A,B)              \tkzGetPoint{M}
   \tkzDefPointWith[orthogonal](I,M) \tkzGetPoint{H}
   \tkzInterLC(I,H)(M,B)             \tkzGetFirstPoint{C}
   \tkzDrawSegment[style=purple](I,C)
   \tkzDrawArc(M,B)(A)
   \tkzDrawSegment[dim={$1$,-16pt,}](A,I)
   \tkzDrawSegment[dim={$(a-1)/2$,-10pt,}](I,M)
   \tkzDrawSegment[dim={$(a+1)/2$,-16pt,}](M,B)   
   \tkzMarkRightAngle[ra](A,I,C)
   \tkzDrawPoints(I,A,B,C,M)  
   \tkzLabelPoint[left](A){$A(0,0)$} 
   \tkzLabelPoints[above right](I,M)
   \tkzLabelPoints[above left](C)
   \tkzLabelPoint[right](B){$B(10,0)$}
   \tkzLabelSegment[right=4pt](I,C){$\sqrt{a^2}=a \ (a>0)$}
\end{tikzpicture}
\end{tkzexample}

\emph{注释}
 
\begin{itemize}
\item 导言区


 让我们先看导言区。如果你需要 \tkzname{xcolor}，请在 \tkzname{tkz-euclide} 之前加载它，即在 \TIKZ 之前。\TIKZ 可能会在活动字符方面引起问题，但它
 提供了 \NameLib{babel} 库，旨在解决这些问题。
 
\begin{tkzltxexample}[]
\documentclass{standalone} % or another class
   % \usepackage{xcolor} % before tikz or tkz-euclide if necessary
\usepackage{tkz-euclide} % no need to load TikZ
   % \usetkzobj{all}  is no longer necessary 
   % \usetikzlibrary{babel}  if there are problems with the active characters
\end{tkzltxexample}

以下代码由几个部分组成：

   \item 固定点的定义：第一部分定义作图所需的点；这些是固定点。在大多数情况下，宏 \tkzcname{tkzInit} 和 \tkzcname{tkzClip} 不是必需的。

\begin{tkzltxexample}[]
  \tkzDefPoint(0,0){A}
  \tkzDefPoint(1,0){I}
\end{tkzltxexample}
 
  \item 第二部分从固定点创建新点：
 点 $B$ 放置在距 $A$ $10$~cm 处。$[AB]$ 的中点是 $M$，然后我们定义在点 $I$ 处与 $(AB)$ 正交的直线。然后我们寻找这条直线与以 $M$ 为圆心、经过 $A$ 的半圆的交点。
  
\begin{tkzltxexample}[]
   \tkzDefPointBy[homothety=center A ratio  10 ](I)
      \tkzGetPoint{B}
   \tkzDefMidPoint(A,B)
      \tkzGetPoint{M}
   \tkzDefPointWith[orthogonal](I,M)
      \tkzGetPoint{H}
   \tkzInterLC(I,H)(M,B)             
   \tkzGetSecondPoint{C}
 \end{tkzltxexample}  
     

 \item 第三部分包含绘图；
 \begin{tkzltxexample}[]
   \tkzDrawSegment[style=purple](I,H)
   \tkzDrawPoints(O,I,A,B,M)
   \tkzDrawArc(M,A)(O)
   \tkzDrawSegment[dim={$1$,-16pt,}](A,I)
   \tkzDrawSegment[dim={$a/2$,-10pt,}](I,M)
   \tkzDrawSegment[dim={$a/2$,-16pt,}](M,B)
 \end{tkzltxexample}
 
\item 标记：第四部分用于标记；


\begin{tkzltxexample}[]
 \tkzMarkRightAngle[ra](A,I,C)
 \end{tkzltxexample}
 
 \item 标注：最后一部分只放置标签。
\begin{tkzltxexample}[]
   \tkzLabelPoint[left](A){$A(0,0)$} 
   \tkzLabelPoint[right](B){$B(10,0)$}
   \tkzLabelSegment[right=4pt](I,C){$\sqrt{a^2}=a \ (a>0)$}
\end{tkzltxexample}

\end{itemize}

\endinput

\section{安装}

\tkzname{\tkznameofpack} 可在 \tkzname{CTAN}\footnote{\tkzname{\tkznameofpack} 是 \NameDist{TeXLive} 的一部分，可由 \tkzname{tlmgr} 安装，在 \NameSys{Windows} 下也是 \NameDist{MiKTeX} 的一部分。} 上获取。如果你想测试 beta 版本，请将文件放在系统能找到的 texmf 文件夹中。
请检查以下几点：

\begin{itemize}\setlength{\itemsep}{5pt}
\item \tkzname{\tkznameofpack} 文件夹必须位于 \tkzname{latex} 搜索的路径中。
\item \tkzname{\tkznameofpack} 使用宏包 \tkzNamePack{xfp}。
\item \PGF 必须安装在你的计算机上。\tkzname{\tkznameofpack} 使用 \TIKZ 的以下几个库：

 \begin{tabular}{l}
    angles,                         \\
    arrows,                         \\
    arrows.meta,                    \\
    backgrounds,                    \\
    calc,                           \\
    decorations,                    \\
    decorations.markings,           \\
    decorations.pathreplacing,      \\
    decorations.shapes,             \\
    decorations.text,               \\
    decorations.pathmorphing,       \\
    intersections,                  \\
    math,                           \\
    plotmarks,                      \\
    positioning,                    \\
    quotes,                         \\
    shapes.misc,                    \\
    through
\end{tabular}

\item 本手册使用 \tkzname{lualatex} 编译，但也可以使用 \tkzname{pdflatex} 和 \tkzname{xelatex}。
\end{itemize}

\endinput

\section{tkz 代码的基本元素}

要使用这个宏包，你需要一些 \LATEX 和 \TIKZ 的知识。

在本节中，我们来看用 \tkzname{\tkznameofpack} 创建图形所使用的 \code{rules} 和 \code{symbols}。

\subsection{对象与语言}

 基本对象是点。你可以在任何时候用定义点时所给的名称来引用它（之后也可以给它另一个名称）。

新的点通过宏获得。\tkzname{\tkznameofpack} 宏的名称以 \code{tkz} 开头。主要有四类，以以下开头：
|\tkzDef...| |\tkzDraw...| |\tkzMark...| 和 |\tkzLabel...|。
使用的点作为参数放在圆括号中传入，而创建的点放在花括号中给出。

图形的代码放在 \tkzimp{tikzpicture} 环境中。


与 \TIKZ 不同，宏不能以“;”结尾。因此我们失去了 \tkzimp{path} 这一重要概念。不过，你可以在 \tkzname{\tkznameofpack} 宏之间放置一些 \TIKZ 代码。


在第一类中，|\tkzDefPoint| 定义固定点；稍后会详细研究。这里我们来看宏 |\tkzDefTriangle|。

这个宏将第三个点与一对点关联，以定义某个三角形：|\tkzDefTriangle(A,B)|。得到的点引用为 |tkzPointResult|，你可以选择另一个引用，例如用 |\tkzGetPoint{C}|。

|\tkzDefTriangle[euclid](A,B) \tkzGetPoint{C}|

圆括号用于传递参数：在 |(A,B)| 中，$A$ 和 $B$ 是用于定义第三个点的点。在 |{C}| 中，花括号用于获取新点。

要从以下类型中选择一种三角形：
  |equilateral|,  |isosceles right|, |half|, |pythagoras|, |school|, |golden or sublime|, |euclid|, |gold|, |cheops|...
 和 |two angles|，你只需在方括号中给出选项，例如：
 
\begin{minipage}{0.5\textwidth}
  \begin{tikzpicture}[scale=.5]
  \tkzDefPoints{0/0/A,8/0/B}
  \foreach \tr in {golden, equilateral}
  {\tkzDefTriangle[\tr](A,B) \tkzGetPoint{C}
  \tkzDrawPoint(C)
  \tkzLabelPoint[right](C){\tr}
  \tkzDrawSegments(A,C C,B)}
  \tkzDrawPoints(A,B)
  \tkzDrawSegments(A,B)
   \tkzLabelPoints(A,B)
  \end{tikzpicture}
\end{minipage}
\begin{minipage}{0.5\textwidth}
  \begin{tkzexample}[code only,small]
    \begin{tikzpicture}[scale=.5]
    \tkzDefPoints{0/0/A,8/0/B}
   \foreach \tr in {golden, equilateral}
    {\tkzDefTriangle[\tr](A,B) \tkzGetPoint{C}
    \tkzDrawPoint(C)
    \tkzLabelPoint[right](C){\tr}
    \tkzDrawSegments(A,C C,B)}
    \tkzDrawPoints(A,B)
    \tkzDrawSegments(A,B)
    \tkzLabelPoints(A,B)
    \end{tikzpicture}
  \end{tkzexample}
\end{minipage}

\subsection{记号与约定}

我特意选择了法国几何约定以及一些个人约定来描述所表示的对象。\tkzname{\tkznameofpack} 定义和绘制的对象是平面中的点、直线和圆。它们是欧几里得几何的基本对象，我们由此构建图形。

根据 \tkzimp{Euclid} 的观点，这些图形仅说明我们心灵所产生的纯粹概念。
因此点没有维度，因而没有真实的存在。同样，直线没有宽度，因而在现实世界中没有存在。我们所考虑的对象只是理想数学对象的表示。\tkzname{\tkznameofpack} 遵循古希腊人的步骤，用直尺和圆规获得几何作图。

以下是使用的记号：

\begin{itemize}
\item 点用小圆点或两条直线的交点（两条直线、一条直线和一个圆，或两个圆）来表示。在第二种情况下，点用十字表示。

\begin{tkzexample}[latex=6cm, small]     
  \begin{tikzpicture}       
    \tkzDefPoints{0/0/A,4/2/B}       
    \tkzDrawPoints(A,B)       
    \tkzLabelPoints(A,B)     
  \end{tikzpicture}    
\end{tkzexample}

或者

\begin{tkzexample}[latex=6cm, small]     
  \begin{tikzpicture}       
    \tkzSetUpPoint[shape=cross, color=red]       
    \tkzDefPoints{0/0/A,4/2/B}       
    \tkzDrawPoints(A,B)       
    \tkzLabelPoints(A,B)     
    \end{tikzpicture}    
    \end{tkzexample}  

一旦一个点存在，我们就可以给它一个标签，即拉丁字母表中的一个大写字母（有一些例外），例如 $A$、$B$ 或 $C$。例如：
\begin{itemize}
\item $O$ 是圆的圆心、旋转中心等；
\item $M$ 是中点；
\item $H$ 是高的垂足；
\item $P'$ 是 $P$ 经过变换后的像。
\end{itemize}

请注意，代码中点的引用名称可能与文本中指定它的标签不同：我们可以定义一个点 A 并给它标签 $P$。样式也不同：点 A 标注为 $A$。

\begin{tkzexample}[latex=6cm, small]     
  \begin{tikzpicture}       
    \tkzDefPoint(0,0){A}       
    \tkzDrawPoints(A)       
    \tkzLabelPoint(A){$P$}     
  \end{tikzpicture}    
\end{tkzexample}

例外：某些点，例如三角形各边的中点，共享一个性质，因此它们的名称也共享一个共同的字母。我们用 $M_a$、$M_b$ 和 $M_c$ 或 $M_A$、$M_B$ 和 $M_C$ 来表示这些点。

在代码中，这些点称为 M\_A、M\_B 和 M\_C。

另一个例外涉及中间作图点，它们不标注。它们在代码中通常用小写字母表示。

\item 线段用其两个端点放在方括号中来表示：$[AB]$。

\item 在欧几里得几何中，一条直线由两个点定义：$A$ 和 $B$ 定义直线 $(AB)$。我们也可以用希腊字母为直线命名，如 $(\delta)$ 或 $(\Delta)$，或用小写字母如 $d$ 或 $d'$。

\item 射线（半直线）写作 $[AB)$。

\item 直线之间的关系：如果 $(AB)$ 和 $(CD)$ 垂直，我们写作 $(AB) \perp (CD)$；如果它们平行，我们写作 $(AB) \parallelslant (CD)$。

\item 三角形 $ABC$ 各边的长度为 $AB$、$AC$ 和 $BC$。数也用小写字母表示，因此我们写作 $AB=c$、$AC=b$ 和 $BC=a$。字母 $a$ 也用于表示角；$r$ 常用于表示半径，$d$ 表示直径或距离，$l$ 表示长度。

\item 多边形用其顶点表示：$ABC$ 是三角形，$EFGH$ 是四边形。

\item 角一般以度为单位度量（例如 $60^\circ$）；在等边三角形 $ABC$ 中，我们写作 $\widehat{ABC}=\widehat{B}=60^\circ$。

\item 弧用其端点表示：如果 $A$ 和 $B$ 是一个圆上的两个点，我们写作 $\widearc{AB}$。

\item 如果不会混淆，圆写作 $\mathcal{C}$；$\mathcal{C}$ $(O~;~A)$ 表示以 $O$ 为圆心、经过点 $A$ 的圆；$\mathcal{C}$ $(O~;~1)$ 表示以 $O$ 为圆心、半径为 1 cm 的圆。

\item 三角形特殊直线的名称：我使用角平分线、外角平分线、中垂线（通常称为垂直平分线）、高、中线和陪中线这些术语。

\item ($x_1$,$y_1$) 是点 $A_1$ 的坐标，($x_A$,$y_A$) 是点 $A$ 的坐标。
\end{itemize}

\subsection{\tkzname{Set, Calculate, Draw, Mark, Label}}
本节也可以称为 \texttt{计算与绘图的分离}。

当用 \LATEX 准备文档时，其源代码分为两部分：导言区和文档正文。
用这种方法，可以以最少的工作量对文档进行结构化、设置样式和排版。
我提出一种类似的方法，用 \tkzname{\tkznameofpack} 创建图形。

第一部分定义固定点，第二部分创建新点：\tkzname{Set and Calculate} 是两个主要部分。然后剩下要做的就是绘制（或填充）、标记和标注。\tkzname{\tkznameofpack} 对于最后这些操作中的某些可能不够用，但你随时可以使用 \TIKZ。

最后一点我认为很重要的意见：尽量避免在代码中使用坐标。坐标只应在开头与固定点一起出现；此后使用点的引用。大多数宏都有 \tkzname{nodes} 或 \tkzname{with nodes} 选项。

最好在开头定义不同对象的样式。
\endinput

\section{关于本手册及其示例}

本手册使用 \code{lualatex} 编译。我使用自己的文档类 \tkzname{tkz-doc.cls}，它基于 \tkzname{scrartcl}。

下面是本手册中使用的样式列表。要了解如何使用样式，请参见第 \ref{custom} 节。

|\tkzSetUpColors[background=white,text=black]  |

|\tkzSetUpCompass[color=orange, line width=.2pt,delta=10]|

|\tkzSetUpArc[color=gray,line width=.2pt,delta=10]|

|\tkzSetUpPoint[size=2,color=teal]|

|\tkzSetUpLine[line width=.2pt,color=teal]|

|\tkzSetUpStyle[color=orange,line width=.2pt]{new}|

|\tikzset{every picture/.style={line width=.2pt}}|

|\tikzset{label angle style/.append style={color=teal,font=\footnotesize}}|

|\tikzset{label style/.append style={below,color=teal,font=\scriptsize}}|

有些示例使用预定义的样式，例如：

|\tikzset{new/.style={color=orange,line width=.2pt}}  |


\part{点}
\section{第一步：固定点}

几何作图的第一步是定义固定点，图形将从这些点构建。

总体思路是避免操作坐标，而是使用第一步中固定的点或通过宏包工具获得的点的引用。即使可以直接使用坐标，我认为这也是个坏主意：最好使用命名点。

\tkzname{\tkznameofpack} 使用专用于几何作图的宏和词汇。你当然可以使用 \TIKZ 的工具，但在我看来不混合这两种语法更合乎逻辑。

对于 \TIKZ，\tkzname{\tkznameofpack} 的点是一个特殊的 \code{node}。在下一节中我们将看到如何用坐标定义点。这里不讨论点的样式（颜色和形状）：你会在示例中找到一些提示，并在第 \ref{custom} 节中找到更多信息。


\section{点的定义：\tkzcname{tkzDefPoint} 或 \tkzcname{tkzDefPoints}}

 点可以通过以下任一方式指定：
\begin{itemize}
\item 笛卡尔坐标；
\item 极坐标；
\item 命名点；
\item 相对点。
\end{itemize}


当一个点的名称与唯一一对十进制数关联时，该点就被定义：
 $(x,y)$（$x$ 横坐标，$y$ 纵坐标）或 $(a:d)$（$a$ 角，$d$ 距离）。
 这之所以可能，是因为平面被赋予了一个单位为 $1$~cm 的标准正交笛卡尔坐标系。

 笛卡尔坐标 $(a,b)$ 指向在 $x$ 方向上 $a$ 厘米、在 $y$ 方向上 $b$ 厘米的点。

 极坐标中的一个点需要一个以度为单位的角 $\alpha$，以及到原点的距离 $d$；默认单位是 \texttt{cm}。

 宏 \tkzNameMacro{tkzDefPoint} 通过给点赋值坐标来定义点。这个宏基于

 \tkzNameMacro{coordinate}，这是 \TIKZ 的一个宏，并接受 \TIKZ 选项如 \tkzname{shift}。计算由宏包 \tkzNamePack{xfp} 完成。我们可以使用笛卡尔坐标或极坐标。

\begin{minipage}[t]{0.48\textwidth}
 笛卡尔坐标
\begin{tkzexample}[code only,small]
\begin{tikzpicture}[scale=1] 
  \tkzInit[xmax=5,ymax=5]
  % necessary to limit
  % the size of the axes
  \tkzDrawX[>=latex]
  \tkzDrawY[>=latex]
  \tkzDefPoints{0/0/O,1/0/I,0/1/J}
  \tkzDefPoint(3,4){A} 
  \tkzDrawPoints(O,A) 
  \tkzLabelPoint[above](A){$A_1(x_1,y_1)$} 
  \tkzShowPointCoord[xlabel=$x_1$,
                     ylabel=$y_1$](A) 
  \tkzLabelPoints(O,I)
  \tkzLabelPoints[left](J)
  \tkzDrawPoints[shape=cross](I,J) 
\end{tikzpicture}
\end{tkzexample}%
\end{minipage}
\begin{minipage}[t]{0.45\textwidth}
 极坐标
\begin{tkzexample}[code only,small]
\begin{tikzpicture}[scale=1]
  \tkzInit[xmax=5,ymax=5]
  \tkzDrawX[>=latex]
  \tkzDrawY[>=latex]  
  \tkzDefPoints{0/0/O,1/0/I,0/1/J}
  \tkzDefPoint(40:4){P} 
  \tkzDrawSegment[dim={$d$,
                 16pt,above=6pt}](O,P)
  \tkzDrawPoints(O,P) 
  \tkzMarkAngle[mark=none,->](I,O,P) 
  \tkzFillAngle[opacity=.5](I,O,P) 
  \tkzLabelAngle[pos=1.25](I,O,P){%
                              $\alpha$}
  \tkzLabelPoint[right](P){$P(\alpha:d)$}
  \tkzDrawPoints[shape=cross](I,J)
  \tkzLabelPoints(O,I)
  \tkzLabelPoints[left](J)
\end{tikzpicture}
\end{tkzexample}
\end{minipage}%

\begin{minipage}[b]{0.45\textwidth}
\begin{tikzpicture}[scale=1]
  \tkzInit[xmax=5,ymax=5]
  \tkzDrawX[>=latex]
  \tkzDrawY[>=latex]
  \tkzDefPoints{0/0/O,1/0/I,0/1/J}
  \tkzDefPoint(3,4){A}
  \tkzDrawPoints(O,A)
  \tkzLabelPoint[above](A){$A_1 (x_1,y_1)$}
  \tkzShowPointCoord[xlabel=$x_1$,ylabel=$y_1$](A)
  \tkzLabelPoints(O,I)
  \tkzLabelPoints[left](J)
  \tkzDrawPoints[shape=cross](I,J)
\end{tikzpicture}
\end{minipage}
\begin{minipage}[b]{0.45\textwidth}
\begin{tikzpicture}[,scale=1]
  \tkzInit[xmax=5,ymax=5]
  \tkzDrawX[>=latex]
  \tkzDrawY[>=latex]  
  \tkzDefPoints{0/0/O,1/0/I,0/1/J}
  \tkzDefPoint(40:4){P} 
  \tkzDrawSegment[dim={$d$,
                 16pt,above=6pt}](O,P)
  \tkzDrawPoints(O,P) 
  \tkzMarkAngle[mark=none,->](I,O,P) 
  \tkzFillAngle[opacity=.5](I,O,P) 
  \tkzLabelAngle[pos=1.25](I,O,P){$\alpha$}  
  \tkzLabelPoint[right](P){$P  (\alpha : d )$} 
  \tkzDrawPoints[shape=cross](I,J) 
  \tkzLabelPoints(O,I)
  \tkzLabelPoints[left](J) 
\end{tikzpicture}
\end{minipage}%

\subsection{定义命名点：\tkzcname{tkzDefPoint}}

\begin{NewMacroBox}{tkzDefPoint}{\oarg{local options}\parg{$x,y$}\marg{ref} or \parg{$\alpha$:$d$}\marg{ref}}%
\begin{tabular}{lll}%
参数 & 默认值 & 定义  \\
\midrule
\TAline{($x,y$)}{no default}{$x$ 和 $y$ 是两个度量，默认为 cm。}
\TAline{($\alpha$:$d$)}{no default}{$\alpha$ 是以度为单位的角，$d$ 是一个度量}
\TAline{\{ref\}}{no default}{赋值给该点的引用：$A$、$T\_a$、$P1$ 或 $P_1$}
\bottomrule
\end{tabular}

\medskip
\emph{这个宏的必选参数是两个十进制数：第一种情况是两个长度，第二种情况是一个以度为单位的角和一个长度。不要混淆点的引用和名称：引用供计算使用，但名称通常与引用相同。}

\medskip
\begin{tabular}{lll}%
\toprule
选项             & 默认值 & 定义  \\
\midrule
\TOline{label} {no default} {允许你在预定义距离处放置标签}
\TOline{shift} {no default} {将 $(x,y)$ 或 $(\alpha:d)$ 加到所有坐标上}
\end{tabular}
\end{NewMacroBox}

\subsubsection{笛卡尔坐标}
 
\begin{tkzexample}[latex=6cm,small]
  \begin{tikzpicture}
  \tkzInit[xmax=5,ymax=5] % limits the size of the axes
  \tkzDrawX[>=latex]
  \tkzDrawY[>=latex]  
  \tkzDefPoint(0,0){A}
  \tkzDefPoint(4,0){B}
  \tkzDefPoint(0,3){C} 
  \tkzDrawPolygon(A,B,C)
  \tkzDrawPoints(A,B,C)
  \end{tikzpicture}
\end{tkzexample}

\subsubsection{用 \tkzNamePack{xfp} 进行计算}

 \begin{tkzexample}[latex=7cm,small]
\begin{tikzpicture}[scale=1]
  \tkzInit[xmax=4,ymax=4]
  \tkzDrawX\tkzDrawY
  \tkzDefPoint(-1+2,sqrt(4)){O}
  \tkzDefPoint({3*ln(exp(1))},{exp(1)}){A}
  \tkzDefPoint({4*sin(pi/6)},{4*cos(pi/6)}){B}
  \tkzDrawPoints(O,B,A)
\end{tikzpicture}
\end{tkzexample}

\subsubsection{极坐标}

\begin{tkzexample}[latex=7cm,small]
  \begin{tikzpicture}
  \foreach \an [count=\i] in {0,60,...,300}
   { \tkzDefPoint(\an:3){A_\i}}
  \tkzDrawPolygon(A_1,A_...,A_6)
  \tkzDrawPoints(A_1,A_...,A_6)
  \end{tikzpicture}
\end{tkzexample}

\subsubsection{相对点}
首先，我们可以使用 \TIKZ 的 \tkzNameEnv{scope} 环境。
在下面的示例中，我们有一种定义等边三角形的方法。

\begin{tkzexample}[latex=7cm,small]
\begin{tikzpicture}[scale=1]
 \begin{scope}[rotate=30]
  \tkzDefPoint(2,3){A}
  \begin{scope}[shift=(A)]
     \tkzDefPoint(90:5){B}
     \tkzDefPoint(30:5){C}
  \end{scope}
 \end{scope}
 \tkzDrawPolygon(A,B,C)
\tkzLabelPoints[above](B,C)
\tkzLabelPoints[below](A)
\tkzDrawPoints(A,B,C)
\end{tikzpicture}
\end{tkzexample}

\subsection{相对于另一个点的点：\tkzcname{tkzDefShiftPoint}}
\begin{NewMacroBox}{tkzDefShiftPoint}{\oarg{Point}\parg{$x,y$}\marg{ref} or \parg{$\alpha$:$d$}\marg{ref}}%
\begin{tabular}{lll}%
参数 & 默认值 & 定义 \\
\midrule
\TAline{($x,y$)}{no default}{$x$ 和 $y$ 是两个度量，默认为 cm。}
\TAline{($\alpha$:$d$)}{no default}{$\alpha$ 是以度为单位的角，$d$ 是一个度量}
\TAline{\{ref\}}{no default}{赋值给该点的引用：$A$、$T\_a$、$P1$ 或 $P_1$}

\midrule
选项 & 默认值 & 定义 \\

\midrule
\TOline{[pt]} {no default} {\tkzcname{tkzDefShiftPoint}[A](0:4)\{B\}}
\end{tabular}
\end{NewMacroBox}

\subsubsection{等腰三角形}
这个宏将一个点相对于另一个点放置；它等价于一次平移。下面是如何构造一个以 $A$ 为主要顶点、该顶点处角为 $30^{\circ}$ 的等腰三角形。

\begin{tkzexample}[latex=7cm,small]
\begin{tikzpicture}[rotate=-30]
 \tkzDefPoint(2,3){A}
 \tkzDefShiftPoint[A](0:4){B}
 \tkzDefShiftPoint[A](30:4){C}
 \tkzDrawSegments(A,B B,C C,A)
 \tkzMarkSegments[mark=|](A,B A,C)
 \tkzDrawPoints(A,B,C)
 \tkzLabelPoints[right](B,C)
 \tkzLabelPoints[above left](A)
\end{tikzpicture}
\end{tkzexample}

\subsubsection{等边三角形}
让我们看看如何得到一个等边三角形（有更简单的方法，见 \tkzcname{tkzDefTriangle}）。

\begin{tkzexample}[latex=7cm,small]
\begin{tikzpicture}[scale=1]
 \tkzDefPoint(2,3){A}
 \tkzDefShiftPoint[A](30:3){B}
 \tkzDefShiftPoint[A](-30:3){C}
 \tkzDrawPolygon(A,B,C)
 \tkzDrawPoints(A,B,C)
 \tkzLabelPoints[right](B,C)
 \tkzLabelPoints[above left](A)
 \tkzMarkSegments[mark=|](A,B A,C B,C)
\end{tikzpicture}
\end{tkzexample}

\subsubsection{平行四边形}
有更简单的方法，见 \tkzcname{tkzDefParallelogram}。
\begin{tkzexample}[latex=7cm,small]
\begin{tikzpicture}
 \tkzDefPoint(0,0){A}
 \tkzDefPoint(30:3){B}
 \tkzDefShiftPointCoord[B](10:2){C}
 \tkzDefShiftPointCoord[A](10:2){D}
 \tkzDrawPolygon(A,...,D)
 \tkzDrawPoints(A,...,D)
\end{tikzpicture}
\end{tkzexample}

\subsection{多个点的定义：\tkzcname{tkzDefPoints}}

\begin{NewMacroBox}{tkzDefPoints}{\oarg{local options}\marg{$x_1/y_1/r_1,x_2/y_2/r_2$, ...}}%
$x_i$ 和 $y_i$ 是引用为 $r_i$ 的点的坐标。

\begin{tabular}{lll}%
\toprule
参数 & 默认值 & 示例  \\
\midrule
\TAline{$x_i/y_i/r_i$}{}{\tkzcname{tkzDefPoints\{0/0/O,2/2/A\}}}
\end{tabular}

\medskip
\begin{tabular}{lll}%
选项             & 默认值 & 定义   \\
\midrule
\TOline{shift} {no default} {将 $(x,y)$ 或 $(\alpha:d)$ 加到所有坐标上}
\end{tabular}
\end{NewMacroBox}

\subsection{创建三角形}
\begin{tkzexample}[latex=6cm,small]
\begin{tikzpicture}[scale=.75]
 \tkzDefPoints{0/0/A,4/0/B,4/3/C}
 \tkzDrawPolygon(A,B,C)
 \tkzDrawPoints(A,B,C)
\end{tikzpicture}
\end{tkzexample}

\subsection{创建正方形}
请注意这里绘制多边形的语法。
\begin{tkzexample}[latex=6cm,small]
\begin{tikzpicture}[scale=1]
 \tkzDefPoints{0/0/A,2/0/B,2/2/C,0/2/D}
 \tkzDrawPolygon(A,...,D)
 \tkzDrawPoints(A,...,D)
\end{tikzpicture}
\end{tkzexample}

\endinput

既然固定点已经定义，我们就可以通过包提供的宏或自己创建的宏，由这些点的引用得到新的点。计算通常由包完成，对用户是透明的。
您可能需要一些数学常数；以下是包中定义的常数列表。


\section{辅助工具}
\subsection{常数}

\tkzname{\tkznameofpack} 已知以下常数：
\begin{tkzltxexample}[]
  \def\tkzPhi{1.618034}
  \def\tkzInvPhi{0.618034}
  \def\tkzSqrtPhi{1.27202}
  \def\tkzSqrTwo{1.414213}
  \def\tkzSqrThree{1.7320508}
  \def\tkzSqrFive{2.2360679}
  \def\tkzSqrTwobyTwo{0.7071065}
  \def\tkzPi{3.1415926}
  \def\tkzEuler{2.71828182}
\end{tkzltxexample}

\subsection{通过计算得到新点}

当 \tkzname{\tkznameofpack} 的某个宏创建了一个新点时，该点在内部以引用 \tkzname{tkzPointResult} 存储。您可以用宏 \tkzcname{tkzGetPoint} 为它指定自己的引用；新引用放在花括号之间。

\begin{NewMacroBox}{tkzGetPoint}{\marg{ref}}%
如果结果在 \tkzname{tkzPointResult} 中，可以用 \tkzcname{tkzGetPoint} 获取它。

 \medskip
\begin{tabular}{lll}%
\toprule
参数 & 默认值 & 示例 \\
\midrule
\TAline{ref}{no default}{ \tkzcname{tkzGetPoint\{M\} } 见下例}
\end{tabular}
\end{NewMacroBox}

有时您需要获取两个点。这可以用下面的宏完成。

\begin{NewMacroBox}{tkzGetPoints}{\marg{ref1}\marg{ref2}}%
结果分别存放在 \tkzname{tkzFirstPointResult} 和 \tkzname{tkzSecondPointResult} 中。

 \medskip
\begin{tabular}{lll}%
\toprule
参数 & 默认值 & 示例 \\
\midrule
\TAline{\{ref1,ref2\}}{no default}{ \tkzcname{tkzGetPoints\{M,N\} } 这是 \tkzcname{tkzInterCC}}
\end{tabular}
\end{NewMacroBox}

如果只需要第一个或第二个点，也可以使用：

\begin{NewMacroBox}{tkzGetFirstPoint}{\marg{ref1}}%

 \medskip
\begin{tabular}{lll}%
\toprule
参数 & 默认值 & 示例 \\
\midrule
\TAline{ref1}{no default}{ \tkzcname{tkzGetFirstPoint\{M\} }}
\end{tabular}
\end{NewMacroBox}

\begin{NewMacroBox}{tkzGetSecondPoint}{\marg{ref2}}%

 \medskip
\begin{tabular}{lll}%
\toprule
参数 & 默认值 & 示例 \\
\midrule
\TAline{ref2}{no default}{ \tkzcname{tkzGetSecondPoint\{M\} }}
\end{tabular}
\end{NewMacroBox}

有时结果由一个点和一个长度组成。点用 \tkzcname{tkzGetPoint} 获取，长度用 \tkzcname{tkzGetLength} 获取。

\begin{NewMacroBox}{tkzGetLength}{\marg{name of a macro}}%

 \medskip
\begin{tabular}{lll}%
\toprule
参数 & 默认值 & 示例 \\
\midrule
\TAline{name of a macro}{no default}{ \tkzcname{tkzGetLength\{rAB\} \tkzcname{rAB} 给出以 cm 为单位的长度}}
\end{tabular}
\end{NewMacroBox}

%\tkzcname{tkzCalcLength}(A,B) After \tkzcname{tkzGetLength\{dAB\}} \tkzcname{dAB} gives $AB$ in cm}


\section{特殊点}
以下是一些特殊点。
%<--------------------------------------------------------------------------->
\subsection{线段的中点 \tkzcname{tkzDefMidPoint}}
该宏定义一条线段的中点。

\begin{NewMacroBox}{tkzDefMidPoint}{\parg{pt1,pt2}}%
结果存放在 \tkzname{tkzPointResult} 中。可以用 \tkzcname{tkzGetPoint} 获取。

 \medskip
\begin{tabular}{lll}%
\toprule
参数 & 默认值 & 定义 \\
\midrule
\TAline{(pt1,pt2)}{no default}{pt1 和 pt2 为两个点}
\end{tabular}
\end{NewMacroBox}

\subsubsection{\tkzcname{tkzDefMidPoint} 的使用}
另请参阅 \tkzcname{tkzDefPoint} 的使用。
\begin{tkzexample}[latex=7cm,small]
\begin{tikzpicture}[scale=1]
 \tkzDefPoint(2,3){A}
 \tkzDefPoint(6,2){B}
 \tkzDefMidPoint(A,B)
 \tkzGetPoint{M}
 \tkzDrawSegment(A,B)
 \tkzDrawPoints(A,B,M)
 \tkzLabelPoints[below](A,B,M)
\end{tikzpicture}
\end{tkzexample}

\subsection{\tkzname{黄金比例} \tkzcname{tkzDefGoldenRatio}}
来自维基百科：在数学中，若两个量的比等于它们的和与较大量的比，则称这两个量成黄金比例。用代数语言表述，对于量 $a$ 和 $b$，若 $a > b > 0$，则 $(a+b):a = a:b$。

$ \frac{a+b}{a} = \frac{a}{b} = \phi = \frac{1 + \sqrt{5}}{2}$


方程 $x^2 - x - 1 = 0$ 的两个解之一
即为黄金比例 $\phi$，$\phi = \frac{1 + \sqrt{5}}{2}$。

\begin{NewMacroBox}{tkzDefGoldenRatio}{\parg{pt1,pt2}}%
\begin{tabular}{lll}%
参数 & 默认值 & 示例 \\
\midrule
\TAline{(pt1,pt2)}{no default}{\tkzcname{tkzDefGoldenRatio(A,C)} \tkzcname{tkzGetPoint}\{B\}}
\bottomrule
\end{tabular}

\medskip
$AB=a$，$BC=b$，且 $\dfrac{AC}{AB} = \dfrac{AB}{BC} =\phi$
\end{NewMacroBox}

\subsubsection{用黄金比例分割线段}
\begin{tkzexample}[latex=7cm,small]
\begin{tikzpicture}
 \tkzDefPoints{0/0/A,6/0/C}
 \tkzDefMidPoint(A,C) \tkzGetPoint{I}
 %\tkzDefPointWith[linear,K=\tkzInvPhi](A,C)
 \tkzDefGoldenRatio(A,C) \tkzGetPoint{B}
 \tkzDrawSegments(A,C)
 \tkzDrawPoints(A,B,C)
 \tkzLabelPoints(A,B,C)
\end{tikzpicture}
\end{tkzexample}

\subsubsection{黄金 arbelos（鞋匠刀形）}
\begin{tkzexample}[latex=7cm,small]
\begin{tikzpicture}[scale=.6]
\tkzDefPoints{0/0/A,10/0/B}
\tkzDefGoldenRatio(A,B)     \tkzGetPoint{C}
\tkzDefMidPoint(A,B)        \tkzGetPoint{O_1}
\tkzDefMidPoint(A,C)        \tkzGetPoint{O_2}
\tkzDefMidPoint(C,B)        \tkzGetPoint{O_3}
\tkzDrawSemiCircles[fill=purple!15](O_1,B)
\tkzDrawSemiCircles[fill=teal!15](O_2,C O_3,B)
\end{tikzpicture}
\end{tkzexample}

也可以使用下面的宏。
\subsection{\tkzname{重心坐标} 与 \tkzcname{tkzDefBarycentricPoint}}

$n$ 个点 $pt_1$、$pt_2$、\dots、$pt_n$ 定义了 $n$ 个以原点为共同起点的向量 $\overrightarrow{v_1}$、$\overrightarrow{v_2}$、\dots、$\overrightarrow{v_n}$。$\alpha_1$、$\alpha_2$、
\dots、$\alpha_n$ 为 $n$ 个数；向量
\begin{align*}
  \frac{\alpha_1 \overrightarrow{v_1} + \alpha_2 \overrightarrow{v_2} + \cdots + \alpha_n \overrightarrow{v_n}}{\alpha_1
    + \alpha_2 + \cdots + \alpha_n}
\end{align*}
确定唯一一个点。

\begin{NewMacroBox}{tkzDefBarycentricPoint}{\parg{pt1=$\alpha_1$,pt2=$\alpha_2$,\dots}}%
\begin{tabular}{lll}%
参数 & 默认值 & 定义 \\
\midrule
\TAline{(pt1=$\alpha_1$,pt2=$\alpha_2$,\dots)}{no default}{每个点都被赋予一个权重}
\bottomrule
\end{tabular}

\medskip
\emph{至少需要两个点。结果存放在 \tkzname{tkzPointResult} 中。}
\end{NewMacroBox}


\subsubsection{两个点的情形}
在下面的示例中，我们以系数 $1$ 和 $2$ 求点 $A$ 和 $B$ 的重心，即：
\[
  \overrightarrow{AI}= \frac{2}{3}\overrightarrow{AB}
\]

\begin{tkzexample}[latex=7cm,small]
\begin{tikzpicture}
  \tkzDefPoint(2,3){A}
  \tkzDefShiftPointCoord[2,3](30:4){B}
  \tkzDefBarycentricPoint(A=1,B=2)
  \tkzGetPoint{G}
  \tkzDrawLine(A,B)
  \tkzDrawPoints(A,B,G)
  \tkzLabelPoints(A,B,G)
\end{tikzpicture}
\end{tkzexample}

\subsubsection{三个点的情形}
此时 $M$ 就是三角形的重心。

为简洁和一致起见，也提供了 \tkzcname{tkzCentroid}。
\begin{tkzexample}[latex=7cm,small]
\begin{tikzpicture}[scale=.8]
  \tkzDefPoints{2/1/A,5/3/B,0/6/C}
  \tkzDefBarycentricPoint(A=1,B=1,C=1)
  \tkzGetPoint{G}
  \tkzDefMidPoint(A,B)  \tkzGetPoint{C'}
  \tkzDefMidPoint(A,C)  \tkzGetPoint{B'}
  \tkzDefMidPoint(C,B)  \tkzGetPoint{A'}
  \tkzDrawPolygon(A,B,C)
  \tkzDrawLines[add=0 and 1,new](A,G B,G C,G)
  \tkzDrawPoints[new](A',B',C',G)
  \tkzDrawPoints(A,B,C)
  \tkzLabelPoint[above right](G){$G$}
  \tkzAutoLabelPoints[center=G](A,B,C)
  \tkzLabelPoints[above right](A')
  \tkzLabelPoints[below](B',C')
\end{tikzpicture}
\end{tkzexample}


\subsection{内、外相似中心：\tkzcname{tkzDefSimilitudeCenter}}
将一个圆映射到另一个圆的两个位似的中心，称为外相似中心和内相似中心。可以使用 \tkzcname{tkzDefIntSimilitudeCenter} 和 \tkzcname{tkzDefExtSimilitudeCenter}，但下面的宏更方便。

\begin{NewMacroBox}{tkzDefSimilitudeCenter}{\oarg{options}\parg{O,A}\parg{O',B}}%

\begin{tabular}{lll}%
参数 & 示例 & 说明                         \\
\midrule
\TAline{\parg{pt1,pt2}\parg{pt3,pt4}}{$(O,A)(O',B)$} {$r=OA,r'=O'B$}
\end{tabular}

\medskip
\begin{tabular}{lll}%
\toprule
选项 & 默认值 & 定义                         \\
\midrule
\TOline{ext}{ext}{外相似中心}
\TOline{int}{ext}{内相似中心}
\end{tabular}
\end{NewMacroBox}

\subsubsection{用 \tkzname{node} 获取内、外相似中心}
\begin{tkzexample}[latex=7.5cm,small]
\begin{tikzpicture}[scale=.7]
 \tkzDefPoints{0/0/O,4/-5/A,3/0/B,5/-5/C}
 \tkzDefSimilitudeCenter[int](O,B)(A,C)    
 \tkzGetPoint{I}
 \tkzDefSimilitudeCenter[ext](O,B)(A,C)   
 \tkzGetPoint{J}
 \tkzDefLine[tangent from = I](O,B)       
 \tkzGetPoints{D}{E}
 \tkzDefLine[tangent from = I](A,C)       
 \tkzGetPoints{D'}{E'}
 \tkzDefLine[tangent from = J](O,B)       
 \tkzGetPoints{F}{G}
 \tkzDefLine[tangent from = J](A,C)
 \tkzGetPoints{F'}{G'}
 \tkzDrawCircles(O,B A,C)
 \tkzDrawSegments[add = .5 and .5,new](D,D' E,E')
 \tkzDrawSegments[add= 0 and 0.25,new](J,F J,G)
 \tkzDrawPoints(O,A,I,J,D,E,F,G,D',E',F',G')
\end{tikzpicture}
\end{tkzexample}

\subsubsection{D'Alembert 定理} % (fold)
\label{ssub:d_alembert_theorem}

\begin{tkzexample}[latex=7cm,small]
 \begin{tikzpicture}[scale=.6,rotate=90]
 \tkzDefPoints{0/0/A,3/0/a,7/-1/B,5.5/-1/b}
 \tkzDefPoints{5/-4/C,4.25/-4/c}
 \tkzDrawCircles(A,a B,b C,c)
 \tkzDefExtSimilitudeCenter(A,a)(B,b) \tkzGetPoint{I}
 \tkzDefExtSimilitudeCenter(A,a)(C,c) \tkzGetPoint{J}
 \tkzDefExtSimilitudeCenter(C,c)(B,b) \tkzGetPoint{K}
 \tkzDefIntSimilitudeCenter(A,a)(B,b) \tkzGetPoint{I'}
 \tkzDefIntSimilitudeCenter(A,a)(C,c) \tkzGetPoint{J'}
 \tkzDefIntSimilitudeCenter(C,c)(B,b) \tkzGetPoint{K'}
 \tkzDrawPoints(A,B,C,I,J,K,I',J',K')
 \tkzDrawSegments[new](I,K A,I A,J B,I B,K C,J C,K)
 \tkzDrawSegments[new](I,J' I',J I',K)
 \end{tikzpicture}
\end{tkzexample}

% subsubsection d_alembert_theorem (end)

也可以使用 \tkzcname{tkzDefBarycentricPoint} 来找位似中心：

|\tkzDefBarycentricPoint(O=\r,A=\R)     \tkzGetPoint{I}| \\
|\tkzDefBarycentricPoint(O={-\r},A=\R)  \tkzGetPoint{J}|

\subsubsection{使用 \tkzname{node} 的示例}
\begin{tkzexample}[latex=7cm,small]
\begin{tikzpicture}[rotate=60,scale=.5]
 \tkzDefPoints{0/0/A,5/0/C}
 \tkzDefGoldenRatio(A,C) \tkzGetPoint{B}
 \tkzDefSimilitudeCenter(A,B)(C,B)\tkzGetPoint{J}
 \tkzDefTangent[from = J](A,B)  \tkzGetPoints{F}{G}
 \tkzDefTangent[from = J](C,B)  \tkzGetPoints{F'}{G'}
 \tkzDrawCircles(A,B C,B)
 \tkzDrawSegments[add= 0 and 0.25,cyan](J,F J,G)
 \tkzDrawPoints(A,J,F,G,F',G')
\end{tikzpicture}
\end{tkzexample}
\newpage
%<---------------------------------------------------------------------->
\subsection{ \tkzname{调和分割} 与 \tkzcname{tkzDefHarmonic}}
%<---------------------------------------------------------------------->

\begin{NewMacroBox}{tkzDefHarmonic}{\oarg{options}\parg{pt1,pt2,pt3} or \parg{pt1,pt2,k}}%

\begin{tabular}{lll}%
选项 & 默认值 & 定义                         \\
\midrule
\TOline{both}{both}{\parg{A,B,2} 求 C 和 D 使得 $(A,B;C,D) = -1$ 且 CA=2CB }
\TOline{ext}{both}{\parg{A,B,C} 求 D 使得 $(A,B;C,D) = -1$}
\TOline{int}{both}{\parg{A,B,D} 求 C 使得 $(A,B;C,D) = -1$}
\end{tabular}
\end{NewMacroBox}

\subsubsection{选项 \tkzname{ext} 与 \tkzname{int}}
\begin{tkzexample}[vbox,small]
  \begin{tikzpicture}
  \tkzDefPoints{0/0/A,6/0/B,4/0/C}
  \tkzDefHarmonic[ext](A,B,C) \tkzGetPoint{J}
  \tkzDefHarmonic[int](A,B,J) \tkzGetPoint{I}
  \tkzDrawPoints(A,B,I,J)
  \tkzDrawLine[add=.5 and 1](A,B)
  \tkzLabelPoints(A,B,I,J)
  \end{tikzpicture}
\end{tkzexample}

\subsubsection{角平分线与调和分割} % (fold)
\label{ssub:bisector_and_harmonic_division}

\begin{tkzexample}[vbox,small]
  \begin{tikzpicture}[scale=1.25]
  \tkzDefPoints{0/0/A,4/0/C,5/3/X}
  \tkzDefLine[bisector](A,X,C) \tkzGetPoint{x}
  \tkzInterLL(X,x)(A,C)        \tkzGetPoint{B}
  \tkzDefHarmonic[ext](A,C,B)  \tkzGetPoint{D}
  \tkzDrawPolygon(A,X,C)
  \tkzDrawSegments(X,B C,D D,X)
  \tkzDrawPoints(A,B,C,D,X)
  \tkzMarkAngles[mark=s|](A,X,B B,X,C)
  \tkzMarkRightAngle[size=.4,
                     fill=gray!20,
                     opacity=.3](B,X,D)
  \tkzLabelPoints(A,B,C,D)
  \tkzLabelPoints[above right](X)
  \end{tikzpicture}
\end{tkzexample}


% subsubsection bisector_and_harmonic_division (end)
\subsubsection{选项 \tkzname{both}}
\tkzname{both} 是默认选项。
\begin{tkzexample}[vbox,small]
\begin{tikzpicture}
 \tkzDefPoints{0/0/A,6/0/B}
 \tkzDefHarmonic(A,B,{1/2})\tkzGetPoints{I}{J}
 \tkzDrawPoints(A,B,I,J)
 \tkzDrawLine[add=1 and .5](A,B)
 \tkzLabelPoints(A,B,I,J)
\end{tikzpicture}
\end{tkzexample}

%<---------------------------------------------------------------------->
\subsection{\tkzname{等距点} 与 \tkzcname{tkzDefEquiPoints}}
%<---------------------------------------------------------------------->

\begin{NewMacroBox}{tkzDefEquiPoints}{\oarg{local options}\parg{pt1,pt2}}%
\begin{tabular}{lll}%
参数 & 默认值 & 定义 \\
\midrule
\TAline{(pt1,pt2)}{no default}{两个元素的无序列表}
\end{tabular}

\begin{tabular}{lll}%
选项 & 默认值 & 定义  \\
\midrule
\TOline{dist} {2 (cm)} {两点之间距离的一半}
\TOline{from=pt} {no default} {参考点}
\TOline{show} {false} {若为 true 则显示圆规作图痕迹}
\TOline{/compass/delta} {0} {圆规标记的大小}
\end{tabular}

\medskip
\emph{该宏给出一条直线上与给定点等距的两个点。}
\end{NewMacroBox}


\subsubsection{带选项使用 \tkzcname{tkzDefEquiPoints}}
\begin{tkzexample}[latex=7cm,small]
\begin{tikzpicture}
  \tkzSetUpCompass[color=purple,line width=1pt]
  \tkzDefPoints{0/1/A,5/2/B,3/4/C}
  \tkzDefEquiPoints[from=C,dist=1,show,
      /tkzcompass/delta=20](A,B)
   \tkzGetPoints{E}{H}
   \tkzDrawLines[color=blue](C,E C,H A,B)
   \tkzDrawPoints[color=blue](A,B,C)
   \tkzDrawPoints[color=red](E,H)
   \tkzLabelPoints(E,H)
   \tkzLabelPoints[color=blue](A,B,C)
\end{tikzpicture}
\end{tkzexample}
%<---------------------------------------------------------------------->
%                          Middle of an arc                             >
%<---------------------------------------------------------------------->
\subsection{弧的中点}
\begin{NewMacroBox}{tkzDefMidArc}{\parg{pt1,pt2,pt3}}%
\begin{tabular}{lll}%
参数 & 默认值 & 定义 \\
\midrule
\TAline{$pt1,pt2,pt3$}{no default}{$pt1$ 为圆心，$\widearc{pt2pt3}$ 为圆弧}
\end{tabular}
\end{NewMacroBox}

\begin{tkzexample}[vbox,small]
  \begin{tikzpicture}[scale=1]
   \tkzDefPoints{0/0/A,10/0/B}
   \tkzDefGoldenRatio(A,B)                              \tkzGetPoint{C}
   \tkzDefMidPoint(A,B)                                 \tkzGetPoint{O_1}
   \tkzDefMidPoint(A,C)                                 \tkzGetPoint{O_2}
   \tkzDefMidPoint(C,B)                                 \tkzGetPoint{O_3}
   \tkzDefMidArc(O_3,B,C)                               \tkzGetPoint{P}
   \tkzDefMidArc(O_2,C,A)                               \tkzGetPoint{Q}
   \tkzDefMidArc(O_1,B,A)                               \tkzGetPoint{L}
   \tkzDefPointBy[rotation=center C angle 90](B)        \tkzGetPoint{c}
   \tkzInterCC[common=B](P,B)(O_1,B)                    \tkzGetFirstPoint{P_1}
   \tkzInterCC[common=C](P,C)(O_2,C)                    \tkzGetFirstPoint{P_2}
   \tkzInterCC[common=C](Q,C)(O_3,C)                    \tkzGetFirstPoint{P_3}
   \tkzInterLC[near](c,C)(O_1,A)                        \tkzGetFirstPoint{D}
   \tkzInterLL(A,P_1)(C,D)                              \tkzGetPoint{P_1'}
   \tkzDefPointBy[inversion = center A through D](P_2)  \tkzGetPoint{P_2'}
   \tkzDefCircle[circum](P_3,P_2,P_1)                   \tkzGetPoint{O_4}
   \tkzInterLL(B,Q)(A,P)                                \tkzGetPoint{S}
   \tkzDefMidPoint(P_2',P_1')                           \tkzGetPoint{o}
   \tkzDefPointBy[inversion = center A through D](S)    \tkzGetPoint{S'}
   \tkzDrawArc[cyan,delta=0](Q,A)(P_1)
   \tkzDrawArc[cyan,delta=0](P,P_1)(B)
   \tkzDrawSemiCircles[teal](O_1,B O_2,C O_3,B)
   \tkzDrawCircles[new](o,P O_4,P_1)
   \tkzDrawSegments(A,B)
   \tkzDrawSegments[cyan](A,P_1 A,S' A,P_2')
   \tkzDrawSegments[purple](B,L C,P_2' B,Q B,L S',P_1')
   \tkzDrawLines[add=0 and .8](B,P_2')
   \tkzDrawLines[add=0 and .4](C,D)
   \tkzDrawPoints(A,B,C,P,Q,P_3,P_2,P_1,P_1',D,P_2',L,S,S')
   \tkzLabelPoints(A,B,C,P_3)
   \tkzLabelPoints[above](P,Q,P_1)
   \tkzLabelPoints[above right](P_2,P_2',D,S')
   \tkzLabelPoints[above left](L,S)
    \tkzLabelPoints[below left](P_1')
  \end{tikzpicture}
\end{tkzexample}

%<---------------------------------------------------------------------->
%                          Point on a line                              >
%<---------------------------------------------------------------------->

\section{直线或圆上的点}
\subsection{直线上的点 与 \tkzcname{tkzDefPointOnLine}}

\begin{NewMacroBox}{tkzDefPointOnLine}{\oarg{local options}\parg{A,B}}%
\begin{tabular}{lll}%
参数 & 默认值 & 定义                 \\
\midrule
\TAline{pt1,pt2} {no default}  {用于确定一条直线的两个点}
\bottomrule
\end{tabular}

\medskip
\begin{tabular}{lll}%
选项 & 默认值 & 定义 \\
\midrule
\TOline{pos=nb}  {}{nb 为小数}
\end{tabular}
\end{NewMacroBox}

\subsubsection{选项 \tkzname{pos} 的使用}
\begin{tkzexample}[latex=7cm,small]
\begin{tikzpicture}
\tkzDefPoints{0/0/A,3/0/B}
\tkzDefPointOnLine[pos=1.2](A,B)\tkzGetPoint{P}
\tkzDefPointOnLine[pos=-0.2](A,B)\tkzGetPoint{R}
\tkzDefPointOnLine[pos=0.5](A,B) \tkzGetPoint{S}
\tkzDrawLine[new](A,B)
\tkzDrawPoints(A,B,P)
\tkzLabelPoints(A,B)
\tkzLabelPoint[above](P){pos=$1.2$}
\tkzLabelPoint[above](R){pos=$-.2$}
\tkzLabelPoint[above](S){pos=$.5$}
\tkzDrawPoints(A,B,P,R,S)
\tkzLabelPoints(A,B)
\end{tikzpicture}
\end{tkzexample}

\subsection{圆上的点 与 \tkzcname{tkzDefPointOnCircle}}
参数顺序已更改：现在为圆心、角度以及点或半径。
有两个选项使用弧度：\tkzname{through in rad} 和 \tkzname{R in rad}。

\begin{NewMacroBox}{tkzDefPointOnCircle}{\oarg{local options}}%
\begin{tabular}{lll}%
选项 & 默认值 & 示例/定义 \\
\midrule
\TOline{through}  {}{through = 圆心 K1 角度 30 点 B]}
\TOline{R} {}{R = 圆心 K1 角度 30 半径 \tkzcname{rAp}}
\TOline{through in rad}  {}{through in rad= 圆心 K1 角度 pi/4 点 B]}
\TOline{R in rad} {}{R in rad = 圆心 K1 角度 pi/6 半径 \tkzcname{rAp}}
\end{tabular}

\medskip
\emph{参数顺序为：圆心、角度、以及点或半径。}
\end{NewMacroBox}

\subsubsection{Altshiller 定理}
将两个正交圆的交点与其中一个圆上一点相连所得的两条直线，在另一个圆上交于一对对径点（Altshiller-Court，第 176 页）。

\begin{tkzexample}[latex=6cm,small]
\begin{tikzpicture}[scale=.4]
\tkzDefPoints{0/0/P,5/0/Q,3/2/I}
\tkzDefCircle[orthogonal from=P](Q,I)
\tkzGetFirstPoint{E}
\tkzDrawCircles(P,E Q,E)
\tkzInterCC[common=E](P,E)(Q,E) \tkzGetFirstPoint{F}
\tkzDefPointOnCircle[through =  center P angle 80 point E]
 \tkzGetPoint{A}
\tkzInterLC[common=E](A,E)(Q,E)  \tkzGetFirstPoint{C}
\tkzInterLL(A,F)(C,Q)  \tkzGetPoint{D}
\tkzDrawLines[add=0 and .75](P,Q)
\tkzDrawLines[add=0 and 2](A,E)
\tkzDrawSegments(P,E E,F F,C A,F C,D)
\tkzDrawPoints(P,Q,E,F,A,C,D)
\tkzLabelPoints(P,Q,F,C,D)
\tkzLabelPoints[above](E,A)
\end{tikzpicture}
\end{tkzexample}

\subsubsection{\tkzcname{tkzDefPointOnCircle} 的使用}

\begin{tkzexample}[latex=6cm,small]
\begin{tikzpicture}
\tkzDefPoints{0/0/A,4/0/B,0.8/3/C}
\tkzDefPointOnCircle[R = center B  angle 90 radius 1]
\tkzGetPoint{I}
\tkzDefCircle[circum](A,B,C)
\tkzGetPoints{G}{g}
\tkzDefPointOnCircle[through = center G angle 30 point g]
\tkzGetPoint{J}
\tkzDefCircle[R](B,1) \tkzGetPoint{b}
\tkzDrawCircle[teal](B,b)
\tkzDrawCircle(G,J)
\tkzDrawPoints(A,B,C,G,I,J)
\tkzAutoLabelPoints[center=G](A,B,C,J)
\tkzLabelPoints[below](G,I)
\end{tikzpicture}
\end{tkzexample}

\newpage
\section{与三角形相关的特殊点}

\subsection{三角形中心：\tkzcname{tkzDefTriangleCenter}}

\begin{NewMacroBox}{tkzDefTriangleCenter}{\oarg{local options}\parg{A,B,C}}%
\tkzHandBomb\ 该宏定义三角形的某种中心。注意：参数是三个点组成的列表。该宏与 \tkzcname{tkzGetPoint} 配合使用以获取所需的中心。

如果不需要保留结果，可以直接使用 \tkzname{tkzPointResult}。

\medskip
\begin{tabular}{lll}%
\toprule
参数 & 默认值 & 示例 \\
\midrule
\TAline{(pt1,pt2,pt3)}{no default}{ \tkzcname{tkzDefTriangleCenter[ortho](B,C,A)}}
\midrule
选项 & 默认值 & 定义                         \\
\midrule
\TOline{ortho}  {circum}{三条高的交点}
\TOline{orthic}  {circum}{\dots}
\TOline{centroid} {circum}{三条中线的交点}
\TOline{median} {circum}{ \dots }
\TOline{circum}{circum}{外接圆的圆心}
\TOline{in}    {circum}{三角形内切圆的圆心}
\TOline{in}    {circum}{角平分线的交点}
\TOline{ex}    {circum}{三角形旁切圆的圆心}
\TOline{euler}{circum}{Euler 圆（九点圆）的圆心}
\TOline{gergonne}{circum}{由切点三角形（Contact triangle）定义}
\TOline{symmedian} {circum}{Lemoine 点，即陪位重心（symmedian center），亦称 Grebe 点}
\TOline{lemoine} {circum}{ \dots}
\TOline{grebe} {circum}{ \dots}
\TOline{spieker} {circum}{Spieker 圆的圆心}
\TOline{nagel}{circum}{Nagel 点（界心）}
\TOline{mittenpunkt} {circum}{亦称 middlespoint}
\TOline{feuerbach}{circum}{费尔巴哈点（Feuerbach Point）}

\end{tabular}
\end{NewMacroBox}

\subsubsection{选项 \tkzname{ortho} 或 \tkzname{orthic}}
三角形三条高的交点 $H$ 称为垂心。

\begin{tkzexample}[latex=6cm,small]
\begin{tikzpicture}
  \tkzDefPoint(0,0){A}
  \tkzDefPoint(5,1){B}
  \tkzDefPoint(1,4){C}
  \tkzDefTriangleCenter[ortho](B,C,A)
  \tkzGetPoint{H}
  \tkzDefSpcTriangle[orthic,name=H](A,B,C){a,b,c}
  \tkzDrawPolygon(A,B,C)
  \tkzDrawSegments[new](A,Ha B,Hb C,Hc)
  \tkzDrawPoints(A,B,C,H)
  \tkzLabelPoint(H){$H$}
  \tkzLabelPoints[below](A,B)
  \tkzLabelPoints[above](C)
 \tkzMarkRightAngles(A,Ha,B B,Hb,C C,Hc,A)
\end{tikzpicture}
\end{tkzexample}

\subsubsection{选项 \tkzname{centroid}}
\begin{tkzexample}[latex=6cm,small]
\begin{tikzpicture}[scale=.75]
  \tkzDefPoints{0/0/A,5/0/B,1/4/C}
  \tkzDefTriangleCenter[centroid](A,B,C)
  \tkzGetPoint{G}
  \tkzDrawPolygon(A,B,C)
  \tkzDrawLines[add = 0 and 2/3,new](A,G B,G C,G)
  \tkzDrawPoints(A,B,C,G)
  \tkzLabelPoint(G){$G$}
\end{tikzpicture}
\end{tkzexample}

\subsubsection{选项 \tkzname{circum}}
\begin{tkzexample}[latex=6cm,small]
 \begin{tikzpicture}
  \tkzDefPoints{0/1/A,3/2/B,1/4/C}
  \tkzDefTriangleCenter[circum](A,B,C)
  \tkzGetPoint{O}
  \tkzDrawPolygon(A,B,C)
  \tkzDrawCircle(O,A)
  \tkzDrawPoints(A,B,C,O)
  \tkzLabelPoint(O){$O$}
\end{tikzpicture}
\end{tkzexample}

\subsubsection{选项 \tkzname{in}}
在几何学中，三角形的内切圆（incircle）是三角形内部最大的圆；它与三角形的三条边都相切。内切圆的圆心是一个三角形中心，称为内心（incenter）。
内切圆的圆心即内心，可以作为三条内角平分线的交点求得。旁切圆的圆心是一个角的内角平分线（例如顶点 $A$ 处）与另外两个角的外角平分线的交点。该旁切圆的圆心称为相对于顶点 $A$ 的旁心（excenter），或 $A$ 所对的旁心。由于一个角的内角平分线与外角平分线互相垂直，因此内切圆圆心与三个旁切圆圆心构成一个垂心组。\\
(Article on \href{https://en.wikipedia.org/wiki/Incircle_and_excircles_of_a_triangle}{Wikipedia})

 \medskip
 我们得到三角形内切圆的圆心。结果当然存放在 \tkzname{tkzPointResult} 中；用 \tkzcname{tkzGetPoint} 取回。

\begin{tkzexample}[latex=7cm,small]
\begin{tikzpicture}
\tkzDefPoints{0/0/A,6/0/B,0.8/4/C}
\tkzDefTriangleCenter[in](A,B,C)
   \tkzGetPoint{I}
\tkzDrawLines(A,B B,C C,A)
\tkzDefCircle[in](A,B,C) \tkzGetPoints{I}{i}
\tkzDrawCircle(I,i)
\tkzDrawPoint[red](I)
\tkzDrawPoints(A,B,C)
\tkzLabelPoint(I){$I$}
\end{tikzpicture}
\end{tkzexample}

\subsubsection{选项 \tkzname{ex}}
三角形的旁切圆（excircle 或 escribed circle）是位于三角形外部、与三角形的一条边相切且与另外两条边的延长线相切的圆。每个三角形都有三个不同的旁切圆，分别与三角形的一条边相切。\\
(Article on \href{https://en.wikipedia.org/wiki/Incircle_and_excircles_of_a_triangle}{Wikipedia})


 我们得到三角形旁切圆的圆心。结果当然存放在 \tkzname{tkzPointResult} 中；用 \tkzcname{tkzGetPoint} 取回。

\begin{tkzexample}[latex=7cm,small]
\begin{tikzpicture}[scale=.5]
  \tkzDefPoints{0/1/A,3/2/B,1/4/C}
  \tkzDefTriangleCenter[ex](B,C,A)
  \tkzGetPoint{J_c}
  \tkzDefPointBy[projection=onto A--B](J_c)
  \tkzGetPoint{Tc}
  \tkzDrawPolygon(A,B,C)
  \tkzDrawCircle[new](J_c,Tc)
  \tkzDrawLines[add=1.5 and 0](A,C B,C)
  \tkzDrawPoints(A,B,C,J_c)
  \tkzLabelPoints(J_c)
\end{tikzpicture}
\end{tkzexample}

\subsubsection{选项 \tkzname{euler}}
该宏给出九点圆（nine-point circle）的圆心，九点圆也称 Euler 圆或 Feuerbach 圆。九点圆是过任意参考三角形 $ABC$ 的三个顶点向对边所作垂线的垂足 $H_A$、$H_B$、$H_C$ 的圆。Euler 于 1765 年证明它还过 $ABC$ 三边的中点 $M_A$、$M_B$、$M_C$。由 Feuerbach 定理，九点圆还过顶点与垂心 $H$ 连线线段的中点 $E_A$、$E_B$、$E_C$。这些点通常称为 Euler 点。\\

\begin{tkzexample}[latex=6cm,small]
\begin{tikzpicture}[scale=1.2,rotate=90]
 \tkzDefPoints{0/0/A,6/0/B,0.8/4/C}
 \tkzDefSpcTriangle[medial,name=M](A,B,C){_A,_B,_C}
 \tkzDefTriangleCenter[euler](A,B,C)\tkzGetPoint{N}
 % I= N nine points
 \tkzDefTriangleCenter[ortho](A,B,C)\tkzGetPoint{H}
 \tkzDefMidPoint(A,H) \tkzGetPoint{E_A}
 \tkzDefMidPoint(C,H) \tkzGetPoint{E_C}
 \tkzDefMidPoint(B,H) \tkzGetPoint{E_B}
 \tkzDefSpcTriangle[ortho,name=H](A,B,C){_A,_B,_C}
 \tkzDrawPolygon(A,B,C)
 \tkzDrawCircle[new](N,E_A)
 \tkzDrawSegments[new](A,H_A B,H_B C,H_C)
 \tkzDrawPoints(A,B,C,N,H)
 \tkzDrawPoints[new](M_A,M_B,M_C)
 \tkzDrawPoints( H_A,H_B,H_C)
 \tkzDrawPoints[green](E_A,E_B,E_C)
 \tkzAutoLabelPoints[center=N,
 font=\scriptsize](A,B,C,M_A,M_B,M_C,H_A,H_B,H_C,%
   E_A,E_B,E_C)
 \tkzLabelPoints[font=\scriptsize](H,N)
 \tkzMarkSegments[mark=s|,size=3pt,
 color=blue,line width=1pt](B,E_B E_B,H)
\end{tikzpicture}
\end{tkzexample}


\subsubsection{选项 \tkzname{symmedian}}

The point of concurrence $K$ of the symmedians, sometimes also called the Lemoine point (in England and France) or the Grebe point (in Germany).\\
\href{https://mathworld.wolfram.com/SymmedianPoint.html}{Weisstein, Eric W. "Symmedian Point." From MathWorld--A Wolfram Web Resource.}

\begin{tkzexample}[latex=7cm,small]
\begin{tikzpicture}
  \tkzDefPoint(0,0){A}
  \tkzDefPoint(5,0){B}
  \tkzDefPoint(1,4){C}
  \tkzDefTriangleCenter[symmedian](A,B,C)
  \tkzGetPoint{K}
  \tkzDefTriangleCenter[median](A,B,C)
  \tkzGetPoint{G}
  \tkzDefTriangleCenter[in](A,B,C)\tkzGetPoint{I}
  \tkzDefSpcTriangle[centroid,name=M](A,B,C){a,b,c}
  \tkzDefSpcTriangle[incentral,name=I](A,B,C){a,b,c}
  \tkzDrawPolygon(A,B,C)
  \tkzDrawLines[add = 0 and 2/3,new](A,K B,K C,K)
  \tkzDrawSegments[color=cyan](A,Ma B,Mb C,Mc)
  \tkzDrawSegments[color=green](A,Ia B,Ib C,Ic)
  \tkzDrawPoints(A,B,C,K,G,I)
  \tkzLabelPoints[font=\scriptsize](A,B,K,G,I)
  \tkzLabelPoints[above,font=\scriptsize](C)
\end{tikzpicture}
\end{tkzexample}

\subsubsection{选项 \tkzname{spieker}}
Spieker 中心是 Spieker 圆的圆心 $Sp$，即参考三角形的中点三角形（medial triangle）的内心。\\
\href{https://mathworld.wolfram.com/SpiekerCenter.html}{Weisstein, Eric W. "Spieker Center." From MathWorld--A Wolfram Web Resource. }

\begin{tkzexample}[latex=8cm,small]
\begin{tikzpicture}
 \tkzDefPoints{0/0/A,6/0/B,5/5/C}
 \tkzDefSpcTriangle[medial](A,B,C){Ma,Mb,Mc}
 \tkzDefTriangleCenter[centroid](A,B,C)
 \tkzGetPoint{G}
 \tkzDefTriangleCenter[spieker](A,B,C)
 \tkzGetPoint{Sp}
 \tkzDrawPolygon[](A,B,C)
 \tkzDrawPolygon[new](Ma,Mb,Mc)
 \tkzDefCircle[in](Ma,Mb,Mc) \tkzGetPoints{I}{i}
 \tkzDrawCircle(I,i)
 \tkzDrawPoints(B,C,A,Sp,Ma,Mb,Mc)
 \tkzAutoLabelPoints[center=G,dist=.3](Ma,Mb)
 \tkzLabelPoints[right](Sp)
 \tkzLabelPoints[below](A,B,Mc)
 \tkzLabelPoints[above](C)
\end{tikzpicture}
\end{tkzexample}

\subsubsection{选项 \tkzname{gergonne}}

Gergonne 点是将三角形各顶点与其内切圆在对边上的切点相连所得三条直线的交点。
（Joseph Gergonne，法国数学家。）

\begin{tkzexample}[latex=8cm,small]
\begin{tikzpicture}
\tkzDefPoints{0/0/B,3.6/0/C,2.8/4/A}
\tkzDefTriangleCenter[gergonne](A,B,C)
\tkzGetPoint{Ge}
\tkzDefSpcTriangle[intouch](A,B,C){C_1,C_2,C_3}
\tkzDefCircle[in](A,B,C) \tkzGetPoints{I}{i}
\tkzDrawLines[add=.25 and .25,teal](A,B A,C B,C)
\tkzDrawSegments[new](A,C_1 B,C_2 C,C_3)
\tkzDrawPoints(A,...,C,C_1,C_2,C_3)
\tkzDrawPoints[red](Ge)
\tkzLabelPoints(B,C,C_1,Ge)
\tkzLabelPoints[above](A,C_2,C_3)
\end{tikzpicture}
\end{tkzexample}

\subsubsection{选项 \tkzname{nagel}}
设 $Ta$ 为圆心在 $Ja$ 的旁切圆与三角形 $ABC$ 的边 $BC$ 的切点，类似地定义 $Tb$ 和 $Tc$。则直线 $ATa$、$BTb$、$CTc$ 交于 Nagel 点 $Na$。\\
\href{https://mathworld.wolfram.com/NagelPoint.html}{Weisstein, Eric W. "Nagel point." From MathWorld--A Wolfram Web Resource. }


\begin{tkzexample}[latex=7cm,small]
  \begin{tikzpicture}[scale=.5]
  \tkzDefPoints{0/0/A,6/0/B,4/6/C}
  \tkzDefSpcTriangle[ex](A,B,C){Ja,Jb,Jc}
  \tkzDefSpcTriangle[extouch](A,B,C){Ta,Tb,Tc}
  \tkzDefTriangleCenter[nagel](A,B,C)
  \tkzGetPoint{Na}
  \tkzDrawPolygon[blue](A,B,C)
  \tkzDrawLines[add=0 and 1](A,Ta B,Tb C,Tc)
  \tkzDrawPoints[new](Ja,Jb,Jc,Ta,Tb,Tc)
  \tkzClipBB
  \tkzDrawLines[add=1 and 1,dashed](A,B B,C C,A)
  \tkzDrawCircles[new](Ja,Ta Jb,Tb Jc,Tc)
  \tkzDrawSegments[new,dashed](Ja,Ta Jb,Tb Jc,Tc)
  \tkzDrawPoints(B,C,A)
  \tkzDrawPoints[new](Na)
  \tkzLabelPoints(B,C,A)
  \tkzLabelPoints[new](Na)
  \tkzLabelPoints[new](Ja,Jb,Jc,Ta,Tb,Tc)
  \tkzMarkRightAngles[fill=gray!20](Ja,Ta,C
              Jb,Tb,A Jc,Tc,B)
  \end{tikzpicture}
\end{tkzexample}


\subsubsection{Option   \tkzname{mittenpunkt}}

The mittenpunkt (also called the middlespoint) of a triangle $ABC$ is the symmedian point of the excentral triangle, i.e., the point of concurrence $M$ of the lines from the excenters through the midpoints of the corresponding sides.\\
\href{https://mathworld.wolfram.com/Mittenpunkt.html}{Weisstein, Eric W. "Mittenpunkt." From MathWorld--A Wolfram Web Resource.}


\begin{tkzexample}[latex=8cm,small]
\begin{tikzpicture}[scale=.5]
 \tkzDefPoints{0/0/A,6/0/B,4/6/C}
 \tkzDefSpcTriangle[centroid](A,B,C){Ma,Mb,Mc}
 \tkzDefSpcTriangle[ex](A,B,C){Ja,Jb,Jc}
 \tkzDefSpcTriangle[extouch](A,B,C){Ta,Tb,Tc}
 \tkzDefTriangleCenter[mittenpunkt](A,B,C)
 \tkzGetPoint{Mi}
 \tkzDrawPoints[new](Ma,Mb,Mc,Ja,Jb,Jc)
 \tkzClipBB
 \tkzDrawPolygon[blue](A,B,C)
 \tkzDrawLines[add=0 and 1](Ja,Ma
               Jb,Mb Jc,Mc)
 \tkzDrawLines[add=1 and 1](A,B A,C B,C)
 \tkzDrawCircles[new](Ja,Ta Jb,Tb Jc,Tc)
 \tkzDrawPoints(B,C,A)
 \tkzDrawPoints[new](Mi)
 \tkzLabelPoints(Mi)
 \tkzLabelPoints[left](Mb)
 \tkzLabelPoints[new](Ma,Mc,Jb,Jc)
 \tkzLabelPoints[above left](Ja,Jc)
\end{tikzpicture}
\end{tkzexample}

\subsubsection{Relation between  \tkzname{gergonne}, \tkzname{centroid} and \tkzname{mittenpunkt}}

Gergonne 点 $Ge$、重心 $G$ 和 mittenpunkt $M$ 三点共线，且 $GeG/GM=2$。

\begin{tkzexample}[vbox,small]
\begin{tikzpicture}
\tkzDefPoints{0/0/A,2/2/B,8/0/C}
\tkzDefTriangleCenter[gergonne](A,B,C) \tkzGetPoint{Ge}
\tkzDefTriangleCenter[centroid](A,B,C)
\tkzGetPoint{G}
\tkzDefTriangleCenter[mittenpunkt](A,B,C)
\tkzGetPoint{M}
\tkzDrawLines[add=.25 and .25,teal](A,B A,C B,C)
\tkzDrawLines[add=.25 and .25,new](Ge,M)
\tkzDrawPoints(A,...,C)
\tkzDrawPoints[red,size=2](G,M,Ge)
\tkzLabelPoints(A,...,C,M,G,Ge)
\tkzMarkSegment[mark=s||](Ge,G)
\tkzMarkSegment[mark=s|](G,M)
\end{tikzpicture}
\end{tkzexample}

\endinput
\section{通过变换定义点}
可用的变换包括：

\begin{itemize}
   \item 平移；
   \item 位似；
   \item 正交反射（即轴对称）；
   \item 中心对称；
   \item 正交投影；
   \item 旋转（角度制或弧度制）；
   \item 关于某个圆的反演。
\end{itemize}

\subsection{\tkzcname{tkzDefPointBy}}
变换通过选项来选择。共有两个宏：\tkzcname{tkzDefPointBy} 用于变换单个点，\tkzcname{tkzDefPointsBy} 用于变换一组点。默认情况下，点 $A$ 的像记为 $A'$。例如：
\begin{tkzltxexample}[]
\tkzDefPointBy[translation= from A to A'](B) 
\end{tkzltxexample}
结果保存在 \tkzname{tkzPointResult} 中。

\medskip
\begin{NewMacroBox}{tkzDefPointBy}{\oarg{local options}\parg{pt}}%
参数为一个已存在的点，其像保存在 \tkzname{tkzPointResult} 中。若要保留该点，可使用 \tkzcname{tkzGetPoint\{M\}} 将其命名为 \tkzname{M}。

\begin{tabular}{lll}%
\toprule
参数 &  含义 & 示例               \\ 
\midrule
\TAline{pt}   {已有点的名称}   {$(A)$}
\bottomrule
\end{tabular}

\begin{tabular}{lll}%
选项     &     & 示例                         \\ 
\midrule
\TOline{translation}{= from \#1 to \#2}{[translation=from A to B](E)}
\TOline{homothety}  {= center \#1 ratio \#2}{[homothety=center A ratio .5](E)}
\TOline{reflection} {= over \#1--\#2}{[reflection=over A--B](E)}
\TOline{symmetry }  {= center \#1}{[symmetry=center A](E)}
\TOline{projection }{= onto \#1--\#2}{[projection=onto A--B](E)}
\TOline{rotation }  {= center \#1 angle \#2}{[rotation=center O angle 30](E)}
\TOline{rotation in rad}{= center \#1 angle \#2}{[rotation in rad=center O angle pi/3](E)} 
\TOline{rotation with nodes}{= center \#1 from \#2 to \#3}{[center O from A to B](E)} 
\TOline{inversion}{= center \#1 through \#2}{[inversion =center O through A](E)} 
\TOline{inversion negative}{= center \#1 through \#2}{...} 
\bottomrule
\end{tabular}

\medskip
\emph{该宏仅定义像点，并不绘制。}
\end{NewMacroBox} 

\subsubsection{\tkzname{translation}（平移）} 

\begin{tkzexample}[latex=7cm,small]
\begin{tikzpicture}[>=latex] 
 \tkzDefPoints{0/0/A,3/1/B,3/0/C}
 \tkzDefPointBy[translation= from B to A](C) 
 \tkzGetPoint{D} 
 \tkzDrawPoints[teal](A,B,C,D)
 \tkzLabelPoints[color=teal](A,B,C,D) 
 \tkzDrawSegments[orange,->](A,B D,C) 
\end{tikzpicture} 
\end{tkzexample}

\subsubsection{\tkzname{reflection}（反射，即轴对称）} 

\begin{tkzexample}[latex=7cm,small] 
\begin{tikzpicture}[scale=.75]
 \tkzDefPoints{-2/-2/A,-1/-1/C,-4/2/D,-4/0/O}    
 \tkzDrawCircle(O,A)
 \tkzDefPointBy[reflection = over C--D](A)
 \tkzGetPoint{A'}
 \tkzDefPointBy[reflection = over C--D](O)
 \tkzGetPoint{O'}
 \tkzDrawCircle(O',A')
 \tkzDrawLine[add= .5 and .5](C,D)
 \tkzDrawPoints(C,D,O,O')
\end{tikzpicture}
\end{tkzexample}


\subsubsection{\tkzname{homothety}（位似）与 \tkzname{projection}（投影）}

\begin{tkzexample}[latex=7cm,small] 
\begin{tikzpicture}
  \tkzDefPoints{0/1/A,5/3/B,3/4/C}
  \tkzDefLine[bisector](B,A,C) \tkzGetPoint{a} 
  \tkzDrawLine[add=0 and 0,color=magenta!50 ](A,a) 
  \tkzDefPointBy[homothety=center A ratio .5](a)   
  \tkzGetPoint{a'} 
  \tkzDefPointBy[projection = onto A--B](a') 
  \tkzGetPoint{k'}  
  \tkzDefPointBy[projection = onto A--B](a) 
  \tkzGetPoint{k} 
  \tkzDrawLines[add= 0 and .3](A,k A,C)   
  \tkzDrawSegments[blue](a',k' a,k) 
  \tkzDrawPoints(a,a',k,k',A)
  \tkzDrawCircles(a',k' a,k)   
  \tkzLabelPoints(a,a',k,A)
\end{tikzpicture}
\end{tkzexample}  


\subsubsection{\tkzname{projection}（投影）}
\begin{tkzexample}[latex=7cm,small] 
\begin{tikzpicture}[scale=1.25]  
 \tkzDefPoints{0/0/A,0/4/B}
 \tkzDefTriangle[pythagore](B,A) \tkzGetPoint{C}
 \tkzDefLine[bisector](B,C,A) \tkzGetPoint{c}
 \tkzInterLL(C,c)(A,B)  \tkzGetPoint{D}
 \tkzDefPointBy[projection=onto B--C](D) 
 \tkzGetPoint{G}
 \tkzInterLC(C,D)(D,A) \tkzGetPoints{E}{F}
 \tkzDrawPolygon(A,B,C)
 \tkzDrawSegment(C,D)
 \tkzDrawCircle(D,A)
 \tkzDrawSegment[new](D,G)
 \tkzMarkRightAngle[fill=orange!10](D,G,B)
 \tkzDrawPoints(A,C,F) \tkzLabelPoints(A,C,F)
 \tkzDrawPoints(B,D,E,G)   
 \tkzLabelPoints[above right](B,D,E)
  \tkzLabelPoints[above](G)
 \end{tikzpicture}
 \end{tkzexample} 

\subsubsection{\tkzname{symmetry}（中心对称）} 
\begin{tkzexample}[latex=6cm,small] 
\begin{tikzpicture}[scale=1]
  \tkzDefPoints{2/-1/A,2/2/B,0/0/O}
  \tkzDefPointsBy[symmetry=center O](B,A){}
  \tkzDrawLine(A,A')
  \tkzDrawLine(B,B')
  \tkzMarkAngle[mark=s,arc=lll,
      size=1.5,mkcolor=red](A,O,B) 
  \tkzLabelAngle[pos=2,circle,draw,
    fill=blue!10,font=\scriptsize](A,O,B){$60^{\circ}$}
  \tkzDrawPoints(A,B,O,A',B') 
  \tkzLabelPoints(B,B')
  \tkzLabelPoints[below](A,O,A')  
\end{tikzpicture}   
\end{tkzexample}

\subsubsection{\tkzname{rotation}（旋转）} 
\begin{tkzexample}[latex=7cm,small] 
\begin{tikzpicture}[scale=0.75] 
 \tkzDefPoints{0/0/A,5/0/B}
 \tkzDrawSegment(A,B)
 \tkzDefPointBy[rotation=center A angle 60](B) 
 \tkzGetPoint{C} 
 \tkzDefPointBy[symmetry=center C](A) 
 \tkzGetPoint{D} 
 \tkzDrawSegment(A,tkzPointResult) 
 \tkzDrawLine(B,D)
 \tkzDrawArc(A,B)(C) \tkzDrawArc(B,C)(A)
 \tkzDrawArc(C,D)(D)  
 \tkzMarkRightAngle(D,B,A) 
 \tkzDrawPoints(A,B) 
 \tkzLabelPoints(A,B)
 \tkzLabelPoints[above](C)
 \tkzLabelPoints[right](D)
\end{tikzpicture}  
\end{tkzexample}  

\subsubsection{\tkzname{rotation in radian}（弧度制旋转）} 
\begin{tkzexample}[latex=6cm,small]
\begin{tikzpicture}
  \tkzDefPoint["$A$" left](1,5){A}
  \tkzDefPoint["$B$" right](4,3){B}
  \tkzDefPointBy[rotation in rad= center A angle pi/3](B)
  \tkzGetPoint{C}  
  \tkzDrawSegment(A,B)
  \tkzDrawPoints(A,B,C) 
  \tkzCompass(A,C)
  \tkzCompass(B,C) 
  \tkzLabelPoints(C)
\end{tikzpicture}
\end{tkzexample} 

\subsubsection{\tkzname{rotation with nodes}（按节点旋转）} 
\begin{tkzexample}[latex=6cm,small]
\begin{tikzpicture}
 \tkzDefPoint(0,0){O}    
 \tkzDefPoint(0:2){A} 
 \tkzDefPoint(40:2){B}  
 \tkzDefPoint(20:4){C}
 \tkzDrawLine(O,A)
 \tkzDefPointBy[rotation with nodes%
             =center O from A to B](C)  
 \tkzGetPoint{D}
\tkzDrawPoints(A,B,C,D)
\tkzDrawCircle(O,A)
\tkzLabelPoints(A,C,D)
\tkzLabelPoints[above](B)
\end{tikzpicture}
\end{tkzexample} 

\subsubsection{\tkzname{inversion}（反演）} 

反演是将点变换为其对应反演点集合的过程。设反演圆的圆心为 $O$、反演半径为 $k$，若 $P'$ 是三角形 $OQP$ 中从 $O$ 引出的高的垂足，其中 $Q$ 是圆上满足 $OQ \perp PQ$ 的一点，则称 $P$ 与 $P'$ 关于该反演圆互为反演点。\\
量 $k^2$ 称为圆幂（Coxeter 1969, p. 81）。
 
(\url{https://mathworld.wolfram.com/Inversion.html})

若干性质：
\begin{itemize}
\item 不过反演中心的圆，其反形仍为圆。
\item 过反演中心的圆，其反形为直线。
\item 不过反演中心的直线，其反形为过反演中心的圆。
\item 与反演圆正交的圆，其反形为自身。
\item 过反演中心的直线，其反形为自身。
\item 反演保持夹角不变。
\end{itemize}

说明：

直接构造：
（圆心 $O$，幂 $k^2={OA}^2=OP \times OP'$）

\begin{tkzexample}[latex=6cm,small]
\begin{tikzpicture}[scale=.5]
  \tkzDefPoints{4/0/A,6/0/P,0/0/O}
  \tkzDefPointBy[inversion = center O through A](P)
  \tkzGetPoint{P'}
  \tkzDrawSegments(O,P)
  \tkzDrawCircle(O,A)
  \tkzLabelPoints[above right,font=\scriptsize](O,A,P,P')
  \tkzDrawPoints(O,A,P,P')
\end{tikzpicture}
\end{tkzexample} 

\begin{tkzexample}[latex=6cm,small]
\begin{tikzpicture}[scale=.5]
  \tkzDefPoints{4/0/A,6/0/P,0/0/O}
  \tkzDefLine[orthogonal=through P](O,P)
  \tkzGetPoint{L}
  \tkzDefLine[tangent from = P](O,A) \tkzGetPoints{R}{Q}
  \tkzDefPointBy[projection=onto O--A](Q) \tkzGetPoint{P'}
  \tkzDrawSegments(O,P O,A)
  \tkzDrawSegments[new](O,P O,Q P,Q Q,P')
  \tkzDrawCircle(O,A)
  \tkzDrawLines[add=1 and 0](P,L)
  \tkzLabelPoints[below,font=\scriptsize](O,P')
  \tkzLabelPoints[above right,font=\scriptsize](P,Q)
  \tkzDrawPoints(O,P) \tkzDrawPoints[new](Q,P')
  \tkzLabelSegment[above](O,Q){$k$}
  \tkzMarkRightAngles(A,P',Q P,Q,O)
  \tkzLabelCircle[above=.5cm,
      font=\scriptsize](O,A)(100){inversion circle}
  \tkzLabelPoint[left,font=\scriptsize](O){inversion center}
  \tkzLabelPoint[left,font=\scriptsize](L){polar}
\end{tikzpicture}
\end{tkzexample} 


\subsubsection{直线的反演：示例 1}
\begin{tkzexample}[latex=6cm,small]  
\begin{tikzpicture}[scale=.5]
\tkzDefPoints{0/0/O,3/0/I,4/3/P,6/-3/Q}
\tkzDrawCircle(O,I)
\tkzDefPointBy[projection= onto P--Q](O) \tkzGetPoint{A}
\tkzDefPointBy[inversion = center O through I](A)
\tkzGetPoint{A'}
\tkzDefPointBy[inversion = center O through I](P)
\tkzGetPoint{P'}
\tkzDefCircle[diameter](O,A')\tkzGetPoint{o}
\tkzDrawCircle[new](o,A')
\tkzDrawLines[add=.25 and .25,red](P,Q)
\tkzDrawLines[add=.25 and .25](O,A)
\tkzDrawSegments(O,P)
\tkzDrawPoints(A,P,O) \tkzDrawPoints[new](A',P')
\end{tikzpicture}
\end{tkzexample} 

\subsubsection{直线的反演：示例 2}
\begin{tkzexample}[latex=6cm,small]  
\begin{tikzpicture}[scale=.8]
\tkzDefPoints{0/0/O,3/0/I,3/2/P,3/-2/Q}
\tkzDrawCircle(O,I)
\tkzDefPointBy[projection= onto P--Q](O) \tkzGetPoint{A}
\tkzDefPointBy[inversion = center O through I](A)
\tkzGetPoint{A'}
\tkzDefPointBy[inversion = center O through I](P)
\tkzGetPoint{P'}
\tkzDefCircle[diameter](O,A')\tkzGetPoint{o}
\tkzDrawCircle[new](o,A')
\tkzDrawLines[add=.25 and .25,red](P,Q)
\tkzDrawLines[add=.25 and .25](O,A)
\tkzDrawSegments(O,P)
\tkzDrawPoints(A,P,O) \tkzDrawPoints[new](A',P')
\end{tikzpicture}
\end{tkzexample}

\subsubsection{直线的反演：示例 3}
\begin{tkzexample}[latex=6cm,small]  
\begin{tikzpicture}[scale=.8]
\tkzDefPoints{0/0/O,3/0/I,2/1/P,2/-2/Q}
\tkzDrawCircle(O,I)
\tkzDefPointBy[projection= onto P--Q](O) \tkzGetPoint{A}
\tkzDefPointBy[inversion = center O through I](A)
\tkzGetPoint{A'}
\tkzDefPointBy[inversion = center O through I](P)
\tkzGetPoint{P'}
\tkzDefCircle[diameter](O,A')
\tkzDrawCircle[new](I,A')
\tkzDrawLines[add=.25 and .75,red](P,Q)
\tkzDrawLines[add=.25 and .25](O,A')
\tkzDrawSegments(O,P')
\tkzDrawPoints(A,P,O) \tkzDrawPoints[new](A',P')
\end{tikzpicture}
\end{tkzexample}

\subsubsection{圆的反演与 \tkzname{homothety}（位似）}
\begin{tkzexample}[latex=7cm,small]
\begin{tikzpicture}[scale=.7]
\tkzDefPoints{0/0/O,3/2/A,2/1/P}
\tkzDefLine[tangent from = O](A,P) \tkzGetPoints{T}{X}
\tkzDefPointsBy[homothety = center O
                ratio 1.25](A,P,T){}
\tkzInterCC(A,P)(A',P') \tkzGetPoints{C}{D}
\tkzCalcLength(A,P)
\tkzGetLength{rAP}
\tkzDefPointOnCircle[R=center A angle 190 radius \rAP]
\tkzGetPoint{M}
\tkzDefPointBy[inversion = center O through C](M)
\tkzGetPoint{M'}
\tkzDrawCircles[new](A,P A',P')
\tkzDrawCircle(O,C)
\tkzDrawLines[add=0 and .5](O,T' O,A' O,M' O,P')
\tkzDrawPoints(A,A',P,P',O,T,T',M,M')
\tkzLabelPoints(O,T,T',M,M')
\tkzLabelPoints[below](P,P')
\end{tikzpicture}
\end{tkzexample}


\subsubsection{三角形关于其内切圆的 \tkzname{inversion}（反演）}
\begin{tkzexample}[latex=6cm,small] 
\begin{tikzpicture}[scale=1]
\tkzDefPoints{0/0/A,5/1/B,3/6/C}
\tkzDefTriangleCenter[in](A,B,C) \tkzGetPoint{O} 
\tkzDefPointBy[projection= onto A--C](O) \tkzGetPoint{b}
\tkzDefPointBy[projection= onto A--C](O) \tkzGetPoint{b}
\tkzDefPointBy[projection= onto B--C](O) \tkzGetPoint{a}
\tkzDefPointBy[projection= onto A--B](O) \tkzGetPoint{c}
\tkzDefPointsBy[inversion = center O through b](a,b,c)%
                                             {Ia,Ib,Ic}
\tkzDefMidPoint(O,Ia) \tkzGetPoint{Ja}
\tkzDefMidPoint(O,Ib) \tkzGetPoint{Jb}
\tkzDefMidPoint(O,Ic) \tkzGetPoint{Jc}
\tkzInterCC(Ja,O)(Jb,O) \tkzGetPoints{O}{x}
\tkzInterCC(Ja,O)(Jc,O) \tkzGetPoints{y}{O}
\tkzInterCC(Jb,O)(Jc,O) \tkzGetPoints{O}{z}
\tkzDrawPolygon(A,B,C)
\tkzDrawCircle(O,b)\tkzDrawPoints(A,B,C,O)
\tkzDrawCircles[dashed,gray](Ja,y Jb,x Jc,z)
\tkzDrawArc[line width=1pt,orange,delta=0](Jb,x)(z)
\tkzDrawArc[line width=1pt,orange,delta=0](Jc,z)(y)
\tkzDrawArc[line width=1pt,orange,delta=0](Ja,y)(x)
\tkzLabelPoint[below](A){$A$}\tkzLabelPoint[above](C){$C$}
\tkzLabelPoint[right](B){$B$}\tkzLabelPoint[below](O){$O$}
\end{tikzpicture}
\end{tkzexample}

\subsubsection{与反演圆正交的圆的 \tkzname{inversion}（反演）} 
反演圆本身、与它正交的圆，以及过反演中心的直线，在反演下均保持不变。若一个圆与反演圆相交，则交点保持不变，因而这些点也位于反形圆上：见下图中的 $I$ 与 $J$。

\begin{tkzexample}[latex=5cm,small]
\begin{tikzpicture}[scale=1]
\tkzDefPoint(0,0){O}\tkzDefPoint(1,0){A}
\tkzDefPoint(-1.5,-1.5){z1} 
\tkzDefPoint(1.5,-1.25){z2} 
\tkzDefCircle[orthogonal through=z1 and z2](O,A)
\tkzGetPoint{c} 
\tkzDrawCircle[new](c,z1) 
\tkzDefPointBy[inversion =  center O through A](z1)
\tkzGetPoint{Z1} 
\tkzInterCC(O,A)(c,z1) \tkzGetPoints{I}{J}
\tkzDefPointBy[inversion =  center O through A](I)
\tkzGetPoint{I'}
\tkzDrawCircle(O,A)
\tkzDrawPoints(O,A,z1,z2) 
\tkzDrawPoints[new](c,Z1,I,J) 
\tkzLabelPoints(O,A,z1,z2,c,Z1,I,J)
\end{tikzpicture}
\end{tkzexample}



更复杂的示例参见 \tkzname{Pappus} 一节（\ref{pappus}）。

\subsubsection{\tkzname{inversion negative}（负反演）}
即先做反演，再以 $O$ 为中心做一次中心对称。
\begin{tkzexample}[latex=5cm,small]  
\begin{tikzpicture}[scale=1.5]
  \tkzDefPoints{1/0/A,0/0/O}
  \tkzDefPoint(-1.5,-1.5){z1}
  \tkzDefPoint(0.35,-2){z2} 
  \tkzDefPointBy[inversion negative = center O through A](z1)
  \tkzGetPoint{Z1} 
  \tkzDefPointBy[inversion negative = center O through A](z2)
  \tkzGetPoint{Z2}
  \tkzDrawCircle(O,A)  
  \tkzDrawPoints[color=black, fill=red,size=4](Z1,Z2)    
  \tkzDrawSegments(z1,Z1 z2,Z2)
  \tkzDrawPoints[color=black, fill=red,size=4](O,z1,z2)
  \tkzLabelPoints[font=\scriptsize](O,A,z1,z2,Z1,Z2)  
\end{tikzpicture}
\end{tkzexample} 


\newpage
\subsection{多个点的变换：\tkzcname{tkzDefPointsBy}} 
前一个宏的变体，用于同时定义多个像点。
你可以将像点的名称作为参数给出，或将该参数留空：此时像点名称由原点名称加撇号构成。

\begin{tkzltxexample}[]
\tkzDefPointsBy[translation= from A to A'](B,C){}
\end{tkzltxexample}
像点为 $B'$ 与 $C'$。

\begin{tkzltxexample}[]
\tkzDefPointsBy[translation= from A to A'](B,C){D,E} 
\end{tkzltxexample}
像点为 $D$ 与 $E$。

\begin{tkzltxexample}[]
\tkzDefPointsBy[translation= from A to A'](B)
\end{tkzltxexample}
像点为 $B'$。
\begin{NewMacroBox}{tkzDefPointsBy}{\oarg{local options}\parg{list of points}\marg{list of points}}%
\begin{tabular}{lll}%
参数 &  示例  &                  \\ 
\midrule
\TAline{\parg{list of points}\marg{list of pts}}{(A,B)\{E,F\}}{$E$、$F$ 分别为 $A$、$B$ 的像}   \\
\bottomrule
\end{tabular}

\medskip
若像点列表为空，则每个像点的名称为原点名称后加“'”。

\medskip
\begin{tabular}{lll}%
\toprule
选项     &     & 示例                         \\ 
\midrule
\TOline{translation = from \#1 to \#2}{}{[translation=from A to B](E)\{\}}
\TOline{homothety = center \#1 ratio \#2}{}{[homothety=center A ratio .5](E)\{F\}}
\TOline{reflection = over \#1--\#2}{}{[reflection=over A--B](E)\{F\}}
\TOline{symmetry = center \#1}{}{[symmetry=center A](E)\{F\}}
\TOline{projection = onto \#1--\#2}{}{[projection=onto A--B](E)\{F\}}
\TOline{rotation = center \#1 angle \#2}{}{[rotation=center  angle 30](E)\{F\}}
\TOline{rotation in rad = center \#1 angle \#2}{}{例如 angle pi/3}
\TOline{rotation with nodes = center \#1 from \#2 to \#3}{}{[center O from A to B](E)\{F\}} 
\TOline{inversion = center \#1 through \#2}{}{[inversion = center O through A](E)\{F\}} 
\TOline{inversion negative = center \#1 through \#2}{}{...} 
\bottomrule
\end{tabular}

\medskip
这些点仅被定义，并不绘制。
\end{NewMacroBox}

\subsubsection{多个点的 \tkzname{translation}（平移）}
\begin{tkzexample}[latex=7cm,small]
\begin{tikzpicture}[>=latex] 
 \tkzDefPoints{0/0/A,3/0/B,3/1/A',1/2/C}
 \tkzDefPointsBy[translation= from A to A'](B,C){} 
 \tkzDrawPolygon(A,B,C)
 \tkzDrawPolygon[new](A',B',C')
 \tkzDrawPoints(A,B,C)
 \tkzDrawPoints[new](A',B',C') 
 \tkzLabelPoints(A,B,A',B')  
 \tkzLabelPoints[above](C,C')
 \tkzDrawSegments[color = gray,->,
              style=dashed](A,A' B,B' C,C') 
\end{tikzpicture}
\end{tkzexample}

\subsubsection{多个点的 \tkzname{symmetry}（中心对称）：一个卵形}

\begin{tkzexample}[latex=7cm,small]
\begin{tikzpicture}[scale=0.4]
  \tkzDefPoint(-4,0){I}
  \tkzDefPoint(4,0){J}
  \tkzDefPoint(0,0){O} 
  \tkzInterCC(J,O)(O,J) \tkzGetPoints{L}{H}
  \tkzInterCC(I,O)(O,I) \tkzGetPoints{K}{G} 
  \tkzInterLL(I,K)(J,H) \tkzGetPoint{M}
  \tkzInterLL(I,G)(J,L) \tkzGetPoint{N}
  \tkzDefPointsBy[symmetry=center J](L,H){D,E} 
  \tkzDefPointsBy[symmetry=center I](G,K){C,F}
  \begin{scope}[line style/.style = {very thin,teal}]
    \tkzDrawLines[add=1.5 and 1.5](I,K I,G J,H J,L) 
    \tkzDrawLines[add=.5 and .5](I,J) 
    \tkzDrawCircles(O,I I,O J,O) 
    \tkzDrawArc[delta=0,orange](N,D)(C) 
    \tkzDrawArc[delta=0,orange](M,F)(E) 
    \tkzDrawArc[delta=0,orange](J,E)(D) 
    \tkzDrawArc[delta=0,orange](I,C)(F) 
  \end{scope}   
\end{tikzpicture} 
\end{tkzexample}

\endinput

\section{用向量定义点}

\subsection{\tkzcname{tkzDefPointWith}}
满足特定向量条件的点可由 \tkzcname{tkzDefPointWith} 创建。基本原理如下：将两个点（即一个向量）作为参数传入。根据所选选项的不同，新点与第一个点（个别情况除外）所构成的向量，与原向量共线或垂直，其长度或与原向量长度成正比，或与单位长度成正比。由于这种点通常只是临时使用，不必立即命名：结果保存在 \tkzname{tkzPointResult} 中，再用 \tkzNameMacro{tkzGetPoint} 以另一个名称取出。

某些选项会设定给定点与所得点之间的距离。一般情况下，该距离为作为参数给出的两点之间的距离；若选用 \tkzname{normed} 类选项，则该距离为 1~cm。选项 $K$ 则给出该距离的倍数。

\begin{NewMacroBox}{tkzDefPointWith}{\parg{pt1,pt2}}%
 本宏用于定义满足向量条件的点。

\medskip
  
\begin{tabular}{lll}%
\toprule
参数             & 含义 & 说明                         \\ 
\midrule
\TAline{(pt1,pt2)} {点对}{结果点保存在 \tkzname{tkzPointResult} 中} \\

\bottomrule
\end{tabular}

\medskip
下文假设所得点均通过 \tkzNameMacro{tkzGetPoint\{C\}} 取出。

\begin{tabular}{lll}%
\toprule
选项             & 示例 & 说明                         \\ 
\midrule
\TOline{orthogonal}{[orthogonal](A,B)}{$AC=AB$ 且 $\overrightarrow{AC} \perp \overrightarrow{AB}$}
\TOline{orthogonal normed}{[orthogonal normed](A,B)}{$AC=1$ 且 $\overrightarrow{AC} \perp \overrightarrow{AB}$} 
\TOline{linear}{[linear](A,B)}{$\overrightarrow{AC}=K \times \overrightarrow{AB}$}
\TOline{linear normed}{[linear normed](A,B)}{$AC=K$ 且 $\overrightarrow{AC}=k\times \overrightarrow{AB}$ }  
\TOline{colinear= at \#1}{[colinear= at C](A,B)}{$\overrightarrow{CD}= \overrightarrow{AB}$ }
\TOline{colinear normed= at \#1}{[colinear normed= at C](A,B)}{$\overrightarrow{CD}= \overrightarrow{AB}$ }
\TOline{K}{[linear](A,B),K=2}{$\overrightarrow{AC}=2\times \overrightarrow{AB}$}
\end{tabular}
\end{NewMacroBox}

\subsubsection{选项 \tkzname{colinear at}：简单示例}
 $(\overrightarrow{AB}=\overrightarrow{CD})$
\begin{tkzexample}[latex=6cm,small]
\begin{tikzpicture}[scale=1.2,
   vect/.style={->,shorten >=1pt,>=latex'}]
  \tkzDefPoint(2,3){A}   \tkzDefPoint(4,2){B}
  \tkzDefPoint(0,1){C}
  \tkzDefPointWith[colinear=at C](A,B)
  \tkzGetPoint{D}
  \tkzDrawPoints[new](A,B,C,D)
  \tkzLabelPoints[above right=3pt](A,B,C,D)
  \tkzDrawSegments[vect](A,B C,D)
\end{tikzpicture}
\end{tkzexample}

\subsubsection{选项 \tkzname{colinear at}：较复杂的示例}
\begin{tkzexample}[vbox,small]
\begin{tikzpicture}[scale=.75]
\tkzDefPoints{0/0/B,3.6/0/C,1.5/4/A}
\tkzDefSpcTriangle[ortho](A,B,C){Ha,Hb,Hc}
\tkzDefTriangleCenter[ortho](A,B,C) \tkzGetPoint{H}
\tkzDefSquare(A,C) \tkzGetPoints{R}{S}
\tkzDefSquare(B,A) \tkzGetPoints{M}{N}
\tkzDefSquare(C,B) \tkzGetPoints{P}{Q}
\tkzDefPointWith[colinear= at M](A,S) \tkzGetPoint{A'}
\tkzDefPointWith[colinear= at P](B,N) \tkzGetPoint{B'}
\tkzDefPointWith[colinear= at Q](C,R) \tkzGetPoint{C'}
\tkzDefPointBy[projection=onto P--Q](Ha) \tkzGetPoint{Pa}
\tkzDrawPolygon[teal,thick](A,C,R,S)\tkzDrawPolygon[teal,thick](A,B,N,M)
\tkzDrawPolygon[teal,thick](C,B,P,Q)
\tkzDrawPoints[teal,size=2](A,B,C,Ha,Hb,Hc,A',B',C')
\tkzDrawSegments[ultra thin,red](M,A' A',S P,B' B',N Q,C' C',R B,S C,M C,N B,R A,P A,Q)
\tkzDrawSegments[ultra thin,teal, dashed](A,Ha B,Hb C,Hc)
\tkzDefPointBy[rotation=center A angle 90](S) \tkzGetPoint{S'}
\tkzDrawSegments[ultra thin,teal,dashed](B,S' A,S' A,A' M,S' B',Q P,C' M,S Ha,Pa)
\tkzDrawArc(A,S)(S')
\end{tikzpicture}
\end{tkzexample}

\subsubsection{选项 \tkzname{colinear at}：如何使用 $K$}
$K$ 的用法
\begin{tkzexample}[latex=7cm,small]
\begin{tikzpicture}[vect/.style={->,
               shorten >=1pt,>=latex'}]
  \tkzDefPoints{0/0/A,5/0/B,1/2/C}
  \tkzDefPointWith[colinear=at C](A,B)
  \tkzGetPoint{G}
  \tkzDefPointWith[colinear=at C, K=0.5](A,B)
  \tkzGetPoint{H}
  \tkzLabelPoints(A,B,C,G,H)
  \tkzDrawPoints(A,B,C,G,H)
  \tkzDrawSegments[vect](A,B C,H)
\end{tikzpicture}
\end{tkzexample}

\subsubsection{选项 \tkzname{colinear at} 配合 $K=\sqrt{2}/2$}
取 $K=\frac{\sqrt{2}}{2}$

\begin{tkzexample}[latex=6cm,small]
\begin{tikzpicture}[vect/.style={->,
            shorten >=1pt,>=latex'}]
 \tkzDefPoints{1/1/A,4/2/B,2/2/C}
 \tkzDefPointWith[colinear=at C,K=sqrt(2)/2](A,B)
 \tkzGetPoint{D}
 \tkzDrawPoints[color=red](A,B,C,D)
 \tkzDrawSegments[vect](A,B C,D)
\end{tikzpicture}
\end{tkzexample}

\subsubsection{选项 \tkzname{orthogonal} 配合 $K=1$}
此时 $AB=AC$，因为 $K=1$。
\begin{tkzexample}[latex=6cm,small]
\begin{tikzpicture}[scale=1.2,
  vect/.style={->,shorten >=1pt,>=latex'}]
  \tkzDefPoints{2/3/A,4/2/B}
   \tkzDefPointWith[orthogonal,K=1](A,B)
     \tkzGetPoint{C}
   \tkzDrawPoints[color=red](A,B,C)
   \tkzLabelPoints[right=3pt](B,C)
   \tkzLabelPoints[below=3pt](A)
   \tkzDrawSegments[vect](A,B A,C)
   \tkzMarkRightAngle(B,A,C)
\end{tikzpicture}
\end{tkzexample}



\subsubsection{选项 \tkzname{orthogonal} 配合 $K=-1$}
 取 $K=-1$。
因 $\lvert K \rvert=1$，故 $OK=OI$，进而 $OI=OJ=OK$。

\begin{tkzexample}[latex=7cm,small]
\begin{tikzpicture}[scale=.75]
  \tkzDefPoints{1/2/O,2/5/I}
  \tkzDefPointWith[orthogonal](O,I)
  \tkzGetPoint{J}
  \tkzDefPointWith[orthogonal,K=-1](O,I)
  \tkzGetPoint{K}
  \tkzDrawSegment(O,I)
  \tkzDrawSegments[->](O,J O,K)
  \tkzMarkRightAngles(I,O,J I,O,K)
  \tkzDrawPoints(O,I,J,K)
  \tkzLabelPoints(O,I,J,K)
\end{tikzpicture}
\end{tkzexample}

\subsubsection{选项 \tkzname{orthogonal}：一个较复杂的示例}
\begin{tkzexample}[latex=7cm,small]
\begin{tikzpicture}[scale=.75]
  \tkzDefPoints{0/0/A,6/0/B}
  \tkzDefMidPoint(A,B)
    \tkzGetPoint{I}
  \tkzDefPointWith[orthogonal,K=-.75](B,A)
  \tkzGetPoint{C}
  \tkzInterLC(B,C)(B,I)
     \tkzGetPoints{D}{F}
  \tkzDuplicateSegment(B,F)(A,F)
  \tkzGetPoint{E}
  \tkzDrawArc[delta=10](F,E)(B)
  \tkzInterLC(A,B)(A,E)
    \tkzGetPoints{N}{M}
  \tkzDrawArc[delta=10](A,M)(E)
  \tkzDrawLines(A,B B,C A,F)
  \tkzCompass(B,F)
  \tkzDrawPoints(A,B,C,F,M,E)
  \tkzLabelPoints(A,B,C,F,M)
  \tkzLabelPoints[above](E)
\end{tikzpicture}
\end{tkzexample}

\subsubsection{选项 \tkzname{colinear} 与 \tkzname{orthogonal}}
\begin{tkzexample}[latex=7cm,small]
\begin{tikzpicture}[scale=1.2,
  vect/.style={->,shorten >=1pt,>=latex'}]
  \tkzDefPoints{2/1/A,6/2/B}
  \tkzDefPointWith[orthogonal,K=.5](A,B)
  \tkzGetPoint{C}
  \tkzDefPointWith[colinear=at C,K=.5](A,B)
  \tkzGetPoint{D}
  \tkzMarkRightAngle[fill=gray!20](B,A,C)
  \tkzDrawSegments[vect](A,B A,C C,D)
  \tkzDrawPoints(A,...,D)
\end{tikzpicture}
\end{tkzexample}

\subsubsection{选项 \tkzname{orthogonal normed}}
 $K=1$，故 $AC=1$。

\begin{tkzexample}[latex=7cm,small]
\begin{tikzpicture}[scale=1.2,
  vect/.style={->,shorten >=1pt,>=latex'}]
  \tkzDefPoints{2/3/A,4/2/B}
  \tkzDefPointWith[orthogonal normed](A,B)
  \tkzGetPoint{C}
  \tkzDrawPoints[color=red](A,B,C)
  \tkzDrawSegments[vect](A,B A,C)
  \tkzMarkRightAngle[fill=gray!20](B,A,C)
\end{tikzpicture}
\end{tkzexample}

\subsubsection{选项 \tkzname{orthogonal normed} 与 K=2}
$K=2$，因此 $AC=2$。

\begin{tkzexample}[latex=7cm,small]
\begin{tikzpicture}[scale=1.2,
   vect/.style={->,shorten >=1pt,>=latex'}]
  \tkzDefPoints{2/3/A,5/1/B}
  \tkzDefPointWith[orthogonal normed,K=2](A,B)
  \tkzGetPoint{C}
  \tkzDrawPoints[color=red](A,B,C)
  \tkzDefCircle[R](A,2) \tkzGetPoint{a}
  \tkzDrawCircle(A,a)
  \tkzDrawSegments[vect](A,B A,C)
  \tkzMarkRightAngle[fill=gray!20](B,A,C)
  \tkzLabelPoints[above=3pt](A,B,C)
\end{tikzpicture}
\end{tkzexample}

\subsubsection{选项 \tkzname{linear}}
此处 $K=0.5$。

这相当于做一次位似，或将一个向量乘以一个实数。下面取 $[AB]$ 的中点。

\begin{tkzexample}[latex=7cm,small]
\begin{tikzpicture}[scale=1.2]
  \tkzDefPoints{1/3/A,4/2/B}
  \tkzDefPointWith[linear,K=0.5](A,B)
  \tkzGetPoint{C}
  \tkzDrawPoints[color=red](A,B,C)
  \tkzDrawSegment(A,B)
  \tkzLabelPoints[above right=3pt](A,B,C)
\end{tikzpicture}
\end{tkzexample}

\subsubsection{选项 \tkzname{linear normed}}
在下面的示例中，$AC=1$，且 $C$ 位于直线 $(AB)$ 上。

\begin{tkzexample}[latex=7cm,small]
\begin{tikzpicture}[scale=1.2]
 \tkzDefPoints{1/3/A,4/2/B}
 \tkzDefPointWith[linear normed](A,B)
 \tkzGetPoint{C}
 \tkzDrawPoints[color=red](A,B,C)
 \tkzDrawSegment(A,B)
 \tkzLabelSegment(A,C){$1$}
 \tkzLabelPoints[above right=3pt](A,B,C)
\end{tikzpicture}
\end{tkzexample}
%<--------------------------------------------------------------------------–>
%         tkzGetVectxy
%<--------------------------------------------------------------------------–>
\subsection{\tkzcname{tkzGetVectxy} }
取出一个向量的坐标。

\begin{NewMacroBox}{tkzGetVectxy}{\parg{$A,B$}\var{text}}%
给出一个向量的坐标。

\medskip
\begin{tabular}{lll}%
\toprule
参数    & 示例 & 说明      \\

\midrule

\TAline{(point)\{name of macro\}} {\tkzcname{tkzGetVectxy}(A,B)\{V\}}{\tkzcname{Vx}、\tkzcname{Vy}：向量 $\overrightarrow{AB}$ 的坐标}
\end{tabular}
\end{NewMacroBox}

\subsubsection{用 \tkzcname{tkzGetVectxy} 传递坐标}

\begin{tkzexample}[latex=7cm,small]
\begin{tikzpicture}
 \tkzDefPoints{0/0/O,1/1/A,4/2/B}
 \tkzGetVectxy(A,B){v}
 \tkzDefPoint(\vx,\vy){V}
 \tkzDrawSegment[->,color=red](O,V)
 \tkzDrawSegment[->,color=blue](A,B)
 \tkzDrawPoints(A,B,O)
 \tkzLabelPoints(A,B,O,V)
\end{tikzpicture}
\end{tkzexample}
\endinput

\section{随机点}
%<--------------------------------------------------------------------------->
%           随机点
%<--------------------------------------------------------------------------->
可以放置随机点的位置包括：
\begin{enumerate}
  \item 矩形内部；
  \item 线段上；
  \item 直线上；
  \item 圆上；
  \item 圆盘内部。
\end{enumerate}

\subsection{获取随机点}
本宏取代了旧版的 \tkzcname{tkzGetRandPointOn}。
\begin{NewMacroBox}{tkzDefRandPointOn}{\oarg{local options}}%
{结果为一个位置随机的点，可用宏 \tkzcname{tkzGetPoint} 为其命名。若无需保留结果，可直接使用 \tkzname{tkzPointResult}。}

\medskip
\begin{tabular}{lll}%
\toprule
选项             & 默认值 & 含义                         \\ 
\midrule
\TOline{rectangle=pt1 and pt2}  {}{[rectangle=A and B]} 
\TOline{segment= pt1--pt2} {}{[segment=A--B]}
\TOline{line=pt1--pt2}{}{[line=A--B]} 
\TOline{circle =center pt1 radius dim}{}{[circle = center A radius 2]} 
\TOline{circle through=center pt1 through pt2}{}{[circle through= center A through B]}
\TOline{disk through=center pt1 through pt2}{}{[disk through=center A through B]}
\end{tabular}
\end{NewMacroBox} 

\subsubsection{矩形内的随机点} 

\begin{tkzexample}[latex=7cm,small]
\begin{tikzpicture}
  \tkzDefPoints{0/0/A,5/3/C}
  \tkzDefRandPointOn[rectangle = A and C]
  \tkzGetPoint{E}
  \tkzDefRectangle(A,C)\tkzGetPoints{B}{D}
  \tkzDrawPolygon[red](A,...,D)
  \tkzDrawPoints(A,...,E) 
  \tkzLabelPoints(A,B)
  \tkzLabelPoints[above](C,D,E) 
\end{tikzpicture} 
\end{tkzexample} 

\subsubsection{线段或直线上的随机点}  
\begin{tkzexample}[latex=7cm,small]
\begin{tikzpicture}  
  \tkzDefPoints{0/0/A,5/2/C}  
  \tkzDefRandPointOn[segment = A--C]\tkzGetPoint{B}
  \tkzDrawLine(A,C)
  \tkzDrawPoints(A,C) \tkzDrawPoint[red](B)
  \tkzLabelPoints(A,C) \tkzLabelPoints[red](B) 
\end{tikzpicture}
\end{tkzexample}


\subsubsection{圆或圆盘内的随机点}
\begin{tkzexample}[latex=7cm,small]
\begin{tikzpicture}
\tkzDefPoints{3/2/A,1/1/B}
\tkzCalcLength(A,B) \tkzGetLength{rAB} 
\tkzDefRandPointOn[circle = center A radius \rAB] 
\tkzGetPoint{C}
\tkzDefRandPointOn[circle through= center A through B]
\tkzGetPoint{D}
\tkzDefRandPointOn[disk through=center A through B]
\tkzGetPoint{E}
\tkzDrawCircle(A,B)
\tkzDrawPoints(A,B)
\tkzLabelPoints(A,B)
\tkzDrawPoints[red](C,D,E)
\tkzLabelPoints[red,right](C,D,E)
\end{tikzpicture}
\end{tkzexample}
   
\endinput


\part{直线、三角形与多边形}
\section{直线}

绘制直线当然是必不可少的，但首先我们必须能够定义一些特殊直线，例如中垂线（垂直平分线）、角平分线、平行线和垂线。基本思路是先确定直线上的两个点。

\subsection{直线的定义}

\begin{NewMacroBox}{tkzDefLine}{\oarg{local options}\parg{pt1,pt2} or \parg{pt1,pt2,pt3}}%
参数为两个或三个点构成的列表。根据情况的不同，本宏会定义得到该直线所需的一个或两个点，再用 \tkzcname{tkzGetPoint} 或 \tkzcname{tkzGetPoints} 取出。
本手册中用“中垂线”一词表示线段的垂直平分线。

\medskip
\begin{tabular}{lll}%
\toprule
参数           & 示例 & 说明                         \\
\midrule
\TAline{\parg{pt1,pt2}}{[mediator]\parg{A,B}}{线段 $[A,B]$ 的中垂线}
\TAline{\parg{pt1,pt2,pt3}}{[bisector]\parg{A,B,C}} {角 $\widehat{ABC}$ 的平分线}
\TAline{\parg{pt1,pt2,pt3}}{[altitude]\parg{A,B,C}} {从 $B$ 引出的高}
\TAline{\parg{pt1}}{[tangent at=A]\parg{O}} {圆心为 $O$ 的圆在点 $A$ 处的切线}
\TAline{\parg{pt1,pt2}}{[tangent from=A]\parg{O,B}} {圆心为 $O$ 且过 $B$ 的圆}
\end{tabular}

\medskip
\begin{tabular}{lll}%
\toprule
选项             & 默认值 & 含义                         \\ 
\TOline{mediator}{}{线段的垂直平分线}  
\TOline{perpendicular=through\dots}{mediator}{过某点且与已知直线垂直的直线} 
\TOline{orthogonal=through\dots}{mediator}{同上}
\TOline{parallel=through\dots}{mediator}{过某点且与已知直线平行的直线}
\TOline{bisector}{mediator}{由三点确定的角的平分线}
\TOline{bisector out}{mediator}{外角平分线}
\TOline{symmedian}{mediator}{从某顶点引出的陪位中线}
\TOline{altitude}{mediator}{从某顶点引出的高}
\TOline{euler}{mediator}{三角形的欧拉线}
\TOline{tangent at}{mediator}{圆上某点处的切线}
\TOline{tangent from}{mediator}{从外部一点引出的切线}
\TOline{K}{1}{垂线的系数}
\TOline{normed}{false}{将所生成的线段单位化}
\end{tabular}
\end{NewMacroBox}  

\subsubsection{使用 \tkzname{mediator}}  
\begin{tkzexample}[latex=5 cm,small]
\begin{tikzpicture}[rotate=25]
 \tkzDefPoints{-2/0/A,1/2/B}
 \tkzDefLine[mediator](A,B)          \tkzGetPoints{C}{D}
 \tkzDefPointWith[linear,K=.75](C,D) \tkzGetPoint{D}
 \tkzDefMidPoint(A,B)                \tkzGetPoint{I}
 \tkzFillPolygon[color=teal!20](A,C,B,D)
 \tkzDrawSegments(A,B C,D)
 \tkzMarkRightAngle(B,I,C) 
 \tkzDrawSegments(D,B D,A)
 \tkzDrawSegments(C,B C,A)
\end{tikzpicture}
\end{tkzexample}  

\subsubsection{用 \tkzname{mediator} 选项绘制包络}
本图基于 O. Reboux 的图形，原图由 D. Rodriguez 用 pst-eucl 制作。

\begin{tkzexample}[latex=7cm,small]
\begin{tikzpicture}[scale=.75]
   % necessary
\tkzInit[xmin=-6,ymin=-4,xmax=6,ymax=6]
\tkzClip
\tkzSetUpLine[thin,color=magenta]
\tkzDefPoint(0,0){O} 
\tkzDefPoint(132:4){A}
\tkzDefPoint(5,0){B}
\foreach \ang in {5,10,...,360}{%
 \tkzDefPoint(\ang:5){M}
 \tkzDefLine[mediator](A,M)
 \tkzGetPoints{x}{y}
 \tkzDrawLine[add= 3 and 3](x,y)}
\end{tikzpicture}
\end{tkzexample}


\subsubsection{用 \tkzname{mediator} 选项绘制抛物线}
本图基于 O. Reboux 的图形，原图由 D. Rodriguez 用 pst-eucl 制作。
无需为确定中垂线的两个点命名。

\begin{tkzexample}[latex=8cm,small]
\begin{tikzpicture}[scale=.6]
\tkzInit[xmin=-6,ymin=-4,xmax=6,ymax=6] 
\tkzClip
\tkzSetUpLine[thin,color=teal]
\tkzDefPoint(0,0){O} 
\tkzDefPoint(132:5){A}
\tkzDefPoint(4,0){B}
\foreach \ang in {5,10,...,360}{%
 \tkzDefPoint(\ang:4){M}
 \tkzDefLine[mediator](A,M) 
 \tkzGetPoints{x}{y}
 \tkzDrawLine[add= 3 and 3](x,y)}
\end{tikzpicture}
\end{tkzexample}

\subsubsection{使用 \tkzname{bisector} 与 \tkzname{normed} 选项} 
\begin{tkzexample}[latex=7 cm,small] 
\begin{tikzpicture}[rotate=25,scale=.75]
 \tkzDefPoints{0/0/C, 2/-3/A, 4/0/B}
 \tkzDefLine[bisector,normed](B,A,C) \tkzGetPoint{a}
 \tkzDrawLines[add= 0 and .5](A,B A,C)
 \tkzShowLine[bisector,gap=4,size=2,color=red](B,A,C)
 \tkzDrawLines[new,dashed,add= 0 and 3](A,a)
\end{tikzpicture}
\end{tkzexample} 

\subsubsection{使用 \tkzname{parallel=through} 选项} % (fold)
\label{ssub:parallel}
阿基米德《引理集》命题 1。

\begin{tkzexample}[latex=7cm,small]
  \begin{tikzpicture}
    \tkzDefPoints{0/0/O_1,0/1/O_2,0/3/A}
    \tkzDefPoint(15:3){F}
    \tkzDefPointBy[symmetry=center O_1](F) 
    \tkzGetPoint{E}
    \tkzDefLine[parallel=through O_2](E,F) 
    \tkzGetPoint{x}   
    \tkzInterLC(x,O_2)(O_2,A) \tkzGetPoints{D}{C} 
    \tkzDrawCircles(O_1,A O_2,A)
    \tkzDrawSegments[new](O_1,A E,F C,D)
    \tkzDrawSegments[purple](A,E A,F)
    \tkzDrawPoints(A,O_1,O_2,E,F,C,D)
    \tkzLabelPoints(A,O_1,O_2,E,F,C,D)
  \end{tikzpicture}
\end{tkzexample}
% subsubsection parallel (end)

\subsubsection{使用 \tkzname{orthogonal} 与 \tkzname{parallel} 选项}    
\begin{tkzexample}[latex=5 cm,small]
\begin{tikzpicture}
   \tkzDefPoints{-1.5/-0.25/A,1/-0.75/B,-0.7/1/C}
   \tkzDrawLine(A,B)
   \tkzLabelLine[pos=1.25,below left](A,B){$(d_1)$}
   \tkzDrawPoints(A,B,C)
   \tkzDefLine[orthogonal=through C](B,A) \tkzGetPoint{c}
   \tkzDrawLine(C,c) 
   \tkzLabelLine[pos=1.25,left](C,c){$(\delta)$}
   \tkzInterLL(A,B)(C,c) \tkzGetPoint{I}
   \tkzMarkRightAngle(C,I,B) 
   \tkzDefLine[parallel=through C](A,B) \tkzGetPoint{c'}
   \tkzDrawLine(C,c') 
   \tkzLabelLine[pos=1.25,below left](C,c'){$(d_2)$}
   \tkzMarkRightAngle(I,C,c')   
\end{tikzpicture}
\end{tkzexample}

\subsubsection{使用 \tkzname{altitude} 选项} % (fold)
\label{sub:altitude}
\begin{tkzexample}[latex=7 cm,small]
\begin{tikzpicture}
\tkzDefPoints{0/0/A,6/0/B,0.8/4/C}	
\tkzDefLine[altitude](A,B,C)     \tkzGetPoint{b}
\tkzDefLine[altitude](B,C,A)     \tkzGetPoint{c}
\tkzDefLine[altitude](B,A,C)     \tkzGetPoint{a}
\tkzDrawPolygon(A,B,C)
\tkzDrawPoints[blue](a,b,c)
\tkzDrawSegments[blue](A,a B,b C,c)
\tkzLabelPoints(A,B,c)
\tkzLabelPoints[above](C,a)
\tkzLabelPoints[above left](b)
\end{tikzpicture}
\end{tkzexample}

% subsection altitude (end)


\subsubsection{使用 \tkzname{euler} 选项} % (fold)
\label{sub:eulerline}
\begin{tkzexample}[latex=7 cm,small]
\begin{tikzpicture}[scale=.75]
\tkzDefPoints{0/0/A,6/0/B,0.8/4/C}			 
\tkzDefLine[euler](A,B,C)             
\tkzGetPoints{h}{e}
\tkzDefTriangleCenter[circum](A,B,C)  
\tkzGetPoint{o}
\tkzDrawPolygon[teal](A,B,C)
\tkzDrawPoints[red](A,B,C,h,e,o)
\tkzDrawLine[add= 2 and 2](h,e)
\tkzLabelPoints(A,B,C,h,e,o)
\tkzLabelPoints[above](C)
\end{tikzpicture}
\end{tkzexample}
% subsection eulerline (end)

\subsubsection{过圆上一点的切线：\tkzname{tangent at}} 
\begin{tkzexample}[latex=7cm,small]
\begin{tikzpicture}[scale=.75]
  \tkzDefPoint(0,0){O}
  \tkzDefPoint(6,6){E}
  \tkzDefRandPointOn[circle=center O radius 3]
  \tkzGetPoint{A}
  \tkzDrawSegment(O,A)
  \tkzDrawCircle(O,A)
  \tkzDefLine[tangent at=A](O)
  \tkzGetPoint{h}
  \tkzDrawLine[add = 4 and 3](A,h)
  \tkzMarkRightAngle[fill=teal!30](O,A,h)
\end{tikzpicture}
\end{tkzexample}

\subsubsection{过圆外一点的切线：切点的选择（\tkzname{tangent from} 选项）} 

切线并不直接绘制。使用 \tkzname{at} 选项时，切线上的一个点由 \tkzname{tkzPointResult} 给出。使用 \tkzname{from} 选项时，得到圆上的两个点，即 \tkzname{tkzFirstPointResult} 与 \tkzname{tkzSecondPointResult}。可通过比较“外点—切点—圆心”所成两个角的方向来选择：两个角方向相反，其中逆时针方向的那一个对应 \tkzname{tkzFirstPointResult}。

\begin{tkzexample}[latex=7cm,small]
\begin{tikzpicture}[scale=1,rotate=-30]
\tkzDefPoints{0/0/Q,0/2/A,6/-1/O}
\tkzDefLine[tangent from = O](Q,A)  
\tkzGetPoints{R}{S} 
\tkzInterLC[near](O,Q)(Q,A)         
\tkzGetPoints{M}{N}
\tkzDrawCircle(Q,M)
\tkzDrawSegments[new,add = 0 and .2](O,R O,S)
\tkzDrawSegments[gray](N,O R,Q S,Q)
\tkzDrawPoints(O,Q,R,S,M,N)
\tkzMarkAngle[gray,-stealth,size=1](O,R,Q)
\tkzFindAngle(O,R,Q)   \tkzGetAngle{an}
\tkzLabelAngle(O,R,Q){%
    $\pgfmathprintnumber{\an}^\circ$}
\tkzMarkAngle[gray,-stealth,size=1](O,S,Q)
\tkzFindAngle(O,S,Q)   \tkzGetAngle{an}
\tkzLabelAngle(O,S,Q){%
    $\pgfmathprintnumber{\an}^\circ$}
\tkzLabelPoints(Q,O,M,N,R)
\tkzLabelPoints[above,text=red](S)
\end{tikzpicture}
\end{tkzexample}

\subsubsection{过圆外一点的切线示例} 
\begin{tkzexample}[latex=7cm,small]
\begin{tikzpicture}[scale=.8] 
\tkzDefPoints{0/0/c,1/0/d,3/0/a0}
\def\tkzRadius{1}
\tkzDrawCircle(c,d) 
 \foreach \an in {0,10,...,350}{
  \tkzDefPointBy[rotation=center c angle \an](a0)  
  \tkzGetPoint{a}
  \tkzDefLine[tangent from = a](c,d) 
  \tkzGetPoints{e}{f}
  \tkzDrawLines(a,f a,e)
  \tkzDrawSegments(c,e c,f)}
\end{tikzpicture} 
\end{tkzexample}

\subsubsection{Andrew Mertz 示例}

\begin{tkzexample}[latex=6cm,small]
 \begin{tikzpicture}[scale=.6] 
 \tkzDefPoint(100:8){A}\tkzDefPoint(50:8){B}  
 \tkzDefPoint(0,0){C} \tkzDefPoint(0,-4){R} 
 \tkzDrawCircle(C,R)
 \tkzDefLine[tangent from = A](C,R)  \tkzGetPoints{D}{E}
\tkzDefLine[tangent from = B](C,R)  \tkzGetPoints{F}{G}
 \tkzDrawSector[fill=teal!20,opacity=0.5](A,E)(D)
 \tkzFillSector[color=teal,opacity=0.5](B,G)(F)
 \end{tikzpicture}
\end{tkzexample}
\url{http://www.texample.net/tikz/examples/}  

\subsubsection{绘制切线：\tkzname{tangent from} 选项}
\begin{tkzexample}[latex=6cm,small]
\begin{tikzpicture}[scale=.6] 
 \tkzDefPoint(0,0){B} 
 \tkzDefPoint(0,8){A} 
 \tkzDefSquare(A,B)
 \tkzGetPoints{C}{D}
 \tkzDrawPolygon(A,B,C,D)
 \tkzClipPolygon(A,B,C,D)
 \tkzDefPoint(4,8){F}
 \tkzDefPoint(4,0){E}
 \tkzDefPoint(4,4){Q}
 \tkzFillPolygon[color = green](A,B,C,D)
 \tkzDrawCircle[fill = orange](B,A)
 \tkzDrawCircle[fill = purple](E,B)  
 \tkzDefLine[tangent from = B](F,A)
 \tkzInterLL(F,tkzSecondPointResult)(C,D)
 \tkzInterLL(A,tkzPointResult)(F,E) 
 \tkzDrawCircle[fill = yellow](tkzPointResult,Q)  
 \tkzDefPointBy[projection= onto B--A](tkzPointResult)
 \tkzDrawCircle[fill = blue!50!black](tkzPointResult,A)
\end{tikzpicture}
\end{tkzexample}

\endinput

\section{三角形}

\subsection{三角形的定义 \tkzcname{tkzDefTriangle}}
下列宏可以由 \tkzname{至少} 两个点出发来定义或构造三角形。

 可定义的三角形包括：
 \begin{itemize}
\item  \tkzname{two angles}：给定两个角确定一个三角形；
\item  \tkzname{equilateral}：确定一个等边三角形；
\item  \tkzname{isosceles right}：确定一个等腰直角三角形；
\item \tkzname{half}：确定一个直角三角形，其两条直角边（直角的两条边）之比为 $2$；
\item \tkzname{pythagore}：确定一个直角三角形，其三边之比为 $3:4:5$；
\item \tkzname{school}：确定一个直角三角形，其三个角分别为 $30^\circ$、$60^\circ$、$90^\circ$；
\item \tkzname{golden}（即 \tkzname{sublime}）：确定黄金三角形，即以 $[AB]$ 为底、底角为 $72^\circ$ 的等腰三角形；

\item \tkzname{euclid}：确定欧几里得三角形，其角与黄金三角形相同，但边 $[AB]$ 是它的一条腰（$A$ 处角为 $36^\circ$）；在旧版本中该选项写作 "euclide"，末尾带字母 e；

\item \tkzname{gold}：确定一个直角三角形，直角在 $B$ 处，两条直角边之比为 $\Phi=1.618034$，即 $AB/BC=\Phi$。该三角形源自黄金矩形。

\item \tkzname{cheops}：确定第三个点，使三角形为等腰三角形，且三边之比为 $2:\Phi:\Phi$。
\end{itemize}

\newpage
\begin{NewMacroBox}{tkzDefTriangle}{\oarg{local options}\parg{A,B}}%
点的顺序很重要，因为三角形按逆时针方向构造。本宏既可配合 \tkzcname{tkzGetPoint} 使用，若无需保留该点，也可直接使用 \tkzname{tkzPointResult}。

\medskip
\begin{tabular}{lll}%
\toprule
选项             & 默认值 & 含义                        \\
\midrule
\TOline{two angles= \#1 and \#2}{无默认}{已知两个角确定三角形}
\TOline{equilateral} {equilateral}{等边三角形}
\TOline{half} {equilateral}{$B$ 为直角，$AB=2BC$，$AC$ 为斜边}
\TOline{isosceles right} {equilateral}{等腰直角三角形}
\TOline{pythagore}{equilateral}{三边之比为 3-4-5 的直角三角形}
\TOline{pythagoras}{equilateral}{同上}
\TOline{egyptian}{equilateral}{同上}
\TOline{school} {equilateral}{三个角分别为 $30^\circ$、$60^\circ$、$90^\circ$}
\TOline{gold}{equilateral}{$B$ 为直角，且 $AB/AC = \Phi$}
\TOline{euclid} {equilateral}{三个角为 $72^\circ$、$72^\circ$、$36^\circ$，$A$ 为顶角}
\TOline{golden} {equilateral}{三个角为 $72^\circ$、$72^\circ$、$36^\circ$，$C$ 为顶角}
\TOline{sublime} {equilateral}{三个角为 $72^\circ$、$72^\circ$、$36^\circ$，$C$ 为顶角}
\TOline{cheops} {equilateral}{$AC=BC$，且 $AC$、$BC$ 与 $2$ 和 $\Phi$ 成比例}
\TOline{swap} {false}{给出关于 $AB$ 对称的点}
\bottomrule
\end{tabular}

\medskip
\emph{用 \tkzcname{tkzGetPoint} 保存该点；否则可直接使用 \tkzname{tkzPointResult}。}
\end{NewMacroBox}

\subsubsection{选项 \tkzname{equilateral}}
\begin{tkzexample}[latex=7 cm,small]
\begin{tikzpicture}
  \tkzDefPoint(0,0){A}
  \tkzDefPoint(4,0){B}
  \tkzDefTriangle[equilateral](A,B)
  \tkzGetPoint{C}
  \tkzDrawPolygons(A,B,C)
  \tkzDefTriangle[equilateral](B,A)
  \tkzGetPoint{D}
  \tkzDrawPolygon(B,A,D)
  \tkzMarkSegments[mark=s|](A,B B,C A,C A,D B,D)
\end{tikzpicture}
\end{tkzexample}


\subsubsection{选项 \tkzname{two angles}}
\begin{tkzexample}[latex=6 cm,small]
\begin{tikzpicture}
\tkzDefPoint(0,0){A} 
\tkzDefPoint(5,0){B} 
\tkzDefTriangle[two angles = 50 and 70](A,B)
\tkzGetPoint{C} 
\tkzDrawSegment(A,B) 
\tkzDrawPoints(A,B) 
\tkzLabelPoints(A,B) 
\tkzDrawSegments[new](A,C B,C) 
\tkzDrawPoints[new](C)
\tkzLabelPoints[above,new](C)
\tkzLabelAngle[pos=1.4](B,A,C){$50^\circ$}
\tkzLabelAngle[pos=0.8](C,B,A){$70^\circ$}
\end{tikzpicture}
\end{tkzexample}

\subsubsection{选项 \tkzname{school}}
三个角分别为 $30^\circ$、$60^\circ$、$90^\circ$。

\begin{tkzexample}[latex=6 cm,small]
\begin{tikzpicture}
  \tkzDefPoints{0/0/A,4/0/B}
  \tkzDefTriangle[school](A,B)  
  \tkzGetPoint{C}
  \tkzMarkRightAngles(C,B,A)
  \tkzLabelAngle[pos=0.8](B,A,C){$30^\circ$}
  \tkzLabelAngle[pos=0.8](C,B,A){$90^\circ$}
  \tkzLabelAngle[pos=0.8](A,C,B){$60^\circ$} 
  \tkzDrawSegments(A,B)
  \tkzDrawSegments[new](A,C B,C)
  \tkzLabelPoints(A,B)
  \tkzLabelPoints[above](C)
\end{tikzpicture}
\end{tkzexample}

\subsubsection{选项 \tkzname{pythagore}}
该三角形的三边长度之比为 $3:4:5$。

\begin{tkzexample}[latex=6 cm,small]
\begin{tikzpicture} 
  \tkzDefPoints{0/0/A,4/0/B} 
  \tkzDefTriangle[pythagore](A,B) 
  \tkzGetPoint{C} 
  \tkzDrawSegments(A,B)
  \tkzDrawSegments[new](A,C B,C)
  \tkzMarkRightAngles(A,B,C)
  \tkzDrawPoints[new](C) 
  \tkzDrawPoints(A,B) 
  \tkzLabelPoints[above](A,B)
  \tkzLabelPoints[new](C)  
\end{tikzpicture}
\end{tkzexample}

\subsubsection{选项 \tkzname{pythagore} 与 \tkzname{swap}}
该三角形的三边长度之比为 $3:4:5$。

\begin{tkzexample}[latex=6 cm,small]
\begin{tikzpicture} 
  \tkzDefPoints{0/0/A,4/0/B} 
  \tkzDefTriangle[pythagore,swap](A,B) 
  \tkzGetPoint{C} 
  \tkzDrawSegments(A,B)
  \tkzDrawSegments[new](A,C B,C)
  \tkzMarkRightAngles(A,B,C)
  \tkzLabelPoint[above,new](C){$C$} 
  \tkzDrawPoints[new](C) 
  \tkzDrawPoints(A,B) 
  \tkzLabelPoints(A,B) 
\end{tikzpicture}
\end{tkzexample}

\subsubsection{选项 \tkzname{golden}}
\begin{tkzexample}[latex=6 cm,small]
\begin{tikzpicture}[scale=.8]
\tkzDefPoint(0,0){A} \tkzDefPoint(4,0){B} 
\tkzDefTriangle[golden](A,B)\tkzGetPoint{C} 
\tkzDefSpcTriangle[in,name=M](A,B,C){a,b,c}
\tkzDrawPolygon(A,B,C) 
\tkzDrawPoints(A,B) 
\tkzDrawSegment(C,Mc) 
\tkzDrawPoints[new](C)
\tkzLabelPoints(A,B) 
\tkzLabelPoints[above,new](C)
\end{tikzpicture}
\end{tkzexample}

\subsubsection{选项 \tkzname{euclid}}
\tkzimp{euclid} 三角形与 \tkzimp{golden} 三角形形状相同，但在前一个中线段 $[AB]$ 是腰，在后一个中 $[AB]$ 是底。

\begin{tkzexample}[latex=7 cm,small]
\begin{tikzpicture}[scale=.75]
 \tkzDefPoint(0,0){A} \tkzDefPoint(4,0){B}
 \tkzDefTriangle[euclid](A,B)\tkzGetPoint{C}
 \tkzDrawPolygon(A,B,C)
 \tkzDrawPoints(A,B,C)
 \tkzLabelPoints(C)
 \tkzLabelPoints[above](A,B)
 \tkzLabelAngle[pos=0.8](A,B,C){$72^\circ$}
 \tkzLabelAngle[pos=0.8](B,C,A){$72^\circ$}
 \tkzLabelAngle[pos=0.8](C,A,B){$36^\circ$}
\end{tikzpicture}
\end{tkzexample}

\subsubsection{选项 \tkzname{isosceles right}}
\begin{tkzexample}[latex=7 cm,small]
\begin{tikzpicture}
  \tkzDefPoint(0,0){A}
  \tkzDefPoint(4,0){B}
  \tkzDefTriangle[isosceles right](A,B)
  \tkzGetPoint{C}
  \tkzDrawPolygons(A,B,C)
  \tkzDrawPoints(A,B,C)
  \tkzMarkRightAngles(A,C,B)
  \tkzLabelPoints(A,B)
  \tkzLabelPoints[above](C)
\end{tikzpicture}
\end{tkzexample}

\subsubsection{选项 \tkzname{gold}}
\begin{tkzexample}[latex=6 cm,small]
\begin{tikzpicture}
 \tkzDefPoints{0/0/A,4/0/B} 
 \tkzDefTriangle[gold](A,B)
 \tkzGetPoint{C}
 \tkzDrawPolygon(A,B,C)
 \tkzDrawPoints(A,B,C)
 \tkzLabelPoints[above](A,B) 
 \tkzLabelPoints[below](C)
 \tkzMarkRightAngle(A,B,C)
 \tkzText(0,-2){$\dfrac{AB}{BC}=\varphi$}
\end{tikzpicture}
\end{tkzexample}

\clearpage
\subsection{用 \tkzcname{tkzDefSpcTriangle} 构造特殊三角形}

某些三角形的重心已在“特殊点”一节中定义；这里我们要确定若干特殊三角形的三个顶点。

\begin{NewMacroBox}{tkzDefSpcTriangle}{\oarg{local options}\parg{p1,p2,p3}\marg{r1,r2,r3}}
点的顺序非常重要！p1p2p3 定义一个三角形；结果是另一个三角形，其顶点的引用名为 \tkzname{name} 后接 r1、r2、r3。若 \tkzname{name} 为空，则引用名即为 r1、r2、r3。

\medskip
\begin{tabular}{lll}%
\toprule
选项             & 默认值 & 含义                        \\
\midrule
\TOline{orthic} {centroid}{由三条高的垂足确定的三角形……}
\TOline{centroid or medial}{centroid}{三条中线的交点所确定的中点三角形}
\TOline{in or incentral}{centroid}{由三条角平分线确定}
\TOline{ex or excentral} {centroid}{由三个旁心确定}
\TOline{extouch}{centroid}{由三角形与三个旁切圆的切点构成}
\TOline{intouch or contact} {centroid}{由内切圆的三个切点构成}
\TOline{} {}{各顶点}
\TOline{euler} {centroid}{由九点圆上的欧拉线中点构成}
\TOline{symmedial} {centroid}{三条陪位中线的交点}
\TOline{tangential}{centroid}{由外接圆在三个顶点处的切线构成}
\TOline{feuerbach} {centroid}{由九点圆与旁切圆的三个切点构成……}
\TOline{} {} {圆与旁切圆}
\TOline{name} {empty}{用于给顶点命名}
\midrule
\end{tabular}
\end{NewMacroBox}

\subsubsection{如何为顶点命名}

使用 \tkzcname{tkzDefSpcTriangle[medial,name=M](A,B,C)\{\_A,\_B,\_C\}} 可得到三个名为 $M_A$、$M_B$、$M_C$ 的顶点。

使用 \tkzcname{tkzDefSpcTriangle[medial](A,B,C)\{a,b,c\}} 可得到三个名为 $a$、$b$、$c$ 的顶点。

使用 \tkzcname{tkzDefSpcTriangle[medial,name=M\_](A,B,C)\{A,B,C\}} 同样可得到三个名为 $M_A$、$M_B$、$M_C$ 的顶点。

\subsubsection{选项 \tkzname{medial} 或 \tkzname{centroid}}
三角形的重心是三条中线的交点 $G$（有时也记为 $M$），因此有时也称为中位点。重心始终位于三角形内部。
\\

\href{http://mathworld.wolfram.com/TriangleCentroid.html}{Weisstein, Eric W. "Centroid triangle" From MathWorld--A Wolfram Web Resource.}

在下面的示例中，我们得到欧拉线圆（九点圆），它经过上面定义的各点。

\begin{tkzexample}[latex=7cm,small]
  \begin{tikzpicture}[rotate=90,scale=.75]
   \tkzDefPoints{0/0/A,6/0/B,0.8/4/C}
   \tkzDefTriangleCenter[centroid](A,B,C)
   \tkzGetPoint{M}
   \tkzDefSpcTriangle[medial,name=M](A,B,C){_A,_B,_C}
   \tkzDrawPolygon(A,B,C)
   \tkzDrawSegments[dashed,new](A,M_A B,M_B C,M_C)
   \tkzDrawPolygon[new](M_A,M_B,M_C)
   \tkzDrawPoints(A,B,C)
   \tkzDrawPoints[new](M,M_A,M_B,M_C)
   \tkzLabelPoints[above](B)
   \tkzLabelPoints[below](A,C,M_B)
   \tkzLabelPoints[right](M_C)
   \tkzLabelPoints[left](M_A)
   \tkzLabelPoints[font=\scriptsize](M)
  \end{tikzpicture}
\end{tkzexample}

\subsubsection{选项 \tkzname{in} 或 \tkzname{incentral}}

内心三角形（incentral triangle）是以参考三角形的三条角平分线与对应对边的交点为顶点的三角形。
\\
\href{http://mathworld.wolfram.com/ContactTriangle.html}{Weisstein, Eric W. "Incentral triangle" From MathWorld--A Wolfram Web Resource.}


\begin{tkzexample}[latex=7cm,small]
\begin{tikzpicture}[scale=1]
  \tkzDefPoints{ 0/0/A,5/0/B,2/3/C}
  \tkzDefSpcTriangle[in,name=I](A,B,C){_a,_b,_c}
  \tkzDefCircle[in](A,B,C) \tkzGetPoints{I}{a}
  \tkzDrawCircle(I,a)
  \tkzDrawPolygon(A,B,C)
  \tkzDrawPolygon[new](I_a,I_b,I_c)
  \tkzDrawSegments[dashed,new](A,I_a B,I_b C,I_c)
  \tkzDrawPoints(A,B,C,I,I_a,I_b,I_c) 
  \tkzLabelPoints[below](A,B,I_c)
  \tkzLabelPoints[above left](I_b)
  \tkzLabelPoints[above right](C,I_a)
\end{tikzpicture}
\end{tkzexample}

\subsubsection{选项 \tkzname{ex} 或 \tkzname{excentral}}

三角形 $ABC$ 的旁心三角形（excentral triangle）是以 $ABC$ 的三个旁心为顶点的三角形 $J_aJ_bJ_c$。

\begin{tkzexample}[latex=7cm,small]
\begin{tikzpicture}[scale=.6]
 \tkzDefPoints{0/0/A,6/0/B,0.8/4/C}
 \tkzDefSpcTriangle[excentral,name=J](A,B,C){_a,_b,_c}
 \tkzDefSpcTriangle[extouch,name=T](A,B,C){_a,_b,_c}
 \tkzDrawPolygon(A,B,C)
 \tkzDrawPolygon[new](J_a,J_b,J_c)
 \tkzClipBB
 \tkzDrawPoints(A,B,C)
 \tkzDrawPoints[new](J_a,J_b,J_c)
 \tkzLabelPoints(A,B,C)
 \tkzLabelPoints[new](J_b,J_c)
 \tkzLabelPoints[new,above](J_a)
 \tkzDrawCircles[gray](J_a,T_a J_b,T_b J_c,T_c) 
\end{tikzpicture}
\end{tkzexample}


\subsubsection{选项 \tkzname{intouch} 或 \tkzname{contact}}
三角形 $ABC$ 的切点三角形（contact triangle），又称内心切点三角形（intouch triangle），是由 $ABC$ 的内切圆与三边的切点构成的三角形。\\
\href{http://mathworld.wolfram.com/ContactTriangle.html}{Weisstein, Eric W. "Contact triangle" From MathWorld--A Wolfram Web Resource.}

下面我们得到内切圆与各边相切的切点。
\begin{tkzexample}[latex=7cm,small]
\begin{tikzpicture}[scale=.75]
 \tkzDefPoints{0/0/A,6/0/B,0.8/4/C}          
 \tkzDefSpcTriangle[intouch,name=X](A,B,C){_a,_b,_c}
 \tkzInCenter(A,B,C)\tkzGetPoint{I}
 \tkzDefCircle[in](A,B,C) \tkzGetPoints{I}{i}
 \tkzDrawCircle(I,i)
 \tkzDrawPolygon(A,B,C)
 \tkzDrawPolygon[new](X_a,X_b,X_c)
 \tkzDrawPoints(A,B,C)
 \tkzDrawPoints[new](X_a,X_b,X_c)
 \tkzLabelPoints[right](X_a)
 \tkzLabelPoints[left](X_b)
 \tkzLabelPoints[above](C)
 \tkzLabelPoints[below](A,B,X_c)
\end{tikzpicture} 
\end{tkzexample}

\subsubsection{选项 \tkzname{extouch}}
旁切切点三角形（extouch triangle）$T_aT_bT_c$ 是由三角形 $ABC$ 的三边与三个旁切圆（圆心分别为 $J_a$、$J_b$、$J_c$）的切点构成的三角形。点 $T_a$、$T_b$、$T_c$ 也可分别看作从 $A$、$B$、$C$ 出发将 $ABC$ 的周长二等分的点。\\
\href{http://mathworld.wolfram.com/ExtouchTriangle.html}{Weisstein, Eric W. "Extouch triangle" From MathWorld--A Wolfram Web Resource.}

下面我们得到旁切圆的切点，以及由三个旁心构成的三角形。

\begin{tkzexample}[latex=8cm,small]
\begin{tikzpicture}[scale=.7]
\tkzDefPoints{0/0/A,6/0/B,0.8/4/C}
\tkzDefSpcTriangle[excentral,
                 name=J](A,B,C){_a,_b,_c}
\tkzDefSpcTriangle[extouch,
                  name=T](A,B,C){_a,_b,_c}
\tkzDefTriangleCenter[nagel](A,B,C)
\tkzGetPoint{N_a}
\tkzDefTriangleCenter[centroid](A,B,C)
\tkzGetPoint{G}
\tkzDrawPoints[new](J_a,J_b,J_c)
\tkzClipBB \tkzShowBB
\tkzDrawCircles[gray](J_a,T_a J_b,T_b J_c,T_c)
\tkzDrawLines[add=1 and 1](A,B B,C C,A)
\tkzDrawSegments[new](A,T_a B,T_b C,T_c)
\tkzDrawSegments[new](J_a,T_a J_b,T_b J_c,T_c)
\tkzDrawPolygon(A,B,C)
\tkzDrawPolygon[new](T_a,T_b,T_c)
\tkzDrawPoints(A,B,C,N_a)
\tkzDrawPoints[new](T_a,T_b,T_c)
\tkzLabelPoints[below left](A)
\tkzLabelPoints[below](N_a,B)
\tkzLabelPoints[above](C)
\tkzLabelPoints[new,below left](T_b)
\tkzLabelPoints[new,below right](T_c)
\tkzLabelPoints[new,right=6pt](T_a)
\tkzMarkRightAngles[fill=gray!15](J_a,T_a,B
 J_b,T_b,C J_c,T_c,A)
\end{tikzpicture}
\end{tkzexample}

\subsubsection{选项 \tkzname{orthic}}

给定三角形 $ABC$，以三条高在各边上的垂足为顶点的三角形 $H_AH_BH_C$ 称为垂足三角形（orthic triangle），有时也叫高三角形。三条直线 $AH_A$、$BH_B$、$CH_C$ 共点于 $ABC$ 的垂心 $H$。

\begin{tkzexample}[latex=7cm,small]
\begin{tikzpicture}[scale=.75]
\tkzDefPoints{1/5/A,0/0/B,7/0/C}
 \tkzDefSpcTriangle[orthic](A,B,C){H_A,H_B,H_C}
 \tkzDefTriangleCenter[ortho](B,C,A)
 \tkzGetPoint{H}
 \tkzDefPointWith[orthogonal,normed](H_A,B)
 \tkzGetPoint{a}
 \tkzDrawSegments[new](A,H_A B,H_B C,H_C)   
 \tkzMarkRightAngles[fill=gray!20,
         opacity=.5](A,H_A,C B,H_B,A C,H_C,A)
 \tkzDrawPolygon[fill=teal!20,opacity=.3](A,B,C)
 \tkzDrawPoints(A,B,C)
 \tkzDrawPoints[new](H_A,H_B,H_C)
 \tkzDrawPolygon[new,fill=orange!20,
                opacity=.3](H_A,H_B,H_C)
 \tkzLabelPoints(C)
 \tkzLabelPoints[left](B)
 \tkzLabelPoints[above](A)
 \tkzLabelPoints[new](H_A)
 \tkzLabelPoints[new,above left](H_C)
 \tkzLabelPoints[new,above right](H_B,H)
\end{tikzpicture}
\end{tkzexample}
    
\subsubsection{选项 \tkzname{feuerbach}}
费尔巴哈三角形（Feuerbach triangle）是由九点圆与三个旁切圆的三个切点构成的三角形。\\
\href{http://mathworld.wolfram.com/FeuerbachTriangle.html}{Weisstein, Eric W. "Feuerbach triangle" From MathWorld--A Wolfram Web Resource.}

 这些切点确定了费尔巴哈三角形。

\begin{tkzexample}[latex=8cm,small]
\begin{tikzpicture}[scale=1]
  \tkzDefPoint(0,0){A}
  \tkzDefPoint(3,0){B}
  \tkzDefPoint(0.5,2.5){C}
  \tkzDefCircle[euler](A,B,C) \tkzGetPoint{N}
  \tkzDefSpcTriangle[feuerbach,
                       name=F](A,B,C){_a,_b,_c}
  \tkzDefSpcTriangle[excentral,
                       name=J](A,B,C){_a,_b,_c}
  \tkzDefSpcTriangle[extouch,
                        name=T](A,B,C){_a,_b,_c}
  \tkzLabelPoints[below left](J_a,J_b,J_c)  
  \tkzClipBB \tkzShowBB
  \tkzDrawCircle[purple](N,F_a)
  \tkzDrawPolygon(A,B,C)
  \tkzDrawPolygon[new](F_a,F_b,F_c)
  \tkzDrawCircles[gray](J_a,F_a J_b,F_b J_c,F_c)
  \tkzDrawPoints[blue](J_a,J_b,J_c,%
          F_a,F_b,F_c,A,B,C)
  \tkzLabelPoints(A,B,F_c)
  \tkzLabelPoints[above](C)    
  \tkzLabelPoints[right](F_a)
  \tkzLabelPoints[left](F_b)        
\end{tikzpicture}
\end{tkzexample}

\subsubsection{选项 \tkzname{tangential}} 
切线三角形（tangential triangle）是由给定三角形 $ABC$ 的外接圆在三个顶点处的切线所围成的三角形 $T_aT_bT_c$。因此它是 $ABC$ 关于外心 $O$ 的反垂足三角形。\\ 
\href{http://mathworld.wolfram.com/TangentialTriangle.html}{Weisstein, Eric W. "Tangential Triangle." From MathWorld--A Wolfram Web Resource. }


\begin{tkzexample}[latex=8cm,small]
\begin{tikzpicture}[scale=.5,rotate=80]
  \tkzDefPoints{0/0/A,6/0/B,1.8/4/C}           
  \tkzDefSpcTriangle[tangential,
    name=T](A,B,C){_a,_b,_c}
  \tkzDrawPolygon(A,B,C)
  \tkzDrawPolygon[new](T_a,T_b,T_c)
  \tkzDrawPoints(A,B,C)
  \tkzDrawPoints[new](T_a,T_b,T_c)
  \tkzDefCircle[circum](A,B,C)  
  \tkzGetPoint{O} 
  \tkzDrawCircle(O,A)
  \tkzLabelPoints(A)
  \tkzLabelPoints[above](B)
  \tkzLabelPoints[left](C)
  \tkzLabelPoints[new](T_b,T_c)
  \tkzLabelPoints[new,left](T_a)
\end{tikzpicture} 
\end{tkzexample} 

\subsubsection{选项 \tkzname{euler}} 
三角形 $ABC$ 的欧拉三角形（Euler triangle）是三角形 $E_AE_BE_C$，其顶点分别是垂心 $H$ 与三个顶点连线的中点。这些顶点称为欧拉点，位于九点圆上。
\\
\href{https://mathworld.wolfram.com/EulerTriangle.html}{Weisstein, Eric W. "Euler Triangle." From MathWorld--A Wolfram Web Resource.} 

\begin{tkzexample}[latex=7cm,small]
\begin{tikzpicture}[rotate=90,scale=1.25]
 \tkzDefPoints{0/0/A,6/0/B,0.8/4/C}
 \tkzDefSpcTriangle[medial,
     name=M](A,B,C){_A,_B,_C}
 \tkzDefTriangleCenter[euler](A,B,C)
     \tkzGetPoint{N} % I= N nine points
 \tkzDefTriangleCenter[ortho](A,B,C)
        \tkzGetPoint{H}
 \tkzDefMidPoint(A,H) \tkzGetPoint{E_A}
 \tkzDefMidPoint(C,H) \tkzGetPoint{E_C}
 \tkzDefMidPoint(B,H) \tkzGetPoint{E_B}
 \tkzDefSpcTriangle[ortho,name=H](A,B,C){_A,_B,_C}
 \tkzDrawPolygon(A,B,C)
 \tkzDrawCircle(N,E_A)
 \tkzDrawSegments[new](A,H_A B,H_B C,H_C)
 \tkzDrawPoints(A,B,C,N,H)
 \tkzDrawPoints[red](M_A,M_B,M_C)
 \tkzDrawPoints[blue]( H_A,H_B,H_C)
 \tkzDrawPoints[green](E_A,E_B,E_C)
 \tkzAutoLabelPoints[center=N,font=\scriptsize]%
(A,B,C,M_A,M_B,M_C,H_A,H_B,H_C,E_A,E_B,E_C)
\tkzLabelPoints[font=\scriptsize](H,N)
\tkzMarkSegments[mark=s|,size=3pt,
  color=blue,line width=1pt](B,E_B E_B,H)
   \tkzDrawPolygon[color=cyan](M_A,M_B,M_C)
\end{tikzpicture}
\end{tkzexample}

\subsubsection{选项 \tkzname{euler} 与选项 \tkzname{orthic}} 
\begin{tkzexample}[vbox,small]
  \begin{tikzpicture}[scale=1.25]
    \tkzDefPoints{0/0/A,6/0/B,0.8/4/C}
    \tkzDefSpcTriangle[euler,name=E](A,B,C){a,b,c}
    \tkzDefSpcTriangle[orthic,name=H](A,B,C){a,b,c}
    \tkzDefExCircle(A,B,C) \tkzGetPoints{I}{i}
    \tkzDefExCircle(C,A,B) \tkzGetPoints{J}{j}
    \tkzDefExCircle(B,C,A) \tkzGetPoints{K}{k}
    \tkzDrawPoints[orange](I,J,K)
    \tkzLabelPoints[font=\scriptsize](A,B,C,I,J,K)
    \tkzClipBB
    \tkzInterLC(I,C)(I,i) \tkzGetSecondPoint{Fc}
    \tkzInterLC(J,B)(J,j) \tkzGetSecondPoint{Fb}
    \tkzInterLC(K,A)(K,k) \tkzGetSecondPoint{Fa}
    \tkzDrawLines[add=1.5 and 1.5](A,B A,C B,C)
    \tkzDefCircle[euler](A,B,C) \tkzGetPoints{E}{e}
    \tkzDrawCircle[orange](E,e)
    \tkzDrawSegments[orange](E,I E,J E,K)
    \tkzDrawSegments[dashed](A,Ha B,Hb C,Hc)
    \tkzDrawCircles(J,j I,i K,k)
    \tkzDrawPoints(A,B,C)
    \tkzDrawPoints[orange](E,I,J,K,Ha,Hb,Hc,Ea,Eb,Ec,Fa,Fb,Fc)
    \tkzLabelPoints[font=\scriptsize](E,Ea,Eb,Ec,Ha,Hb,Hc,Fa,Fb,Fc)  
  \end{tikzpicture}
\end{tkzexample}

\subsubsection{选项 \tkzname{symmedial}}
陪位中线三角形（symmedial triangle）$K_AK_BK_C$ 是以三条陪位中线与参考三角形 $ABC$ 的交点为顶点的三角形。

\begin{tkzexample}[latex=7cm,small]
\begin{tikzpicture}
\tkzDefPoint(0,0){A}
\tkzDefPoint(5,0){B}
\tkzDefPoint(.75,4){C}
\tkzDefTriangleCenter[symmedian](A,B,C)\tkzGetPoint{K} 
\tkzDefSpcTriangle[symmedial,name=K_](A,B,C){A,B,C}
\tkzDrawPolygon(A,B,C)
\tkzDrawSegments[new](A,K_A B,K_B C,K_C)
\tkzDrawPoints(A,B,C,K,K_A,K_B,K_C)
\tkzLabelPoints(A,B,K,K_C)
\tkzLabelPoints[above](C)
\tkzLabelPoints[right](K_A)
\tkzLabelPoints[left](K_B)
\end{tikzpicture}
\end{tkzexample}

\subsection{交换三角形中的两个点}

\begin{NewMacroBox}{tkzPermute}{\parg{$pt1$,$pt2$,$pt3$}}%
\begin{tabular}{lll}%
参数             & 示例 & 说明                         \\
\midrule
\TAline{(pt1,pt2,pt3)} {\tkzcname{tkzPermute}(A,B,C)}{$A$ 与角 $\widehat{B,A,C}$ 保持不变，$B$ 与 $C$ 互换}
\midrule
\end{tabular}

\medskip
\emph{三角形本身不变。}
\end{NewMacroBox}

\subsubsection{对 \tkzname{school} 三角形的修改}
该三角形是在射线 $[Ax)$ 上由线段 $[AB]$ 构造出来的。

若希望线段 $[AC]$ 位于 $[Ax)$ 上，只需交换 $B$ 与 $C$。

\begin{tkzexample}[latex=7cm,small]
\begin{tikzpicture}
  \tkzDefPoints{0/0/A,4/0/B,6/0/x}
  \tkzDefTriangle[school](A,B)  
  \tkzGetPoint{C}
  \tkzPermute(A,B,C)
  \tkzDrawSegments(A,B C,x)
  \tkzDrawSegments(A,C B,C)
  \tkzDrawPoints(A,B,C)
  \tkzLabelPoints(A,C,x)
  \tkzLabelPoints[above](B)
  \tkzMarkRightAngles(C,B,A)
\end{tikzpicture}
\end{tkzexample}

说明：只有第一个点保持不变。后两个参数的顺序无关紧要。

\endinput

\section{多边形的定义}
\subsection{正方形顶点的定义} \label{def_square}
前面我们已经看到如何定义一些三角形。下面来看一些四边形和正多边形。

\begin{NewMacroBox}{tkzDefSquare}{\parg{pt1,pt2}}%
给定两个点，可再得到两个点，使这四个点按顺序构成一个正方形；正方形按逆时针方向定义。\\结果分别保存在 \tkzname{tkzFirstPointResult} 和 \tkzname{tkzSecondPointResult} 中。\\
可用 \tkzcname{tkzGetPoints} 为它们重新命名。

\medskip
\begin{tabular}{lll}%
\toprule
参数             & 示例 & 说明                         \\ 
\midrule
\TAline{\parg{pt1,pt2}}{\tkzcname{tkzDefSquare}\parg{A,B}}{正方形按逆时针方向定义。}
\end{tabular}
\end{NewMacroBox}

\subsubsection{用两个点调用 \tkzcname{tkzDefSquare}}
注意交换前两个点会得到不同结果。

\begin{tkzexample}[latex=4cm,small]
\begin{tikzpicture}[scale=.5]
  \tkzDefPoint(0,0){A} \tkzDefPoint(3,0){B}
  \tkzDefSquare(A,B)
  \tkzDrawPolygon[new](A,B,tkzFirstPointResult,%
               tkzSecondPointResult)
  \tkzDefSquare(B,A)
  \tkzDrawPolygon(B,A,tkzFirstPointResult,%
               tkzSecondPointResult) 
\end{tikzpicture} 
\end{tkzexample}

有时我们可能只需要其中一个点，用来绘制等腰直角三角形，这时可使用 \\ \tkzcname{tkzGetFirstPoint} 或 \tkzcname{tkzGetSecondPoint}。

\subsubsection{用 \tkzcname{tkzDefSquare} 得到等腰直角三角形}
\begin{tkzexample}[latex=7cm,small]
\begin{tikzpicture}[scale=1]
  \tkzDefPoint(0,0){A}
  \tkzDefPoint(3,0){B}
  \tkzDefSquare(A,B) \tkzGetFirstPoint{C}
  \tkzDrawSegment(A,B)
  \tkzDrawSegments[new](A,C B,C)
  \tkzMarkRightAngles(A,B,C)
  \tkzDrawPoints(A,B) \tkzDrawPoint[new](C)
  \tkzLabelPoints(A,B)
  \tkzLabelPoints[new,above](C)
\end{tikzpicture}
\end{tkzexample}

\subsubsection{勾股定理与 \tkzcname{tkzDefSquare}}
\begin{tkzexample}[latex=8cm,small]
\begin{tikzpicture}[scale=.5]
\tkzDefPoint(0,0){C}
\tkzDefPoint(4,0){A}
\tkzDefPoint(0,3){B} 
\tkzDefSquare(B,A)\tkzGetPoints{E}{F} 
\tkzDefSquare(A,C)\tkzGetPoints{G}{H} 
\tkzDefSquare(C,B)\tkzGetPoints{I}{J} 
\tkzDrawPolygon(A,B,C) 
\tkzDrawPolygon(A,C,G,H) 
\tkzDrawPolygon(C,B,I,J) 
\tkzDrawPolygon(B,A,E,F) 
\tkzLabelSegment(A,C){$a$} 
\tkzLabelSegment[right](C,B){$b$} 
\tkzLabelSegment[swap](A,B){$c$} 
\end{tikzpicture}
\end{tkzexample}

\subsection{矩形顶点的定义}
。

\begin{NewMacroBox}{tkzDefRectangle}{\parg{pt1,pt2}}%
给定两个点，可再得到两个点，使这四个点按顺序构成一个矩形；矩形按逆时针方向定义。作为参数给出的两个点是一条对角线的两个端点，矩形的各边分别平行于坐标轴。\\
 结果保存在 \tkzname{tkzFirstPointResult} 和 \tkzname{tkzSecondPointResult} 中。\\
可用 \tkzcname{tkzGetPoints} 为它们重新命名。

\medskip
\begin{tabular}{lll}%
\toprule
参数             & 示例 & 说明                         \\ 
\midrule
\TAline{\parg{pt1,pt2}}{\tkzcname{tkzDefRectangle}\parg{A,B}}{矩形按逆时针方向定义。}
\end{tabular}
\end{NewMacroBox}

\subsubsection{矩形定义示例}
\begin{tkzexample}[latex=7 cm,small]
\begin{tikzpicture}
\tkzDefPoints{0/0/A,5/2/C}
\tkzDefRectangle(A,C) \tkzGetPoints{B}{D}
\tkzDrawPolygon[fill=teal!15](A,...,D)
\end{tikzpicture}
\end{tkzexample}

\subsection{平行四边形的定义} 

定义平行四边形的顶点：由给定的三个点补出第四个点。
\begin{NewMacroBox}{tkzDefParallelogram}{\parg{pt1,pt2,pt3}}%
\begin{tabular}{lll}%
\toprule
参数 &  默认值 & 含义  \\ 
\midrule
\TAline{\parg{pt1,pt2,pt3}}{无默认}{需要给出三个点}
\bottomrule
\end{tabular}
\end{NewMacroBox}

由三个点可得到另一个点，使这四个点按顺序构成一个平行四边形。
\\ 结果保存在 \tkzname{tkzPointResult} 中。\\
可用 \tkzcname{tkzGetPoint} 为它重新命名。


\subsubsection{平行四边形定义示例}

\begin{tkzexample}[latex=7 cm,small]
\begin{tikzpicture}[scale=1]
 \tkzDefPoints{0/0/A,3/0/B,4/2/C} 
 \tkzDefParallelogram(A,B,C) 
 % or   \tkzDefPointWith[colinear= at C](B,A) 
 \tkzGetPoint{D}
 \tkzDrawPolygon(A,B,C,D)
 \tkzLabelPoints(A,B) 
 \tkzLabelPoints[above right](C,D)
 \tkzDrawPoints(A,...,D)
\end{tikzpicture}
\end{tkzexample}


\subsection{黄金矩形} 
 \begin{NewMacroBox}{tkzDefGoldenRectangle}{\parg{point,point}}%
本宏确定一个矩形，使其邻边之比为 $\Phi$。\\
 新生成的点保存在 \tkzname{tkzFirstPointResult} 和 \tkzname{tkzSecondPointResult} 中。\\
 可用宏 \tkzcname{tkzGetPoints} 取出它们。

\begin{tabular}{lll}%
\toprule
参数             & 示例 & 说明                         \\
\midrule
\TAline{\parg{pt1,pt2}}{\parg{A,B}}{若新生成 C、D，则 $AB/BC=\Phi$。}
 \end{tabular}
 
 \tkzcname{tkzDefGoldenRectangle} 或 \tkzcname{tkzDefGoldRectangle}
\end{NewMacroBox}


\subsubsection{黄金矩形}
\begin{tkzexample}[latex=6 cm,small]
\begin{tikzpicture}[scale=.6]
 \tkzDefPoint(0,0){A}      \tkzDefPoint(8,0){B}
 \tkzDefGoldRectangle(A,B) \tkzGetPoints{C}{D}
 \tkzDefGoldRectangle(B,C) \tkzGetPoints{E}{F}
 \tkzDefGoldRectangle(C,E) \tkzGetPoints{G}{H}
 \tkzDrawPolygon(A,B,C,D)
 \tkzDrawSegments(E,F G,H)
\end{tikzpicture}
\end{tkzexample}

\subsubsection{黄金矩形的尺规作法}
下面给出不使用上述宏来构造黄金矩形的方法。

\begin{tkzexample}[latex=8cm,small]
\begin{tikzpicture}[scale=.5] 
\tkzDefPoint(0,0){A}
\tkzDefPoint(8,0){B} 
\tkzDefMidPoint(A,B)
\tkzGetPoint{I} 
\tkzDefSquare(A,B)\tkzGetPoints{C}{D} 
\tkzInterLC(A,B)(I,C)\tkzGetPoints{G}{E} 
\tkzDefPointWith[colinear= at C](E,B) 
 \tkzGetPoint{F}
\tkzDefPointBy[projection=onto D--C ](E) 
 \tkzGetPoint{H}
\tkzDrawArc[style=dashed](I,E)(D)
\tkzDrawPolygon(A,B,C,D) 
\tkzDrawPoints(C,D,E,F,H) 
\tkzLabelPoints(A,B,C,D,E,F,H)
\tkzLabelPoints[above](C,D,F,H)  
\tkzDrawSegments[style=dashed,color=gray]%
(E,F C,F B,E F,H H,C E,H) 
\end{tikzpicture}
\end{tkzexample}




\subsection{正多边形} 
 \begin{NewMacroBox}{tkzDefRegPolygon}{\oarg{local options}\parg{pt1,pt2}}%
给定边数后，本宏根据所选选项，从中心或从一条边出发确定一个正多边形。

\begin{tabular}{lll}%
\toprule
参数             & 示例 & 说明                         \\
\midrule
\TAline{\parg{pt1,pt2}}{\parg{O,A}}{配合选项 \code{center}，$O$ 为正多边形的中心。}
\TAline{\parg{pt1,pt2}}{\parg{A,B}}{配合选项 \code{side}，$[AB]$ 为正多边形的一条边。}
 \end{tabular}

\medskip
\begin{tabular}{lll}%
\toprule
选项             & 默认值 & 示例                         \\
\midrule
\TOline{name}{P}{顶点命名为 $P1$、$P2$、……}
\TOline{sides}{5}{边数。}
\TOline{center}{center}{第一个点为中心。}
\TOline{side}{center}{给出的两个点为顶点。}
\TOline{Options TikZ}{...}{}
\end{tabular} 
\end{NewMacroBox}

\subsubsection{选项 \tkzname{center}}
\begin{tkzexample}[latex=7cm, small]   
\begin{tikzpicture}
  \tkzDefPoints{0/0/P0,0/0/Q0,2/0/P1}
  \tkzDefMidPoint(P0,P1) \tkzGetPoint{Q1}
  \tkzDefRegPolygon[center,sides=7](P0,P1)
  \tkzDefMidPoint(P1,P2) \tkzGetPoint{Q1}
  \tkzDefRegPolygon[center,sides=7,name=Q](P0,Q1)
  \tkzFillPolygon[teal!20](Q0,Q1,P2,Q2)
  \tkzDrawPolygon(P1,P...,P7)
  \foreach \j in {1,...,7} {%
  \tkzDrawSegment[black](P0,Q\j)}
\end{tikzpicture}
\end{tkzexample}

\subsubsection{选项 \tkzname{side}}
\begin{tkzexample}[latex=7cm, small]   
\begin{tikzpicture}[scale=1]
    \tkzDefPoints{-4/0/A, -1/0/B}
    \tkzDefRegPolygon[side,sides=5,name=P](A,B)
    \tkzDrawPolygon[thick](P1,P...,P5)
\end{tikzpicture}
\end{tkzexample}
\endinput


\part{圆}
\newpage
\section{圆}

在以下宏命令中，有一个用于绘制圆，这并非什么了不起的功能：你需要圆心以及半径或圆上一点。对我来说最常见的用法是给定圆心并过给定点作圆；这是默认方式，否则使用选项 \tkzname{R}。还有许多特殊的圆，例如三角形的外接圆。

\begin{itemize}
  \item 第一个宏命令 \tkzcname{tkzDefCircle}，给出某个圆的圆心和半径（单位为 cm），你可以用宏命令 \tkzcname{tkzGetPoint} 和 \tkzcname{tkzGetLength} 获取；
 
 \item 用于绘制圆的宏命令 \tkzcname{tkzDrawCircle}；
 
 \item 用于填充圆盘而不绘制圆的宏命令 \tkzcname{tkzFillCircle}；
 
 \item 有时图形必须包含在一个圆盘内：这就是 \tkzcname{tkzClipCircle} 的作用；
  
 \item 最后，\tkzcname{tkzLabelCircle} 用于给圆添加标签。
\end{itemize} 

\subsection{圆的特征：\tkzcname{tkzDefCircle}}
 
本宏命令给出某些圆的特征（圆心和半径）。

\begin{NewMacroBox}{tkzDefCircle}{\oarg{local options}\parg{A,B} or \parg{A,B,C}}%
\tkzHandBomb\ 注意参数是由两个或三个点组成的列表。本宏命令既可与 \\ \tkzcname{tkzGetPoints} 配合使用，以获取圆心和圆上一点，也可在不需要保留结果时直接使用 \\ \tkzname{tkzFirstPointResult} 和 \tkzname{tkzSecondPointResult}。\tkzcname{tkzGetLength} 用于获取半径。

\medskip
\begin{tabular}{lll}%
\toprule
参数           & 示例 & 说明                         \\
\midrule
\TAline{\parg{pt1,pt2} or \parg{pt1,pt2,pt3}}{\parg{A,B}} {$[AB]$ 为半径，$A$ 为圆心}
\bottomrule
\end{tabular} 

\medskip
\begin{tabular}{lll}%
\toprule
选项             & 默认值 & 定义                         \\ 
\midrule
\TOline{R}       {circum}{由圆心和半径确定的圆} 
\TOline{diameter}{circum}{由确定直径的两个点确定的圆}
\TOline{circum}       {circum}{三角形的外接圆} 
\TOline{in}           {circum}{三角形的内切圆} 
\TOline{ex}           {circum}{三角形的旁切圆}
\TOline{euler or nine}{circum}{欧拉圆（九点圆）}
\TOline{spieker}      {circum}{斯俾克圆}
\TOline{apollonius}   {circum}{阿波罗尼斯圆}
\TOline{orthogonal from} {circum}{[orthogonal from = A ](O,M)}
\TOline{orthogonal through}{circum}{[orthogonal through = A and B](O,M)} 
\TOline{K} {1}{阿波罗尼斯圆所使用的系数} 
 \bottomrule
\end{tabular}

\medskip
\emph{在以下示例中，我用一个尚未介绍的宏命令绘制圆。你可能只需要圆心和圆上一点。}
\end{NewMacroBox} 

\subsubsection{使用选项 \tkzname{R} 的示例}  
借助宏命令 \tkzcname{tkzGetPoint}，我们得到圆上一点，即其"东极"。

\begin{tkzexample}[latex=7cm,small]  
\begin{tikzpicture}[scale=1]
  \tkzDefPoint(3,3){C}
  \tkzDefPoint(5,5){A}
   \tkzCalcLength(A,C) \tkzGetLength{rAC}
  \tkzDefCircle[R](C,\rAC) \tkzGetPoint{B}
  \tkzDrawCircle(C,B)
  \tkzDrawSegment(C,A)
  \tkzLabelSegment[above left](C,A){$2\sqrt{2}$}
  \tkzDrawPoints(A,B,C)
  \tkzLabelPoints(A,C,B)
\end{tikzpicture} 
\end{tkzexample}     


\subsubsection{使用选项 \tkzname{diameter} 的示例}  
 这里更简单的做法是直接求 $[AB]$ 的中点。结果即为圆心，如有必要则
\begin{tkzexample}[latex=7cm,small]  
\begin{tikzpicture}
  \tkzDefPoint(0,0){O}
  \tkzDefPoint(2,2){B}
  \tkzDefCircle[diameter](O,B) \tkzGetPoint{A}
  \tkzDrawCircle(A,B)
  \tkzDrawPoints(O,A,B)
  \tkzDrawSegment(O,B)
  \tkzLabelPoints(O,A,B)
  \tkzLabelSegment[above left](O,A){$\sqrt{2}$}
  \tkzLabelSegment[above left](A,B){$\sqrt{2}$}
  \tkzMarkSegments[mark=s||](O,A A,B)
\end{tikzpicture}
\end{tkzexample}    

\subsubsection{给定三角形的内切圆与外接圆} 
 
\begin{tkzexample}[latex=7cm,small]  
\begin{tikzpicture}[scale=.75]
 \tkzDefPoint(2,2){A}  \tkzDefPoint(5,-2){B}
 \tkzDefPoint(1,-2){C}
 \tkzDefCircle[in](A,B,C)
 \tkzGetPoints{I}{x}   
 \tkzDefCircle[circum](A,B,C)
 \tkzGetPoint{K}  
 \tkzDrawCircles[new](I,x K,A) 
 \tkzLabelPoints[below](B,C)
 \tkzLabelPoints[above left](A,I,K)
 \tkzDrawPolygon(A,B,C)
 \tkzDrawPoints(A,B,C,I,K) 
\end{tikzpicture} 
\end{tkzexample}

\subsubsection{使用选项 \tkzname{ex} 的示例}
我们希望定义三角形中与顶点 $C$ 相对的旁切圆。

\begin{tkzexample}[latex=8cm,small]
\begin{tikzpicture}[scale=.75]
  \tkzDefPoints{ 0/0/A,4/0/B,0.8/4/C}
  \tkzDefCircle[ex](B,C,A)                   
  \tkzGetPoints{J_c}{h}
  \tkzDefPointBy[projection=onto A--C ](J_c)   
  \tkzGetPoint{X_c}
  \tkzDefPointBy[projection=onto A--B ](J_c)   
  \tkzGetPoint{Y_c}     
  \tkzDefCircle[in](A,B,C)    
  \tkzGetPoints{I}{y}
  \tkzDrawCircles[color=lightgray](J_c,h I,y)
  \tkzDefPointBy[projection=onto A--C ](I)
  \tkzGetPoint{F}
  \tkzDefPointBy[projection=onto A--B ](I)
  \tkzGetPoint{D}
  \tkzDrawPolygon(A,B,C)
  \tkzDrawLines[add=0 and 1.5](C,A C,B)
  \tkzDrawSegments(J_c,X_c I,D  I,F J_c,Y_c)
  \tkzMarkRightAngles(A,F,I B,D,I J_c,X_c,A
     J_c,Y_c,B)
  \tkzDrawPoints(B,C,A,I,D,F,X_c,J_c,Y_c)
  \tkzLabelPoints(B,A,J_c,I,D)
  \tkzLabelPoints[above](Y_c)
  \tkzLabelPoints[left](X_c)
  \tkzLabelPoints[above left](C)
  \tkzLabelPoints[left](F)
\end{tikzpicture}  
\end{tkzexample}

\subsubsection{给定三角形使用选项 \tkzname{euler} 的欧拉圆}
 
我们验证该圆经过每条边的中点。
\begin{tkzexample}[latex=6cm,small]  
\begin{tikzpicture}[scale=.75]
   \tkzDefPoint(5,3.5){A} 
   \tkzDefPoint(0,0){B} \tkzDefPoint(7,0){C}
   \tkzDefCircle[euler](A,B,C)
   \tkzGetPoints{E}{e}
   \tkzDefSpcTriangle[medial](A,B,C){M_a,M_b,M_c}
   \tkzDrawCircle[new](E,e)
   \tkzDrawPoints(A,B,C,E,M_a,M_b,M_c)    
   \tkzDrawPolygon(A,B,C)    
   \tkzLabelPoints[below](B,C)  
   \tkzLabelPoints[left](A,E)   
\end{tikzpicture}
\end{tkzexample}

\subsubsection{给定线段的阿波罗尼斯圆：选项 \tkzname{apollonius}} 
 
\begin{tkzexample}[latex=9cm,small]    
\begin{tikzpicture}[scale=0.75]
  \tkzDefPoint(0,0){A} 
  \tkzDefPoint(4,0){B}
  \tkzDefCircle[apollonius,K=2](A,B)
  \tkzGetPoints{K1}{x}
  \tkzDrawCircle[color = teal!50!black,
      fill=teal!20,opacity=.4](K1,x)
  \tkzDefCircle[apollonius,K=3](A,B)
  \tkzGetPoints{K2}{y}
  \tkzDrawCircle[color=orange!50,
      fill=orange!20,opacity=.4](K2,y) 
  \tkzLabelPoints[below](A,B,K1,K2)
  \tkzDrawPoints(A,B,K1,K2) 
  \tkzDrawLine[add=.2 and 1](A,B)  
\end{tikzpicture}
\end{tkzexample}  

\subsubsection{给定三角形的旁切圆：选项 \tkzname{ex}}
 你也可以获取圆心及其在三角形某条边上的投影。
 
 可使用 \tkzcname{tkzGetFirstPoint\{Jb\}} 和 \tkzcname{tkzGetSecondPoint\{Tb\}}。
 
\begin{tkzexample}[latex=8cm,small]  
\begin{tikzpicture}[scale=.6]
  \tkzDefPoint(0,0){A}
  \tkzDefPoint(3,0){B}
  \tkzDefPoint(1,2.5){C}
  \tkzDefCircle[ex](A,B,C) \tkzGetPoints{I}{i}
  \tkzDefCircle[ex](C,A,B) \tkzGetPoints{J}{j}
  \tkzDefCircle[ex](B,C,A) \tkzGetPoints{K}{k}
  \tkzDefCircle[in](B,C,A) \tkzGetPoints{O}{o}
  \tkzDrawCircles[new](J,j I,i K,k O,o) 
  \tkzDrawLines[add=1.5 and 1.5](A,B A,C B,C)
  \tkzDrawPolygon[purple](I,J,K)
  \tkzDrawSegments[new](A,K B,J C,I)
  \tkzDrawPoints(A,B,C)
  \tkzDrawPoints[new](I,J,K)   
  \tkzLabelPoints(A,B,C,I,J,K)
\end{tikzpicture}
\end{tkzexample}
 
\subsubsection{使用选项 \tkzname{spieker} 的斯俾克圆}   
中点三角形 $M_aM_bM_c$ 的内切圆即为斯俾克圆：

\begin{tkzexample}[latex=6cm, small]
\begin{tikzpicture}[scale=1.25]
  \tkzDefPoints{ 0/0/A,4/0/B,0.8/4/C}
  \tkzDefSpcTriangle[medial](A,B,C){M_a,M_b,M_c}
  \tkzDefTriangleCenter[spieker](A,B,C) 
  \tkzGetPoint{S_p}
  \tkzDrawPolygon(A,B,C)
  \tkzDrawPolygon[cyan](M_a,M_b,M_c)
  \tkzDrawPoints(B,C,A)
  \tkzDefCircle[spieker](A,B,C)
  \tkzDrawPoints[new](M_a,M_b,M_c,S_p)
  \tkzDrawCircle[new](tkzFirstPointResult,%
    tkzSecondPointResult)
  \tkzLabelPoints[right](M_a)
  \tkzLabelPoints[left](M_b)
  \tkzLabelPoints[below](A,B,M_c,S_p)
  \tkzLabelPoints[above](C)
\end{tikzpicture}
\end{tkzexample}
 
\subsection{旁心的投影}

\begin{NewMacroBox}{tkzDefProjExcenter}{\oarg{local options}\parg{A,B,C}\parg{a,b,c}\marg{X,Y,Z}}%
每个旁心在三角形 $ABC$ 的三边上各有三个投影。我们可以用一个宏命令获取全部：\\ \tkzcname{tkzDefProjExcenter[name=J](A,B,C)(a,b,c)\{Y,Z,X\}}。

\medskip
\begin{tabular}{lll}%
\toprule
选项             & 默认值 & 定义                        \\
\midrule
\TOline{name} {无默认}{用于命名顶点}
\bottomrule
\end{tabular}

\begin{tabular}{lll}%
参数 & 默认值 & 定义 \\
\midrule
\TAline{(pt1=$\alpha_1$,pt2=$\alpha_2$,\dots)}{无默认}{每个点都被赋予一个权重}
\end{tabular}

\medskip
\end{NewMacroBox}

\subsubsection{\tkzname{旁切圆}}

\begin{tkzexample}[vbox,small]
\begin{tikzpicture}[scale=.6]
\tikzset{line style/.append style={line width=.2pt}}
\tikzset{label style/.append style={color=teal,font=\footnotesize}}
\tkzDefPoints{0/0/A,5/0/B,0.8/4/C}
\tkzDefSpcTriangle[excentral,name=J](A,B,C){a,b,c} 
\tkzDefSpcTriangle[intouch,name=I](A,B,C){a,b,c}
\tkzDefProjExcenter[name=J](A,B,C)(a,b,c){X,Y,Z}
\tkzDefCircle[in](A,B,C)   \tkzGetPoint{I} \tkzGetSecondPoint{T}  
\tkzDrawCircles[red](Ja,Xa Jb,Yb Jc,Zc)
\tkzDrawCircle(I,T) 
\tkzDrawPolygon[dashed,color=blue](Ja,Jb,Jc)
\tkzDrawLines[add=1.5 and 1.5](A,C A,B B,C)
\tkzDrawSegments(Ja,Xa Ja,Ya Ja,Za
                 Jb,Xb Jb,Yb Jb,Zb
                 Jc,Xc Jc,Yc Jc,Zc
                 I,Ia I,Ib I,Ic)
\tkzMarkRightAngles[size=.2,fill=gray!15](Ja,Za,B Ja,Xa,B Ja,Ya,C Jb,Yb,C)
\tkzMarkRightAngles[size=.2,fill=gray!15](Jb,Zb,B Jb,Xb,C Jc,Yc,A Jc,Zc,B)
\tkzMarkRightAngles[size=.2,fill=gray!15](Jc,Xc,C I,Ia,B I,Ib,C I,Ic,A)
\tkzDrawSegments[blue](Jc,C Ja,A Jb,B)
\tkzDrawPoints(A,B,C,Xa,Xb,Xc,Ja,Jb,Jc,Ia,Ib,Ic,Ya,Yb,Yc,Za,Zb,Zc)
\tkzLabelPoints(A,Ya,Yb,Ja,I)
\tkzLabelPoints[left](Jb,Ib,Yc)
\tkzLabelPoints[below](Zb,Ic,Jc,B,Za,Xa)
\tkzLabelPoints[above right](C,Zc,Yb)
\tkzLabelPoints[right](Xb,Ia,Xc)
\end{tikzpicture}
\end{tkzexample}
 
\subsubsection{\tkzname{Orthogonal from}}
给定圆心且与给定圆正交的圆。\tkzcname{tkzGetPoints\{z1\}\{z2\}} 给出该圆上的两个点。

\begin{tkzexample}[latex=7cm,small]
\begin{tikzpicture}[scale=.75]
  \tkzDefPoints{0/0/O,1/0/A}
  \tkzDefPoints{1.5/1.25/B,-2/-3/C}
  \tkzDefCircle[orthogonal from=B](O,A)
  \tkzGetPoints{z1}{z2}
  \tkzDefCircle[orthogonal from=C](O,A)
  \tkzGetPoints{t1}{t2}
  \tkzDrawCircle(O,A)
  \tkzDrawCircles[new](B,z1 C,t1)
  \tkzDrawPoints(t1,t2,C)
  \tkzDrawPoints(z1,z2,O,A,B)
  \tkzLabelPoints[right](O,A,B,C)
\end{tikzpicture}
\end{tkzexample}

\subsubsection{\tkzname{Orthogonal through}}
与给定圆正交且经过两个给定点的圆。
\begin{tkzexample}[latex=6cm,small]
\begin{tikzpicture}[scale=1]
  \tkzDefPoint(0,0){O}
  \tkzDefPoint(1,0){A}
  \tkzDrawCircle(O,A)
  \tkzDefPoint(-1.5,-1.5){z1}
  \tkzDefPoint(1.5,-1.25){z2}
  \tkzDefCircle[orthogonal through=z1 and z2](O,A)
   \tkzGetPoint{c}
  \tkzDrawCircle[new](tkzPointResult,z1)
  \tkzDrawPoints[new](O,A,z1,z2,c)
  \tkzLabelPoints[right](O,A,z1,z2,c)
\end{tikzpicture}
\end{tkzexample}

\endinput

\section{由变换得到的圆}
\subsection{宏命令 \tkzcname{tkzDefCircleBy}}
可用的变换有：

\begin{itemize}
   \item 平移；
   \item 位似；
   \item 正交反射或轴对称；
   \item 中心对称；
   \item 正交投影；
   \item 旋转（角度制）；
   \item 反演。
\end{itemize}

变换通过选项来选择。例如，我们写：
\begin{tkzltxexample}[]
\tkzDefCircleBy[translation= from A to A'](O,M) 
\end{tkzltxexample}
$O$ 是圆心，$M$ 是圆上一点。
像仍是一个圆：新圆心为 |tkzFirstPointResult|，|tkzSecondPointResult| 是新圆上的一点。你可以用宏命令 \tkzcname{tkzGetPoints} 获取结果。使用反演时，过反演中心的圆的像是一条直线：此时两个点即为该直线上的两个点。
\medskip
\begin{NewMacroBox}{tkzDefCircleBy}{\oarg{local options}\parg{pt1,pt2}}%
参数为一对点，结果同样是一对点。要保留这些点，\tkzcname{tkzGetPoints\{O'\}\{M'\}} 将圆心命名为 \tkzname{O'}，将圆上的点命名为 \tkzname{M'}。

\begin{tabular}{lll}%
\toprule
参数 &  定义 & 示例               \\ 
\midrule
\TAline{pt1,pt2}   {已有点}   {$(O,M)$}
\bottomrule
\end{tabular}

\begin{tabular}{lll}%
选项     &     & 示例                         \\ 
\midrule
\TOline{translation}{= from \#1 to \#2}{[translation=from A to B](O,M)}
\TOline{homothety}  {= center \#1 ratio \#2}{[homothety=center A ratio .5](O,M)}
\TOline{reflection} {= over \#1--\#2}{[reflection=over A--B](O,M)}
\TOline{symmetry }  {= center \#1}{[symmetry=center A](O,M)}
\TOline{projection }{= onto \#1--\#2}{[projection=onto A--B](O,M)}
\TOline{rotation }  {= center \#1 angle \#2}{[rotation=center O angle 30](O,M)}
\TOline{inversion}{= center \#1 through \#2}{[inversion =center O through A](O,M)} 
% \TOline{inversion negative}{= center \#1 through \#2}{[inversion negative =center O through A](O,M)}
\bottomrule
\end{tabular}

\medskip
\emph{像仅被定义，并未绘制。}
\end{NewMacroBox} 

\subsubsection{\tkzname{平移}}
\begin{tkzexample}[latex=7cm,small]
\begin{tikzpicture}[>=latex] 
 \tkzDefPoint(0,0){A}  \tkzDefPoint(3,1){B}
 \tkzDefPoint(3,2){C}   \tkzDefPoint(4,3){D}
 \tkzDefCircleBy[translation= from B to A](C,D) 
 \tkzGetPoints{C'}{D'} 
 \tkzDrawPoints[teal](A,B,C,D,C',D')
 \tkzDrawSegments[orange,->](A,B)
 \tkzDrawCircles(C,D C',D')
 \tkzLabelPoints[color=teal](A,B,C,C') 
 \tkzLabelPoints[color=teal,above](D,D') 
\end{tikzpicture} 
\end{tkzexample}

\subsubsection{\tkzname{反射}（轴对称）}

\begin{tkzexample}[latex=7cm,small]
\begin{tikzpicture}[>=latex] 
 \tkzDefPoint(0,0){A}  \tkzDefPoint(3,1){B}
 \tkzDefPoint(3,2){C}   \tkzDefPoint(4,3){D}
 \tkzDefCircleBy[reflection = over A--B](C,D)
 \tkzGetPoints{C'}{D'} 
 \tkzDrawPoints[teal](A,B,C,D,C',D')
 \tkzDrawLine[add =0 and 1][orange](A,B)
 \tkzDrawCircles(C,D C',D')
 \tkzLabelPoints[color=teal](A,B,C,C') 
 \tkzLabelPoints[color=teal,right](D,D') 
\end{tikzpicture} 
\end{tkzexample}

\subsubsection{\tkzname{位似}}

\begin{tkzexample}[latex=7cm,small]
\begin{tikzpicture}[scale=1.2]
 \tkzDefPoint(0,0){A}   \tkzDefPoint(3,1){B}
 \tkzDefPoint(3,2){C}   \tkzDefPoint(4,3){D}
 \tkzDefCircleBy[homothety=center A ratio .5](C,D)
 \tkzGetPoints{C'}{D'}
 \tkzDrawPoints[teal](A,C,D,C',D')
 \tkzDrawCircles(C,D C',D')
 \tkzLabelPoints[color=teal](A,C,C')
 \tkzLabelPoints[color=teal,right](D,D') 
\end{tikzpicture}
\end{tkzexample}

\subsubsection{\tkzname{中心对称}}
\begin{tkzexample}[latex=7cm,small]
\begin{tikzpicture}[scale=1]
 \tkzDefPoint(3,1){B}
 \tkzDefPoint(3,2){C}   \tkzDefPoint(4,3){D}
 \tkzDefCircleBy[symmetry=center B](C,D)
 \tkzGetPoints{C'}{D'}
 \tkzDrawPoints[teal](B,C,D,C',D')
 \tkzDrawLines[orange](C,C' D,D')
 \tkzDrawCircles(C,D C',D')
 \tkzLabelPoints[color=teal](C,C')
 \tkzLabelPoints[color=teal,above](D)
 \tkzLabelPoints[color=teal,below](D')
\end{tikzpicture}
\end{tkzexample}

\subsubsection{\tkzname{旋转}}
\begin{tkzexample}[latex=7cm,small]
\begin{tikzpicture}[scale=0.5]
 \tkzDefPoint(3,-1){B}
 \tkzDefPoint(3,2){C}   \tkzDefPoint(4,3){D}
 \tkzDefCircleBy[rotation=center B angle 90](C,D)
 \tkzGetPoints{C'}{D'}
 \tkzDrawPoints[teal](B,C,D,C',D')
 \tkzLabelPoints[color=teal](B,C,D,C',D')
 \tkzDrawCircles(C,D C',D')
\end{tikzpicture}
\end{tkzexample}

\subsubsection{\tkzname{反演}}

\begin{tkzexample}[latex=7cm,small]
\begin{tikzpicture}[scale=1.5]
\tkzSetUpPoint[size=3,color=red,fill=red!20]
\tkzSetUpStyle[color=purple,ultra thin]{st1}
\tkzSetUpStyle[color=cyan,ultra thin]{st2}
\tkzDefPoint(2,0){A} \tkzDefPoint(3,0){B} 
\tkzDefPoint(3,2){C} \tkzDefPoint(4,2){D} 
\tkzDefCircleBy[inversion = center B through A](C,D) 
\tkzGetPoints{C'}{D'}
\tkzDrawPoints(A,B,C,D,C',D') 
\tkzLabelPoints(A,B,C,D,C',D')
\tkzDrawCircles(B,A)
\tkzDrawCircles[st1](C,D)
\tkzDrawCircles[st2](C',D')
\end{tikzpicture}
\end{tkzexample}

\endinput


\part{交点与角}
\section{交点}

你可以求两条直线、一条直线与一个圆、以及两个圆的交点。

相关命令不会检查交点是否存在：这由用户自行负责（参见测试命令 \tkzcname{tkzTestInterLC} 和 \tkzcname{tkzTestInterCC}）。

\subsection{两条直线的交点 \tkzcname{tkzInterLL}}
\begin{NewMacroBox}{tkzInterLL}{\parg{$A,B$}\parg{$C,D$}}%
定义两条直线 $(AB)$ 与 $(CD)$ 的交点 \tkzname{tkzPointResult}。已知点以成对形式给出（每条直线两个），放在圆括号中，所得结果点用宏命令 \tkzcname{tkzGetPoint} 获取。
\end{NewMacroBox}

\subsubsection{两条直线相交的示例}

\begin{tkzexample}[latex=7cm,small]
\begin{tikzpicture}[rotate=-45,scale=.75]
  \tkzDefPoint(2,1){A}   
  \tkzDefPoint(6,5){B}
  \tkzDefPoint(3,6){C}   
  \tkzDefPoint(5,2){D}
  \tkzDrawLines(A,B C,D)
  \tkzInterLL(A,B)(C,D)  
     \tkzGetPoint{I}
  \tkzDrawPoints[color=blue](A,B,C,D)
   \tkzDrawPoint[color=red](I)
\end{tikzpicture}
\end{tkzexample}

\subsection{直线与圆的交点  \tkzcname{tkzInterLC}}

与前面一样，直线由一对点确定。圆
 也由一对点确定：
\begin{itemize}
\item $(O,C)$，一对点：第一个是圆心，第二个是圆上任意一点；
\item $(O,r)$，其中 $r$ 是半径。
\end{itemize}

\begin{NewMacroBox}{tkzInterLC}{\oarg{options}\parg{$A,B$}\parg{$O,C$} or \parg{$O,r$} or \parg{$O,C,D$}}%
因此参数为两对。

\medskip
\begin{tabular}{lll}%
\toprule
选项            & 默认值 & 定义                         \\ 
\midrule
\TOline{N}         {N}    {$(O,C)$ 确定圆}
\TOline{R}         {N}    {$(O, 1 )$ 单位为 1 cm}  
\TOline{with nodes}{N}    {$(O,C,D)$，$CD$ 为半径}  
\TOline{common=pt} {}     {pt 为公共点；tkzFirstPoint 给出另一个点}
\TOline{near}      {}     {tkzFirstPoint 是最靠近直线第一个点的点}
\bottomrule
\end{tabular}

\medskip   
本宏命令定义直线 $(AB)$ 与以 $O$ 为圆心、$r$ 为半径的圆的交点 $I$ 和 $J$（若存在）；否则在 |.log| 文件中报告错误。使用 \tkzname{with nodes} 时无需计算半径，半径即为 $[CD]$ 的长度。
若未使用 \tkzname{common} 和 \tkzname{near}，\tkzname{tkzFirstPoint} 是给出最小角度的点（该角由 \tkzname{tkzSecondPoint} 与圆心构成）。
\end{NewMacroBox}

\begin{NewMacroBox}{tkzTestInterLC}{\parg{$O,A$}\parg{$O',B$}}%
参数为两对，分别定义一条直线和一个圆（各由圆心和圆上一点给出）。若直线与圆相交，则测试命令 \tkzcname{iftkzFlagLC} 为真。
\end{NewMacroBox}

\subsubsection{测试直线与圆是否相交：\tkzcname{tkzTestInterLC}}

\begin{tkzexample}[latex=7cm,small]
  \begin{tikzpicture}[scale=1]
    \tkzDefPoints{% x   y   name
                    3    /4    /I,
                    3    /2    /P,
                    0    /2    /La,
                    8    /3    /Lb}
  \tkzDrawCircle(I,P)
  \foreach \i in {1,...,3}{%
     \coordinate  (Lb) at (8,\i);
     \tkzDrawLine(La,Lb)
     \tkzTestInterLC(La,Lb)(I,P)
      \iftkzFlagLC
      \tkzInterLC(La,Lb)(I,P)  
      \tkzGetPoints{a}{b}
      \tkzDrawPoints(a,b)
      \fi
     }
  \end{tikzpicture}
\end{tkzexample}


\subsubsection{直线与圆相交}

在下面的示例中，圆由两个点绘制，直线与圆的交点使用两对参数给出。我们将比较角 $\widehat{D,E,O}$ 和 $\widehat{E,D,O}$。这两个角的方向相反。\tkzname{tkzFirstPoint} 被赋予形成最小角度（逆时针方向）的那个点。逆时针角 $\widehat{D,E,O}$ 的大小为 $360^\circ$ 减去 $\widehat{O,E,D}$ 的大小。

\begin{tkzexample}[latex=7cm,small]
\begin{tikzpicture}[scale=.75]
 \tkzInit[xmax=5,ymax=4]
 \tkzDefPoint(1,1){O} 
 \tkzDefPoint(-2,4){La} 
 \tkzDefPoint(5,0){Lb} 
 \tkzDefPoint(3,3){C}
 \tkzInterLC(La,Lb)(O,C)  \tkzGetPoints{D}{E}  
 \tkzMarkAngle[->,size=1.5](E,D,O)
 \tkzDrawPolygons[new](O,D,E)
 \tkzMarkAngle[->,size=1.5](D,E,O)
 \tkzDrawCircle(O,C)
 \tkzDrawPoints[color=teal](O,La,Lb,C)
 \tkzDrawPoints[color=red](D,E)
 \tkzDrawLine(La,Lb)
 \tkzLabelPoints[above right](O,La,Lb,C,D,E)
\end{tikzpicture} 
\end{tkzexample}

\subsubsection{经过圆心的直线：选项 \tkzname{common}}
这是一种特殊情况：角度无法比较，因此必须使用选项 \tkzname{near}。\tkzname{tkzFirstPoint} 被赋予最靠近直线所给第一个点的那个点。这里我们希望 $A$ 最靠近 $Lb$。

\begin{tkzexample}[latex=8cm,small]
\begin{tikzpicture}
\tkzDefPoints{% x   y   name
             0    /1    /D,
             6    /0    /B,
             3    /3    /O,
             2    /2    /La,
             5    /5    /Lb}
  \tkzDrawCircle(O,D)
  \tkzDrawLine(La,Lb)
  \tkzInterLC[near](Lb,La)(O,D)  
  \tkzGetFirstPoint{A}
  \tkzDrawSegments(O,A)
  \tkzDrawPoints(O,D,A,La,Lb)
  \tkzLabelPoints(O,D,A,La,Lb)
\end{tikzpicture}
\end{tkzexample}

\subsubsection{使用选项 \tkzname{common} 的直线与圆相交}
一种常见的特殊情况：直线上有一点在圆上，我们需要求另一个公共点。
\begin{tkzexample}[latex=7cm,small]
\begin{tikzpicture}[scale=.5]
 \tkzDefPoints{0/0/O,-5/0/A,2/-2/B,0/5/D}
 \tkzInterLC[common=A](B,A)(O,D)
 \tkzGetFirstPoint{C}
 \tkzDrawPoints(O,A,B)
 \tkzDrawCircle(O,A)
 \tkzDrawLine(A,C)
 \tkzDrawPoint(C)
 \tkzLabelPoints(A,B,C)
\end{tikzpicture}
\end{tkzexample}


\subsubsection{直线与圆相交：点的顺序}
思路是比较由直线的第一个点、某个结果点以及圆心所构成的角。第一个点对应最小角。

如你所见，$\widehat{BCO} < \widehat{BEO}$；事实上，$\widehat{BEO}$ 是逆时针方向。

\begin{tkzexample}[latex=6cm,small]
\begin{tikzpicture}[scale=.5]
  \tkzDefPoints{0/0/O,5/1/A,2/2/B,3/1/D}
  \tkzInterLC[common=A](B,D)(O,A) \tkzGetPoints{C}{E}
  \tkzDrawPoints(O,A,B,D)
  \tkzDrawCircle(O,A) \tkzDrawLine(E,C)
  \tkzDrawSegments[dashed](B,O O,C)
  \tkzMarkAngle[->,size=1.5](B,C,O)
  \tkzDrawSegments[dashed](O,E)
  \tkzMarkAngle[->,size=1.5](B,E,O)
  \tkzDrawPoints(C,E)
  \tkzLabelPoints[above](O,E)
  \tkzLabelPoints[right](A,B,C,D)
\end{tikzpicture}
\end{tkzexample}

\subsubsection{使用 \tkzcname{foreach} 的示例}
\begin{tkzexample}[latex=7cm,small]
\begin{tikzpicture}[scale=3,rotate=180]
\tkzDefPoint(0,1){J} 
\tkzDefPoint(0,0){O}
\foreach \i in {0,-5,-10,...,-90}{
 \tkzDefPoint({2.5*cos(\i*pi/180)},%
   {1+2.5*sin(\i*pi/180)}){P}
 \tkzInterLC[R](P,J)(O,1)\tkzGetPoints{N}{M}
 \tkzDrawSegment[color=orange](J,N)
 \tkzDrawPoints[red](N)} 
\foreach \i in {-90,-95,...,-175,-180}{
 \tkzDefPoint({2.5*cos(\i*pi/180)},%
   {1+2.5*sin(\i*pi/180)}){P} 
 \tkzInterLC[R](P,J)(O,1)\tkzGetPoints{N}{M}
 \tkzDrawSegment[color=orange](J,M)
 \tkzDrawPoints[red](M)}   
\end{tikzpicture} 
\end{tkzexample}

\subsubsection{使用选项 \tkzname{near} 的直线与圆相交}
$D$ 是最靠近 $b$ 的点。

\begin{tkzexample}[vbox,small]
  \begin{tikzpicture}
    \tkzDefPoints{0/0/A,12/0/C}
    \tkzDefGoldenRatio(A,C)                          \tkzGetPoint{B}
    \tkzDefMidPoint(A,C)                             \tkzGetPoint{O}
    \tkzDefMidPoint(A,B)                             \tkzGetPoint{O_1}
    \tkzDefMidPoint(B,C)                             \tkzGetPoint{O_2}
    \tkzDefPointBy[rotation=center O_2 angle 90](C)  \tkzGetPoint{P}
    \tkzDefPointBy[rotation=center O_1 angle 90](B)  \tkzGetPoint{Q}
    \tkzDefPointBy[rotation=center B angle 90](C)    \tkzGetPoint{b}
    \tkzInterLC[near](b,B)(O,A)                      \tkzGetFirstPoint{D}
    \tkzInterCC(D,B)(O,C)                            \tkzGetPoints{V}{U}
    \tkzDefPointBy[projection=onto U--V](O_1)        \tkzGetPoint{M}
    \tkzDefPointBy[projection=onto U--V](O_2)        \tkzGetPoint{N}  
    \tkzDrawPoints(A,B,C,O,O_1,O_2,D,U,V,M,N,b)
    \tkzDrawSemiCircles[teal](O,C O_1,B O_2,C)
    \tkzDrawSegments(A,C B,D U,V A,D C,D M,B B,N)
    \tkzDrawArc(D,U)(V)
    \tkzLabelPoints(A,B,C,O,O_1,O_2)
    \tkzLabelPoints[above](D,U,V,M,N)
  \end{tikzpicture}
\end{tkzexample}


\subsubsection{更复杂的直线与圆相交示例}
图形取自 \url{http://gogeometry.com/problem/p190_tangent_circle}

\begin{tkzexample}[latex=6.5cm,small]
\begin{tikzpicture}[scale=.75]
 \tkzDefPoint(0,0){A}  
 \tkzDefPoint(8,0){B}
 \tkzDefMidPoint(A,B)             \tkzGetPoint{O}
 \tkzDefMidPoint(O,B)             \tkzGetPoint{O'}
 \tkzDefLine[tangent from=A](O',B)\tkzGetFirstPoint{E}
 \tkzInterLC(A,E)(O,B)            \tkzGetFirstPoint{D}
 \tkzDefPointBy[projection=onto A--B](D)  
 \tkzGetPoint{F}
 \tkzDrawCircles(O,B O',B)
 \tkzDrawSegments(A,D A,B D,F) 
 \tkzDrawSegments[color=red,line width=1pt,
     opacity=.4](A,O F,B)
 \tkzDrawPoints(A,B,O,O',E,D) 
 \tkzMarkRightAngle(D,F,B)
 \tkzLabelPoints[below right](A,B,O,O',E,D) 
\end{tikzpicture}
\end{tkzexample}

\subsubsection{由圆心和度量值确定的圆及特殊情况}
我们来看一些特殊情况，例如与圆相切的直线。

\begin{tkzexample}[latex=7cm,small]
\begin{tikzpicture}[scale=.5]
 \tkzDefPoint(0,8){A}      \tkzDefPoint(8,0){B}
 \tkzDefPoint(8,8){C}      \tkzDefPoint(4,4){D}
 \tkzDefPoint(2,4){E}      \tkzDefPoint(4,2){F}
 \tkzDefPoint(8,4){G}
 \tkzInterLC(A,C)(D,G)     \tkzGetPoints{I1}{I2}
 \tkzInterLC(B,C)(D,G)     \tkzGetPoints{J1}{J2}
 \tkzInterLC[near](A,B)(D,G)  \tkzGetPoints{K1}{K2}
 \tkzInterLC(E,F)(D,G)     \tkzGetPoints{E1}{E2}
 \tkzDrawCircle(D,G)
 \tkzDrawPoints[color=red](I1,J1,K1,K2,E1,E2)
 \tkzDrawLines(A,B B,C A,C I2,J2 E1,E2)
 \tkzDrawPoints[color=blue](A,...,F)
 \tkzDrawPoints[color=red](I2,J2)
 \tkzLabelPoints[left](B,D,E,F)
 \tkzLabelPoints[below left](A,C)
 \tkzLabelPoints[below=4pt](I1,K1,K2,E2)
 \tkzLabelPoints[left](J1,E1)
\end{tikzpicture}

\end{tkzexample}

\subsubsection{半径的计算}
 使用 \tkzname{pgfmath} 和 \tkzcname{pgfmathsetmacro}

半径可以是计算的结果，该计算在求交宏命令之前完成，而不是在其内部进行。
长度可以通过多种方式计算。当然你可以
 使用 \tkzname{pgfmath} 模块和宏命令 \tkzcname{pgfmathsetmacro}。在某些情况下，结果精度不够：计算 $0.0002 \div 0.0001$ 用 pgfmath 得到 $1.98$，而用 xfp 得到 $2$。

使用 \tkzname{xfp} 和 \tkzcname{fpeval}：

\begin{tkzexample}[latex=7cm,small]
  \begin{tikzpicture}
  \tkzDefPoint(2,2){A}
  \tkzDefPoint(5,4){B}
  \tkzDefPoint(4,4){O}
  \pgfmathsetmacro\tkzLen{\fpeval{0.0002/0.0001}}
 % or \edef\tkzLen{\fpeval{0.0002/0.0001}}
  \tkzInterLC[R](A,B)(O, \tkzLen)
  \tkzGetPoints{I}{J}
  \tkzDrawCircle(O,I)
  \tkzDrawPoints[color=blue](A,B)
  \tkzDrawPoints[color=red](I,J)
  \tkzDrawLine(I,J)
\end{tikzpicture}
  \end{tkzexample}


\subsubsection{选项 \code{with nodes}}
\begin{tkzexample}[latex=8cm,small]
\begin{tikzpicture}[scale=.75]
\tkzDefPoints{0/0/A,4/0/B,1/1/D,2/0/E}
\tkzDefTriangle[equilateral](A,B)
\tkzGetPoint{C}
\tkzInterLC[with nodes](D,E)(C,A,B)
\tkzGetPoints{F}{G}
\tkzDrawCircle(C,A)
\tkzDrawPolygon(A,B,C)
\tkzDrawPoints(A,...,G)
\tkzDrawLine(F,G)
\end{tikzpicture}
\end{tkzexample}

\subsection{两个圆的交点  \tkzcname{tkzInterCC}}

最常见的情况是由圆心和圆上一点确定两个圆，但与前面一样，选项 \tkzname{R} 允许你用度量值给出半径。

\begin{NewMacroBox}{tkzInterCC}{\oarg{options}\parg{$O,A$}\parg{$O',A'$} or \parg{$O,r$}\parg{$O',r'$} or   \parg{$O,A,B$} \parg{$O',C,D$}}%
\begin{tabular}{lll}%
选项       & 默认值 & 定义                         \\
\midrule
\TOline{N}   {N}    {$OA$ 和 $O'A'$ 为半径，$O$ 和 $O'$ 为圆心。}
\TOline{R}   {N}    {$r$ 和 $r'$ 为度量值，表示半径。}
\TOline{with nodes} {N}  {在 (A,A,C)(C,B,F) 中，AC 和 BF 给出半径。 }
\TOline{common=pt}  {}   {pt 为公共点；tkzFirstPoint 给出另一个点。}
\bottomrule
\end{tabular}

\medskip
本宏命令定义以 $O$ 和 $O'$ 为圆心的两个圆的交点 $I$ 和 $J$。若两圆没有公共点，宏命令将以未处理的错误结束。若圆心为 $O$ 和 $O'$，交点为 $A$ 和 $B$，则角 $\widehat{O,A,O'}$ 和 $\widehat{O,B,O'}$ 方向相反；\tkzname{tkzFirstPoint} 被赋予形成 \code{clockwise}（顺时针）角的那个点。
\end{NewMacroBox}

\begin{NewMacroBox}{tkzTestInterCC}{\parg{$O,A$}\parg{$O',B$}}%
参数为两对，分别定义两个圆（每个圆由圆心和圆上一点给出）。若两圆相交，则测试命令 \tkzcname{iftkzFlagCC} 为真。
\end{NewMacroBox}

\subsubsection{测试两圆是否相交：\tkzcname{tkzTestInterCC}}

\begin{tkzexample}[latex=7cm,small]
\begin{tikzpicture}[scale=.75]
  \tkzDefPoints{% x   y   name
                   0    /0    /A,
                   2    /0    /B,
                   4    /0    /I,
                   1    /0    /P}
\tkzDrawCircle(A,B)
\foreach \i in {1,...,3}{%
   \coordinate  (P) at (\i,0);
\tkzDrawCircle[new](I,P)
   \tkzTestInterCC(A,B)(I,P)
    \iftkzFlagCC
    \tkzInterCC(A,B)(I,P)  \tkzGetPoints{a}{b}
    \tkzDrawPoints(a,b)
    \fi}
  \end{tikzpicture}
\end{tkzexample}

\subsubsection{带有 \tkzname{common} 公共点的两圆相交}

\begin{tkzexample}[latex=7cm,small]
\begin{tikzpicture}[scale=.5]
  \tkzDefPoints{0/0/O,5/-1/A,2/2/B}
  \tkzDrawPoints(O,A,B)
  \tkzDrawCircles(O,B A,B)
  \tkzInterCC[common=B](O,B)(A,B)
  \tkzGetFirstPoint{C}
  \tkzDrawPoint(C)
  \tkzLabelPoints[above](O,A,B,C)
\end{tikzpicture}
\end{tkzexample}

\subsubsection{两圆相交：点的顺序}
思路是比较由第一个圆心、某个结果点以及第二个圆的圆心所构成的角。第一个点对应最小角。

如你所见，$\widehat{ODB} < \widehat{OBE}$。

\begin{tkzexample}[latex=6cm,small]
\begin{tikzpicture}[scale=.5]
  \pgfkeys{/pgf/number format/.cd,fixed relative,
     precision=4}
  \tkzDefPoints{0/0/O,5/-1/A,2/2/B,2/-1/C}
  \tkzDrawPoints(O,A,B)
  \tkzDrawCircles(O,A B,C)
  \tkzInterCC(O,A)(B,C)\tkzGetPoints{D}{E}
  \tkzDrawPoints(C,D,E)
  \tkzLabelPoints(O,A,B,C)
  \tkzLabelPoints[above](D,E) 
  \tkzDrawSegments[cyan](D,O D,B)
  \tkzMarkAngle[red,->,size=1.5](O,D,B)
  \tkzFindAngle(O,D,B)   \tkzGetAngle{an}
  \tkzLabelAngle(O,D,B){$ \pgfmathprintnumber{\an}$}
  \tkzDrawSegments[cyan](E,O E,B)
  \tkzMarkAngle[red,->,size=1.5](O,E,B)  
  \tkzFindAngle(O,E,B)   \tkzGetAngle{an}
  \tkzLabelAngle(O,E,B){$ \pgfmathprintnumber{\an}$}  
\end{tikzpicture}
\end{tkzexample}

\subsubsection{等边三角形的构造}
$\widehat{A,C,B}$ 是一个顺时针角。
\begin{tkzexample}[latex=7cm,small]
\begin{tikzpicture}[trim left=-1cm,scale=.5]
 \tkzDefPoint(1,1){A}
 \tkzDefPoint(5,1){B}
 \tkzInterCC(A,B)(B,A)\tkzGetPoints{C}{D}
 \tkzDrawPoint[color=black](C)
 \tkzDrawCircles(A,B B,A)
 \tkzCompass[color=red](A,C)
 \tkzCompass[color=red](B,C)
 \tkzDrawPolygon(A,B,C)
 \tkzMarkSegments[mark=s|](A,C B,C)
 \tkzLabelPoints[](A,B)
 \tkzLabelPoint[above](C){$C$}
\end{tikzpicture}
\end{tkzexample}


\subsubsection{线段三等分}
 这里的思路是用尺规将线段三等分。

\begin{tkzexample}[latex=7cm,small]
\begin{tikzpicture}[scale=.6]
 \tkzDefPoint(0,0){A}
 \tkzDefPoint(3,2){B}
 \tkzInterCC(A,B)(B,A)          \tkzGetSecondPoint{D}
 \tkzInterCC(D,B)(B,A)          \tkzGetPoints{A}{C}
 \tkzInterCC(D,B)(A,B)          \tkzGetPoints{E}{B}
 \tkzInterLC[common=D](C,D)(E,D)\tkzGetFirstPoint{F}
 \tkzInterLL(A,F)(B,C)          \tkzGetPoint{O}
 \tkzInterLL(O,D)(A,B)          \tkzGetPoint{H}
 \tkzInterLL(O,E)(A,B)          \tkzGetPoint{G}
 \tkzDrawCircles(D,E A,B B,A E,A)
 \tkzDrawSegments[](O,F O,B O,D O,E)
 \tkzDrawPoints(A,...,H)
 \tkzDrawSegments(A,B B,D A,D A,E E,F C,F B,C)
 \tkzMarkSegments[mark=s|](A,G G,H H,B)
\end{tikzpicture}
\end{tkzexample}

\subsubsection{使用选项 \code{with nodes}}
\begin{tkzexample}[latex=6cm,small]
\begin{tikzpicture}[scale=.5]
 \tkzDefPoints{0/0/A,0/5/B,5/0/C}
 \tkzDefPoint(54:5){F}
 \tkzInterCC[with nodes](A,A,C)(C,B,F)
 \tkzGetPoints{a}{e}
 \tkzInterCC(A,C)(a,e) \tkzGetFirstPoint{b}
 \tkzInterCC(A,C)(b,a) \tkzGetFirstPoint{c}
 \tkzInterCC(A,C)(c,b) \tkzGetFirstPoint{d}
 \tkzDrawCircle[new](A,C)
 \tkzDrawPoints(a,b,c,d,e)
 \tkzDrawPolygon(a,b,c,d,e)
 \foreach \vertex/\num in {a/36,b/108,c/180,
                          d/252,e/324}{%
 \tkzDrawPoint(\vertex)
 \tkzLabelPoint[label=\num:$\vertex$](\vertex){}
 \tkzDrawSegment(A,\vertex)
 }
\end{tikzpicture}
\end{tkzexample}

\subsubsection{求交组合示例}
\begin{tkzexample}[latex=8cm,small]
\begin{tikzpicture}[scale = .7]
  \tkzDefPoint(2,2){A}
  \tkzDefPoint(0,0){B}
  \tkzDefPoint(-2,2){C}
  \tkzDefPoint(0,4){D}
  \tkzDefPoint(4,2){E}
  \tkzCircumCenter(A,B,C)\tkzGetPoint{O}
  \tkzInterCC[R](O,2)(D,2)\tkzGetPoints{M1}{M2}
  \tkzInterCC(O,A)(D,O) \tkzGetPoints{1}{2}
  \tkzInterLC(A,E)(B,M1)\tkzGetSecondPoint{M3}
  \tkzInterLC(O,C)(M3,D)\tkzGetSecondPoint{L}
  \tkzDrawSegments(C,L)
  \tkzDrawPoints(A,B,C,D,E,M1,M2,M3,O,L)
  \tkzDrawSegments(O,E)
  \tkzDrawSegments[new](C,A D,B)
  \tkzDrawPoint(O)
  \tkzDrawCircles[new](M3,D B,M2 D,O)
  \tkzDrawCircle(O,A)
  \tkzLabelPoints[below right](A,B,C,D,E,M1,M2,M3,O,L)
\end{tikzpicture}
\end{tkzexample}


\subsubsection{Altshiller-Court 定理}
 过两个正交圆的交点与其中一个圆上一点的两条直线，与另一个圆交于一对直径的两个端点（Altshiller-Court，第 176 页）。


\begin{tkzexample}[vbox,small]
\begin{tikzpicture}
  \tkzDefPoints{0/0/P,5/0/Q,3/2/I}
  \tkzDefCircle[orthogonal from=P](Q,I) 
  \tkzGetFirstPoint{E}
  \tkzDrawCircles(P,E Q,E)
  \tkzInterCC[common=E](P,E)(Q,E) \tkzGetFirstPoint{F}
  \tkzDefPointOnCircle[through =  center P angle 80 point E] 
  \tkzGetPoint{A}
  \tkzInterLC[common=E](A,E)(Q,E)  \tkzGetFirstPoint{C}
  \tkzInterLL(A,F)(C,Q)  \tkzGetPoint{D}
  \tkzDrawLines[add=0 and 1](P,Q)
  \tkzDrawLines[add=0 and 2](A,E)
  \tkzDrawSegments(P,E E,F F,C A,F C,D)
  \tkzDrawPoints(P,Q,E,F,A,C,D)
  \tkzLabelPoints(P,Q,F)
  \tkzLabelPoints[above](E,A)
  \tkzLabelPoints[left](D)
  \tkzLabelPoints[above right](C)
\end{tikzpicture}
\end{tkzexample}


\endinput

\section{角} 
\subsection{在 \tkzname{tkz-euclide} 中的定义与用法}
在欧几里得几何中，角是由两条射线（称为角的边）共享一个公共端点（称为角的顶点）所形成的图形 [Wikipedia]。在 \tkzname{tkz-euclide} 中，一条射线由两个点确定，因此角由三个点确定，如 $\widehat{AOB}$；顶点 $O$ 是第二个点。它们的顺序很重要，因为角默认按逆时针方向给出。
在三角学以及一般数学中，平面角约定以逆时针方向度量，从指向正右方（或东方）的 $0^\circ$ 开始，到指向正上方（或北方）的 $90^\circ$ [Wikipedia]。
我们约定：逆时针度量的角为正。

  \begin{center}
    \begin{tikzpicture}[scale=.75]
      \node{顺时针};
      \tkzDefPoint(0,0){O} \tkzDefPoint(90:2){A}\tkzDefPoint(180:2){B}   
      \tkzDrawArc[black,line width=2pt,arrows = {Stealth-}](O,B)(A)
    \end{tikzpicture}
    \begin{tikzpicture}[scale=.75]
          \node{逆时针};
      \tkzDefPoint(0,0){O} \tkzDefPoint(90:2){A}\tkzDefPoint(0:2){B}   
      \tkzDrawArc[black,line width=2pt,arrows = {-Stealth}](O,A)(B)
    \end{tikzpicture}  
  \end{center}
  
 角涉及多个宏命令，例如 \tkzcname{tkzDefPoint}、\tkzcname{tkzDefPointBy[rotation = \dots]}、\tkzcname{tkzDrawArc}
 以及下文中介绍的 \tkzcname{tkzGetAngle}。除最后一个外，所有这些宏命令都接受负角。
 
 \begin{figure}[!ht]
 \centering
 \begin{tabular}{|c|c|}
 \hline
 \tkzsubf{\begin{tikzpicture}
 \tkzDefPoint(0,0){O}    \tkzDefPoint(0:2){A}
 \tkzDefPointBy[rotation=center O angle 80](A)  \tkzGetPoint{B}
 \tkzDrawSegments[-{Stealth}](O,A O,B)
 \tkzMarkAngles[size=2,-Stealth,teal](A,O,B)
 \tkzFindAngle(A,O,B)   \tkzGetAngle{an}
 \tkzLabelAngle[pos=1,teal](A,O,B){$ \pgfmathprintnumber{\an}^\circ$}
 \tkzLabelPoints(A)  \tkzLabelPoints[above](B)
 \end{tikzpicture}}
      {从 $(O,A)$ 旋转 $80^\circ$ 到 $(O,B)$\\ 
    {\textbackslash}tkzDefPointBy[rotation=center O angle 80]}
 &
 \tkzsubf{\begin{tikzpicture}
 \tkzDefPoint(0,0){O}    \tkzDefPoint(0:2){A}
 \tkzDefPointBy[rotation=center O angle -80](A)  \tkzGetPoint{B}
 \tkzDrawSegments[-{Stealth}](O,A O,B)
 \tkzMarkAngles[size=2,Stealth-,red](B,O,A)
 \tkzFindAngle(B,O,A)   \tkzGetAngle{an}
 \tkzLabelAngle[pos=1,red](B,O,A){$-\pgfmathprintnumber{\an}^\circ$}
\tkzLabelPoints[right](A)  \tkzLabelPoints[below](B)
 \end{tikzpicture}}
  {从 $(O,A)$ 旋转 $-80^\circ$ 到 $(O,B)$\\ 
     {\textbackslash}tkzDefPointBy[rotation=center O angle -80]}
 \\ \hline
 \tkzsubf{\begin{tikzpicture}
 \tkzDefPoint(0,0){O}    \tkzDefPoint(0:2){A}
 \tkzDefPointBy[rotation=center O angle 80](A)  \tkzGetPoint{B}
 \tkzDrawSegments[-{Stealth}](O,A O,B)
 \tkzMarkAngles[size=1.5,-Stealth,teal](A,O,B)
 \tkzFindAngle(A,O,B)   \tkzGetAngle{an}
 \tkzLabelAngle[pos=1,teal](A,O,B){$ \pgfmathprintnumber{\an}^\circ$}
\tkzLabelPoints(A)  \tkzLabelPoints[above](B)
 \end{tikzpicture}}
      { {\textbackslash}tkzFindAngle(A,O,B) 给出 $80$}
 &
 \tkzsubf{\begin{tikzpicture}
 \tkzDefPoint(0,0){O}    \tkzDefPoint(0:2){A}
 \tkzDefPointBy[rotation=center O angle -80](A)  \tkzGetPoint{B}
 \tkzDrawSegments[-{Stealth}](O,A O,B)
 \tkzMarkAngles[size=1,-Stealth,red](A,O,B)
 \tkzFindAngle(A,O,B)   \tkzGetAngle{an}
 \tkzLabelAngle[pos=.75,red](A,O,B){$\pgfmathprintnumber{\an}^\circ$}
\tkzLabelPoints[right](A)  \tkzLabelPoints[below](B)
 \end{tikzpicture}}
  {{\textbackslash}tkzFindAngle(A,O,B) 给出 $\pgfmathprintnumber{\an}^\circ$}
 \\\hline
 \end{tabular}
 \end{figure}

如我们所见，$-80^\circ$ 旋转定义了一个顺时针角，但宏命令
\tkzcname{tkzFindAngle} 恢复出的是一个逆时针角。

\subsection{获取角度 \tkzcname{tkzGetAngle}}
\begin{NewMacroBox}{tkzGetAngle}{\parg{name of macro}}%
将一个角的度数赋值给某个宏。该值为正数，介于 $0^\circ$ 与 $360^\circ$ 之间。本宏命令获取 \tkzcname{tkzAngleResult} 并存入一个新宏。

\medskip

\begin{tabular}{lll}%
\toprule
参数             & 示例 & 说明             \\
\midrule
\TAline{宏名称} {\tkzcname{tkzGetAngle}\{ang\}}{\tkzcname{ang} 中保存角度值。}
\end{tabular}
\end{NewMacroBox}

这个辅助宏获取下一个宏命令 \tkzcname{tkzFindAngle} 的结果。

\subsection{由三个点构成的角}

\begin{NewMacroBox}{tkzFindAngle}{\parg{pt1,pt2,pt3}}%
结果保存在宏 \tkzcname{tkzAngleResult} 中。

\medskip

\begin{tabular}{lll}%
\toprule
参数     & 示例 & 说明     \\
\midrule
\TAline{(pt1,pt2,pt3)} {\tkzcname{tkzFindAngle}(A,B,C)}{\tkzcname{tkzAngleResult} 给出角（$\overrightarrow{BA},\overrightarrow{BC}$）}
\bottomrule
\end{tabular}

\medskip
角度大小恒为正，介于 $0^\circ$ 与 $360^\circ$ 之间。按照通常约定，小于平角的逆时针角大小介于 $0^\circ$ 与 $180^\circ$ 之间，而小于平角的顺时针角大小大于 $180^\circ$。\tkzcname{tkzGetAngle} 用于获取该角。
\end{NewMacroBox}
 
\subsubsection{角度度量的验证}
    
\begin{tkzexample}[latex=7cm,small]
\begin{tikzpicture}[scale=.75]
  \tkzDefPoint(-1,1){A}
  \tkzDefPoint(5,2){B}
  \tkzDefEquilateral(A,B)
  \tkzGetPoint{C}
  \tkzDrawPolygon(A,B,C)
  \tkzFindAngle(B,A,C) \tkzGetAngle{angleBAC}
  \edef\angleBAC{\fpeval{round(\angleBAC)}}
  \tkzDrawPoints(A,B,C) 
  \tkzLabelPoints(A,B)
  \tkzLabelPoint[right](C){$C$}
  \tkzLabelAngle(B,A,C){\angleBAC$^\circ$}
  \tkzMarkAngle[size=1.5](B,A,C)
\end{tikzpicture}
\end{tkzexample}

\subsubsection{确定三角形的三个角}

\begin{tkzexample}[latex=6cm,small]
\begin{tikzpicture}
\tikzset{label angle style/.append style={pos=1.4}}
\tkzDefPoints{0/0/a,5/3/b,3/6/c}
\tkzDrawPolygon(a,b,c)
\tkzFindAngle(c,b,a)\tkzGetAngle{angleCBA}
\pgfmathparse{round(1+\angleCBA)}
\let\angleCBA\pgfmathresult
\tkzFindAngle(a,c,b)\tkzGetAngle{angleACB}
\pgfmathparse{round(\angleACB)}
\let\angleACB\pgfmathresult
\tkzFindAngle(b,a,c)\tkzGetAngle{angleBAC}
\pgfmathparse{round(\angleBAC)}
\let\angleBAC\pgfmathresult
\tkzMarkAngle(c,b,a)
\tkzLabelAngle(c,b,a){\tiny $\angleCBA^\circ$}
\tkzMarkAngle(a,c,b)
\tkzLabelAngle(a,c,b){\tiny $\angleACB^\circ$}
\tkzMarkAngle(b,a,c)
\tkzLabelAngle(b,a,c){\tiny $\angleBAC^\circ$}
\end{tikzpicture}
\end{tkzexample}

\subsubsection{两圆之间的夹角}
我们求两圆在交点处由切线所构成的角。

\begin{tkzexample}[latex=7cm,small]
\begin{tikzpicture}[scale=.4]
\pgfkeys{/pgf/number format/.cd,%
          fixed,precision=1}
\tkzDefPoints{0/0/A,6/0/B,4/2/C}
\tkzDrawCircles(A,C B,C)
\tkzDefLine[tangent at=C](A) \tkzGetPoint{a}
\tkzDefPointsBy[symmetry = center C](a){d}
\tkzDefLine[tangent at=C](B) \tkzGetPoint{b}
\tkzDrawLines[add=1 and 4](a,C  C,b)
\tkzFillAngle[fill=teal,opacity=.2%
                        ,size=2](b,C,d)
\tkzFindAngle(b,C,d)\tkzGetAngle{bcd}
\tkzLabelAngle[pos=1.25](b,C,d){%
  \tiny $\pgfmathprintnumber{\bcd}^\circ$}
\end{tikzpicture}
\end{tkzexample}

\subsection{直线与水平轴的夹角 \tkzcname{tkzFindSlopeAngle}}
比上一个更有趣：结果介于 $-180^\circ$ 与 $+180^\circ$ 之间。

\begin{NewMacroBox}{tkzFindSlopeAngle}{\parg{A,B}}%
确定直线 $(AB)$ 的斜率（以角度表示）。结果保存在宏 \tkzcname{tkzAngleResult} 中。

\medskip
\begin{tabular}{lll}%
\toprule
参数  & 示例 & 说明     \\
\midrule
\TAline{(pt1,pt2)} {\tkzcname{tkzFindSlopeAngle}(A,B)}{}
\bottomrule
\end{tabular}

\medskip
\tkzcname{tkzGetAngle} 用于获取结果；若不需要保留，可直接使用 \tkzcname{tkzAngleResult}。
\end{NewMacroBox}

\subsubsection{如何使用 \tkzcname{tkzFindSlopeAngle}}

 这里 $(AB)$ 是 $\widehat{CAD}$ 的角平分线，$(AD)$ 的斜率为零。我们先求 $(AB)$ 的斜率，然后旋转两次。

\begin{tkzexample}[latex=7cm,small]
\begin{tikzpicture}
 \tkzDefPoint(1,5){A} \tkzDefPoint(5,2){B}  
 \tkzFindSlopeAngle(A,B)\tkzGetAngle{tkzang}
 \tkzDefPointBy[rotation= center A angle \tkzang ](B)
 \tkzGetPoint{C}
 \tkzDefPointBy[rotation= center A angle -\tkzang ](B) 
 \tkzGetPoint{D}
 \tkzDrawSegment(A,B)
 \tkzDrawSegments[new](A,C A,D)
 \tkzDrawPoints(A,B,C,D)
 \tkzCompass[length=1](A,C)
 \tkzCompass[delta=10,brown](B,C)  
 \tkzLabelPoints(B,C,D)  
 \tkzLabelPoints[above left](A)
\end{tikzpicture}
\end{tkzexample}

\subsubsection{\tkzcname{tkzFindSlopeAngle} 与 \tkzcname{tkzGetAngle} 的使用}
这是另一种构造垂直平分线（中垂线）的方法。

\begin{tkzexample}[latex=6cm,small]
\begin{tikzpicture}
 \tkzInit
 \tkzDefPoint(0,0){A}        \tkzDefPoint(3,2){B}
 \tkzDefLine[mediator](A,B)  \tkzGetPoints{I}{J}
 \tkzCalcLength(A,B)         \tkzGetLength{dAB}
 \tkzFindSlopeAngle(A,B)     \tkzGetAngle{tkzangle}
 \begin{scope}[rotate=\tkzangle]
   \tkzSetUpArc[color=gray,line width=0.2pt,%
     /tkzcompass/delta=10]
   \tkzDrawArc[R,arc](B,3/4*\dAB)(120,240)
   \tkzDrawArc[R,arc](A,3/4*\dAB)(-45,60)
   \tkzDrawLine(I,J)         \tkzDrawSegment(A,B)
  \end{scope}
  \tkzDrawPoints(A,B,I,J)    \tkzLabelPoints(A,B)
   \tkzLabelPoints[right](I,J)
\end{tikzpicture}
\end{tkzexample}

\subsubsection{\tkzcname{tkzFindSlopeAngle} 的另一用法}

\begin{tkzexample}[latex=7cm,small]
\begin{tikzpicture}[scale=1.5]
  \tkzDefPoint(1,2){A}    \tkzDefPoint(3,4){B}
  \tkzDefPoint(3,2){C}    \tkzDefPoint(3,1){D}
  \tkzDrawSegments(A,B A,C A,D)
  \tkzDrawPoints[color=red](A,B,C,D)  
  \tkzLabelPoints(A,B,C,D)
  \tkzFindSlopeAngle(A,B)\tkzGetAngle{SAB} 
  \tkzFindSlopeAngle(A,C)\tkzGetAngle{SAC}
  \tkzFindSlopeAngle(A,D)\tkzGetAngle{SAD}
  \pgfkeys{/pgf/number format/.cd,fixed,precision=2}
  \tkzText(1,5){The slope of (AB) is : 
     $\pgfmathprintnumber{\SAB}^\circ$}     
  \tkzText(1,4.5){The slope of (AC) is : 
     $\pgfmathprintnumber{\SAC}^\circ$}    
  \tkzText(1,4){The slope of (AD) is : 
     $\pgfmathprintnumber{\SAD}^\circ$}
\end{tikzpicture}
\end{tkzexample}

\endinput


\part{绘图、填充与裁剪}
\section{绘图}
\tkzname{\tkznameofpack} 可以绘制五种类型的对象：点、直线或线段、圆、圆弧和扇形。

%<---------------------------------------------------------------------------->
%    POINT(S)
%<---------------------------------------------------------------------------->
\subsection{绘制一个点或多个点}
有两种选择：用 \tkzcname{tkzDrawPoint} 绘制单个点，或用 \tkzcname{tkzDrawPoints} 绘制一个或多个点。

\subsubsection{绘制点 \tkzcname{tkzDrawPoint}} \hypertarget{tdrp}{}

\begin{NewMacroBox}{tkzDrawPoint}{\oarg{local options}\parg{name}}%
\begin{tabular}{lll}%
参数 &  默认值 & 定义                 \\
\midrule
\TAline{点名称} {无默认}  {只接受一个点名}
\bottomrule
\end{tabular}

\medskip
参数为必填。圆盘颜色取圆的颜色但更浅；所有样式都可以修改。点是一个节点，因此在图形缩放时不会受影响。

\medskip
\begin{tabular}{lll}%
\toprule
选项             & 默认值 & 定义 \\
\midrule
\TOline{\TIKZ\ 选项}{}{所有 \TIKZ\ 选项均有效。}
\TOline{shape}  {circle}{可选 \tkzname{cross} 或 \tkzname{cross out}}
\TOline{size}   {6}{$6 \times$ \tkzcname{pgflinewidth}}
\TOline{color}  {black}{可修改默认颜色 }
\bottomrule
\end{tabular}

\medskip
{我们还可以创建其他形状，例如 \tkzname{cross}}
\end{NewMacroBox}

默认情况下，\tkzname{point style } 定义如下：

\begin{tkzltxexample}[]
  \tikzset{point style/.style = {%
           draw         = black,
           inner sep    = 0pt,
           shape        = circle,
           minimum size = 3 pt,
           fill         = black
                               }
        } 
\end{tkzltxexample}

\subsubsection{点绘制示例}
注意 \tkzname{scale} 不影响点的形状，这是正常的。大多数情况下，一种点的形状就足够了；可以在开头定义，用宏命令或在配置文件中定义。

\begin{tkzexample}[latex=5cm,small]
  \begin{tikzpicture}[scale=.5]
   \tkzDefPoint(1,3){A}
   \tkzDefPoint(4,1){B}
   \tkzDefPoint(0,0){O}
   \tkzDrawPoint[color=red](A)
   \tkzDrawPoint[fill=blue!20,draw=blue](B)
   \tkzDrawPoint[shape=cross,size=8pt,color=teal](O)
  \end{tikzpicture}
\end{tkzexample}

你可以一次绘制多个点，但这个宏命令比前一个稍慢，并且所有点都使用相同的选项。
\newpage
\hypertarget{tdrps}{}
\begin{NewMacroBox}{tkzDrawPoints}{\oarg{local options}\parg{liste}}%
\begin{tabular}{lll}%
参数 &  默认值  & 定义 \\
\midrule
\TAline{点列表}{无默认}{例如 \tkzcname{tkzDrawPoints(A,B,C)}}
\bottomrule
\end{tabular}

\medskip
\begin{tabular}{lll}%
选项             & 默认值 & 定义 \\
\midrule
\TOline{shape}  {circle}{可选 \tkzname{cross} 或 \tkzname{cross out}}
\TOline{size}  {6}{$6 \times$ \tkzcname{pgflinewidth}}
\TOline{color}  {black}{可修改默认颜色 }
\bottomrule
\end{tabular}

\medskip
\tkzHandBomb\ 注意结尾的 \code{s}：漏掉它会在绘制多个点时引发连锁错误。选项与前一个宏命令相同。
\end{NewMacroBox}

\subsubsection{带大小、颜色和填充的点}

\begin{tkzexample}[latex=7cm,small]
\begin{tikzpicture}
\tkzDefPoints{1/3/A,4/1/B,0/0/C}
\tkzDrawPoints[size=3,color=red,fill=red!50](A,B,C)
\end{tikzpicture}
\end{tkzexample}
%<---------------------------------------------------------------------------->
%    LINE(S)
%<---------------------------------------------------------------------------->
\section{绘制直线}
以下宏命令用于绘制和标注直线。
\subsection{绘制直线}
要绘制直线，只需给出一对点。你可以使用选项 \tkzname{add} 来延长直线（此选项由 \tkzimp{Mark Wibrow} 提供，参见下面的代码）。

直线的默认样式为：

\begin{tkzltxexample}[]
  \tikzset{line style/.style = {%
    line width = 0.6pt,
    color      = black,
    style      = solid,
    add        = {.2} and  {.2}%
   }}
\end{tkzltxexample}
其中
   
\begin{tkzltxexample}[]
  \tikzset{%
    add/.style args={#1 and #2}{
        to path={%
 ($(\tikztostart)!-#1!(\tikztotarget)$)--($(\tikztotarget)!-#2!(\tikztostart)$)%
  \tikztonodes}}}
\end{tkzltxexample}

你可以用 \tkzcname{tkzSetUpLine} 修改此样式，参见 \ref{tkzsetupline}。

\newpage
\begin{NewMacroBox}{tkzDrawLine}{\oarg{local options}\parg{pt1,pt2} }%
参数是由两个或三个点组成的列表。对于射线，你可以用 \tkzname{add} 创建样式。

\begin{tabular}{lll}%
\toprule
选项             & 默认值 & 定义                         \\ 
\midrule
\TOline{\TIKZ\ 选项}{}{所有 \TIKZ\ 选项均有效。}
\TOline{add}{0.2 and 0.2}{add = $kl$ and $kr$，……}
\TOline{\dots}{\dots}{允许线段向左右两侧延长。 }
\bottomrule
\end{tabular}

\tkzname{add} 定义经过点 pt1 和 pt2 的直线长度。两个数字均为百分比。绘图时可使用 \TIKZ\ 样式。
\end{NewMacroBox}

\subsubsection{使用 \tkzname{add} 的示例}
\begin{tkzexample}[latex=5cm,small]
\begin{tikzpicture}
 \tkzInit[xmin=-2,xmax=3,ymin=-2.25,ymax=2.25]
 \tkzClip[space=.25]
 \tkzDefPoint(0,0){A} \tkzDefPoint(2,0.5){B}
 \tkzDefPoint(0,-1){C}\tkzDefPoint(2,-0.5){D} 
 \tkzDefPoint(0,1){E} \tkzDefPoint(2,1.5){F} 
 \tkzDefPoint(0,-2){G} \tkzDefPoint(2,-1.5){H}
 \tkzDrawLine(A,B)    \tkzDrawLine[add = 0 and .5](C,D) 
 \tkzDrawLine[add = 1 and 0](E,F)
 \tkzDrawLine[add = 0 and 0](G,H) 
 \tkzDrawPoints(A,B,C,D,E,F,G,H)    
 \tkzLabelPoints(A,B,C,D,E,F,G,H)  
\end{tikzpicture}
\end{tkzexample} 

你可以用相同的选项一次绘制多条直线。
\begin{NewMacroBox}{tkzDrawLines}{\oarg{local options}\parg{pt1,pt2 pt3,pt4 ...}}% 
参数为由空格分隔的点对列表。绘图时可使用 \TIKZ\ 样式。
\end{NewMacroBox}      

\subsubsection{使用 \tkzcname{tkzDrawLines} 的示例}    

\begin{tkzexample}[latex=8cm,small]
\begin{tikzpicture}
  \tkzDefPoint(0,0){A}
  \tkzDefPoint(2,0){B}
  \tkzDefPoint(1,2){C}
  \tkzDefPoint(3,2){D}   
  \tkzDrawLines(A,B C,D A,C B,D)
  \tkzLabelPoints(A,B,C,D)
\end{tikzpicture}
\end{tkzexample}
%<---------------------------------------------------------------------------->
%    SEGMENT(S)
%<---------------------------------------------------------------------------->
\section{绘制线段}
当然，我们也有绘制线段的宏命令。

\subsection{绘制线段 \tkzcname{tkzDrawSegment}} 
\begin{NewMacroBox}{tkzDrawSegment}{\oarg{local options}\parg{pt1,pt2}}%
参数是由两个点组成的列表。绘图时可使用 \TIKZ\ 样式。
 
\medskip
\begin{tabular}{lll}%
参数    & 示例 & 定义    \\
\midrule
\TAline{(pt1,pt2)}{(A,B)}{绘制线段 $[A,B]$}
\bottomrule 
\end{tabular}
 
\medskip
\begin{tabular}{lll}%
选项    & 示例 & 定义    \\
\midrule
\TOline{\TIKZ\ 选项}{}{所有 \TIKZ\ 选项均有效。}
\TOline{dim}{无默认}{dim = \{label,dim,option\}，……}
\TOline{\dots}{\dots}{允许给图形添加尺寸标注。}
\bottomrule 
\end{tabular}

这当然等价于 \tkzcname{draw (A)--(B);}。你也可以使用选项 \tkzname{add}。
\end{NewMacroBox}

\subsubsection{使用点引用的示例}     

\begin{tkzexample}[latex=6cm,small]
\begin{tikzpicture}[scale=1.5]
  \tkzDefPoint(0,0){A}
  \tkzDefPoint(2,1){B}
  \tkzDrawSegment[color=red,thin](A,B)
  \tkzDrawPoints(A,B)    
  \tkzLabelPoints(A,B)  
\end{tikzpicture}
\end{tkzexample}

\subsubsection{用选项 \tkzname{add} 延长线段} 

\begin{tkzexample}[latex=7cm,small]
\begin{tikzpicture}
  \tkzDefPoints{0/0/A,6/0/B,0.8/4/C}
  \tkzDefTriangleCenter[euler](A,B,C) 
  \tkzGetPoint{E}
  \tkzDefCircle[euler](A,B,C)\tkzGetPoints{E}{e}
  \tkzDrawCircle[red](E,e)
  \tkzDrawLines[add=.5 and .5](A,B A,C B,C)
  \tkzDrawPoints(A,B,C,E)
  \tkzLabelPoints(A,B,C,E)
  \end{tikzpicture}
\end{tkzexample}

\subsubsection{用选项 \tkzname{dim} 添加尺寸标注（代码由 Muzimuzhi Z 提供）} 
这段代码来自 tex.stackexchange.com 上一个问题的回答
(change-color-and-style-of-dimension-lines-in-tkz-euclide )。
\tkzname{dim} 的代码基于 \TIKZ\ 选项，因此你必须给出单位。
现在你可以用多种方式使用两种样式 |dim style| 和 |dim fence style|。
查看示例以了解用这些样式可以做什么。

\begin{verbatim}
   \tikzset{dim style/.append style={dashed}} % append if you want to keep precedent style.
   or 
   \begin{scope}[ dim style/.append style={orange},
       dim fence style/.style={dashed}]
\end{verbatim}

\begin{tkzexample}[latex=7cm]
\begin{tikzpicture}[scale=.75]
  \tkzDefPoints{0/3/A, 1/-3/B}
  \tkzDrawPoints(A,B)
  \tkzDrawSegment[dim={\(l_0\),1cm,right=2mm}, 
    dim style/.append style={red, 
    dash pattern={on 2pt off 2pt}}](A,B)
  \tkzDrawSegment[dim={\(l_1\),2cm,right=2mm}, 
    dim style/.append style={blue}](A,B)
  \begin{scope}[ dim style/.style={orange},
      dim fence style/.style={dashed}]
    \tkzDrawSegment[dim={\(l_2\),3cm,right=2mm}](A,B)  
    \tkzDrawSegment[dim={\(l_3\),-2cm,right=2mm}](A,B)   
  \end{scope}  
  \tkzLabelPoints[left](A,B)
\end{tikzpicture}
\end{tkzexample}

\subsubsection{用选项 \tkzname{dim} 添加尺寸标注：第一部分} 
\begin{tkzexample}[latex=7cm,small]
\begin{tikzpicture}[scale=2]
\pgfkeys{/pgf/number format/.cd,fixed,precision=2}
\tkzDefPoint(0,0){A}
\tkzDefPoint(3.07,0){B}
\tkzInterCC[R](A,2.37)(B,1.82)
\tkzGetPoints{C}{C'}
\tkzDefCircle[in](A,B,C) \tkzGetPoints{G}{g}
\tkzDrawCircle(G,g)
\tkzDrawPolygon(A,B,C)
\tkzDrawPoints(A,B,C)
\tkzCalcLength(A,B)\tkzGetLength{ABl}
\tkzCalcLength(B,C)\tkzGetLength{BCl}
\tkzCalcLength(A,C)\tkzGetLength{ACl}
\begin{scope}[dim style/.style={dashed,sloped,teal}]
  \tkzDrawSegment[dim={\pgfmathprintnumber\BCl,6pt,%
                                      text=red}](C,B)
  \tkzDrawSegment[dim={\pgfmathprintnumber\ACl,%
                                        6pt,}](A,C)
  \tkzDrawSegment[dim={\pgfmathprintnumber\ABl,%
                                      -6pt,}](A,B)
\end{scope}
\tkzLabelPoints(A,B) \tkzLabelPoints[above](C)
\end{tikzpicture}
\end{tkzexample}

\subsubsection{用选项 \tkzname{dim} 添加尺寸标注：第二部分} 
\begin{tkzexample}[latex=6cm,small]
\begin{tikzpicture}[scale=.5]
  \tkzDefPoints{0/0/O,-2/0/A,2/0/B,
                -2/4/C,2/4/D,2/-4/E,-2/-4/F}
  \tkzDrawPolygon(C,...,F)
  \tkzDrawSegments(A,B)
  \tkzDrawPoints(A,...,F,O)
  \tkzLabelPoints[below left](A,...,F,O)
  \tkzDrawSegment[dim={ $\sqrt{5}$,2cm,}](C,E)
  \tkzDrawSegment[dim={ $\frac{\sqrt{5}}{2}$,1cm,}](O,E)
  \tkzDrawSegment[dim={ $2$,2cm,left=8pt}](F,C)
  \tkzDrawSegment[dim={ $1$,1cm,left=8pt}](F,A)
\end{tikzpicture}
\end{tkzexample}

\subsection{绘制线段 \tkzcname{tkzDrawSegments}} 
如果选项相同，可以用一个宏命令绘制多条线段。

\begin{NewMacroBox}{tkzDrawSegments}{\oarg{local options}\parg{pt1,pt2 pt3,pt4 ...}}%
参数是由点对组成的列表。绘图时可使用 \TIKZ\ 样式。
\end{NewMacroBox}

\begin{tkzexample}[latex=6cm,small]
\begin{tikzpicture}
  \tkzInit[xmin=-1,xmax=3,ymin=-1,ymax=2]
  \tkzClip[space=1]
  \tkzDefPoint(0,0){A}
  \tkzDefPoint(2,1){B} 
  \tkzDefPoint(3,0){C} 
  \tkzDrawSegments(A,B B,C)
  \tkzDrawPoints(A,B,C)    
  \tkzLabelPoints(A,C) 
  \tkzLabelPoints[above](B)  
\end{tikzpicture}
\end{tkzexample}

\subsubsection{在线段上放置箭头}
\begin{tkzexample}[latex=6cm,small]
\begin{tikzpicture}
\tkzSetUpStyle[postaction=decorate,
    decoration={markings, 
    mark=at position .5 with {\arrow[thick]{#1}}
      }]{myarrow}
  \tkzDefPoint(0,0){A}
  \tkzDefPoint(4,-4){B}
  \tkzDrawSegments[myarrow=stealth](A,B)
  \tkzDrawPoints(A,B) 
\end{tikzpicture}
\end{tkzexample}

\subsection{绘制三角形的边}

\subsubsection{如何绘制高} 
\begin{tkzexample}[latex=7cm,small]
  \begin{tikzpicture}[rotate=-90]
  \tkzDefPoint(0,1){A}
  \tkzDefPoint(2,4){C}
  \tkzDefPointWith[orthogonal normed,K=7](C,A)
  \tkzGetPoint{B}
  \tkzDefSpcTriangle[orthic,name=H](A,B,C){a,b,c}
  \tkzDrawLine[dashed,color=magenta](C,Hc)
  \tkzDrawSegment[green!60!black](A,C)
  \tkzDrawSegment[green!60!black](C,B)
  \tkzDrawSegment[green!60!black](B,A)
  \tkzLabelPoint[left](A){$A$}
  \tkzLabelPoint[right](B){$B$}
  \tkzLabelPoint[above](C){$C$}
  \tkzLabelPoint[left](Hc){$Hc$}
  \tkzLabelSegment[auto](B,A){$c$}
  \tkzLabelSegment[auto,swap](B,C){$a$}
  \tkzLabelSegment[auto,swap](C,A){$b$}
  \tkzMarkAngle[size=1,color=cyan,mark=|](C,B,A)
  \tkzMarkAngle[size=1,color=cyan,mark=|](A,C,Hc)
  \tkzMarkAngle[size=0.75,
                color=orange,mark=||](Hc,C,B)
  \tkzMarkAngle[size=0.75,
                color=orange,mark=||](B,A,C)
  \tkzMarkRightAngle(A,C,B)
  \tkzMarkRightAngle(B,Hc,C)
  \end{tikzpicture} 
\end{tkzexample}

\subsection{绘制多边形} 
 \begin{NewMacroBox}{tkzDrawPolygon}{\oarg{local options}\parg{points list}}%
只需给出点列表：宏命令用所给的 \TIKZ\ 选项绘制多边形。你可以用 $(A,...,E)$ 代替 $(A,B,C,D,E)$，用 $(P_1,P...,P_5)$ 代替 $(P_1,P_2,P_3,P_4,P_5)$。

\begin{tabular}{lll}%
\toprule
参数             & 示例 & 说明                         \\
\midrule
\TAline{\parg{pt1,pt2,pt3,...}}{|\BS tkzDrawPolygon[gray,dashed](A,B,C)|}{绘制三角形}
\end{tabular}

\medskip
\begin{tabular}{lll}%
\toprule
选项             & 默认值 & 示例                         \\
\midrule
\TOline{Options TikZ}{...}{|\BS tkzDrawPolygon[red,line width=2pt](A,B,C)|}
 \end{tabular} 
\end{NewMacroBox}

\subsubsection{\tkzcname{tkzDrawPolygon}}

\begin{tkzexample}[latex=7cm, small]  
\begin{tikzpicture} [rotate=18,scale=1]
 \tkzDefPoints{0/0/A,2.25/0.2/B,2.5/2.75/C,-0.75/2/D}
 \tkzDrawPolygon(A,B,C,D)
 \tkzDrawSegments[style=dashed](A,C B,D) 
\end{tikzpicture}
\end{tkzexample}

\subsubsection{选项 \tkzname{two angles}}
\begin{tkzexample}[latex=6 cm,small]
\begin{tikzpicture}
\tkzDefPoint(0,0){A} 
\tkzDefPoint(6,0){B} 
\tkzDefTriangle[two angles = 50 and 70](A,B) \tkzGetPoint{C}
\tkzDrawPolygon(A,B,C)
\tkzLabelAngle[pos=1.4](B,A,C){$50^\circ$}
\tkzLabelAngle[pos=0.8](C,B,A){$70^\circ$}
\end{tikzpicture}
\end{tkzexample}

\subsubsection{直线样式}
\begin{tkzexample}[latex=8 cm,small]
\begin{tikzpicture}[scale=.6]
\tkzSetUpLine[line width=5mm,color=teal]
\tkzDefPoint(0,0){O}
\foreach \i in {0,...,5}{%
 \tkzDefPoint({30+60*\i}:4){p\i}}
\tkzDefMidPoint(p1,p3) \tkzGetPoint{m1}
\tkzDefMidPoint(p3,p5) \tkzGetPoint{m3}
\tkzDefMidPoint(p5,p1) \tkzGetPoint{m5}
\tkzDrawPolygon[line join=round](p1,p3,p5)
\tkzDrawPolygon[teal!80,
line join=round](p0,p2,p4)
\tkzDrawSegments(m1,p3 m3,p5 m5,p1)
\tkzDefCircle[R](O,4.8)\tkzGetPoint{o}
\tkzDrawCircle[teal](O,o)
\end{tikzpicture}
\end{tkzexample}

\subsection{绘制折线} 
 \begin{NewMacroBox}{tkzDrawPolySeg}{\oarg{local options}\parg{points list}}%
只需给出点列表：宏命令用所给的 \TIKZ\ 选项绘制折线。

\begin{tabular}{lll}%
\toprule
参数             & 示例 & 说明                         \\
\midrule
\TAline{\parg{pt1,pt2,pt3,...}}{|\BS tkzDrawPolySeg[gray,dashed](A,B,C)|}{绘制三角形}
\end{tabular}

\medskip
\begin{tabular}{lll}%
\toprule
选项             & 默认值 & 示例                         \\
\midrule
\TOline{Options TikZ}{...}{|\BS tkzDrawPolySeg[red,line width=2pt](A,B,C)|}
 \end{tabular} 
\end{NewMacroBox}

\subsubsection{折线}

\begin{tkzexample}[latex=7cm, small]  
\begin{tikzpicture}
 \tkzDefPoints{0/0/A,6/0/B,3/4/C,2/2/D}          
 \tkzDrawPolySeg(A,...,D)
 \tkzDrawPoints(A,...,D)
\end{tikzpicture}
\end{tkzexample}

\subsubsection{在半圆内作两个正方形。}
一道算额（Sangaku）！问题是证明可以在半圆内作两个正方形，并根据半径确定它们的边长。

\begin{tkzexample}[latex=7 cm,small]
\begin{tikzpicture}[scale=.75] 
  \tkzDefPoints{0/0/A,8/0/B,4/0/I}
  \tkzDefSquare(A,B)    \tkzGetPoints{C}{D} 
  \tkzInterLC(I,C)(I,B) \tkzGetPoints{E'}{E} 
  \tkzInterLC(I,D)(I,B) \tkzGetPoints{F'}{F} 
  \tkzDefPointsBy[projection=onto A--B](E,F){H,G} 
  \tkzDefPointsBy[symmetry = center H](I){J} 
  \tkzDefSquare(H,J)     \tkzGetPoints{K}{L} 
  \tkzDrawSector(I,B)(A) 
  \tkzDrawPolySeg(H,E,F,G) 
  \tkzDrawPolySeg(J,K,L) 
  \tkzDrawPoints(E,G,H,F,J,K,L)
\end{tikzpicture}
\end{tkzexample}

\subsubsection{折线：索引记法}

\begin{tkzexample}[latex=7cm, small]  
\begin{tikzpicture}
\foreach \pt in {1,2,...,8} {%
\tkzDefPoint(\pt*20:3){P_\pt}}     
\tkzDrawPolySeg(P_1,P_...,P_8)
\tkzDrawPoints(P_1,P_...,P_8)
\end{tikzpicture}
\end{tkzexample}
%<---------------------------------------------------------------------------->
%    CIRCLE
%<---------------------------------------------------------------------------->
\section{用 \tkzcname{tkzDrawCircle} 绘制圆}

\subsection{绘制一个圆}
\begin{NewMacroBox}{tkzDrawCircle}{\oarg{local options}\parg{A,B}}%
\tkzHandBomb\ 注意只需两个点即可确定半径。选项 \tkzname{R} 允许你直接以度量值给出半径。

\medskip
\begin{tabular}{lll}%
\toprule
参数           & 示例 & 说明                         \\
\midrule
\TAline{\parg{pt1,pt2}}{\parg{A,B}} {以 A 为圆心过 B}
 \bottomrule
\end{tabular}   

\medskip
当然，绘图时你可以添加所有 \TIKZ\ 样式。
\end{NewMacroBox}
 
 \subsubsection{圆与样式：绘制圆并给圆盘上色}
 我们将看到，可以在绘制圆的同时给圆盘上色。
 
\begin{tkzexample}[latex=7cm,small]
\begin{tikzpicture}
  \tkzDefPoint(0,0){O} 
  \tkzDefPoint(3,0){A}
 % circle with center O and passing through A
  \tkzDrawCircle(O,A) 
 % diameter circle $[OA]$
 \tkzDefCircle[diameter](O,A) \tkzGetPoint{I}
 \tkzDrawCircle[new,fill=orange!10,opacity=.5](I,A)
 % circle with center O and radius = exp(1) cm
  \edef\rayon{\fpeval{0.25*exp(1)}}
  \tkzDefCircle[R](O,\rayon) \tkzGetPoint{o}
   \tkzDrawCircle[color=orange](O,o) 
\end{tikzpicture} 
\end{tkzexample}  

\subsection{绘制多个圆}  
\begin{NewMacroBox}{tkzDrawCircles}{\oarg{local options}\parg{A,B C,D \dots}}%
\tkzHandBomb\ 注意参数是点对列表。可绘制的圆与前一个宏命令相同。选项 \tkzname{R} 允许你直接以度量值给出半径。

\medskip
\begin{tabular}{lll}%
\toprule
参数           & 示例 & 说明                         \\
\midrule
\TAline{\parg{pt1,pt2 pt3,pt4 ...}}{\parg{A,B C,D}} {两点列表}
\bottomrule
\end{tabular}   

\medskip
\begin{tabular}{lll}%
\toprule
选项             & 默认值 & 定义                         \\ 
\midrule
\TOline{through}{through}{由两点确定半径的圆}
 \bottomrule
\end{tabular}

\medskip
你无需使用默认选项 \tkzname{through}。
当然，绘图时你可以添加所有 \TIKZ\ 样式。
\end{NewMacroBox}

 \subsubsection{由三角形确定的圆。} 
 
\begin{tkzexample}[latex=9cm,small]
\begin{tikzpicture}
  \tkzDefPoints{0/0/A,2/0/B,3/2/C}
  \tkzDrawPolygon(A,B,C)
  \tkzDrawCircles(A,B B,C C,A)
  \tkzDrawPoints(A,B,C)
  \tkzLabelPoints(A,B,C) 
\end{tikzpicture} 
\end{tkzexample}

\subsubsection{同心圆。} 
 
\begin{tkzexample}[latex=7cm,small]
\begin{tikzpicture}
   \tkzDefPoints{0/0/A,1/0/a,2/0/b,3/0/c}
   \tkzDrawCircles(A,a A,b A,c)
   \tkzDrawPoint(A)
   \tkzLabelPoints(A)
\end{tikzpicture}
\end{tkzexample}

\subsubsection{旁切圆。} 

\begin{tkzexample}[latex=8cm,small] 
\begin{tikzpicture}[scale=1] 
\tkzDefPoints{0/0/A,4/0/B,1/2.5/C}
\tkzDrawPolygon(A,B,C)
\tkzDefCircle[ex](B,C,A) 
\tkzGetPoint{J_c} \tkzGetSecondPoint{T_c}
\tkzDrawCircle(J_c,T_c)
\tkzDrawLines[add=0 and 1](C,A C,B)
\tkzDrawSegment(J_c,T_c)
\tkzMarkRightAngle(J_c,T_c,B)
\tkzDrawPoints(A,B,C,J_c,T_c)
\end{tikzpicture}
\end{tkzexample}
 
\subsubsection{心形线}  
基于 O. Reboux 的构想，由 D. Rodriguez 用 pst-eucl（一个 PSTricks 模块）实现。

 它的名字源自希腊语 \textit{kardia}（心脏），因其形状而得名；由 Johan Castillon 命名（Wikipedia）。     
 
\begin{tkzexample}[latex=7cm,small]
\begin{tikzpicture}[scale=.5]
  \tkzDefPoint(0,0){O} 
  \tkzDefPoint(2,0){A}
  \foreach \ang in {5,10,...,360}{%
     \tkzDefPoint(\ang:2){M}
     \tkzDrawCircle(M,A) 
   }  
\end{tikzpicture} 
\end{tkzexample}

\newpage

\subsection{绘制半圆}
\begin{NewMacroBox}{tkzDrawSemiCircle}{\oarg{local options}\parg{O,A}}%

\medskip
\begin{tabular}{lll}%
\toprule
参数           & 示例 & 说明                         \\
\midrule
\TAline{\parg{pt1,pt2}}{\parg{O,A}} {$OA$ 为半径}
\bottomrule
\end{tabular} 
    
$O$ 为圆心，$A$ 为半圆的一个端点。
\end{NewMacroBox}  

\subsubsection{\tkzcname{tkzDrawSemiCircle} 的用法}   

\begin{tkzexample}[latex=7cm,small]
\begin{tikzpicture}
   \tkzDefPoint(0,0){A} \tkzDefPoint(6,0){B}
   \tkzDefMidPoint(A,B)  \tkzGetPoint{O}
   \tkzDrawSemiCircle[blue](O,B)
   \tkzDrawSemiCircle[red](O,A)
   \tkzDrawPoints(O,A,B)
   \tkzLabelPoints[below right](O,A,B)
 \end{tikzpicture}
\end{tkzexample}

\subsection{绘制多个半圆}

\begin{NewMacroBox}{tkzDrawSemiCircles}{\oarg{local options}\parg{A,B C,D \dots}}%

\medskip
\begin{tabular}{lll}%
\toprule
参数           & 示例 & 说明                         \\
\midrule
\TAline{\parg{pt1,pt2 pt3,pt4 ...}}{\parg{A,B C,D}} {两点列表}
\bottomrule
\end{tabular} 
    
\end{NewMacroBox}  

\subsubsection{\tkzcname{tkzDrawSemiCircles} 的用法：黄金 arbelos}  

\begin{tkzexample}[vbox,small]
\begin{tikzpicture}[scale=.75]
\tkzDefPoints{0/0/A,10/0/B}
\tkzDefGoldenRatio(A,B) \tkzGetPoint{C}
\tkzDefMidPoint(A,B)                     \tkzGetPoint{O_0}
\tkzDefMidPoint(A,C)                     \tkzGetPoint{O_1}
\tkzDefMidPoint(C,B)                     \tkzGetPoint{O_2}
\tkzLabelPoints(A,B,C)
\tkzDrawSegment(A,B)
\tkzDrawPoints(A,B,C)
\begin{scope}[local bounding box = graph]
  \tkzDrawSemiCircles[color=black](O_0,B)
\end{scope}
\useasboundingbox (graph.south west) rectangle (graph.north east);
\tkzClipCircle[out](O_1,C)\tkzClipCircle[out](O_2,B)
\tkzDrawSemiCircles[draw=none,fill=teal!15](O_0,B)
\tkzDrawSemiCircles[color=black](O_1,C O_2,B)
\end{tikzpicture}
\end{tkzexample}

%<---------------------------------------------------------------------------->
%    Ellipse
%<---------------------------------------------------------------------------->
\section{用 \tkzcname{tkzDrawEllipse} 绘制椭圆}

\subsection{绘制椭圆}
\begin{NewMacroBox}{tkzDrawEllipse}{\oarg{local options}\parg{C,a,b,An}}%


\medskip
\begin{tabular}{lll}%
\toprule
参数           & 示例 & 说明                         \\
\midrule
\TAline{\parg{C,a,b,An}}{\parg{C,4,2,45}} {C 为圆心；4 和 2 为半轴长} \\
 & & 45 为长轴斜率  \\
 \bottomrule
\end{tabular}  
 
\medskip
当然，绘图时你可以添加所有 \TIKZ\ 样式。
\end{NewMacroBox}

\subsubsection{绘制椭圆示例}
\begin{tkzexample}[latex=6cm,small]
   \begin{tikzpicture}[scale=.75]
      \tkzDefPoint(0,4){C}
      \tkzDrawEllipse[blue](C,4,2,45)
      \tkzLabelPoints(C)
   \end{tikzpicture}
\end{tkzexample}

%<---------------------------------------------------------------------------->
%    ARC
%<---------------------------------------------------------------------------->
\section{绘制圆弧} 
\subsection{宏命令 \tkzcname{tkzDrawArc}}
\begin{NewMacroBox}{tkzDrawArc}{\oarg{local options}\parg{O,\dots}\parg{\dots}}%
本宏命令绘制以 $O$ 为圆心的圆弧；参数取决于所选项。我们需要一个起点和一个终点。可以直接给起点（最简单），也可以给圆弧半径；后一种情况需要两个角度。角度可以直接给出，也可以由点与圆心共同确定。角度以度为单位。

\medskip
\begin{tabular}{lll}%
\toprule
选项             & 默认值 & 定义                        \\ 
\midrule
\TOline{towards}{towards}{$O$ 为圆心，弧从 $A$ 到 $(OB)$} 
\TOline{rotate} {towards}{弧从 $A$ 开始，角度决定其长度} 
\TOline{R}{towards}{给出半径和两个角度} 
\TOline{R with nodes}{towards}{给出半径和两个点}
\TOline{angles}{towards}{给出半径和两个点}
\TOline{delta}{0}{两侧各添加的角度 }
\TOline{reverse}{false}{反转弧的路径，便于反转箭头} 
\bottomrule
\end{tabular}

\medskip
当然，绘图时你可以添加所有 \TIKZ\ 样式。

\medskip

\begin{tabular}{lll}%
\toprule
选项             & 参数 & 示例                         \\ 
\midrule
\TOline{towards}{\parg{pt,pt}\parg{pt}}{\tkzcname{tkzDrawArc[delta=10](O,A)(B)}} 
\TOline{rotate} {\parg{pt,pt}\parg{an}}{\tkzcname{tkzDrawArc[rotate,color=red](O,A)(90)}}
\TOline{R}{\parg{pt,$r$}\parg{an,an}}{\tkzcname{tkzDrawArc[R](O,2)(30,90)}}
\TOline{R with nodes}{\parg{pt,$r$}\parg{pt,pt}}{\tkzcname{tkzDrawArc[R with nodes](O,2)(A,B)}}
\TOline{angles}{\parg{pt,pt}\parg{an,an}}{\tkzcname{tkzDrawArc[angles](O,A)(0,90)}}
\end{tabular}
\end{NewMacroBox}

下面是几个示例： 

\subsubsection{选项 \tkzname{towards}：从 $A$ 到 $B$ 的弧}
你不需要写 \tkzname{towards}。在第一个示例中，弧从 $A$ 开始到 $B$ 结束。从 $B$ 到 $A$ 的弧不同：弧总是按逆时针绘制。
\begin{tkzexample}[latex=6cm,small]
\begin{tikzpicture}[scale=.75]
  \tkzDefPoint(0,0){O}
  \tkzDefPoint(2,-1){A}
  \tkzDefPointBy[rotation= center O angle 90](A)
  \tkzGetPoint{B}
  \tkzDrawArc[color=orange,<->](O,A)(B) 
  \tkzDrawArc(O,B)(A)
  \tkzDrawLines[add = 0 and .5](O,A O,B)
  \tkzDrawPoints(O,A,B)
  \tkzLabelPoints[below](O,A,B)  
\end{tikzpicture}
\end{tkzexample}

\subsubsection{选项 \tkzname{towards}：终止于直线上的弧}
在这个示例中，弧从 $A$ 开始但终止于直线 $(OB)$。
 
\begin{tkzexample}[latex=6cm,small]
\begin{tikzpicture}[scale=0.75] 
  \tkzDefPoint(0,0){O}
  \tkzDefPoint(2,-1){A}
  \tkzDefPoint(1,1){B} 
  \tkzDrawArc[color=blue,->](O,A)(B)
  \tkzDrawArc[color=gray](O,B)(A)
  \tkzDrawArc(O,B)(A)
  \tkzDrawLines[add = 0 and .5](O,A O,B) 
  \tkzDrawPoints(O,A,B)
  \tkzLabelPoints[below](O,A,B)  
\end{tikzpicture}
\end{tkzexample}

\subsubsection{选项 \tkzname{rotate}}
\begin{tkzexample}[latex=6cm,small] 
\begin{tikzpicture}[scale=0.75] 
  \tkzDefPoint(0,0){O}
  \tkzDefPoint(2,-2){A}
  \tkzDefPoint(60:2){B}
  \tkzDrawLines[add = 0 and .5](O,A O,B)
  \tkzDrawArc[rotate,color=red](O,A)(180)
  \tkzDrawPoints(O,A,B)
  \tkzLabelPoints[below](O,A,B) 
\end{tikzpicture}
\end{tkzexample} 

\subsubsection{选项 \tkzname{R}} 
\begin{tkzexample}[latex=6cm,small]   
\begin{tikzpicture}[scale=0.75] 
  \tkzDefPoints{0/0/O}
  \tkzSetUpCompass[<->]
  \tkzDrawArc[R,color=teal,double](O,3)(270,360)
  \tkzDrawArc[R,color=orange,double](O,2)(0,270) 
  \tkzDrawPoint(O)
  \tkzLabelPoint[below](O){$O$}  
\end{tikzpicture} 
\end{tkzexample}

\subsubsection{选项 \tkzname{R with nodes}} 
\begin{tkzexample}[latex=6cm,small]
\begin{tikzpicture}[scale=0.75] 
  \tkzDefPoint(0,0){O}
  \tkzDefPoint(2,-1){A}
  \tkzDefPoint(1,1){B}
  \tkzCalcLength(B,A)\tkzGetLength{radius}
  \tkzDrawArc[R with nodes](B,\radius)(A,O)
\end{tikzpicture}
\end{tkzexample}

\subsubsection{选项 \tkzname{delta}}
与 \tkzcname{tkzCompass} 一样，此选项在两侧延长弧；\tkzname{delta} 是以度为单位的角度。也可以用 \tkzcname{tkzSetUpArc} 为所有弧设置该值。

\begin{tkzexample}[latex=7cm,small] 
\begin{tikzpicture} 
 \tkzDefPoint(0,0){A}
 \tkzDefPoint(3,0){B}
 \tkzDefPointBy[rotation= center A angle 60](B)
 \tkzGetPoint{C} 
 \begin{scope}% style only local
   \tkzDefPointBy[symmetry= center C](A)
   \tkzGetPoint{D} 
   \tkzDrawSegments(A,B A,D)
   \tkzDrawLine(B,D)
   \tkzSetUpCompass[color=orange]
   \tkzDrawArc[orange,delta=10](A,B)(C)
   \tkzDrawArc[orange,delta=10](B,C)(A)
   \tkzDrawArc[orange,delta=10](C,D)(D)
 \end{scope}

 \tkzDrawPoints(A,B,C,D)
 \tkzLabelPoints[below right](A,B,C,D)
 \tkzMarkRightAngle(D,B,A)
\end{tikzpicture}
\end{tkzexample} 

\subsubsection{选项 \tkzname{angles}：示例 1}

\begin{tkzexample}[latex=6cm,small]
\begin{tikzpicture}[scale=.75]
  \tkzDefPoint(0,0){A}
  \tkzDefPoint(5,0){B}  
  \tkzDefPoint(2.5,0){O} 
  \tkzDefPointBy[rotation=center O angle 60](B)
  \tkzGetPoint{D}
  \tkzDefPointBy[symmetry=center D](O)
  \tkzGetPoint{E}
  \begin{scope}
    \tkzDrawArc[angles](O,B)(0,180)
    \tkzDrawArc[angles,](B,O)(100,180)  
    \tkzCompass[delta=20](D,E) 
    \tkzDrawLines(A,B O,E B,E)
    \tkzDrawPoints(A,B,O,D,E)
  \end{scope}
  \tkzLabelPoints[below right](A,B,O,D,E)
  \tkzMarkRightAngle(O,B,E) 
\end{tikzpicture} 
\end{tkzexample}

\subsubsection{选项 \tkzname{angles}：示例 2}

\begin{tkzexample}[latex=7cm,small]
  \begin{tikzpicture}
   \tkzDefPoint(0,0){O}
   \tkzDefPoint(5,0){I} 
   \tkzDefPoint(0,5){J}
   \tkzInterCC(O,I)(I,O)\tkzGetPoints{B}{C}  
   \tkzInterCC(O,I)(J,O)\tkzGetPoints{D}{A}
   \tkzInterCC(I,O)(J,O)\tkzGetPoints{L}{K}
   \tkzDrawArc[angles](O,I)(0,90)
   \tkzDrawArc[angles,color=gray,
               style=dashed](I,O)(90,180)
   \tkzDrawArc[angles,color=gray,
               style=dashed](J,O)(-90,0)
   \tkzDrawPoints(A,B,K)
   \foreach \point in {I,A,B,J,K}{%
               \tkzDrawSegment(O,\point)} 
  \end{tikzpicture} 
\end{tkzexample}

\subsubsection{选项 \tkzname{reverse}：反转箭头}

\begin{tkzexample}[latex=6cm,small]
  \begin{tikzpicture}
    \tkzDefPoints{0/0/O,3/0/U}
    \tkzDefPoint(10:1){A}
    \tkzDefPoint(90:1){B}
    \tkzLabelPoints(A,B)
    \tkzDrawArc[reverse,tkz arrow={Stealth}](O,A)(B)
    \tkzDrawPoints(A,B,O)
  \end{tikzpicture}
\end{tkzexample}
%<---------------------------------------------------------------------------->
%    SECTOR
%<---------------------------------------------------------------------------->
\section{绘制扇形}
\subsection{\tkzcname{tkzDrawSector}} 
\tkzHandBomb\ 注意参数取决于所选项。
\begin{NewMacroBox}{tkzDrawSector}{\oarg{local options}\parg{O,\dots}\parg{\dots}}%
\begin{tabular}{SlSlSl}%
选项             & 默认值 & 定义                         \\ 
\midrule
\TOline{towards}{towards}{$O$ 为圆心，弧从 $A$ 到 $(OB)$}
\TOline{rotate} {towards}{弧从 $A$ 开始，角度决定其长度 } 
\TOline{R}{towards}{给出半径和两个角度}
\TOline{R with nodes}{towards}{给出半径和两个点}

\end{tabular} 

\medskip
\emph{当然，绘图时你可以添加所有 \TIKZ\ 样式。}

\begin{tabular}{lll}%

选项             & 参数 & 示例                         \\ 
\midrule
\TOline{towards}{\parg{pt,pt}\parg{pt}}{\tkzcname{tkzDrawSector(O,A)(B)}}
\TOline{rotate} {\parg{pt,pt}\parg{an}}{\tkzcname{tkzDrawSector[rotate,color=red](O,A)(90)}} 
\TOline{R}{\parg{pt,$r$}\parg{an,an}}{\tkzcname{tkzDrawSector[R,color=teal](O,2)(30,90)}}
\TOline{R with nodes}{\parg{pt,$r$}\parg{pt,pt}}{\tkzcname{tkzDrawSector[R with nodes](O,2)(A,B)}}
\end{tabular}
\end{NewMacroBox}

下面是几个示例： 

\subsubsection{\tkzcname{tkzDrawSector} 与 \tkzname{towards}} 
你不需要写 \tkzname{towards}。你可以使用 \tkzname{fill} 作为选项。

\begin{tkzexample}[latex=7cm,small]
\begin{tikzpicture}
  \tkzDefPoint(0,0){O}
  \tkzDefPoint(-30:1){A} 
  \tkzDefPointBy[rotation = center O angle -60](A) 
  \tkzDrawSector[teal](O,A)(tkzPointResult)
 \begin{scope}[shift={(-60:1)}]
  \tkzDefPoint(0,0){O}
  \tkzDefPoint(-30:1){A} 
  \tkzDefPointBy[rotation = center O angle -60](A) 
  \tkzDrawSector[red](O,tkzPointResult)(A)
  \end{scope}
\end{tikzpicture}   
\end{tkzexample}

\subsubsection{\tkzcname{tkzDrawSector} 与 \tkzname{rotate}}  
\begin{tkzexample}[latex=7cm,small]  
\begin{tikzpicture}[scale=2]
 \tkzDefPoints{0/0/O,2/2/A,2/1/B}
 \tkzDrawSector[rotate,orange](O,A)(20)
 \tkzDrawSector[rotate,teal](O,B)(-20)
\end{tikzpicture} 
\end{tkzexample}  

\subsubsection{\tkzcname{tkzDrawSector} 与 \tkzname{R}}  
\begin{tkzexample}[latex=7cm,small]
\begin{tikzpicture}[scale=1.25]
 \tkzDefPoint(0,0){O}
 \tkzDefPoint(2,-1){A}
 \tkzDrawSector[R](O,1)(30,90)
 \tkzDrawSector[R](O,1)(90,180)
 \tkzDrawSector[R](O,1)(180,270)
 \tkzDrawSector[R](O,1)(270,360) 
\end{tikzpicture}
\end{tkzexample}

\subsubsection{\tkzcname{tkzDrawSector} 与 \tkzname{R with nodes}：两个扇形}
在这个示例中我使用了选项 \tkzname{fill}，但 \tkzcname{tkzFillSector} 同样可行。
\begin{tkzexample}[latex=7cm,small]
\begin{tikzpicture}[scale=1.25]
 \tkzDefPoint(0,0){O}
 \tkzDefPoint(4,-2){A}
 \tkzDefPoint(4,1){B}
 \tkzDefPoint(3,3){C}
 \tkzDrawSector[R with nodes,%
                fill=teal!20](O,1)(B,C)
 \tkzDrawSector[R with nodes,%
                fill=orange!20](O,1.25)(A,B)  
\tkzDrawSegments(O,A O,B O,C)
\tkzDrawPoints(O,A,B,C) 
\tkzLabelPoints(A,B,C) 
\tkzLabelPoints[left](O) 
\end{tikzpicture}
\end{tkzexample}

\subsubsection{\tkzcname{tkzDrawSector} 与 \tkzname{R with nodes}：六边形中的扇形}
\begin{tkzexample}[latex=6cm,small]
\begin{tikzpicture} [scale=.4]
 \tkzDefPoints{-1/-2/A,1/3/B}
 \tkzDefRegPolygon[side,sides=6](A,B) 
 \tkzGetPoint{O} 
 \tkzDrawPolygon[fill=black!10, draw=blue](P1,P...,P6) 
 \tkzLabelRegPolygon[sep=1.05](O){A,...,F}
 \tkzDrawCircle[dashed](O,A)
 \tkzLabelSegment[above,sloped,
                  midway](A,B){\(A B = 16m\)}
 \foreach \i  [count=\xi from 1]  in {2,...,6,1}
   {%
    \tkzDefMidPoint(P\xi,P\i)
    \path (O) to [pos=1.1] node {\xi} (tkzPointResult) ;
    }
  \tkzDefRandPointOn[segment = P3--P5] 
  \tkzGetPoint{S}
  \tkzDrawSegments[thick,dashed,red](A,S S,B)
  \tkzDrawPoints(P1,P...,P6,S)
  \tkzLabelPoint[left,above](S){$S$}
  \tkzDrawSector[R with nodes,fill=red!20](S,2)(A,B)
  \tkzLabelAngle[pos=1.5](A,S,B){$\alpha$}
\end{tikzpicture}
\end{tkzexample}

\endinput

\section{填充}
\subsection{给圆盘上色}
以前可以用宏命令 \tkzcname{tkzDrawCircle} 实现，但必须同时绘制圆；现在不再需要这样做。
  
\begin{NewMacroBox}{tkzFillCircle}{\oarg{local options}\parg{A,B}}%
\begin{tabular}{lll}%
选项             & 默认值 & 定义                         \\ 
\midrule
\TOline{radius}  {radius}{两点确定半径}
\TOline{R} {radius}{一个点和半径的度量值 }
\bottomrule
\end{tabular}

\medskip
你不需要写 \tkzname{radius}，因为它是默认选项。当然，你可以添加所有 \TIKZ\ 样式。
\end{NewMacroBox}  


\subsubsection{阴阳} 
\begin{tkzexample}[latex=8cm,small]
  \begin{tikzpicture}[scale=.75]
    \tkzDefPoint(0,0){O}
    \tkzDefPoint(-4,0){A}
    \tkzDefPoint(4,0){B}
    \tkzDefPoint(-2,0){I}
    \tkzDefPoint(2,0){J}
    \tkzDrawSector[fill=teal](O,A)(B)
    \tkzFillCircle[fill=white](J,B) 
    \tkzFillCircle[fill=teal](I,A) 
    \tkzDrawCircle(O,A) 
  \end{tikzpicture}
\end{tkzexample}


\subsubsection{取自一道算额} 

\begin{tkzexample}[latex=7cm,small]
\begin{tikzpicture}
   \tkzDefPoint(0,0){B}  \tkzDefPoint(6,0){C}%
   \tkzDefSquare(B,C)    \tkzGetPoints{D}{A} 
   \tkzClipPolygon(B,C,D,A) 
   \tkzDefMidPoint(A,D)  \tkzGetPoint{F}
   \tkzDefMidPoint(B,C)  \tkzGetPoint{E}
   \tkzDefMidPoint(B,D)  \tkzGetPoint{Q}           
   \tkzDefLine[tangent from = B](F,A) 
   \tkzGetPoints{H}{G} 
   \tkzInterLL(F,G)(C,D) \tkzGetPoint{J}
   \tkzInterLL(A,J)(F,E) \tkzGetPoint{K}
   \tkzDefPointBy[projection=onto B--A](K)   
   \tkzGetPoint{M}  
   \tkzDrawPolygon(A,B,C,D)
   \tkzFillCircle[red!20](E,B)
   \tkzFillCircle[blue!20](M,A)
   \tkzFillCircle[green!20](K,Q)
  \tkzDrawCircles(B,A M,A E,B K,Q) 
\end{tikzpicture}
\end{tkzexample} 

\subsubsection{裁剪与填充 第一部分} 
\begin{tkzexample}[latex=7cm,small]
\begin{tikzpicture}
\tkzDefPoints{0/0/A,4/0/B,2/2/O,3/4/X,4/1/Y,1/0/Z,
              0/3/W,3/0/R,4/3/S,1/4/T,0/1/U}
\tkzDefSquare(A,B)\tkzGetPoints{C}{D}
\tkzDefPointWith[colinear normed=at X,K=1](O,X)
 \tkzGetPoint{F}
\begin{scope}
  \tkzFillCircle[fill=teal!20](O,F)
  \tkzFillPolygon[white](A,...,D)
  \tkzClipPolygon(A,...,D)
  \foreach \c/\t in {S/C,R/B,U/A,T/D}
  {\tkzFillCircle[teal!20](\c,\t)}
\end{scope}
\foreach \c/\t in {X/C,Y/B,Z/A,W/D}
{\tkzFillCircle[white](\c,\t)}
  \foreach \c/\t in {S/C,R/B,U/A,T/D}
  {\tkzFillCircle[teal!20](\c,\t)}
\end{tikzpicture}
\end{tkzexample}

\subsubsection{裁剪与填充 第二部分} 
\begin{tkzexample}[latex=7cm, small]  
\begin{tikzpicture}[scale=.75]
\tkzDefPoints{0/0/A,8/0/B,8/8/C,0/8/D}
\tkzDefMidPoint(A,B) \tkzGetPoint{F}
\tkzDefMidPoint(B,C) \tkzGetPoint{E}
\tkzDefMidPoint(D,B) \tkzGetPoint{I}
\tkzDefMidPoint(I,B) \tkzGetPoint{a}
\tkzInterLC(B,I)(B,C) \tkzGetSecondPoint{K}
\tkzDefMidPoint(I,K) \tkzGetPoint{b}
\begin{scope}
  \tkzFillSector[fill=blue!10](B,C)(A)
  \tkzDefMidPoint(A,B) \tkzGetPoint{x}
  \tkzDrawSemiCircle[fill=white](x,B)
  \tkzDefMidPoint(B,C) \tkzGetPoint{y}
  \tkzDrawSemiCircle[fill=white](y,C)
  \tkzClipCircle(E,B)
  \tkzClipCircle(F,B)
  \tkzFillCircle[fill=blue!10](B,A)
\end{scope}
\tkzDrawSemiCircle[thick](F,B)
\tkzDrawSemiCircle[thick](E,C)
\tkzDrawArc[thick](B,C)(A)
\tkzDrawSegments[thick](A,B B,C)
\tkzDrawPoints(A,B,C,E,F)
\tkzLabelPoints[centered](a,b)
\tkzLabelPoints(A,B,C,E,F)
\end{tikzpicture}
\end{tkzexample}

\subsubsection{裁剪与填充 第三部分} 

\begin{tkzexample}[latex=7cm, small] 
\begin{tikzpicture}
  \tkzDefPoint(0,0){A} \tkzDefPoint(1,0){B}
  \tkzDefPoint(2,0){C} \tkzDefPoint(-3,0){a} 
  \tkzDefPoint(3,0){b}  \tkzDefPoint(0,3){c} 
  \tkzDefPoint(0,-3){d}
\begin{scope}
 \tkzClipPolygon(a,b,c,d)
 \tkzFillCircle[teal!20](A,C)
\end{scope}
 \tkzFillCircle[white](A,B)
 \tkzDrawCircle[color=red](A,C)
 \tkzDrawCircle[color=red](A,B)
\end{tikzpicture}
\end{tkzexample}

\subsection{给多边形上色} 
 \begin{NewMacroBox}{tkzFillPolygon}{\oarg{local options}\parg{points list}}%
你可以在绘制多边形的同时给它上色；本宏命令给多边形内部上色而不绘制其边框。

\medskip
\begin{tabular}{lll}%
\toprule
参数                & 示例 & 说明                         \\ 
\midrule
\TAline{\parg{pt1,pt2,\dots}}{\parg{A,B,\dots}}{}
%\bottomrule
 \end{tabular}
\end{NewMacroBox} 

\subsubsection{\tkzcname{tkzFillPolygon}} 
\begin{tkzexample}[latex=7cm, small]  
\begin{tikzpicture}[scale=.5]
   \tkzDefPoint(0,0){C} \tkzDefPoint(4,0){A}
   \tkzDefPoint(0,3){B}
   \tkzDefSquare(B,A)      \tkzGetPoints{E}{F}
   \tkzDefSquare(A,C)      \tkzGetPoints{G}{H}
   \tkzDefSquare(C,B)       \tkzGetPoints{I}{J}
   \tkzFillPolygon[color  =  orange!30   ](A,C,G,H)
   \tkzFillPolygon[color  =  teal!40  ](C,B,I,J)
   \tkzFillPolygon[color  =  purple!20](B,A,E,F)
   \tkzDrawPolygon[line width  =  1pt](A,B,C)
   \tkzDrawPolygon[line width  =  1pt](A,C,G,H)
   \tkzDrawPolygon[line width  =  1pt](C,B,I,J)
   \tkzDrawPolygon[line width  =  1pt](B,A,E,F)
   \tkzLabelSegment[above](C,A){$a$}
   \tkzLabelSegment[right](B,C){$b$}
   \tkzLabelSegment[below left](B,A){$c$}
\end{tikzpicture}
\end{tkzexample}

\subsection{\tkzcname{tkzFillSector}}
\tkzHandBomb\ 注意参数取决于所选项。
\begin{NewMacroBox}{tkzFillSector}{\oarg{local options}\parg{O,\dots}\parg{\dots}}%
\begin{tabular}{lll}%
选项          & 默认值 & 定义      \\ 
\midrule
\TOline{towards}{towards}{$O$ 为圆心，弧从 $A$ 到 $(OB)$}
\TOline{rotate} {towards}{弧从 A 开始，角度决定其长度 } 
\TOline{R}{towards}{给出半径和两个角度}
\TOline{R with nodes}{towards}{给出半径和两个点}
\bottomrule
\end{tabular} 

\medskip
当然，绘图时你可以添加所有 \TIKZ\ 样式。

\medskip
\begin{tabular}{lll}%
\toprule
选项             & 参数 & 示例                         \\ 
\midrule
\TOline{towards}{\parg{pt,pt}\parg{pt}}{\tkzcname{tkzFillSector(O,A)(B)}}
\TOline{rotate} {\parg{pt,pt}\parg{an}}{\tkzcname{tkzFillSector[rotate,color=red](O,A)(90)}}
\TOline{R}{\parg{pt,$r$}\parg{an,an}}{\tkzcname{tkzFillSector[R,color=blue](O,2)(30,90)}} 
\TOline{R with nodes}{\parg{pt,$r$}\parg{pt,pt}}{\tkzcname{tkzFillSector[R with nodes](O,2)(A,B)}}
\end{tabular}   
\end{NewMacroBox} 

\subsubsection{\tkzcname{tkzFillSector} 与 \tkzname{towards}} 
你不需要写 \tkzname{towards}。注意轮廓不会被绘制：只有面被上色。
\begin{tkzexample}[latex=5.75cm,small]
  \begin{tikzpicture}[scale=.6] 
  \tkzDefPoint(0,0){O}
  \tkzDefPoint(-30:3){A}
  \tkzDefPointBy[rotation = center O angle -60](A)   
  \tkzFillSector[fill=purple!20](O,A)(tkzPointResult) 
    \begin{scope}[shift={(-60:1)}]
    \tkzDefPoint(0,0){O}
    \tkzDefPoint(-30:3){A}
    \tkzDefPointBy[rotation = center O angle -60](A)
    \tkzGetPoint{A'}
    \tkzFillSector[color=teal!40](O,A')(A)
      \end{scope}
  \end{tikzpicture}
\end{tkzexample}


\subsubsection{\tkzcname{tkzFillSector} 与 \tkzname{rotate}} 
\begin{tkzexample}[latex=5.75cm,small]
\begin{tikzpicture}[scale=1.5]
 \tkzDefPoint(0,0){O} \tkzDefPoint(2,2){A}
 \tkzFillSector[rotate,color=purple!20](O,A)(30)
 \tkzFillSector[rotate,color=teal!40](O,A)(-30)
\end{tikzpicture}    
\end{tkzexample} 

\subsection{给角上色：\tkzcname{tkzFillAngle}}

最简单的情况：
\begin{NewMacroBox}{tkzFillAngle}{\oarg{local options}\parg{A,O,B}}%
$O$ 是角的顶点，$OA$ 和 $OB$ 是它的两边。注意角取决于点的顺序。

\medskip

\begin{tabular}{lll}%
\toprule
选项             & 默认值 & 定义                        \\ 
\midrule
\TOline{size}{1}{此选项决定被上色角扇形的半径。}

\bottomrule
\end{tabular} 

\medskip
当然，你可以添加所有 \TIKZ\ 样式，例如填充和渐变。
\end{NewMacroBox}  

\subsubsection{使用 \tkzname{size} 的示例}  
\begin{tkzexample}[latex=7cm,small]
\begin{tikzpicture} 
   \tkzInit 
   \tkzDefPoints{0/0/O,2.5/0/A,1.5/2/B}
   \tkzFillAngle[size=2, fill=gray!10](A,O,B)
   \tkzDrawLines(O,A O,B)
   \tkzDrawPoints(O,A,B)
\end{tikzpicture}
\end{tkzexample}


\subsubsection{改变点的顺序} 
\begin{tkzexample}[latex=7cm,small]
\begin{tikzpicture} 
   \tkzInit 
   \tkzDefPoints{0/0/O,2.5/0/A,1.5/2/B}
   \tkzFillAngle[size=2,fill=gray!10](B,O,A)
   \tkzDrawLines(O,A O,B)
   \tkzDrawPoints(O,A,B)
\end{tikzpicture}
\end{tkzexample}

\begin{tkzexample}[latex=7cm,small]
\begin{tikzpicture} 
   \tkzInit 
   \tkzDefPoints{0/0/O,5/0/A,3/4/B}
   % Don't forget {} to get, () to use
   \tkzFillAngle[size=4,left color=white, 
                 right color=red!50](A,O,B)
   \tkzDrawLines(O,A O,B)
   \tkzDrawPoints(O,A,B)
\end{tikzpicture}
\end{tkzexample}

\begin{NewMacroBox}{tkzFillAngles}{\oarg{local options}\parg{A,O,B}\parg{A',O',B'}etc.}%
对于具有相同选项的多个角，有一个对应的宏命令。
  \end{NewMacroBox}  
  
\subsubsection{多个角}  
\begin{tkzexample}[latex=5cm,small]
\begin{tikzpicture}[scale=0.5]
  \tkzDefPoints{0/0/B,8/0/C,0/8/A,8/8/D}
  \tkzDrawPolygon(B,C,D,A)
  \tkzDefTriangle[equilateral](B,C) \tkzGetPoint{M}
  \tkzInterLL(D,M)(A,B) \tkzGetPoint{N}
  \tkzDefPointBy[rotation=center N angle -60](D) 
  \tkzGetPoint{L}
  \tkzInterLL(N,L)(M,B)     \tkzGetPoint{P}
  \tkzInterLL(M,C)(D,L)     \tkzGetPoint{Q}
  \tkzDrawSegments(D,N N,L L,D B,M M,C)
  \tkzDrawPoints(L,N,P,Q,M,A,D) 
  \tkzLabelPoints[left](N,P,Q)
  \tkzLabelPoints[above](M,A,D)
  \tkzLabelPoints(L,B,C)
  \tkzMarkAngles(C,B,M B,M,C M,C,B D,L,N L,N,D N,D,L)
  \tkzFillAngles[fill=red!20,opacity=.2](C,B,M
      B,M,C M,C,B D,L,N L,N,D N,D,L)
\end{tikzpicture}
\end{tkzexample} 
\endinput

\section{控制边界框}
引自 \tkzimp{PgfManual}：

"当添加 clip 选项时，当前路径将用于裁剪后续的图形。裁剪绝不会扩大裁剪区域。因此，当你针对某一路径裁剪后再针对另一路径裁剪时，实际裁剪区域是两者的交集。
唯一能扩大裁剪路径的方法是结束进行裁剪的 {pgfscope}。在 {pgfscope} 结束时，进入该作用域时生效的裁剪路径会被重新恢复。"


首先，使用 \TIKZ\ 时你无需处理边界框的大小。早期版本的 \tkzNamePack{tkz-euclide} 并不控制它；自第 4 版起，边界框的大小受到了限制。

使用宏 \tkzcname{tkzInit} 后，初始边界框由以 $(0,0)$ 和 $(10,10)$ 为对角顶点的矩形定义。宏 \tkzcname{tkzInit} 可以通过参数 \tkzname{xmin}、\tkzname{xmax}、\tkzname{ymin} 和 \tkzname{ymax} 修改这个初始边界框。当然，任何画在外部的图形都会修改由 \TIKZ\ 维护的边界框。你可以通过 \TIKZ\ 命令或选项来影响这一行为，例如 \tkzcname{useasboundingbox} 或选项 \tkzname{use as bounding box}。这样你可以为图形预留一个框，但图形可能溢出该框并延伸到正文中。
命令 \tkzcname{pgfresetboundingbox} 会清除边界框并重新开始一个新的边界框。

\subsection{\tkzcname{tkzInit} 的作用}
 然而，有时需要控制所显示内容的大小。
 为此，你需要先准备好将要使用的边界框：这正是宏 \tkzNameMacro{tkzInit} 的作用。对于某些图形，固定极值（xmin、xmax、ymin 和 ymax）并对该矩形执行 \code{clip}，以便尽可能精确地控制图形大小，是很有用的。

有两个宏可用于控制边界框：
\begin{itemize}
   \item \tkzcname{tkzInit}
   \item \tkzcname{tkzClip}
\end{itemize}
\vspace{20pt}

我还添加了一些与边界框直接相关的宏：你可以显示它或对其进行裁剪（见下文）。

\subsection{\tkzcname{tkzInit}}

\begin{NewMacroBox}{tkzInit}{\oarg{local options}}\hypertarget{init}{}%
\begin{tabular}{lll}%    
选项  & 默认值 & 说明             \\
\midrule    
\TOline{xmin} {0} {横坐标的最小值，单位为 cm}
\TOline{xmax} {10} {横坐标的最大值，单位为 cm}
\TOline{xstep}{1} {$x$ 方向两个刻度之间的间距}
\TOline{ymin} {0} {$y$ 轴的最小值，单位为 cm}
\TOline{ymax} {10} {$y$ 轴的最大值，单位为 cm}
\TOline{ystep}{1} {$y$ 方向两个刻度之间的间距}  
\bottomrule    
\end{tabular}

\medskip 

\tkzcname{tkzInit} 的作用是定义一个 \textcolor{red}{正交}坐标系和平面上的一个矩形区域，你可以在该区域内用笛卡尔坐标放置图形。
该宏定义了你的工作区域，就像在图形计算器上一样。自 \tkzname{\tkznameofpack} 第 4 版起，\tkzcname{xstep} 和 \tkzcname{ystep} 始终为 1，因此 \tkzcname{tkzInit} 主要用于 \code{clipping out} 之类的操作。
\end{NewMacroBox}


\subsection{\tkzcname{tkzClip}}

\begin{NewMacroBox}{tkzClip}{\oarg{local options}}
该宏会隐藏位于由 (xmin~;~ymin) 与 (xmax~;~ymax) 所定义矩形之外的所有内容。

\medskip
\begin{tabular}{lll}
\hline
选项  & 默认值 & 说明             \\
\midrule
\TOline{space} {1} {在背景矩形的右、左、下、上四个方向上增加的数值}
\bottomrule
\end{tabular}

\medskip

选项 \tkzname{space} 会扩大图形的可见部分，使其变为由 (xmin-space~;~ymin-space) 与 (xmax+space~;~ymax+space) 所定义的矩形。\tkzname{space} 可以为负值！单位为 cm，且无需写出。
\end{NewMacroBox}



该宏对初始矩形进行 \code{clip}，使得只有位于该矩形内部的路径才会被绘制。

\begin{tkzexample}[latex=8cm,small]
\begin{tikzpicture}
 \tkzInit[xmax=4, ymax=3]
 \tkzDefPoints{-1/-1/A,5/2/B}
 \tkzDrawX \tkzDrawY 
 \tkzGrid
 \tkzClip
 \tkzDrawSegment(A,B)
\end{tikzpicture}
\end{tkzexample} 

你可以增加一些空间：
\begin{tkzltxexample}[]
  \tkzClip[space=1]
\end{tkzltxexample} 

\subsection{\tkzcname{tkzClip} 与选项 \tkzname{space}} 
该选项在 \code{clipped} 矩形周围增加一些空间。
\begin{tkzexample}[latex=8cm,small]
\begin{tikzpicture}
 \tkzInit[xmax=4, ymax=3]
 \tkzDefPoints{-1/-1/A,5/2/B}
 \tkzDrawX \tkzDrawY 
 \tkzGrid
 \tkzClip[space=1]
 \tkzDrawSegment(A,B)
\end{tikzpicture}
\end{tkzexample}   
\code{clipped} 矩形的尺寸为 \tkzname{xmin-1}、\tkzname{ymin-1}、\tkzname{xmax+1} 和 \tkzname{ymax+1}。 

%<--------------------------------------------------------------------------->
%              tkzShowBB
%<--------------------------------------------------------------------------->
\subsection{\tkzcname{tkzShowBB}}
最简单的宏。
\begin{NewMacroBox}{tkzShowBB}{\oarg{local options}}% 
该宏以矩形框的形式显示边界框。它接受 \TIKZ\ 选项。
\end{NewMacroBox} 


\subsubsection{\tkzcname{tkzShowBB} 示例}
\begin{tkzexample}[latex=8cm,small]
\begin{tikzpicture}[scale=.5]
  \tkzInit[ymax=5,xmax=8]
  \tkzGrid  
  \tkzDefPoint(3,0){A}
   \begin{scope}
    \tkzClipBB
    \tkzDefCircle[R](A,5) \tkzGetPoint{a}
    \tkzDrawCircle(A,a)
    \tkzShowBB[line width = 4pt,fill=teal!10,%
              opacity=.4]
   \end{scope}
\tkzDefCircle[R](A,4) \tkzGetPoint{b}
\tkzDrawCircle[red](A,b)
\end{tikzpicture}
\end{tkzexample}
%<--------------------------------------------------------------------------->
%         tkzClipBB
%<--------------------------------------------------------------------------->
\subsection{\tkzcname{tkzClipBB}}
\begin{NewMacroBox}{tkzClipBB}{}%
其思想是将后续的构造限制在当前边界框之内。
\end{NewMacroBox}

\subsubsection{\tkzcname{tkzClipBB} 与角平分线示例}

\begin{tkzexample}[latex=6cm,small]
  \begin{tikzpicture}
  \tkzInit[xmin=-3,xmax=6, ymin=-1,ymax=6]
  \tkzDefPoint(0,0){O}\tkzDefPoint(3,1){I}
  \tkzDefPoint(1,4){J}
  \tkzDefLine[bisector](I,O,J) \tkzGetPoint{i}
  \tkzDefLine[bisector out](I,O,J) \tkzGetPoint{j}
  \tkzDrawPoints(O,I,J,i,j)
  \tkzClipBB
  \tkzDrawLines[add = 1 and 2,color=orange](O,I O,J)
  \tkzDrawLines[add = 1 and 2](O,i O,j)
  \tkzShowBB
  \end{tikzpicture}
\end{tkzexample}


\newpage

\section{裁剪不同对象}

\subsection{裁剪多边形} 
 \begin{NewMacroBox}{tkzClipPolygon}{\oarg{local options}\parg{points list}}%
该宏将图形限制在给定的多边形之内。

\medskip
\begin{tabular}{lll}%
\toprule
参数       & 示例 & 说明     \\ 
\midrule
\TAline{\parg{pt1,pt2,pt3,\dots}}{\parg{A,B,C}}{}
\midrule
选项  & 默认值 & 说明             \\
\midrule    
\TOline{out} {} {裁剪对象的外部}
 \end{tabular}
\end{NewMacroBox}

\subsubsection{\tkzcname{tkzClipPolygon}}

\begin{tkzexample}[latex=7cm,small]
  \begin{tikzpicture}[scale=1.25] 
  \tkzDefPoint(0,0){A} 
  \tkzDefPoint(4,0){B} 
  \tkzDefPoint(1,3){C} 
  \tkzDrawPolygon(A,B,C) 
  \tkzDefPoint(0,2){D} 
  \tkzDefPoint(2,0){E} 
  \tkzDrawPoints(D,E) 
  \tkzLabelPoints(D,E) 
  \tkzClipPolygon(A,B,C) 
  \tkzDrawLine[new](D,E)
\end{tikzpicture}
\end{tkzexample}

\subsubsection{\tkzcname{tkzClipPolygon[out]}}

\begin{tkzexample}[latex=7cm,small]
  \begin{tikzpicture}[scale=1]
  \tkzDefPoint(0,0){P1}
  \tkzDefPoint(4,0){P2}
  \tkzDefPoint(4,4){P3}
  \tkzDefPoint(0,4){P4}
  \tkzDefPoint(1,1){Q1}
  \tkzDefPoint(3,1){Q2}
  \tkzDefPoint(3,3){Q3}
  \tkzDefPoint(1,3){Q4}
  \tkzDrawPolygon(P1,P2,P3,P4)
  \begin{scope}
     \tkzClipPolygon[out](Q1,Q2,Q3,Q4)
    \tkzFillPolygon[teal!20](P1,P2,P3,P4)
  \end{scope}
  \tkzDrawPolygon(Q1,Q2,Q3,Q4)
  \end{tikzpicture}
\end{tkzexample}

\subsubsection{示例：在正方形中使用 \code{Clip} 绘制 Sangaku} 
\begin{tkzexample}[latex=7cm, small]  
\begin{tikzpicture}[scale=.75]
 \tkzDefPoint(0,0){A} \tkzDefPoint(8,0){B}
 \tkzDefSquare(A,B)   \tkzGetPoints{C}{D}
 \tkzDefPoint(4,8){F}
 \tkzDefTriangle[equilateral](C,D) 
 \tkzGetPoint{I}
 \tkzDefPointBy[projection=onto B--C](I) 
 \tkzGetPoint{J}
 \tkzInterLL(D,B)(I,J)  \tkzGetPoint{K}
 \tkzDefPointBy[symmetry=center K](B)  
 \tkzGetPoint{M}
 \tkzClipPolygon(B,C,D,A)
 \tkzFillPolygon[color = orange](A,B,C,D)
 \tkzFillCircle[color = yellow](M,I)
 \tkzFillCircle[color = blue!50!black](F,D)
\end{tikzpicture}
\end{tkzexample}

\subsection{裁剪圆盘}

\begin{NewMacroBox}{tkzClipCircle}{\oarg{local options}\parg{A,B}}%
\begin{tabular}{lll}%
\toprule
参数           & 示例 & 说明                         \\
\midrule
\TAline{\parg{A,B}}{\parg{A,B}} {AB 为半径}
\bottomrule
\end{tabular}  
 
\medskip
\begin{tabular}{lll}%
选项             & 默认值 & 说明                         \\ 
\midrule
\TOline{out} {} {裁剪对象的外部}
 \bottomrule
\end{tabular}

\medskip
无需写出 \tkzname{radius}，因为它是默认选项。
\end{NewMacroBox}

 \subsubsection{简单裁剪} 
\begin{tkzexample}[latex=6cm,small] 
\begin{tikzpicture}[scale=.5]
  \tkzDefPoint(0,0){A} \tkzDefPoint(2,2){O}
  \tkzDefPoint(4,4){B} \tkzDefPoint(5,5){C}
  \tkzDrawPoints(O,A,B,C) 
  \tkzLabelPoints(O,A,B,C)
  \tkzDrawCircle(O,A) 
  \tkzClipCircle(O,A)
  \tkzDrawLine(A,C)
  \tkzDrawCircle[fill=teal!10,opacity=.5](C,O)
\end{tikzpicture} 
\end{tkzexample}

\subsection{向外裁剪}

\begin{tkzexample}[latex=6cm,small]
\begin{tikzpicture}
  \tkzInit[xmin=-3,ymin=-2,xmax=4,ymax=3]
   \tkzDefPoint(0,0){O}
   \tkzDefPoint(-4,-2){A}
   \tkzDefPoint(3,1){B}
   \tkzDefCircle[R](O,2) \tkzGetPoint{o}
   \tkzDrawPoints(A,B) % to have a good bounding box
   \begin{scope}
    \tkzClipCircle[out](O,o)
    \tkzDrawLines(A,B)
   \end{scope}
\end{tikzpicture}  
\end{tkzexample} 

\subsection{圆盘的交集}

\begin{tkzexample}[latex=6cm,small]
\begin{tikzpicture}
\tkzDefPoints{0/0/O,4/0/A,0/4/B}
\tkzDrawPolygon[fill=teal](O,A,B)
\tkzClipPolygon(O,A,B)
\tkzClipCircle(A,O)
\tkzClipCircle(B,O)
\tkzFillPolygon[white](O,A,B)
\end{tikzpicture}
\end{tkzexample} 

关于裁剪圆的更复杂示例，请参见 \ref{About clipping circles} 节。

\subsection{裁剪扇形}
\tkzHandBomb\  注意，参数取决于所使用的选项。
\begin{NewMacroBox}{tkzClipSector}{\oarg{local options}\parg{O,\dots}\parg{\dots}}%
\begin{tabular}{lll}%
选项             & 默认值 & 说明                         \\ 
\midrule
\TOline{towards}{towards}{$O$ 为圆心，扇形从 $A$ 起沿 $(OB)$ 方向}
\TOline{rotate} {towards}{扇形从 $A$ 起，由角度确定其张角} 
\TOline{R}{towards}{给定半径和两个角度} 
\bottomrule
\end{tabular} 

\medskip
当然，你可以为图形添加所有 \TIKZ\ 样式。

\medskip   
\begin{tabular}{lll}%
\toprule
选项             & 参数 & 示例                         \\ 
\midrule
\TOline{towards}{\parg{pt,pt}\parg{pt}}{\tkzcname{tkzClipSector(O,A)(B)}}
\TOline{rotate} {\parg{pt,pt}\parg{angle}}{\tkzcname{tkzClipSector[rotate](O,A)(90)}} 
\TOline{R}{\parg{pt,$r$}\parg{angle 1,angle 2}}{\tkzcname{tkzClipSector[R](O,2)(30,90)}}
\end{tabular}
\end{NewMacroBox}

\subsubsection{裁剪扇形：第一个示例}
\begin{tkzexample}[latex=7cm,small] 
\begin{tikzpicture}[scale=0.5]
\tkzDefPoint(0,0){a}
\tkzDefPoint(12,0){b}
\tkzDefPoint(4,10){c}
\tkzInterCC[R](a,6)(b,8) 
\tkzGetFirstPoint{AB1} \tkzGetSecondPoint{AB2}
\tkzInterCC[R](a,6)(c,6) 
\tkzGetFirstPoint{AC1} \tkzGetSecondPoint{AC2}
\tkzInterCC[R](b,8)(c,6) 
\tkzGetFirstPoint{BC1} \tkzGetSecondPoint{BC2}
\tkzDrawArc(a,AB2)(AB1)
\tkzDrawArc(b,AB1)(AB2)
\tkzDrawArc(a,AC2)(AC1)
\tkzDrawArc(c,AC1)(AC2)
\tkzDrawArc(b,BC2)(BC1)
\tkzDrawArc(c,BC1)(BC2)
\begin{scope}
\tkzClipSector(b,BC2)(BC1)
\tkzFillSector[teal!40!white](c,BC1)(BC2)
\end{scope}
\begin{scope}
\tkzClipSector(a,AB2)(AB1)
\tkzFillSector[teal!40!white](b,AB1)(AB2)
\end{scope}
\begin{scope}
\tkzClipSector(a,AC2)(AC1)
\tkzFillSector[teal!40!white](c,AC1)(AC2)
\end{scope}
\end{tikzpicture}
\end{tkzexample}

\subsubsection{裁剪扇形：第二个示例}
\begin{tkzexample}[latex=7cm,small] 
\begin{tikzpicture}[scale=1.5] 
  \tkzDefPoint(0,0){O}
  \tkzDefPoint(2,-1){A}
  \tkzDefPoint(1,1){B} 
  \tkzDrawSector[new,dashed](O,A)(B)
  \tkzDrawSector[new](O,B)(A)
\begin{scope}
\tkzClipSector(O,B)(A)
\tkzDefSquare(O,B) \tkzGetPoints{B'}{O'}
\tkzDrawPolygon[color=teal,fill=teal!20](O,B,B',O')
\end{scope}
\tkzDrawPoints(A,B,O) 
\end{tikzpicture} 
\end{tkzexample}


\subsection{\TIKZ\ 选项：trim left 或 trim right}
参见 \tkzimp{pgfmanual}。

\subsection{\TIKZ\ 控制 \tkzcname{pgfinterruptboundingbox} 与 \tkzcname{endpgfinterruptboundingbox}}
该命令会临时中断边界框的计算并开始一个新的边界框。
参见 \tkzimp{pgfmanual}

\subsubsection{控制边界框} 
\begin{tkzexample}[latex=7cm,small] 
\begin{tikzpicture}
\tkzDefPoint(0,5){A}\tkzDefPoint(5,4){B}
\tkzDefPoint(0,0){C}\tkzDefPoint(5,1){D}
\tkzDrawSegments(A,B C,D A,C)
\pgfinterruptboundingbox
   \tkzInterLL(A,B)(C,D)\tkzGetPoint{I}
\endpgfinterruptboundingbox
\tkzClipBB
\tkzDrawCircle(I,B)
\end{tikzpicture}
\end{tkzexample}

\subsection{反向裁剪：\tkzname{tkzreverseclip}}

要使用该选项，必须先定义一个边界框。

\begin{tkzltxexample}[]
  \tikzset{tkzreverseclip/.style={insert path={
     (current bounding box.south west) --(current bounding box.north west)
   --(current bounding box.north east) --  (current bounding box.south east)
   -- cycle} }}
\end{tkzltxexample}

\subsubsection{\tkzcname{tkzClipPolygon[out]} 示例}
\tkzcname{tkzClipPolygon[out]}、\tkzcname{tkzClipCircle[out]} 使用了该选项。
\begin{tkzexample}[vbox,small]
\begin{tikzpicture}[scale=1]
\tkzInit[xmin=-5,xmax=5,ymin=-4,ymax=6]
\tkzClip
\tkzDefPoints{-.5/0/P1,.5/0/P2}
\foreach \i [count=\j from 3] in {2,...,7}{%
    \tkzDefShiftPoint[P\i]({45*(\i-1)}:1){P\j}}  
\tkzClipPolygon[out](P1,P...,P8)
\tkzCalcLength(P1,P5)\tkzGetLength{r}
\begin{scope}[blend group=screen]
  \foreach \i in {1,...,8}{%
   \tkzDefCircle[R](P\i,\r) \tkzGetPoint{x}
   \tkzFillCircle[color=teal](P\i,x)}
  \end{scope}
\end{tikzpicture}
\end{tkzexample}

\endinput


\part{标记与标注}
\section{标记}
\subsection{标记线段 \tkzcname{tkzMarkSegment}}
\hypertarget{tms}{}  
  
 \begin{NewMacroBox}{tkzMarkSegment}{\oarg{local options}\parg{pt1,pt2}}% 
该宏用于在线段上放置一个标记。

\medskip
\begin{tabular}{lll}%
\toprule
选项             & 默认值 & 说明   \\
\midrule
\TOline{pos}{.5}{标记的位置} 
\TOline{color}{black}{标记的颜色} 
\TOline{mark}{none}{标记的选择} 
\TOline{size}{4pt}{标记的大小}
\bottomrule
\end{tabular}

可选的标记包括 \TIKZ\ 提供的那些，但也有基于 Yves Combe 的思路所创建的其他标记。
\end{NewMacroBox} 

\subsubsection{多个标记}
\begin{tkzexample}[latex=5cm,small] 
\begin{tikzpicture}
  \tkzDefPoint(2,1){A}
  \tkzDefPoint(6,4){B}
  \tkzDrawSegment(A,B)
  \tkzMarkSegment[color=brown,size=2pt,pos=0.4, mark=z](A,B) 
  \tkzMarkSegment[color=blue,pos=0.2, mark=oo](A,B)
  \tkzMarkSegment[pos=0.8,mark=s,color=red](A,B) 
\end{tikzpicture}
\end{tkzexample}

\subsubsection{\tkzname{mark} 的用法}      
\begin{tkzexample}[latex=5cm,small] 
\begin{tikzpicture}
  \tkzDefPoint(2,1){A} 
  \tkzDefPoint(6,4){B}
  \tkzDrawSegment(A,B)
  \tkzMarkSegment[color=gray,pos=0.2,mark=s|](A,B)
  \tkzMarkSegment[color=gray,pos=0.4,mark=s||](A,B)
  \tkzMarkSegment[color=brown,pos=0.6,mark=||](A,B)
  \tkzMarkSegment[color=red,pos=0.8,mark=|||](A,B)
\end{tikzpicture}
\end{tkzexample}


\subsection{标记多条线段 \tkzcname{tkzMarkSegments}}
\hypertarget{tmss}{} 
 
\begin{NewMacroBox}{tkzMarkSegments}{\oarg{local options}\parg{pt1,pt2 pt3,pt4 ...}}%
参数为由空格分隔的点对列表。可以使用 \TIKZ\ 样式。
\end{NewMacroBox} 

\subsubsection{等腰三角形的标记}      
\begin{tkzexample}[latex=6cm,small]
\begin{tikzpicture}[scale=1]
 \tkzDefPoints{0/0/O,2/2/A,4/0/B,6/2/C}
 \tkzDrawSegments(O,A A,B)
 \tkzDrawPoints(O,A,B)
 \tkzDrawLine(O,B)   
 \tkzMarkSegments[mark=||,size=6pt](O,A A,B)
\end{tikzpicture}
\end{tkzexample} 

\subsection{带构造的标记：等边三角形}   
\begin{tkzexample}[latex=5cm,small] 
 \begin{tikzpicture}[scale=1]
  \tkzDefPoint(0,0){A}\tkzDefPoint(3,2){B} 
  \tkzDefPoint(4,0){C}\tkzDefPoint(2.5,1){P}
  \tkzDrawPolygon(A,B,C)
  \tkzDefEquilateral(A,P) \tkzGetPoint{P'}
  \tkzDefPointsBy[rotation=center A angle 60](P,B){P',C'}
  \tkzDrawPolygon(A,P,P')
  \tkzDrawPolySeg(P',C',A,P,B)
  \tkzDrawSegment(C,P)
  \tkzDrawPoints(A,B,C,C',P,P')
  \tkzMarkSegments[mark=s|,size=6pt,
  color=blue](A,P P,P' P',A) 
  \tkzMarkSegments[mark=||,color=orange](B,P P',C')
  \tkzLabelPoints(A,C) \tkzLabelPoints[below](P) 
  \tkzLabelPoints[above right](P',C',B) 
\end{tikzpicture} 
\end{tkzexample}  

\subsection{标记圆弧 \tkzcname{tkzMarkArc}}
\hypertarget{tma}{}  
  
 \begin{NewMacroBox}{tkzMarkArc}{\oarg{local options}\parg{pt1,pt2,pt3}}% 
该宏用于在圆弧上放置标记：pt1 为圆心，pt2 和 pt3 为圆弧的两个端点。

\medskip
\begin{tabular}{lll}%
\toprule
选项             & 默认值 & 说明   \\
\midrule
\TOline{pos}{.5}{标记的位置} 
\TOline{color}{black}{标记的颜色} 
\TOline{mark}{none}{标记的选择} 
\TOline{size}{4pt}{标记的大小}
\bottomrule
\end{tabular}

可选的标记包括 \TIKZ\ 提供的那些，但也有基于 Yves Combe 的思路所创建的其他标记。
\begin{tkzltxexample}[]
|, ||,|||, z, s, x, o, oo 
\end{tkzltxexample}
\end{NewMacroBox} 

\subsubsection{多个标记}
\begin{tkzexample}[latex=7cm,small] 
\begin{tikzpicture}
\tkzDefPoint(0,0){O}
\pgfmathsetmacro\r{2}
\tkzDefPoint(30:\r){A}
\tkzDefPoint(85:\r){B}
\tkzDrawCircle(O,A)
\tkzMarkArc[color=red,mark=||](O,A,B)
\tkzDrawPoints(B,A,O)
\end{tikzpicture}
\end{tkzexample}

 
\subsection{标记角：\tkzcname{tkzMarkAngle}}
这是一个更为精细的操作，因为选项很多。除 \TIKZ\ 自带的标记符号外，用于标记的符号定义在文件 |tkz-lib-marks.tex| 中，由以下字符指定：\begin{tkzltxexample}[]
|, ||,|||, z, s, x, o, oo 
\end{tkzltxexample}




%                \tkzMarkAngle(B, A, C)
%
% 角标记
% 圆弧（单/双/三）以及相等标记。
%
% 默认值：
%                 arc       = simple
%                 mksize  = 1（圆弧半径）
%                 style traits pleins
%                 mkpos ?  位置：0.5（标记位置）
%                 mark   none
%
% 参数（可选）
%             arc     : l, ll, lll
%             mksize  : 1
%             gap     : 3pt
%             dist    : 1?
%             style   : 线型
%             mkpos   : 0.5
%             mark    : none  , |, ||,|||, z, s, x, o, oo ，不过所有
%  % tikz 符号都是允许的

\begin{NewMacroBox}{tkzMarkAngle}{\oarg{local options}\parg{A,O,B}}%
$O$ 为顶点。可以有多种标记方式：你可以简单地画一段圆弧，或者在该圆弧上添加一个标记。圆弧的样式通过选项 \tkzname{arc} 选择，其半径由 \tkzname{size} 给出，标记的大小由 \tkzname{mksize} 给出；圆弧当然也可以着色。

\medskip

\begin{tabular}{lll}%
\toprule
选项             & 默认值 & 说明                        \\ 
\midrule
\TOline{arc}{l}{选择 l、ll 和 lll（单、双或三重）。}
\TOline{size}{1 (cm)}{圆弧半径。}
\TOline{mark}{none}{标记的选择。}
\TOline{mksize}{4pt}{符号（标记）大小。}
\TOline{mkcolor}{black}{符号（标记）颜色。}
\TOline{mkpos}{0.5}{符号在圆弧上的位置。}
\end{tabular} 
\end{NewMacroBox}  

\DeleteShortVerb{\|}
\subsubsection{\tkzname{mark = x} 与 \tkzname{mark =||} 示例}

\begin{tkzexample}[latex=6cm,small]
\begin{tikzpicture}[scale=.75]
    \tkzDefPoints{0/0/O,5/0/A,3/4/B}
    \tkzMarkAngle[size = 4,mark = x,
                  arc=ll,mkcolor = red,mkpos=.33](A,O,B)
    \tkzMarkAngle[size = 2,mark = ||,
                arc=ll,mkcolor = blue,mkpos=.66](A,O,B)
    \tkzDrawLines(O,A O,B)
    \tkzDrawPoints(O,A,B)
\end{tikzpicture}
\end{tkzexample}

\MakeShortVerb{\|}
\begin{NewMacroBox}{tkzMarkAngles}{\oarg{local options}\parg{A,O,B}\parg{A',O',B'}etc.}%
对于多个角，提供了一个带公共选项的宏。
  \end{NewMacroBox}  

\subsection{标记小角：选项 \tkzname{veclen}}\label{opt-veclen}
  问题来自 "decorate" 操作以及 \tkzcname{tkzMarkAngle} 中 \tkzname{size} 的取值。
  解决方法是将 \tkzcname{tkzMarkAngle} 放置在一个带 \tkzname{veclen} 选项的 scope 中。
  在下例中，若不使用 scope，结果会是：LaTeX Error: Dimension too large。
  
  \begin{tkzexample}[latex=6cm,small]
    \begin{tikzpicture}[scale=1]
      \tkzDefPoint(0,0){O}
      \tkzDefPoint(2.5,0){N}
      \tkzDefPoint(-4.2,0.5){M}
      \tkzDefPointBy[rotation=center O angle 30](N)
      \tkzGetPoint{B}
      \tkzDefPointBy[rotation=center O angle -50](N)
      \tkzGetPoint{A}
      \tkzInterLC[common=B](M,B)(O,B) \tkzGetFirstPoint{C}
      \tkzInterLC[common=A](M,A)(O,A) \tkzGetFirstPoint{A'}
      \tkzDrawSegments(A,C M,A M,B A,B)
      \tkzDrawCircle(O,N)
      \begin{scope}[veclen]
        \tkzMarkAngle[mkpos=.2, size=1.2](C,A,M)
      \end{scope}
      \tkzDrawPoints(O, A, B, M, B, C, A')
      \tkzLabelPoints[right](O,A,B)
      \tkzLabelPoints[above left](M,C)
      \tkzLabelPoint[below left](A'){$A'$}
    \end{tikzpicture}
  \end{tkzexample}

  
\subsection{标记直角：\tkzcname{tkzMarkRightAngle}}

\begin{NewMacroBox}{tkzMarkRightAngle}{\oarg{local options}\parg{A,O,B}}%
选项 \tkzname{german} 会改变标记的样式，选项 \tkzname{size} 会改变其大小。

\medskip
\begin{tabular}{lll}%
\toprule
选项             & 默认值 & 说明         \\ 
\midrule
\TOline{german}{normal}{带内点的德式圆弧。}
\TOline{size}{0.2}{边长。}
\end{tabular} 
\end{NewMacroBox}  

\subsubsection{标记直角} 
\begin{tkzexample}[latex=6cm,small]
\begin{tikzpicture}
  \tkzDefPoints{0/0/A,3/1/B,0.9/-1.2/P}
  \tkzDefPointBy[projection = onto B--A](P)  \tkzGetPoint{H}
  \tkzDrawLines[add=.5 and .5](P,H)
  \tkzMarkRightAngle[fill=blue!20,size=.5,draw](A,H,P) 
  \tkzDrawLines[add=.5 and .5](A,B)
  \tkzMarkRightAngle[fill=red!20,size=.8](B,H,P)
  \tkzDrawPoints[](A,B,P,H)  
\end{tikzpicture}
\end{tkzexample}

\subsubsection{德式直角标记} 
\begin{tkzexample}[latex=6cm,small]
\begin{tikzpicture}
  \tkzDefPoints{0/0/A,3/1/B,0.9/-1.2/P}
  \tkzDefPointBy[projection = onto B--A](P)  \tkzGetPoint{H}
  \tkzDrawLines[add=.5 and .5](P,H)
  \tkzMarkRightAngle[german,size=.5,draw](A,H,P) 
  \tkzDrawPoints[](A,B,P,H) 
  \tkzDrawLines[add=.5 and .5](A,B)
  \tkzMarkRightAngle[german,size=.8](P,H,B) 
\end{tikzpicture}
\end{tkzexample}

\subsubsection{样式混合} 
\begin{tkzexample}[latex=6cm,small]
\begin{tikzpicture}[scale=.75]
  \tkzDefPoint(0,0){A}
  \tkzDefPoint(4,1){B}
  \tkzDefPoint(2,5){C}
  \tkzDefPointBy[projection=onto B--A](C) 
      \tkzGetPoint{H}
  \tkzDrawLine(A,B)
  \tkzDrawLine[add = .5 and .2,color=red](C,H)
  \tkzMarkRightAngle[,size=1,color=red](C,H,A)
  \tkzMarkRightAngle[german,size=.8,color=blue](B,H,C)
  \tkzFillAngle[opacity=.2,fill=blue!20,size=.8](B,H,C)
  \tkzLabelPoints(A,B,C,H)
  \tkzDrawPoints(A,B,C,H)
\end{tikzpicture}
\end{tkzexample}

\subsubsection{完整示例} 

\begin{tkzexample}[latex=6cm,small]
\begin{tikzpicture}[rotate=-90]
\tkzDefPoint(0,1){A}
\tkzDefPoint(2,4){C}
\tkzDefPointWith[orthogonal normed,K=7](C,A)
\tkzGetPoint{B}
\tkzDrawSegment[green!60!black](A,C)
\tkzDrawSegment[green!60!black](C,B)
\tkzDrawSegment[green!60!black](B,A)
\tkzDefSpcTriangle[orthic](A,B,C){N,O,P}
\tkzDrawLine[dashed,color=magenta](C,P)
\tkzLabelPoint[left](A){$A$}
\tkzLabelPoint[right](B){$B$}
\tkzLabelPoint[above](C){$C$}
\tkzLabelPoint[left](P){$P$}
\tkzLabelSegment[auto](B,A){$c$}
\tkzLabelSegment[auto,swap](B,C){$a$}
\tkzLabelSegment[auto,swap](C,A){$b$}
\tkzMarkAngle[size=1,color=cyan,mark=|](C,B,A)
\tkzMarkAngle[size=1,color=cyan,mark=|](A,C,P)
\tkzMarkAngle[size=0.75,color=orange,
    mark=||](P,C,B)
\tkzMarkAngle[size=0.75,color=orange,
   mark=||](B,A,C)
\tkzMarkRightAngle[german](A,C,B)
\tkzMarkRightAngle[german](B,P,C)
\end{tikzpicture} 
\end{tkzexample} 

\subsection{\tkzcname{tkzMarkRightAngles}}
\begin{NewMacroBox}{tkzMarkRightAngles}{\oarg{local options}\parg{A,O,B}\parg{A',O',B'}etc.}%
对于多个角，提供了一个带公共选项的宏。
\end{NewMacroBox}

\subsection{\TIKZ\ 库 \tkzname{angles}} % (fold)
\label{sub:angles_library}

如果你偏好使用 \TIKZ\ 库 \tkzname{angles}，可以通过宏 \tkzcname{tkzPicAngle} 和 \tkzcname{tkzPicRightAngle} 来标记角。

\begin{NewMacroBox}{tkzPicAngle}{\oarg{tikz options}\parg{A,O,B}}%
  
\medskip
\begin{tabular}{lll}%
\toprule
选项             & 示例 & 说明         \\ 
\midrule
\TOline{tikz option}{见下文}{绘制角 $\widehat{AOB}$。}
\end{tabular} 
\end{NewMacroBox}  

\begin{NewMacroBox}{tkzPicRightAngle}{\oarg{tikz options}\parg{A,O,B}}%
  
\medskip
\begin{tabular}{lll}%
\toprule
选项             & 示例 & 说明         \\ 
\midrule
\TOline{tikz option}{见下文}{绘制直角 $\widehat{AOB}$。}
\end{tabular} 

\medskip
\emph{你需要了解库 \tkzname{angles} 的各个选项。}
\end{NewMacroBox} 

\subsubsection{使用 \TIKZ\ 标记角} % (fold)
\label{ssub:angle_with_tikz}


\begin{tkzexample}[latex=7cm,small]
  \begin{tikzpicture}
  \tkzDefPoints{0/0/A,4/0/B}
  \tkzDefTriangle[right,swap](A,B) \tkzGetPoint{C}
  \tkzDrawPolygon(A,B,C)
  \tkzDrawPoints(A,B,C)
  \tkzLabelPoints[below](B,A)
  \tkzLabelPoints[above right](C)
  \tkzPicAngle["$\alpha$",draw=orange,
               <->,angle eccentricity=1.2,
               angle radius=1cm](B,A,C)
  \tkzPicRightAngle[draw,red,thick,
                angle eccentricity=.5,
                pic text=.](C,B,A)
  \end{tikzpicture}
\end{tkzexample}

% subsubsection angle_with_tikz (end)
% subsection angles_library (end)
\endinput

\section{点的标注}
\subsection{单点标注：\tkzcname{tkzLabelPoint}}
\hypertarget{tlp}{}
通过多次使用该宏，你可以为同一点添加多个标签。

\begin{NewMacroBox}{tkzLabelPoint}{\oarg{local options}\parg{point}\var{label}}%
\begin{tabular}{lll}%
参数 &  示例  &                  \\ 
\midrule
\TAline{point}{\tkzcname{tkzLabelPoint(A)\{\$A\_1\$\}}}{}
选项  & 默认值 & 说明\\
\midrule
\TOline{TikZ options}{}{颜色、位置等}
\bottomrule
\end{tabular}

\medskip
你可以任选 \TIKZ\ 样式，特别是诸如 above、right、\dots 之类的放置位置。
\end{NewMacroBox}

\subsubsection{\tkzcname{tkzLabelPoint} 示例} 
\begin{tkzexample}[latex=7cm,small]  
\begin{tikzpicture}
  \tkzDefPoint(0,0){A}
  \tkzDefPoint(4,0){B}
  \tkzDefPoint(0,3){C}
  \tkzDrawSegments(A,B B,C C,A)
  \tkzDrawPoints(A,B,C)
  \tkzLabelPoint[left,red](A){$A$}
  \tkzLabelPoint[right,blue](B){$B$}
  \tkzLabelPoint[above,purple](C){$C$}  
\end{tikzpicture} 
\end{tkzexample} 

\subsubsection{引用与标签}
 点的引用名是你在代码中使用该点时所依据的名称；标签则是显示出来的点名称。
 
\begin{tkzexample}[latex=6cm,small]
 \begin{tikzpicture}
    \tkzDefPoint(2,0){A} 
    \tkzDrawPoint(A)
    \tkzLabelPoint[above](A){$A_1$}  
  \end{tikzpicture}
 \end{tkzexample}
 
\subsection{为多点添加标签 \tkzcname{tkzLabelPoints}}
当点的引用名与标签相同，且所有标签相对于各自点的位置一致时，你可以快速放置多个标签。默认使用 \tkzname{below right}。
\hypertarget{tlps}{}  

\begin{NewMacroBox}{tkzLabelPoints}{\oarg{local options}\parg{$A_1,A_2,...$}}%
\begin{tabular}{lll}
参数 &  示例 & 结果                 \\ 
\midrule
\TAline{点列表}{\tkzcname{tkzLabelPoints(A,B,C)}}{显示 $A$、$B$ 和 $C$}
\bottomrule
\end{tabular}

\medskip
该宏可以减少代码行数，但实际上很少有所有点都需要相同标签位置的情况。
\end{NewMacroBox}

\subsubsection{\tkzcname{tkzLabelPoints} 示例}   
\begin{tkzexample}[latex = 6cm,small]  
\begin{tikzpicture}
  \tkzDefPoint(2,3){A}
  \tkzDefShiftPoint[A](30:2){B}
  \tkzDefShiftPoint[A](30:5){C}
  \tkzDrawPoints(A,B,C)
  \tkzLabelPoints(A,B,C) 
\end{tikzpicture} 
\end{tkzexample}
%<--------------------------------------------------------------------------->
%                       tkzAutoLabelPoints
%<--------------------------------------------------------------------------->
\subsection{标签的自动定位 \tkzcname{tkzAutoLabelPoints}}
点的标签沿由选项 \tkzname{center} 给出的中心与该点所确定的方向放置。标签与点之间的距离为中心到该点距离的百分比，由 \tkzname{dist} 给出。
\begin{NewMacroBox}{tkzAutoLabelPoints}{\oarg{local options}\parg{$A_1,A_2,...$}}%
\begin{tabular}{lll}
参数 &  示例 & 结果      \\ 
\midrule
\TAline{点列表}{\tkzcname{tkzAutoLabelPoint(A,B,C)}}{显示 $A$、$B$ 和 $C$}
\end{tabular}
\end{NewMacroBox}

\subsubsection{使用 \tkzcname{tkzAutoLabelPoints} 为点标注} 
这里标签是相对于 $A$、$B$、$C$ 和 $O$ 的重心进行定位的。
\begin{tkzexample}[latex=4cm,small]
\begin{tikzpicture}[scale=1]
 \tkzDefPoint(2,1){O}
 \tkzDefRandPointOn[circle=center O radius 1.5]\tkzGetPoint{A}
 \tkzDefPointBy[rotation=center O angle 100](A)\tkzGetPoint{C}
 \tkzDefPointBy[rotation=center O angle 78](A)\tkzGetPoint{B}
 \tkzDrawCircle(O,A) 
 \tkzDrawPoints(O,A,B,C) 
 \tkzDrawSegments(C,B B,A A,O O,C)
 \tkzDefTriangleCenter[centroid](A,B,C) \tkzGetPoint{G}
 \tkzDrawPoint(tkzPointResult)
 \tkzAutoLabelPoints[center=G,dist=.5](O,A,B,C)
\end{tikzpicture}
\end{tkzexample}




\section{线段的标注}
\subsection{宏 \tkzcname{tkzLabelSegment} 与 \tkzcname{tkzLabelSegments}}
\hypertarget{tls}{}  
\begin{NewMacroBox}{tkzLabelSegment}{\oarg{local options}\parg{pt1,pt2}\marg{label}}
该宏沿线段或直线放置标签。选项为 \TIKZ\ 的选项，例如 \tkzname{pos}。

\medskip
\begin{tabular}{lll}%%
参数    & 示例 & 说明    \\
\midrule
\TAline{label}{\tkzcname{tkzLabelSegment(A,B)\{$5$\}}}{标签文本} 
\TAline{(pt1,pt2)}{(A,B)}{沿 $[AB]$ 放置标签} 
\bottomrule
\end{tabular}

\medskip
\begin{tabular}{lll}%
选项  & 默认值 & 说明    \\
\midrule
\TOline{pos}{.5}{标签的位置} 
\end{tabular}
\end{NewMacroBox}  

\subsubsection{线段上方与下方的标签}      
\begin{tkzexample}[latex=7 cm,small]
\begin{tikzpicture}
\tkzDefPoint(0,0){A}
\tkzDefPoint(6,0){B}
\tkzDrawSegment(A,B)
\tkzLabelSegment[above,pos=.8](A,B){$a$}
\tkzLabelSegment[below,pos=.2](A,B){$4$}
\end{tikzpicture} 
\end{tkzexample}  

\subsubsection{示例：黑板}  
\begin{tkzexample}[latex=6cm,small]
\tikzstyle{background rectangle}=[fill=black]
\begin{tikzpicture}[show background rectangle,scale=.4]
  \tkzDefPoint(0,0){O}
  \tkzDefPoint(1,0){I}
  \tkzDefPoint(10,0){A}
  \tkzDefPointWith[orthogonal normed,K=4](I,A)
   \tkzGetPoint{H}
  \tkzDefMidPoint(O,A) \tkzGetPoint{M}
  \tkzInterLC(I,H)(M,A)\tkzGetPoints{B}{C}   
  \tkzDrawSegments[color=white,line width=1pt](I,H O,A)
  \tkzDrawPoints[color=white](O,I,A,B,M) 
  \tkzMarkRightAngle[color=white,line width=1pt](A,I,B) 
  \tkzDrawArc[color=white,line width=1pt,
              style=dashed](M,A)(O)    
  \tkzLabelSegment[white,right=1ex,pos=.5](I,B){$\sqrt{a}$} 
  \tkzLabelSegment[white,below=1ex,pos=.5](O,I){$1$}   
  \tkzLabelSegment[pos=.6,white,below=1ex](I,A){$a$} 
\end{tikzpicture} 
\end{tkzexample}

\subsubsection{标签与选项 \tkzname{swap}}
\begin{tkzexample}[latex=7cm,small]
\begin{tikzpicture}[rotate=-60]
\tkzSetUpStyle[red,auto]{label style}
\tkzDefPoint(0,1){A}
\tkzDefPoint(2,4){C}
\tkzDefPointWith[orthogonal normed,K=7](C,A)
\tkzGetPoint{B}
\tkzDefSpcTriangle[orthic](A,B,C){N,O,P}
\tkzDefTriangleCenter[circum](A,B,C)
\tkzGetPoint{O}
\tkzDrawPolygon[green!60!black](A,B,C)
\tkzDrawLine[dashed,color=magenta](C,P)
\tkzLabelSegment(B,A){$c$}
\tkzLabelSegment[swap](B,C){$a$}
\tkzLabelSegment[swap](C,A){$b$}
\tkzMarkAngles[size=1,
     color=cyan,mark=|](C,B,A A,C,P)
\tkzMarkAngle[size=0.75,
     color=orange,mark=||](P,C,B)
\tkzMarkAngle[size=0.75,
      color=orange,mark=||](B,A,C)
\tkzMarkRightAngles[german](A,C,B B,P,C)
\tkzAutoLabelPoints[center = O,dist= .1](A,B,C)
 \tkzLabelPoint[below left](P){$P$}
 \end{tikzpicture} 
\end{tkzexample}

\hypertarget{tlss}{} 
 \begin{NewMacroBox}{tkzLabelSegments}{\oarg{local options}\parg{pt1,pt2 pt3,pt4 ...}}%
参数为点对列表。可以使用 \TIKZ\ 样式。
\end{NewMacroBox} 
 
\subsubsection{等腰三角形的标签}      
\begin{tkzexample}[latex=6cm,small]
\begin{tikzpicture}[scale=1]
 \tkzDefPoints{0/0/O,2/2/A,4/0/B,6/2/C}
 \tkzDrawSegments(O,A A,B)
 \tkzDrawPoints(O,A,B)
 \tkzDrawLine(O,B)   
 \tkzLabelSegments[color=red,above=4pt](O,A A,B){$a$}
\end{tikzpicture}
\end{tkzexample}  

\section{直线的标注}
\subsection{宏 \tkzcname{tkzLabelLine}}

\begin{NewMacroBox}{tkzLabelLine}{\oarg{local options}\parg{pt1,pt2}\marg{label}}
\begin{tabular}{lll}%
参数 &  默认值 & 说明   \\ 
\midrule
\TAline{label}{}{\tkzcname{tkzLabelLine(A,B)}\{\$\tkzcname{Delta}\$\}}
\bottomrule
\end{tabular}

\begin{tabular}{lll}%
选项             & 默认值 & 说明   \\ 
\midrule
\TOline{pos}{.5}{\tkzname{pos} 是 \TIKZ\ 的选项，但在此处至关重要\dots} 
\end{tabular}

作为选项，除 \tkzname{pos} 外，你可以使用所有 \TIKZ\ 样式，特别是诸如 \tkzname{above}、\tkzname{right}、\dots 之类的放置位置。
\end{NewMacroBox}

\subsubsection{\tkzcname{tkzLabelLine} 示例}
一个重要的选项是 \tkzname{pos}：它用于沿直线放置标签。\tkzname{pos} 的值可以大于 1，也可以为负。

\begin{tkzexample}[latex=6cm,small]
\begin{tikzpicture}
   \tkzDefPoints{0/0/A,3/0/B,1/1/C}
   \tkzDefLine[perpendicular=through C,K=-1](A,B)
   \tkzGetPoint{c}
   \tkzDrawLines(A,B C,c)
   \tkzLabelLine[pos=1.25,blue,right](C,c){$(\delta)$} 
   \tkzLabelLine[pos=-0.25,red,left](C,c){again $(\delta)$} 
\end{tikzpicture}
\end{tkzexample}

\section{角的标注}
\subsection{宏 \tkzcname{tkzLabelAngle} 与 \tkzcname{tkzLabelAngles}}

\begin{NewMacroBox}{tkzLabelAngle}{\oarg{local options}\parg{A,O,B}}%
只有一个选项 \tkzname{dist}（可带或不带单位），它可以由 \TIKZ\ 选项 \tkzname{pos}（不带单位）替代。默认情况下，数值以厘米为单位。

\begin{tabular}{lll}%
  \toprule
选项             & 默认值 & 说明                        \\ 
\midrule
\TOline{pos}{1}{或 dist，控制从顶点到标签的距离。}
\bottomrule
\end{tabular} 

\medskip 
你可以使用所有 \TIKZ\ 选项来移动标签：rotate、shift、below 等。
\end{NewMacroBox}  

\subsubsection{示例（由 J. S. Bibra 提供，\TeX.SE）} 
\begin{tkzexample}[latex=7cm,small]
\begin{tikzpicture}[scale=.75]
  \tkzDefPoint(0,0){C}
  \tkzDefPoint(20:9){B}
  \tkzDefPoint(80:5){A}
  \tkzDefPointsBy[projection=onto B--C](A){a}
  \tkzDrawPolygon[thick,fill=yellow!15](A,B,C)
  \tkzDrawSegment[dashed, red](A,a)   
  \tkzDrawSegment[style=red, dashed, 
  dim={$10$,15pt,midway,font=\scriptsize,
   rotate=90}](A,a) 
  \tkzMarkAngle(B,C,A)
  \tkzMarkRightAngle(A,a,C)    
  \tkzMarkRightAngle(C,A,B)
  \tkzFillAngle[fill=blue!20, opacity=0.5](B,C,A)
  \tkzFillAngle[fill=red!20, opacity=0.5](A,B,C)
  \tkzLabelAngle[pos=1.25](A,B,C){$\beta$}
  \tkzLabelAngle[pos=1.25](B,C,A){$\alpha$}
  \tkzMarkAngle(A,B,C)
  \tkzDrawPoints(A,B,C)
  \tkzLabelPoints(B,C)
  \tkzLabelPoints[above](A)
\end{tikzpicture}
\end{tkzexample}

\subsubsection{使用 \tkzname{pos}} 
\begin{tkzexample}[latex=7cm,small]
\begin{tikzpicture}[scale=.75]
  \tkzDefPoints{0/0/O,5/0/A,3/4/B}
  \tkzMarkAngle[size = 4,mark = ||,
      arc=ll,color = red](A,O,B)%     
  \tkzDrawLines(O,A O,B)
  \tkzDrawPoints(O,A,B)
  \tkzLabelAngle[pos=2,draw,circle,
      fill=blue!10](A,O,B){$\alpha$} 
\end{tikzpicture}
\end{tkzexample}

\subsubsection{\tkzname{pos} 与 \tkzcname{tkzLabelAngles}} 
\begin{tkzexample}[latex=7cm,small]
\begin{tikzpicture}[rotate=30]
  \tkzDefPoint(2,1){S} 
  \tkzDefPoint(7,3){T}
  \tkzDefPointBy[rotation=center S angle 60](T)
  \tkzGetPoint{P} 
  \tkzDefLine[bisector,normed](T,S,P)
  \tkzGetPoint{s}
  \tkzDrawPoints(S,T,P)   
  \tkzDrawPolygon[color=blue](S,T,P) 
  \tkzDrawLine[dashed,color=blue,add=0 and 3](S,s)  
  \tkzLabelPoint[above right](P){$P$}
  \tkzLabelPoints(S,T)
  \tkzMarkAngle[size = 1.8,mark = |,arc=ll,
                    color = blue](T,S,P)
  \tkzMarkAngle[size = 2.1,mark = |,arc=l,
                    color = blue](T,S,s)
  \tkzMarkAngle[size = 2.3,mark = |,arc=l,
                    color = blue](s,S,P)  
 \tkzLabelAngle[pos = 1.5](T,S,P){$60^{\circ}$}%    
 \tkzLabelAngles[pos = 2.7](T,S,s s,S,P){%
                            $30^{\circ}$}%   
\end{tikzpicture}
\end{tkzexample}


\begin{NewMacroBox}{tkzLabelAngles}{\oarg{local options}\parg{A,O,B}\parg{A',O',B'}etc.}%
对于多个角，提供了一个带公共选项的宏。
\end{NewMacroBox}  

最后，我们需要为圆添加标签。有多种方法；这里我们介绍 \tkzcname{tkzLabelCircle}。

\section{圆的标注}
\subsection{宏 \tkzcname{tkzLabelCircle}}
\begin{NewMacroBox}{tkzLabelCircle}{\oarg{tikz options}\parg{O,A}\parg{angle}\marg{label}}%
\begin{tabular}{lll}%
选项             & 默认值 & 说明                         \\
\midrule
\TOline{tikz options} {}{圆心为 $O$、过 $A$ 的圆}
\bottomrule
\end{tabular} 

\medskip
\emph{可以使用 \TIKZ\ 样式。标签是被创建出来的，因此需要在花括号之间 \code{passed}。}
\end{NewMacroBox} 

\subsubsection{圆周围的标签}
\begin{tkzexample}[latex=5cm,small] 
\begin{tikzpicture}
 \tkzDefPoint(0,0){O} \tkzDefPoint(2,0){N}
 \tkzDefPointBy[rotation=center O angle 50](N) 
     \tkzGetPoint{M}
 \tkzDefPointBy[rotation=center O angle -20](N) 
      \tkzGetPoint{P}
 \tkzDefPointBy[rotation=center O angle 125](N) 
      \tkzGetPoint{P'}
 \tkzLabelCircle[above=4pt](O,N)(120){$\mathcal{C}$}
 \tkzDrawCircle(O,M) 
 \tkzFillCircle[color=blue!10,opacity=.4](O,M) 
 \tkzLabelCircle[draw,
       text width=2cm,text centered,left=24pt](O,M)(-120)%
          {The circle\\ $\mathcal{C}$}  
 \tkzDrawPoints(M,P)\tkzLabelPoints[right](M,P)   
\end{tikzpicture}      
\end{tkzexample} 

\section{圆弧的标注}
\subsection{宏 \tkzcname{tkzLabelArc}}
\hypertarget{tla}{}  
\begin{NewMacroBox}{tkzLabelArc}{\oarg{local options}\parg{pt1,pt2,pt3}\marg{label}}
该宏沿圆弧放置标签。选项为 \TIKZ\ 的选项，例如 \tkzname{pos}。

\medskip
\begin{tabular}{lll}%%
参数    & 示例 & 说明    \\
\midrule
\TAline{label}{\tkzcname{tkzLabelArc(A,B)\{$5$\}}}{标签文本} 
\TAline{(pt1,pt2,pt3)}{(O,A,B)}{沿圆弧 $\widearc{AB}$ 放置标签} 
\bottomrule
\end{tabular}

\medskip
\begin{tabular}{lll}%
选项  & 默认值 & 说明    \\
\midrule
\TOline{pos}{.5}{标签的位置} 
\end{tabular}
\end{NewMacroBox}  

\subsubsection{圆弧上的标签}      
\begin{tkzexample}[latex=7 cm,small]
\begin{tikzpicture}
\tkzDefPoint(0,0){O}
\pgfmathsetmacro\r{2}
\tkzDefPoint(30:\r){A}
\tkzDefPoint(85:\r){B}
\tkzDrawCircle(O,A)
\tkzDrawPoints(B,A,O)
\tkzLabelArc[right=2pt](O,A,B){$\widearc{AB}$}
\tkzLabelPoints(A,B,O)
\end{tikzpicture}
\end{tkzexample}

\endinput


\part{补充}
\section{使用圆规}    

\subsection{主要宏 \tkzcname{tkzCompass}} 
\begin{NewMacroBox}{tkzCompass}{\oarg{local options}\parg{A,B}}%
该宏留下一个圆规记号，即一段过给定点的圆弧；圆心必须给出。若干特定选项可以改变圆弧的外观，此外还可以使用 \TIKZ\ 选项，如 style、color、line width 等。

你可以通过选项 |length| 或选项 |delta| 来设置圆弧的长度。

\medskip
\begin{tabular}{lll}%
\toprule
选项             & 默认值 & 说明                        \\ 
\midrule
\TOline{delta} {0 (deg)}{对称地增大圆弧的角度} 
\TOline{length}{1 (cm)}{改变长度（单位 cm）}
\end{tabular}
\end{NewMacroBox} 

\subsubsection{选项 \tkzname{length}} 
\begin{tkzexample}[latex=7cm,small]
\begin{tikzpicture}
  \tkzDefPoint(1,1){A}
  \tkzDefPoint(6,1){B}
  \tkzInterCC[R](A,4)(B,3)
  \tkzGetPoints{C}{D}
  \tkzDrawPoint(C)
  \tkzCompass[length=1.5](A,C)
  \tkzCompass(B,C)
  \tkzDrawSegments(A,B A,C B,C)
\end{tikzpicture}
\end{tkzexample}

\subsubsection{选项 \tkzname{delta}} 
\tkzname{delta} 是以度为单位的角度：标记从该点周围的 $-\delta$ 延伸到 $+\delta$。接受小数值。当取默认值 $0$ 时，标记的长度由选项 \tkzname{length} 给出。

\begin{tkzexample}[latex=7cm,small]
\begin{tikzpicture} 
  \tkzDefPoint(0,0){A} 
  \tkzDefPoint(5,0){B}
  \tkzInterCC[R](A,4)(B,3)
  \tkzGetPoints{C}{D}
  \tkzDrawPoints(A,B,C) 
  \tkzCompass[delta=20](A,C)
  \tkzCompass[delta=20](B,C) 
  \tkzDrawPolygon(A,B,C)  
  \tkzMarkAngle(A,C,B)
\end{tikzpicture}
\end{tkzexample} 

\subsection{多重构造 \tkzcname{tkzCompasss}} 
\begin{NewMacroBox}{tkzCompasss}{\oarg{local options}\parg{pt1,pt2 pt3,pt4,\dots}}%
\tkzHandBomb\  注意，参数为点对列表。这样可以节省几行代码。

\medskip
\begin{tabular}{lll}%
\toprule
选项             & 默认值 & 说明                        \\ 
\midrule
\TOline{delta} {0}{通过对称增大来修改圆弧的角度} 
\TOline{length}{1}{改变长度} 
\end{tabular}
\end{NewMacroBox} 

\subsubsection{\tkzcname{tkzCompasss} 的用法} % (fold)
\label{ssub:use_tkzcname_tkzcompasss}

% subsubsection use_tkzcname_tkzcompasss (end)
\begin{tkzexample}[latex=7cm,small]
\begin{tikzpicture}[scale=.6]
 \tkzDefPoint(2,2){A}  \tkzDefPoint(5,-2){B}
 \tkzDefPoint(3,4){C}  \tkzDrawPoints(A,B) 
 \tkzDrawPoint[shape=cross out](C)    
 \tkzCompasss[new](A,B A,C B,C C,B) 
 \tkzShowLine[mediator,new,dashed,length = 2](A,B)
 \tkzShowLine[parallel = through C,
                     color=purple,length=2](A,B)
 \tkzDefLine[mediator](A,B)  
  \tkzGetPoints{i}{j}
 \tkzDefLine[parallel=through C](A,B) 
   \tkzGetPoint{D}
 \tkzDrawLines[add=.6 and .6](C,D A,C B,D)
 \tkzDrawLines(i,j) \tkzDrawPoints(A,B,C,i,j,D)  
 \tkzLabelPoints(A,B,C,i,j,D)
\end{tikzpicture}
\end{tkzexample} 

\endinput

\section{Show}

\subsection{显示某些直线的构造 \tkzcname{tkzShowLine}}

 \begin{NewMacroBox}{tkzShowLine}{\oarg{local options}\parg{pt1,pt2} or \parg{pt1,pt2,pt3}}%
这些构造涉及中垂线（垂直平分线）、过给定点的垂线或平行线，以及角平分线。因此参数为两个或三个点的列表。若干选项可以调整构造。该宏的思路来自 \tkzimp{Yves Combe}。
  

\medskip 
\begin{tabular}{lll}%
\toprule
选项      & 默认值 & 说明    \\ 
\midrule
\TOline{mediator}{mediator}{显示中垂线的构造} 
\TOline{perpendicular}{mediator}{垂线的构造} 
\TOline{orthogonal}{mediator}{同上}
\TOline{bisector}{mediator}{角平分线的构造}
\TOline{K}{1}{三角形内的圆}
\TOline{length}{1}{单位 cm，圆弧的长度}
\TOline{ratio} {.5}{圆弧长度比}
\TOline{gap}{2}{放置构造点}
\TOline{size}{1}{圆弧半径（见角平分线）}
 \bottomrule
\end{tabular}

\medskip
当然，你可以添加所有 \TIKZ\ 样式。
\end{NewMacroBox}

\subsubsection{\tkzcname{tkzShowLine} 与 \tkzname{parallel} 示例} 
\begin{tkzexample}[latex=7cm,small]
\begin{tikzpicture}
 \tkzDefPoints{-1.5/-0.25/A,1/-0.75/B,-1.5/2/C}
 \tkzDrawLine(A,B)
 \tkzDefLine[parallel=through C](A,B) \tkzGetPoint{c} 
 \tkzShowLine[parallel=through C](A,B)
 \tkzDrawLine(C,c) \tkzDrawPoints(A,B,C,c)
\end{tikzpicture}
\end{tkzexample}

\subsubsection{\tkzcname{tkzShowLine} 与 \tkzname{perpendicular} 示例} 
\begin{tkzexample}[latex=6cm,small]
\begin{tikzpicture}
\tkzDefPoints{0/0/A, 3/2/B, 2/2/C} 
\tkzDefLine[perpendicular=through C,K=-.5](A,B) 
\tkzGetPoint{c}
\tkzShowLine[perpendicular=through C,K=-.5,gap=3](A,B)
\tkzDefPointBy[projection=onto A--B](c)
\tkzGetPoint{h} 
\tkzMarkRightAngle[fill=lightgray](A,h,C) 
\tkzDrawLines[add=.5 and .5](A,B C,c) 
\tkzDrawPoints(A,B,C,h,c)
\end{tikzpicture}
\end{tkzexample}

\subsubsection{\tkzcname{tkzShowLine} 与 \tkzname{bisector} 示例} 
\begin{tkzexample}[latex=7 cm,small]
\begin{tikzpicture}[scale=1.25]
 \tkzDefPoints{0/0/A, 4/2/B, 1/4/C}
 \tkzDrawPolygon(A,B,C) 
 \tkzSetUpCompass[color=brown,line width=.1 pt]
 \tkzDefLine[bisector](B,A,C)  \tkzGetPoint{a}
 \tkzDefLine[bisector](C,B,A)  \tkzGetPoint{b}
 \tkzInterLL(A,a)(B,b) \tkzGetPoint{I}
 \tkzDefPointBy[projection = onto A--B](I) 
   \tkzGetPoint{H}
 \tkzShowLine[bisector,size=2,gap=3,blue](B,A,C)
 \tkzShowLine[bisector,size=2,gap=3,blue](C,B,A)   
 \tkzDrawCircle[color=blue,%
 line width=.2pt](I,H) 
 \tkzDrawSegments[color=red!50](I,tkzPointResult)
 \tkzDrawLines[add=0 and -0.3,color=red!50](A,a B,b) 
\end{tikzpicture}
\end{tkzexample}

\subsubsection{\tkzcname{tkzShowLine} 与 \tkzname{mediator} 示例} 
\begin{tkzexample}[latex=7 cm,small]
\begin{tikzpicture}
\tkzDefPoint(2,2){A}
\tkzDefPoint(5,4){B}
\tkzDrawPoints(A,B) 
\tkzShowLine[mediator,color=orange,length=1](A,B) 
\tkzGetPoints{i}{j}
\tkzDrawLines[add=-0.1 and -0.1](i,j)
\tkzDrawLines(A,B)
\tkzLabelPoints[below =3pt](A,B)
\end{tikzpicture}
\end{tkzexample}

\subsection{某些变换的构造 \addbs{tkzShowTransformation}}
\begin{NewMacroBox}{tkzShowTransformation}{\oarg{local options}\parg{pt1,pt2} or \parg{pt1,pt2,pt3}}%
这些构造涉及轴对称、中心对称、正射投影和平移。若干选项可以调整构造。该宏的思路来自 \tkzimp{Yves Combe}。
  
\medskip 
\begin{tabular}{lll}%
\toprule
选项             & 默认值 & 说明                        \\ 
\midrule
\TOline{reflection= over pt1--pt2}{reflection}{轴对称的构造} 
\TOline{symmetry=center pt}{reflection}{中心对称的构造} 
\TOline{projection=onto pt1--pt2}{reflection}{投影的构造}
\TOline{translation=from pt1 to pt2}{reflection}{平移的构造}
\TOline{K}{1}{三角形内的圆}
\TOline{length}{1}{圆弧长度}
\TOline{ratio} {.5}{圆弧长度比}
\TOline{gap}{2}{放置构造点}
\TOline{size}{1}{圆弧半径（见角平分线）}
\end{tabular}
\end{NewMacroBox}

\subsubsection{\tkzcname{tkzShowTransformation} 用法示例} 

\begin{tkzexample}[latex=6cm,small]
\begin{tikzpicture}[scale=.5]
  \tkzDefPoint(0,0){O} \tkzDefPoint(2,-2){A}
  \tkzDefPoint(70:4){B} \tkzDrawPoints(A,O,B)
  \tkzLabelPoints(A,O,B)
  \tkzDrawLine[add= 2 and 2](O,A)
  \tkzDefPointBy[translation=from O to A](B) 
  \tkzGetPoint{C}
  \tkzDrawPoint[color=orange](C)  \tkzLabelPoints(C)
  \tkzShowTransformation[translation=from O to A,%
             length=2](B) 
  \tkzDrawSegments[->,color=orange](O,A B,C)  
  \tkzDefPointBy[reflection=over O--A](B) \tkzGetPoint{E}
  \tkzDrawSegment[blue](B,E)
  \tkzDrawPoint[color=blue](E)\tkzLabelPoints(E) 
  \tkzShowTransformation[reflection=over O--A,size=2](B)   
  \tkzDefPointBy[symmetry=center O](B) \tkzGetPoint{F} 
  \tkzDrawSegment[color=green](B,F)
  \tkzDrawPoint[color=green](F)\tkzLabelPoints(F)
  \tkzShowTransformation[symmetry=center O,%
                      length=2](B) 
  \tkzDefPointBy[projection=onto O--A](C) 
  \tkzGetPoint{H}    
  \tkzDrawSegments[color=magenta](C,H)
  \tkzDrawPoint[color=magenta](H)\tkzLabelPoints(H)
  \tkzShowTransformation[projection=onto O--A,%
                         color=red,size=3,gap=-2](C)   
\end{tikzpicture}
\end{tkzexample}

\subsubsection{\tkzcname{tkzShowTransformation} 用法另一示例} 

稍后你还会看到这个图形，但不包含构造线。
\begin{tkzexample}[latex=7cm,small]  
\begin{tikzpicture}[scale=.6]
  \tkzDefPoints{0/0/A,8/0/B,3.5/10/I}
  \tkzDefMidPoint(A,B) \tkzGetPoint{O} 
  \tkzDefPointBy[projection=onto A--B](I) 
     \tkzGetPoint{J}
  \tkzInterLC(I,A)(O,A)  \tkzGetPoints{M}{M'}
  \tkzInterLC(I,B)(O,A)  \tkzGetPoints{N}{N'}
  \tkzDefMidPoint(A,B) \tkzGetPoint{M}    
  \tkzDrawSemiCircle(M,B)
  \tkzDrawSegments(I,A I,B A,B B,M A,N) 
  \tkzMarkRightAngles(A,M,B A,N,B)  
  \tkzDrawSegment[style=dashed,color=blue](I,J)
  \tkzShowTransformation[projection=onto A--B,
                  color=red,size=3,gap=-3](I)
  \tkzDrawPoints[color=red](M,N)
  \tkzDrawPoints[color=blue](O,A,B,I,M) 
  \tkzLabelPoints(O)  
  \tkzLabelPoints[above right](N,I) 
  \tkzLabelPoints[below left](M,A) 
\end{tikzpicture} 
\end{tkzexample} 

\endinput

\section{量角器}
基于 Yves Combe 的思路，以下宏用于绘制量角器。
原理很简单：只需指定一条射线。量角器放置在其原点 $O$ 处，射线的方向由 $A$ 给出。角度按逆时针方向测量；使用选项 \tkzname{return} 时，量角器被反转，角度按顺时针方向测量。
\subsection{宏 \tkzcname{tkzProtractor}} % (fold)
\label{sub:the_macro_tkzcname_tkzprotractor}

% subsection the_macro_tkzcname_tkzprotractor (end)
\begin{NewMacroBox}{tkzProtractor}{\oarg{local options}\parg{$O,A$}}%
\begin{tabular}{lll}%
选项    & 默认值 & 说明     \\ 
\midrule
\TOline{lw}  {0.4 pt} {线宽}
\TOline{scale}  {1} {比例：调整量角器的大小} 
\TOline{return} {false} {三角函数圆反向}
\end{tabular}
\end{NewMacroBox}

\subsubsection{圆形量角器} 
逆时针测量：

\begin{tkzexample}[latex=7cm,small] 
\begin{tikzpicture}[scale=.5]
\tkzDefPoint(2,0){A}\tkzDefPoint(0,0){O}
\tkzDefShiftPoint[A](31:5){B}
\tkzDefShiftPoint[A](158:5){C}
\tkzDrawPoints(A,B,C)
\tkzDrawSegments[color = red,
    line width = 1pt](A,B A,C)
  \tkzProtractor[scale = 1](A,B) 
\end{tikzpicture}
\end{tkzexample}  

\subsubsection{圆形量角器，透明且反向}

\begin{tkzexample}[latex=7cm,small] 
\begin{tikzpicture}[scale=.5]
  \tkzDefPoint(2,3){A}
  \tkzDefShiftPoint[A](31:5){B}  
   \tkzDefShiftPoint[A](158:5){C}   
  \tkzDrawSegments[color=red,line width=1pt](A,B A,C)
  \tkzProtractor[return](A,C) 
\end{tikzpicture}
\end{tkzexample}
\endinput

\section{杂项与数学工具}
\subsection{复制线段} 
目标是在给定射线上，构造出与给定线段等长的线段。

\begin{NewMacroBox}{tkzDuplicateSegment}{\parg{pt1,pt2}\parg{pt3,pt4}\marg{pt5}}%
该宏在给定射线上构造出与给定线段等长的线段。实际上它定义了一个点。
\tkzcname{tkzDuplicateSegment} 是 \tkzcname{tkzDuplicateLen} 的新名称。

\medskip  
\begin{tabular}{lll}%
\toprule
参数             & 示例 & 说明                         \\ 

\midrule
\TAline{(pt1,pt2)(pt3,pt4)\{pt5\}} {\tkzcname{tkzDuplicateSegment}(A,B)(E,F)\{C\}}{AC=EF 且 $C \in [AB)$} \\  
\bottomrule
\end{tabular}

\medskip
\emph{宏 \tkzcname{tkzDuplicateLength} 与本宏完全相同。}
\end{NewMacroBox}

\subsubsection{\tkzcname{tkzDuplicateSegment} 的用法} 

\begin{tkzexample}[latex=6cm,small]
\begin{tikzpicture}[scale=.5]
 \tkzDefPoints{0/0/A,2/-3/B,2/5/C}
 \tkzDuplicateSegment(A,B)(A,C)  
 \tkzGetPoint{D}
 \tkzDrawSegments[new](A,B A,C)
 \tkzDrawSegment[teal](A,D)
 \tkzDrawPoints[new](A,B,C,D) 
 \tkzLabelPoints[above right=3pt](A,B,C,D)
\end{tikzpicture} 
\end{tkzexample} 

\subsubsection{用 \tkzcname{tkzDuplicateSegment} 求黄金比例} 
\begin{tkzexample}[latex=7cm,small]
\begin{tikzpicture}[rotate=-90,scale=.4]
 \tkzDefPoints{0/0/A,10/0/B}
 \tkzDefMidPoint(A,B)   
 \tkzGetPoint{I}
 \tkzDefPointWith[orthogonal,K=-.75](B,A)
 \tkzGetPoint{C}
 \tkzInterLC(B,C)(B,I)  \tkzGetSecondPoint{D}
 \tkzDuplicateSegment(B,D)(D,A) \tkzGetPoint{E}
 \tkzInterLC(A,B)(A,E)   \tkzGetPoints{N}{M}
 \tkzDrawArc[orange,delta=10](D,E)(B)
 \tkzDrawArc[orange,delta=10](A,M)(E)
 \tkzDrawLines(A,B B,C A,D)
 \tkzDrawArc[orange,delta=10](B,D)(I)
 \tkzDrawPoints(A,B,D,C,M,I)
 \tkzLabelPoints[below left](A,B,D,C,M,I)
\end{tikzpicture}
\end{tkzexample}

\subsubsection{黄金三角形或崇高三角形} 
\begin{tkzexample}[latex=7cm,small]
\begin{tikzpicture}[scale=.75] 
  \tkzDefPoints{0/0/A,5/0/C,0/5/B} 
  \tkzDefMidPoint(A,C)\tkzGetPoint{H} 
  \tkzDuplicateSegment(H,B)(H,A)\tkzGetPoint{D} 
  \tkzDuplicateSegment(A,D)(A,B)\tkzGetPoint{E} 
  \tkzDuplicateSegment(A,D)(B,A)\tkzGetPoint{G} 
  \tkzInterCC(A,C)(B,G)\tkzGetSecondPoint{F}
  \tkzDrawLine(A,C)
  \tkzDrawArc(A,C)(B) 
  \begin{scope}[arc style/.style={color=gray,%
                               style=dashed}]
    \tkzDrawArc(H,B)(D) 
    \tkzDrawArc(A,D)(B) 
    \tkzDrawArc(B,G)(F) 
  \end{scope}
  \tkzDrawSegment[dashed](H,B) 
  \tkzCompass(B,F) 
  \tkzDrawPolygon[new](A,B,F)
  \tkzDrawPoints(A,...,H)
  \tkzLabelPoints[below left](A,...,H)
\end{tikzpicture}
\end{tkzexample}

\subsection{线段长度 \tkzcname{tkzCalcLength}}
\TIKZ\ 有一个名为 \tkzname{veclen} 的函数。它
 可以计算 $AB$，其中 $A$ 和 $B$ 为两个点。

唯一的问题是，\TIKZ\ 版本在某些情况下精度不够。我的版本使用宏包 \tkzNamePack{xfp}：速度较慢，但更精确。

\begin{NewMacroBox}{tkzCalcLength}{\oarg{local options}\parg{pt1,pt2}}%
你可以用宏 \tkzcname{tkzGetLength} 保存结果；例如 \tkzcname{tkzGetLength\{dAB\}} \\
会定义宏 \tkzcname{dAB}。

\medskip
\begin{tabular}{lll}%
\toprule
参数    & 示例 & 说明       \\
\midrule
\TAline{(pt1,pt2)\{name of macro\}} {\tkzcname{tkzCalcLength}(A,B)}{\tkzcname{dAB} 以 cm 为单位给出 $AB$}
\bottomrule
\end{tabular}

\medskip
只有一个选项。

\begin{tabular}{lll}%
   
\toprule
 选项    & 默认值 & 示例       \\
\midrule
\TOline{cm}  {true}{\tkzcname{tkzCalcLength}(A,B) 在 \tkzcname{tkzGetLength\{dAB\}} 之后，\tkzcname{dAB} 以 cm 为单位给出 $AB$}
\end{tabular}
\end{NewMacroBox}

\subsubsection{圆规与正方形的构造}

\begin{tkzexample}[latex=7cm,small]
\begin{tikzpicture}[scale=1]
  \tkzDefPoint(0,0){A} \tkzDefPoint(4,0){B}
  \tkzCalcLength(A,B)\tkzGetLength{dAB}
  \tkzDefLine[perpendicular=through A](A,B)
  \tkzGetPoint{D}
  \tkzDefPointWith[orthogonal,K=-1](B,A)    
    \tkzGetPoint{F}
  \tkzGetPoint{C}
  \tkzDrawLine[add= .6 and .2](A,B)
  \tkzDrawLine(A,D) 
  \tkzShowLine[orthogonal=through A,gap=2](A,B)
  \tkzMarkRightAngle(B,A,D)
  \tkzCompasss(A,D D,C)
  \tkzDrawArc[R](B,\dAB)(80,110)
  \tkzDrawPoints(A,B,C,D)
  \tkzDrawSegments[color=gray,style=dashed](B,C C,D)
  \tkzLabelPoints[below left](A,B,C,D)
\end{tikzpicture}
\end{tkzexample}


\subsubsection{由直径确定的圆的半径}
宏 \tkzcname{tkzDefCircle[radius](A,B)} 定义一个半径，我们用 \tkzcname{tkzGetLength} 取出；结果单位为 \tkzname{cm}。

\begin{tkzexample}[latex=6cm,small]
\begin{tikzpicture}[scale=.5]
 \tkzDefPoint(0,0){A}
 \tkzDefPoint(3,-4){B}
 \tkzDefMidPoint(A,B) \tkzGetPoint{M}
 \tkzCalcLength(M,B)\tkzGetLength{rAB}
 \tkzDrawCircle(A,B)
 \tkzDrawPoints(A,B)
 \tkzLabelPoints(A,B)
 \tkzDrawSegment[dashed](A,B)
 \tkzLabelSegment(A,B){$\pgfmathprintnumber{\rAB}$}
\end{tikzpicture}
\end{tkzexample}


\subsection{从 pt 到 cm 或从 cm 到 pt 的换算}
这些宏很简单：除以或乘以 28.45274。它们是：

\begin{NewMacroBox}{tkzpttocm}{\parg{number}\marg{name of macro}}%
结果保存在一个宏中。

\medskip
\begin{tabular}{lll}%
\toprule
参数             & 示例 & 说明                         \\
\midrule
\TAline{(number)\{name of macro\}} {\tkzcname{tkzpttocm}(120)\{len\}}{\tkzcname{len} 给出以 \tkzname{cm} 为单位的数值}
\bottomrule
\end{tabular}

\medskip
你需要将 \tkzcname{len} 与单位 \tkzname{cm} 一起使用。
\end{NewMacroBox}

\subsection{单位换算} 
\begin{NewMacroBox}{tkzcmtopt}{\parg{number}\marg{name of macro}}%
结果保存在一个宏中。

\medskip
\begin{tabular}{lll}
\toprule
参数             & 示例 & 说明                         \\
\midrule
\TAline{(number)\{name of macro\}}{\tkzcname{tkzcmtopt}(5)\{len\}}{\tkzcname{len} 的长度以 \tkzname{pts} 为单位}
\bottomrule
\end{tabular}

\medskip
\emph{结果可以与 \tkzcname{len}\ \tkzname{pt} 一起使用}
\end{NewMacroBox}


\subsection{获取点的坐标}
%<--------------------------------------------------------------------------–>
%                    点的坐标 
%    结果保存在 #2x 和 #2y 中，#1 为点，我们取其坐标
% 使用 $A$ 点：\tkzGetPointCoord(A){V}，然后 \Vx = xA，\Vy = yA
% 单位 cm 
% tkzGetPointCoord 是否支持 [#1] cm 或 pt ？待办
%<--------------------------------------------------------------------------–>
\begin{NewMacroBox}{tkzGetPointCoord}{\parg{$A$}\marg{name of macro}}%
\begin{tabular}{lll}%
参数             & 示例 & 说明                         \\
\midrule
\TAline{(point)\{name of macro\}} {\tkzcname{tkzGetPointCoord}(A)\{A\}}{\tkzcname{Ax} 和 \tkzcname{Ay} 给出 $A$ 的坐标}
\end{tabular}

\medskip
\emph{将一个点的坐标保存在两个宏中。如果给定的名称为 \tkzname{p}，则 \tkzcname{px} 和 \tkzcname{py} 以 cm 为单位给出该点的坐标。}
\end{NewMacroBox}

\subsubsection{用 \tkzcname{tkzGetPointCoord} 传递坐标}

\begin{tkzexample}[width=8cm,small]
\begin{tikzpicture}
 \tkzInit[xmax=5,ymax=3]
 \tkzGrid[sub,orange]
 \tkzDrawX \tkzDrawY
 \tkzDefPoint(1,0){A}
 \tkzDefPoint(4,2){B}
 \tkzGetPointCoord(A){a}
 \tkzGetPointCoord(B){b}
 \tkzDefPoint(\ax,\ay){C}
 \tkzDefPoint(\bx,\by){D}
 \tkzDrawPoints[color=red](C,D)
\end{tikzpicture}
\end{tkzexample}

\subsubsection{用 \tkzcname{tkzGetPointCoord} 求向量之和}
\begin{tkzexample}[width=6cm,small]
\begin{tikzpicture}[>=latex]
  \tkzDefPoint(1,4){a}
  \tkzDefPoint(3,2){b}
  \tkzDefPoint(1,1){c}
  \tkzDrawSegment[->,red](a,b)
  \tkzGetPointCoord(c){c}
  \draw[color=blue,->](a) -- ([shift=(b)]\cx,\cy) ;
  \draw[color=purple,->](b) -- ([shift=(b)]\cx,\cy) ;
  \tkzDrawSegment[->,blue](a,c)
  \tkzDrawSegment[->,purple](b,c)
\end{tikzpicture}
\end{tkzexample}

\subsection{交换点的标签}

\begin{NewMacroBox}{tkzSwapPoints}{\parg{$pt1$,$pt2$}}%
\begin{tabular}{lll}%
参数             & 示例 & 说明                         \\
\midrule
\TAline{(pt1,pt2)} {\tkzcname{tkzSwapPoints}(A,B)}{现在 $A$ 具有 $B$ 的坐标}
\end{tabular}

\emph{两点交换了各自的坐标。}
\end{NewMacroBox}

\subsubsection{\tkzcname{tkzSwapPoints} 的用法}

\begin{tkzexample}[width=6cm,small]
\begin{tikzpicture}
  \tkzDefPoints{0/0/O,5/-1/A,2/2/B}
   \tkzSwapPoints(A,B)
   \tkzDrawPoints(O,A,B)
   \tkzLabelPoints(O,A,B)
\end{tikzpicture}
\end{tkzexample}

\subsection{点积}
在欧几里得几何中，由笛卡尔坐标计算的两个向量的点积被广泛使用。

\begin{NewMacroBox}{tkzDotProduct}{\parg{$pt1$,$pt2$,$pt3$}}%
  两个向量 $\overrightarrow{u} = [a,b]$ 和  $\overrightarrow{v} = [a',b']$ 的点积定义为：$\overrightarrow{u}\cdot \overrightarrow{v} = aa' + bb'$

$\overrightarrow{u} = \overrightarrow{pt1pt2}$ $\overrightarrow{v} = \overrightarrow{pt1pt3}$
  
\begin{tabular}{lll}%
参数             & 示例 & 说明                         \\
\midrule
\TAline{(pt1,pt2,pt3)} {\tkzcname{tkzDotProduct}(A,B,C)}{结果为 $\overrightarrow{AB}\cdot \overrightarrow{AC}$}
\end{tabular}

\emph{结果是一个数值，可以用 \tkzcname{tkzGetResult} 取出。}
\end{NewMacroBox}

\subsubsection{简单示例} % (fold)
\label{ssub:simple_example}

\begin{tkzexample}[small,latex=7cm]
\begin{tikzpicture}
  \tkzDefPoints{-2/-3/A,4/0/B,1/3/C}
  \tkzDefPointBy[projection= onto A--B](C)  
  \tkzGetPoint{H}
  \tkzDrawSegment(C,H)
  \tkzMarkRightAngle(C,H,A)
  \tkzDrawSegments[vector style](A,B A,C)
  \tkzDrawPoints(A,H) \tkzLabelPoints(A,B,H)
  \tkzLabelPoints[above](C)
  \tkzDotProduct(A,B,C) \tkzGetResult{pabc}
 % \pgfmathparse{round(10*\pabc)/10}
  \let\pabc\pgfmathresult
  \node at (1,-3) {$\overrightarrow{PA}\cdot \overrightarrow{PB}=\pabc$};
  \tkzDotProduct(A,H,B) \tkzGetResult{phab}
 % \pgfmathparse{round(10*\phab)/10}
  \let\phab\pgfmathresult
  \node at (1,-4) {$PA \times PH = \phab $};
\end{tikzpicture}
\end{tkzexample}
% subsubsection simple_example (end)


\subsubsection{共圆点} % (fold)
\label{ssub:cocyclicpts}

\begin{tkzexample}[small,latex=7cm]
\begin{tikzpicture}[scale=.75]
  \tkzDefPoints{1/2/O,5/2/B,2/2/P,3/3/Q}
  \tkzInterLC[common=B](O,B)(O,B) \tkzGetFirstPoint{A}
  \tkzInterLC[common=B](P,Q)(O,B) \tkzGetPoints{C}{D}
  \tkzDrawCircle(O,B)
  \tkzDrawSegments(A,B C,D)
  \tkzDrawPoints(A,B,C,D,P)
  \tkzLabelPoints(P)
  \tkzLabelPoints[below left](A,C)
  \tkzLabelPoints[above right](B,D)
  \tkzDotProduct(P,A,B) \tkzGetResult{pab}
  \pgfmathparse{round(10*\pab)/10}
  \let\pab\pgfmathresult
  \tkzDotProduct(P,C,D) \tkzGetResult{pcd}
  \pgfmathparse{round(10*\pcd)/10}
  \let\pcd\pgfmathresult
  \node at (1,-3) {%
  $\overrightarrow{PA}\cdot \overrightarrow{PB} =
   \overrightarrow{PC}\cdot \overrightarrow{PD}$};
  \node at (1,-4)%
  {$\overrightarrow{PA}\cdot \overrightarrow{PB}=\pab$};
 \node at (1,-5){%
 $\overrightarrow{PC}\cdot \overrightarrow{PD} =\pcd$};
\end{tikzpicture}
\end{tkzexample}
% subsubsection cocyclicpts (end)

\newpage
\subsection{点关于圆的幂}

\begin{NewMacroBox}{tkzPowerCircle}{\parg{$pt1$}\parg{$pt2$,$pt3$}}%
\begin{tabular}{lll}%
参数             & 示例 & 说明                         \\
\midrule
\TAline{(pt1)(pt2,pt3)} {\tkzcname{tkzPowerCircle}(A)(O,M)}{$A$ 关于圆 (O,A) 的幂}
\end{tabular}

\emph{结果是一个数值，表示一个点关于一个圆的幂。}
\end{NewMacroBox}

\subsubsection{由根轴求幂} % (fold)
\label{ssub:power}

在本例中，已画出根轴 $(EF)$。在 $(EF)$ 上选取一点 $H$，并计算了 $H$ 关于圆心为 $A$ 的圆的幂，以及 $HS^2$。你可以验证，$H$ 关于圆心为 $C$ 的圆的幂，以及 $HS'^2$、$HT^2$、$HT'^2$ 给出相同的结果。

\begin{tkzexample}[small,latex=7cm]
\begin{tikzpicture}[scale=.5]
  \tkzDefPoints{-1/0/A,0/5/B,5/-1/C,7/1/D}
  \tkzDrawCircles(A,B C,D)
  \tkzDefRadicalAxis(A,B)(C,D) \tkzGetPoints{E}{F}
  \tkzDrawLine[add=1 and 2](E,F)
  \tkzDefPointOnLine[pos=1.5](E,F) 
  \tkzGetPoint{H}
  \tkzDefLine[tangent from = H](A,B)
  \tkzGetPoints{T}{T'}
  \tkzDefLine[tangent from = H](C,D)
  \tkzGetPoints{S}{S'}
  \tkzDrawSegments(H,T H,T' H,S H,S')
  \tkzDrawPoints(A,B,C,D,E,F,H,T,T',S,S')
  \tkzPowerCircle(H)(A,B) \tkzGetResult{pw}
  \tkzDotProduct(H,S,S) \tkzGetResult{phtt}
  \node {Power $\approx \pw \approx \phtt$};
\end{tikzpicture}
\end{tkzexample}
% subsubsection power (end)

\subsection{根轴}

在几何中，两个不同心圆的根轴是关于两个圆具有相同幂的点的集合。这里 |\tkzDefRadicalAxis(A,B)(C,D)| 给出两个圆 $\mathcal{C}(A,B)$ 和 $\mathcal{C}(C,D)$ 的根轴。 

\begin{NewMacroBox}{tkzDefRadicalAxis}{\parg{$pt1$,$pt2$}\parg{$pt3$,$pt4$}}%
\begin{tabular}{lll}%
参数             & 示例 & 说明                         \\
\midrule
\TAline{(pt1,pt2)(pt3,pt4)} {\tkzcname{tkzDefRadicalAxis}(A,B)(C,D)}{圆心为 $A$ 和 $C$ 的两个圆}
\midrule
\end{tabular}


\emph{结果为根轴上的两个点。}
\end{NewMacroBox}

\subsubsection{两个相离的圆} % (fold)
\label{ssub:two_circles_disjointed}


\begin{tkzexample}[small,latex=8cm]
\begin{tikzpicture}[scale=.75]
    \tkzDefPoints{-1/0/A,0/2/B,4/-1/C,4/0/D}
    \tkzDrawCircles(A,B C,D)
    \tkzDefRadicalAxis(A,B)(C,D)
     \tkzGetPoints{E}{F}
    \tkzDrawLine[add=1 and 2](E,F)
    \tkzDrawLine[add=.5 and .5](A,C)
\end{tikzpicture}
\end{tkzexample}
% subsubsection two_circles_disjointed (end)

\subsubsection{两个相交的圆} % (fold)
\label{sub:two_intersecting_circles}
\begin{tkzexample}[small,latex=8cm]
\begin{tikzpicture}[scale=.5]
  \tkzDefPoints{-1/0/A,0/2/B,3/-1/C,3/-2/D}
  \tkzDrawCircles(A,C B,D)
  \tkzDefRadicalAxis(A,C)(B,D)
  \tkzGetPoints{E}{F}
  \tkzDrawPoints(A,B,C,D,E,F)
    \tkzLabelPoints(A,B,C,D,E,F)
  \tkzDrawLine[add=.25 and .5](E,F)
  \tkzDrawLine[add=.25 and .25](A,B)
\end{tikzpicture}
\end{tkzexample}
% subsection two_intersecting_circles (end)


\subsubsection{两个外切圆} % (fold)
\label{sub:two_externally_tangent_circles}
\begin{tkzexample}[small,latex=8cm]
\begin{tikzpicture}[scale=.5]
  \tkzDefPoints{0/0/A,4/0/B,6/0/C}
  \tkzDrawCircles(A,B C,B)
  \tkzDefRadicalAxis(A,B)(C,B)
  \tkzGetPoints{E}{F}
  \tkzDrawPoints(A,B,C,E,F)
    \tkzLabelPoints(A,B,C,E,F)
  \tkzDrawLine[add=1 and 1](E,F)
  \tkzDrawLine[add=.5 and .5](A,B)
\end{tikzpicture}
\end{tkzexample}
% subsection two_externally_tangent_circles (end)


\subsubsection{两个内切圆} % (fold)
\label{sub:deux_cercles_tangents_interieurement}
\begin{tkzexample}[small,latex=8cm]
\begin{tikzpicture}[scale=.5]
  \tkzDefPoints{0/0/A,3/0/B,5/0/C}
  \tkzDrawCircles(A,C B,C)
  \tkzDefRadicalAxis(A,C)(B,C)
  \tkzGetPoints{E}{F}
  \tkzDrawPoints(A,B,C,E,F)
  \tkzLabelPoints[below right](A,B,C,E,F)
  \tkzDrawLine[add=1 and 1](E,F)
  \tkzDrawLine[add=.5 and .5](A,B)
\end{tikzpicture}
\end{tkzexample}
% subsection deux_cercles_tangents_interieurement (end)

\subsubsection{三个圆} % (fold)
\label{ssub:threecircles}



\begin{tkzexample}[small,latex=7cm]
\begin{tikzpicture}[scale=.4]
  \tkzDefPoints{0/0/A,5/0/a,7/-1/B,3/-1/b,5/-4/C,2/-4/c}
  \tkzDrawCircles(A,a B,b C,c)
  \tkzDefRadicalAxis(A,a)(B,b) \tkzGetPoints{i}{j}
  \tkzDefRadicalAxis(A,a)(C,c) \tkzGetPoints{k}{l}
  \tkzDefRadicalAxis(C,c)(B,b) \tkzGetPoints{m}{n}
  \tkzDrawLines[new](i,j k,l m,n)
\end{tikzpicture}
\end{tkzexample}
% subsubsection threecircles (end)
  
\subsection{\tkzcname{tkzIsLinear}、\tkzcname{tkzIsOrtho}}
 \begin{NewMacroBox}{tkzIsLinear}{\parg{$pt1$,$pt2$,$pt3$}}%
 \begin{tabular}{lll}%
 参数             & 示例 & 说明                         \\
 \midrule
 \TAline{(pt1,pt2,pt3)} {\tkzcname{tkzIsLinear}(A,B,C)}{$A,B,C$ 共线？}
 \midrule
 \end{tabular}
 
 \emph{\tkzcname{tkzIsLinear} 测试三个点 $pt1$、$pt2$、$pt3$ 是否共线。测试允许约 $0.06^\circ$ 的角度误差。}
 \end{NewMacroBox}
 
 \begin{NewMacroBox}{tkzIsOrtho}{\parg{$pt1$,$pt2$,$pt3$}}%
 \begin{tabular}{lll}%
 参数             & 示例 & 说明                         \\
 \midrule
 \TAline{(pt1,pt2,pt3)} {\tkzcname{tkzIsOrtho}(A,B,C)}{$(AB)\perp (AC)$ ？}
 \midrule
 \end{tabular}
 
 \emph{\tkzcname{tkzIsOrtho} 测试直线 $(pt1pt2)$ 与 $(pt1pt3)$ 是否正交，容差相同。}
 \end{NewMacroBox}
 
 \subsubsection{\tkzcname{tkzIsOrtho} 与 \tkzcname{tkzIsLinear} 的用法}
   
\begin{tkzexample}[small,latex=7cm]
  \begin{tikzpicture}
  \tkzDefPoints{1/-2/A,5/0/B}
  \tkzDefCircle[diameter](A,B) \tkzGetPoint{O}
  \tkzDrawCircle(O,A)
  \tkzDefPointBy[rotation= center O angle 60](B) 
  \tkzGetPoint{C}
  \tkzDefPointBy[rotation= center O angle 60](A) 
  \tkzGetPoint{D}
  \tkzDrawCircle(O,A)
  \tkzDrawPoints(A,B,C,D,O)
  \tkzIsOrtho(C,A,B)
  \iftkzOrtho
    \tkzDrawPolygon[blue](A,B,C)
  \tkzDrawPoints[blue](A,B,C,D)
  \else
  \tkzDrawPoints[red](A,B,C,D)
  \fi
   \tkzIsLinear(O,C,D)
    \iftkzLinear
    \tkzDrawSegment[orange](C,D)
    \fi
\end{tikzpicture}
  
\end{tkzexample}

  
\endinput


\part{样式}
\section{预定义样式}\label{custom}
如何进行取决于你如何使用该宏包。一个好方法是尽可能使用预定义样式，将内容与形式分离。如果你计划创建一个包含大量图形的文档（就像本手册一样），这是正确的方法。我们将看到如何为文档定义全局样式，以及如何在局部使用样式。

文件 \tkzname{tkz-euclide.cfg} 包含了主要对象的预定义样式，其中最重要的是点、直线、线段、圆、圆弧和圆规记号。
 如果你总是使用相同的样式并创建大量图形，那么值得创建自己的样式。为此，你需要知道哪些特性可以修改，以及 \TIKZ\ 的一些概念。
 
 预定义样式是全局的：它们在图形创建之前就已存在，最好不要在两个图形之间修改它们。另一方面，它们可以在一个图形内部被临时修改；此时它们被局部定义，不会影响其他图形。
 
 下面是我如何定义你正在阅读的本文档的样式。

\begin{tkzltxexample}[]
  \tkzSetUpColors[background=white,text=black]  
  \tkzSetUpPoint[size=2,color=teal]
  \tkzSetUpLine[line width=.4pt,color=teal]
  \tkzSetUpCompass[color=orange, line width=.4pt,delta=10]
  \tkzSetUpArc[color=gray,line width=.4pt]
  \tkzSetUpStyle[orange]{new}
\end{tkzltxexample}

宏 \tkzcname{tkzSetUpColors} 设置背景色和文本颜色。如果你不使用它，则会使用文档的颜色和字体。下面我们来看看如何定义主要对象的样式。

\section{点样式}
点的定义方式如下：
\begin{tkzltxexample}[]
\tikzset{point style/.style = {%
       draw         = \tkz@euc@pointcolor,
       inner sep    = 0pt,
       shape        = \tkz@euc@pointshape,
       minimum size = \tkz@euc@pointsize,
       fill         = \tkz@euc@pointcolor}}  
\end{tkzltxexample}

你当然可以使用 \tkzcname{tikzset}，但宏包为此提供了一个宏：\tkzcname{tkzSetUpPoint}，它可以全局或局部使用。

\subsection{\tkzcname{tkzSetUpPoint} 的用法}

\begin{NewMacroBox}{tkzSetUpPoint}{\oarg{local options}}%
\begin{tabular}{lll}%
选项 &  默认值 & 说明                 \\ 
\midrule
\TOline{color}{black}{点的颜色} 
\TOline{size}{3}{点的大小} 
\TOline{fill}{black!50}{点的内部颜色} 
\TOline{shape}{circle}{点的形状：circle、cross 或 cross out} 
\end{tabular}
\end{NewMacroBox}



\subsubsection{全局样式或局部样式}
首先，下面是用本手册的样式创建的一个图形；然后在 \tkzNameEnv{tikzpicture} 环境内部修改了点的样式。

你可以全局或局部使用宏 \tkzcname{tkzSetUpPoint}。如果你将其放在导言区或第一个图形之前，点的样式将对文档中所有图形有效。你可以通过将该命令放在 \tkzNameEnv{tikzpicture} 环境内部来局部使用另一种样式。


\begin{tkzexample}[latex=7cm,small]
\begin{tikzpicture}
  \tkzDefPoints{0/0/A,5/0/B,3/2/C,3/1/D}
  \tkzDrawPolygon(A,B,C)
  \tkzDrawPoints(A,B,C)
  \tkzLabelPoints(A,B) 
  \tkzLabelPoints[above right](C) 
\end{tikzpicture}
\end{tkzexample}

\subsubsection{局部样式}
在第二个图形中局部修改了点的样式。
\begin{tkzexample}[latex=7cm,small]
\begin{tikzpicture}
  \tkzSetUpPoint[size=4,color=red,fill=red!20]
  \tkzDefPoints{0/0/A,5/0/B,3/2/C,3/1/D}
  \tkzDrawPolygon(A,B,C)
  \tkzDrawPoints(A,B,C)
  \tkzDrawPoint[shape=cross out,thick](D)
  \tkzLabelPoints(A,B) 
  \tkzLabelPoints[above right](C) 
\end{tikzpicture}
\end{tkzexample}

\subsubsection{\tkzname{Style} 与 \tkzname{scope}}
点恢复了初始样式。点 D 具有一个受环境 \tkzNameEnv{scope} 限制的新样式。也可以使用 |{...}| 或点恢复了初始样式。点 $D$ 具有一个受环境 \tkzNameEnv{scope} 限制的新样式。也可以使用 |{...}| 或 |\begingoup  ... \endgroup|。

\begin{tkzexample}[latex=7cm,small]
\begin{tikzpicture}
  \tkzDefPoints{0/0/A,5/0/B,3/2/C,3/1/D}
  \tkzDrawPolygon(A,B,C)
  \tkzDrawPoints(A,B,C)
  \begin{scope}
    \tkzSetUpPoint[size=4,color=red,fill=red!20]
    \tkzDrawPoint(D)
  \end{scope}
  \tkzLabelPoints(A,B) 
  \tkzLabelPoints[above right](C) 
\end{tikzpicture}
\end{tkzexample}

\subsubsection{\tkzcname{tkzSetUpPoint} 简单示例}

\begin{tkzexample}[latex=5cm,small]
\begin{tikzpicture}
  \tkzSetUpPoint[shape = cross out,color=blue]
  \tkzDefPoint(2,1){A}
  \tkzDefPoint(4,0){B}
  \tkzDrawLine(A,B)
  \tkzDrawPoints(A,B)
\end{tikzpicture}
\end{tkzexample}

\subsubsection{在 group 内部使用 \tkzcname{tkzSetUpPoint}}
\begin{tkzexample}[latex=8cm,small]
\begin{tikzpicture}
  \tkzDefPoints{0/0/A,2/4/B,4/0/C,3/2/D}
  \tkzDrawSegments(A,B A,C A,D)
  {\tkzSetUpPoint[shape=cross out,
            fill= teal!50,
            size=4,color=teal]
  \tkzDrawPoints(A,B)}
  \tkzSetUpPoint[fill= teal!50,size=4,
               color=teal]
   \tkzDrawPoints(C,D)
  \tkzLabelPoints(A,B,C,D)
\end{tikzpicture}
\end{tkzexample}

\section{直线样式}

有几种方法可以改变直线的样式：使用 \tkzcname{tkzSetUpLine}，或者直接使用\\ |\tikzset{line style/.style = ... }|。

关于 \tkzname{line width} 的提醒：有许多预定义样式可以更"自然"地设置线宽。你也可以重新定义这些样式。


\medskip
\begin{tabular}{cc}
预定义样式 & \tkzname{line width} 的值 \cr
\midrule
ultra thin    &  0.1 pt \cr
very thin     &  0.2 pt \cr
thin          &  0.4 pt \cr
semithick     &  0.6 pt \cr
thick         &  0.8 pt \cr
very thick    &  1.2 pt \cr
ultra thick   &  1.6 pt \cr
\midrule
\end{tabular}


\subsection{\tkzcname{tkzSetUpLine} 的用法} \label{tkzsetupline}
该宏定义所有直线的样式。

\begin{NewMacroBox}{tkzSetUpLine}{\oarg{local options}}%
\begin{tabular}{lll}%
选项 &  默认值 & 说明                 \\
\midrule
\TOline{color}{black}{构造线的颜色}
\TOline{line width}{0.4pt}{构造线的粗细}
\TOline{style}{solid}{构造线的样式}
\TOline{add}{.2 and .2}{改变线段的长度}
\end{tabular}
\end{NewMacroBox}

\subsubsection{改变线宽}
\begin{tkzexample}[latex=8cm,small]
\begin{tikzpicture}[scale=.75]
\tkzSetUpLine[line width=1pt]
\begin{scope}[rotate=-90]
    \tkzDefPoints{0/6/A,10/0/B,10/6/C}
    \tkzDefPointBy[projection = onto B--A](C)
    \tkzGetPoint{H}
    \tkzMarkRightAngle[size=.4,
                       fill=teal!20](B,C,A)
    \tkzMarkRightAngle[size=.4,
                       fill=orange!20](B,H,C)
    \tkzDrawPolygon(A,B,C)
    \tkzDrawSegment[new](C,H)
\end{scope}
 \tkzLabelSegment[below](C,B){$a$}
 \tkzLabelSegment[right](A,C){$b$}
 \tkzLabelSegment[left](A,B){$c$}
 \tkzLabelSegment[color=red](C,H){$h$}
 \tkzDrawPoints(A,B,C)
 \tkzLabelPoints[above left](H)
 \tkzLabelPoints(B,C)
 \tkzLabelPoints[above](A)
\end{tikzpicture}
\end{tkzexample}

\subsubsection{改变直线样式}

\begin{tkzexample}[latex=7cm,small]
\begin{tikzpicture}[scale=.5]
\tikzset{line style/.style = {color = gray,
                             style=dashed}}
\tkzDefPoints{1/0/A,4/0/B,1/1/C,5/1/D}
\tkzDefPoints{1/2/E,6/2/F,0/4/A',3/4/B'}
\tkzCalcLength(C,D)
\tkzGetLength{rCD}
\tkzCalcLength(E,F)
\tkzGetLength{rEF}
\tkzInterCC[R](A',\rCD)(B',\rEF)
\tkzGetPoints{I}{J}
\tkzDrawLine(A',B')
\tkzCompass(A',B')
\tkzDrawSegments(A,B C,D E,F)
\tkzDefCircle[R](A',\rCD) \tkzGetPoint{a'}
\tkzDefCircle[R](B',\rEF)\tkzGetPoint{b'}
\tkzDrawCircles(A',a' B',b')
\begin{scope}
  \tkzSetUpLine[color=red]
  \tkzDrawSegments(A',I B',I)
\end{scope}
\tkzDrawPoints(A,B,C,D,E,F,A',B',I,J)
\tkzLabelPoints(A,B,C,D,E,F,A',B',I,J)
\end{tikzpicture}
\end{tkzexample}

\subsubsection{延长直线}
\begin{tkzexample}[latex=7cm,small]
  \begin{tikzpicture}[scale=.75]
  \tkzSetUpLine[add=.5 and .5]
  \tkzDefPoints{0/0/A,4/0/B,1/3/C}
  \tkzDrawLines(A,B B,C A,C)
  \tkzDrawPolygon[red,thick](A,B,C)
  \tkzSetUpPoint[size=4,circle,color=red,fill=red!20]
  \tkzDrawPoints(A,B,C)
  \end{tikzpicture}
\end{tkzexample}

\section{圆弧样式}

\subsection{宏 \tkzcname{tkzSetUpArc}}
\begin{NewMacroBox}{tkzSetUpArc}{\oarg{local options}}%
\begin{tabular}{lll}%
选项 &  默认值 & 说明                 \\
\midrule
\TOline{color}{black}{直线的颜色}
\TOline{line width}{0.4pt}{直线的粗细}
\TOline{style}{solid}{构造线的样式}
\TOline{delta}{0}{在每端将圆弧延长 \tkzname{delta} 度}
\end{tabular}

\medskip
自第 5.14 版起，用本宏设置的 \tkzname{delta} 仅作用于 \tkzcname{tkzDrawArc}、\tkzcname{tkzMarkArc} 和 \tkzcname{tkzLabelArc} 所绘制的圆弧；圆规记号有自己的 \tkzname{delta}，由 \tkzcname{tkzSetUpCompass} 设置。
\end{NewMacroBox}

\subsubsection{\tkzcname{tkzSetUpArc} 的用法}
\begin{tkzexample}[latex=7cm,small]
\begin{tikzpicture}
\def\r{3} \def\angle{200}
\tkzSetUpArc[delta=10,color=purple,line width=.2pt]
\tkzSetUpLabel[font=\scriptsize,red]
\tkzDefPoint(0,0){O}
\tkzDefPoint(\angle:\r){A}
\tkzInterCC(O,A)(A,O) \tkzGetPoints{C'}{C}
\tkzInterCC(O,A)(C,O) \tkzGetPoints{D'}{D}
\tkzInterCC(O,A)(D,O) \tkzGetPoints{X'}{X}
\tkzDrawCircle(O,A)
\tkzDrawArc(A,C')(C)
\tkzDrawArc(C,O)(D)
\tkzDrawArc(D,O)(X)
\tkzDrawLine[add=.1 and .1](A,X)
\tkzDrawPoints(O,A)
\tkzSetUpPoint[size=3,color=purple,fill=purple!10]
\tkzDrawPoints(C,C',D,X)
\tkzLabelPoints[below left](O,A)
\tkzLabelPoints[below](C')
\tkzLabelPoints[below right](X)
\tkzLabelPoints[above](C,D)
\end{tikzpicture}
\end{tkzexample}

\section{圆规样式}
圆规记号有助于理解图形的构造，通过显示获得某些点所需的圆弧。

\subsection{宏 \tkzcname{tkzSetUpCompass} }
\begin{NewMacroBox}{tkzSetUpCompass}{\oarg{local options}}%
\begin{tabular}{lll}%
选项 &  默认值 & 说明                 \\
\midrule
\TOline{color}{black}{构造线的颜色}
\TOline{line width}{0.4pt}{构造线的粗细}
\TOline{style}{solid}{直线样式：solid、dashed、dotted、...}
\TOline{delta}{0}{改变圆规记号的长度}
\end{tabular}
\end{NewMacroBox}

\subsubsection{\tkzcname{tkzSetUpCompass} 的用法}

\begin{tkzexample}[latex=7cm,small]
\begin{tikzpicture}
  \tkzSetUpCompass[color=red,delta=15]
  \tkzDefPoint(1,1){A}
  \tkzDefPoint(6,1){B}
  \tkzInterCC[R](A,4)(B,4) \tkzGetPoints{C}{D}
  \tkzCompass(A,C)
  \tkzCompass(B,C)
  \tkzDrawPolygon(A,B,C)
  \tkzDrawPoints(A,B,C)
\end{tikzpicture} 
\end{tkzexample}

\subsubsection{\tkzcname{tkzSetUpCompass} 与 \tkzcname{tkzShowLine} 配合使用}

\begin{tkzexample}[latex=7cm,small]
\begin{tikzpicture}[scale=.75]
\tkzSetUpStyle[bisector,size=2,gap=3]{showbi}
\tkzSetUpCompass[color=teal,line width=.3 pt] 
\tkzDefPoints{0/1/A, 8/3/B, 3/6/C} 
\tkzDrawPolygon(A,B,C) 
\tkzDefLine[bisector](B,A,C) \tkzGetPoint{a} 
\tkzDefLine[bisector](C,B,A) \tkzGetPoint{b} 
\tkzShowLine[showbi](B,A,C) 
\tkzShowLine[showbi](C,B,A) 
\tkzInterLL(A,a)(B,b) \tkzGetPoint{I} 
\tkzDefPointBy[projection= onto A--B](I)
\tkzGetPoint{H} 
\tkzDrawCircle[new](I,H) 
\tkzDrawSegments[new](I,H) 
\tkzDrawLines[add=0 and .2,new](A,I B,I)
\end{tikzpicture}
\end{tkzexample} 





\section{标签样式}

\subsection{宏 \tkzcname{tkzSetUpLabel} }
宏 \tkzcname{tkzSetUpLabel} 定义点标签的样式。
\begin{NewMacroBox}{tkzSetUpStyle}{\oarg{local options}}%
  选项与 \TIKZ\ 的相同
\end{NewMacroBox}

\subsubsection{\tkzcname{tkzSetUpLabel} 的用法}
\begin{tkzexample}[latex=6cm,small]
\begin{tikzpicture}[scale=.75]
  \tkzSetUpLabel[font=\scriptsize,red]
  \tkzSetUpStyle[line width=1pt,teal]{XY}
  \tkzInit[xmin=-3,xmax=3,ymin=-3,ymax=3]
  \tkzDrawX[noticks,XY]
  \tkzDrawY[noticks,XY]
  \tkzDefPoints{1/0/A,0/1/B,-1/0/C,0/-1/D}
  \tkzDrawPoints[teal,fill=teal!30,size=6](A,...,D)
  \tkzLabelPoint[above right](A){$A(1,0)$}
  \tkzLabelPoint[above right](B){$B(0,1)$}
  \tkzLabelPoint[above left](C){$C(-1,0)$}
  \tkzLabelPoint[below left](D){$D(0,-1)$}
\end{tikzpicture}
\end{tkzexample}


\section{自定义样式}
你可以用 \tkzcname{tkzSetUpStyle} 定义自己的样式。

\subsection{宏 \tkzcname{tkzSetUpStyle} }
\begin{NewMacroBox}{tkzSetUpStyle}{\oarg{local options}}%
  选项与 \TIKZ\ 的相同
\end{NewMacroBox}

\subsubsection{\tkzcname{tkzSetUpStyle} 的用法}
\begin{tkzexample}[latex=2cm,small]
\begin{tikzpicture}
  \tkzSetUpStyle[color=blue!20!black,fill=blue!20]{mystyle}
  \tkzDefPoint(0,0){O}
  \tkzDefPoint(0,1){A}
  \tkzDrawPoints(O) % general style
  \tkzDrawPoints[mystyle,size=4](A) % my style
  \tkzLabelPoints(O,A)
\end{tikzpicture}
\end{tkzexample}

\section{如何使用 \tkzname{arrows}}

在某些国家，箭头用于表示直线平行，
或表示射线，例如角的边。

下面是一些如何放置这些箭头的示例。
\tkzname{tkz-euclide} 会加载库 \tkzname{arrows.meta}。
 
|\usetikzlibrary{arrows.meta}|

该库提供不同样式的箭头；下例使用其中的一些。

\subsection{在直线、射线或线段端点处的箭头}
\tkzname{Stealth}、\tkzname{Triangle}、\tkzname{To}、\tkzname{Latex}、\dots，可以与 \tkzname{reversed} 组合。在一个或两个端点处放置箭头非常容易。

\begin{enumerate}
\item \tkzname{-Triangle} 与 \tkzname{Segment}
 \begin{tkzexample}[latex=6cm,small]
    \begin{tikzpicture}
      \tkzDefPoints{0/0/A,4/0/B}
      \tkzDrawSegment[-Triangle](A,B)
    \end{tikzpicture}
  \end{tkzexample}
\item \tkzname{Stealth-Stealth} 与 \tkzname{Segment}
  \begin{tkzexample}[latex=6cm,small]
    \begin{tikzpicture}
      \tkzDefPoints{0/0/A,4/0/B}
      \tkzDrawSegment[Stealth-Stealth](A,B)
    \end{tikzpicture}
  \end{tkzexample}
\item \tkzname{Latex-Latex}   与 \tkzname{Line}
  \begin{tkzexample}[latex=6cm,small]
    \begin{tikzpicture}
      \tkzDefPoints{0/0/A,4/0/B}
      \tkzDrawLine[red,Latex-Latex](A,B)
      \tkzDrawPoints(A,B)
    \end{tikzpicture}
  \end{tkzexample}
\item \tkzname{To-To} 与 \tkzname{Segment}
  \begin{tkzexample}[latex=6cm,small]
    \begin{tikzpicture}
      \tkzDefPoints{0/0/A,4/0/B}
      \tkzDrawSegment[To-To](A,B)
    \end{tikzpicture}
  \end{tkzexample}
\item \tkzname{Latex-Latex} 与 \tkzname{Segment}
  \begin{tkzexample}[latex=6cm,small]
    \begin{tikzpicture}
      \tkzDefPoints{0/0/A,4/0/B}
      \tkzDrawSegment[Latex-Latex](A,B)
    \end{tikzpicture}
  \end{tkzexample}
\item \tkzname{Latex-}  与 \tkzname{Segment}
  \begin{tkzexample}[latex=6cm,small]
    \begin{tikzpicture}
      \tkzDefPoints{0/0/A,4/0/B}
      \tkzDrawSegment[Latex-](A,B)
    \end{tikzpicture}
  \end{tkzexample}
\item \tkzname{-Latex}  与 \tkzname{Segments}
  \begin{tkzexample}[latex=6cm,small]
\begin{tikzpicture}
 \tkzDefPoints{0/0/A,4/0/B,5/-2/C}
 \tkzDrawSegments[-Latex](A,B A,C)
\end{tikzpicture}
\end{tkzexample}
\end{enumerate}

\subsubsection{缩放箭头头部}

\begin{tkzexample}[latex=6cm,small]
\begin{tikzpicture}
 \tkzDefPoints{0/0/A,4/0/B}
 \tkzDrawSegment[{Latex[scale=2]}-{Latex[scale=2]}](A,B)
\end{tikzpicture}
\end{tkzexample}

\subsubsection{使用 vector 样式}
|\tikzset{vector style/.style={>=Latex,->}}|

你可以重新定义这个样式。
\begin{tkzexample}[latex=6cm,small]
\begin{tikzpicture}
 \tkzDefPoints{0/0/A,4/0/B}
\tkzDrawSegment[vector style](A,B)
\end{tikzpicture}
\end{tkzexample}
  
\subsection{在线段中点处的箭头}

直线上的箭头表示这些直线平行。这取决于国家：在法国，我们更倾向于在图形外部写出 $(AB) \parallel (DC)$。以下代码改编自 \tkzname{Muzimuzhi Z} 在 \href{https://tex.stackexchange.com/questions/632596/how-to-manage-argument-pattern-keys-and-subways}{tex.stackexchange.com} 上的回答。

\medskip
 语法：\\

 \begin{itemize}
\item |tkz arrow|（默认为 \tkzname{Latex}）
\item |tkz arrow=<arrow end tip>|
\item |tkz arrow=<arrow end tip> at <pos> (<pos> = .5 by default)|
\item |tkz arrow={<arrow end tip>[<arrow options>] at <pos>}| 可使用选项 \tkzname{scale}
 \end{itemize}

用法示例：\\

|\tkzDrawSegment[tkz arrow=Stealth] (A,B)|\\
|\tkzDrawSegment[tkz arrow={To[scale=3] at .4}](A,B)|\\
|\tkzDrawSegment[tkz arrow={Latex[scale=5,blue] at .6}](A,B)|

\subsubsection{在平行四边形中}
\begin{tkzexample}[latex=7cm,small]
\begin{tikzpicture}
 \tkzDefPoints{0/0/A,3/0/B,4/2/C} 
 \tkzDefParallelogram(A,B,C) 
 \tkzGetPoint{D}
 \tkzDrawSegments[tkz arrow](A,B D,C)
 \tkzDrawSegments(B,C D,A)
 \tkzLabelPoints(A,B) 
 \tkzLabelPoints[above right](C,D)
 \tkzDrawPoints(A,...,D)
\end{tikzpicture}
\end{tkzexample}

\subsubsection{与另一条直线平行的直线}
\begin{tkzexample}[latex=7cm,small]
\begin{tikzpicture}
 \tkzDefPoints{0/0/A,3/0/B,1/2/C} 
 \tkzDefPointWith[colinear= at C](A,B) 
 \tkzGetPoint{D}
 \tkzDrawSegments[tkz arrow=Triangle](A,B C,D)
 \tkzLabelPoints(A,B,C) 
 \tkzDrawPoints(A,...,C)
\end{tikzpicture}
\end{tkzexample}

\subsubsection{圆上的箭头}
你可以在圆的第一象限处放置一个箭头；旋转可以移动该箭头。
\begin{tkzexample}[latex=7cm,small]
\begin{tikzpicture}
\tkzDefPoints{0/0/A,3/0/B} 
\begin{scope}[rotate=150]
 \tkzDrawCircle[tkz arrow={Latex[scale=2,red]}](A,B)
\end{scope}
\end{tikzpicture}
\end{tkzexample}

\subsection{在多边形的所有边上放置箭头}
一些用户向我询问一种在多边形每条边上放置箭头的方法。我使用了 Paul Gaborit 在 \href{https://tex.stackexchange.com/questions/3161/tikz-how-to-draw-an-arrow-in-the-middle-of-the-line}{tex.stackexchange.com} 上提出的样式。

|\tikzset{tkz arrows/.style=|\\
|{postaction={on each path={tkz arrow={Latex[scale=2,color=black]}}}}}|   

你可以改变这个样式。使用 \tkzname{tkz arrows} 你可以在多边形的每条线段上放置一个箭头。

\subsubsection{用 \tkzname{tkz arrows} 在每条线段上放置箭头} 
\begin{tkzexample}[latex=7cm,small]
\begin{tikzpicture}
 \tkzDefPoints{0/0/A,3/0/B}  
 \tkzDefSquare(A,B) \tkzGetPoints{C}{D}
 \tkzDrawPolygon[tkz arrows](A,...,D)
\end{tikzpicture}
\end{tkzexample}

\subsubsection{将 \tkzname{tkz arrows} 用于圆}
\begin{tkzexample}[latex=7cm,small]
\begin{tikzpicture}
 \tkzDefPoints{0/0/A,3/0/B} 
 \tkzDrawCircle[tkz arrows](A,B)
\end{tikzpicture}
\end{tkzexample}

\endinput


\part{示例}
\section{不同作者}

\subsection{Andrew Swan 的代码}

\begin{tkzexample}[latex=7cm]
\begin{tikzpicture}[scale=1.25]
\def\radius{4}
\def\angle{40}
\pgfmathsetmacro{\htan}{tan(\angle)}
\tkzDefPoint(0,0){A} \tkzDefPoint(0,\radius){F}
\tkzDefPoint(\radius,0){B}
\tkzDefPointBy[rotation= center A angle \angle](B)
\tkzGetPoint{C}
\tkzDefLine[perpendicular=through B,K=1](A,B)
\tkzGetPoint{b}
\tkzInterLL(A,C)(B,b) \tkzGetPoint{D}
\tkzDefLine[perpendicular=through C,K=-1](A,B)
\tkzGetPoint{c}
\tkzInterLL(C,c)(A,B) \tkzGetPoint{E}
\tkzDrawSector[fill=blue,opacity=0.1](A,B)(C)
\tkzDrawArc[thin](A,B)(F)
\tkzMarkAngle(B,A,C)
\tkzLabelAngle[pos=0.8](B,A,C){$x$}
\tkzDrawPolygon(A,B,D)
\tkzDrawSegments(C,B)
\tkzDrawSegments[dashed,thin](C,E)
\tkzLabelPoints[below left](A)
\tkzLabelPoints[below right](B)
\tkzLabelPoints[above](C)
\tkzLabelPoints[above right](D)
\begin{scope}[/pgf/decoration/raise=5pt]
\draw [decorate,decoration={brace,mirror,
   amplitude=10pt},xshift=0pt,yshift=-4pt]
(A) -- (B) node [black,midway,yshift=-20pt]
{\footnotesize $1$};
\draw [decorate,decoration={brace,amplitude=10pt},
       xshift=4pt,yshift=0pt]
(D) -- (B) node [black,midway,xshift=27pt]
{\footnotesize $\tan x$};
\draw [decorate,decoration={brace,amplitude=10pt},
       xshift=4pt,yshift=0pt]
(E) -- (C) node [black,midway,xshift=-27pt]
{\footnotesize $\sin x$};
\end{scope}
\end{tikzpicture}
\end{tkzexample}


\subsection{Dimitris Kapetas 的示例}

在这个示例中，你需要为 \tkzcname{tkzMarkAngle} 设置 \tkzname{mkpos=.2}，因为角 $\widehat{CAM}$ 太小了。
另一种方法是使用 \tkzcname{tkzFillAngle}。


\begin{tkzexample}[latex=7cm,small]
\begin{tikzpicture}[scale=1]
  \tkzDefPoint(0,0){O}
  \tkzDefPoint(2.5,0){N}
  \tkzDefPoint(-4.2,0.5){M}
  \tkzDefPointBy[rotation=center O angle 30](N)
  \tkzGetPoint{B}
  \tkzDefPointBy[rotation=center O angle -50](N)
  \tkzGetPoint{A}
  \tkzInterLC[common=B](M,B)(O,B) \tkzGetFirstPoint{C}
  \tkzInterLC[common=A](M,A)(O,A) \tkzGetFirstPoint{A'}
  \tkzMarkAngle[mkpos=.2, size=0.5](A,C,B)
  \tkzMarkAngle[mkpos=.2, size=0.5](A,M,C)
  \tkzDrawSegments(A,C M,A M,B A,B)
  \tkzDrawCircle(O,N)
  \tkzLabelCircle[above left](O,N)(120){%
                 $\mathcal{C}$}
  \begin{scope}[veclen] 
  \tkzMarkAngle[mkpos=.2, size=1.2](C,A,M)
  \end{scope}
  \tkzDrawPoints(O, A, B, M, B, C)
  \tkzLabelPoints[right](O,A,B)
  \tkzLabelPoints[above left](M,C)
  \tkzLabelPoint[below left](A'){$A'$}
\end{tikzpicture}
\end{tkzexample}


\subsection{John Kitzmiller 的示例}
证明 $\dfrac{AC}{CE}=\dfrac{BD}{DF}$。

John 提供的另一个有趣示例，展示了如何将 \TIKZ\ 选项（如 \tkzname{decoration} 和 \tkzname{postaction}）与 \tkzname{tkz-euclide} 配合使用。

\begin{tkzexample}[vbox,small]
\begin{tikzpicture}[scale=1.5,decoration={markings,
  mark=at position 3cm with {\arrow[scale=2]{>}}}]
  \tkzDefPoints{0/0/E, 6/0/F, 0/1.8/P, 6/1.8/Q, 0/3/R, 6/3/S}
  \tkzDrawLines[postaction={decorate}](E,F P,Q R,S)
  \tkzDefPoints{3.5/3/A, 5/3/B}
  \tkzDrawSegments(E,A F,B)
  \tkzInterLL(E,A)(P,Q) \tkzGetPoint{C}
  \tkzInterLL(B,F)(P,Q) \tkzGetPoint{D}
  \tkzLabelPoints[above right](A,B)
  \tkzLabelPoints[below](E,F)
  \tkzLabelPoints[above left](C)
  \tkzDrawSegments[style=dashed](A,F)
  \tkzInterLL(A,F)(P,Q) \tkzGetPoint{G}
  \tkzLabelPoints[above right](D,G)
  \tkzDrawSegments[color=teal, line width=3pt, opacity=0.4](A,C A,G)
  \tkzDrawSegments[color=magenta, line width=3pt, opacity=0.4](C,E G,F)
  \tkzDrawSegments[color=teal, line width=3pt, opacity=0.4](B,D)
  \tkzDrawSegments[color=magenta, line width=3pt, opacity=0.4](D,F)
\end{tikzpicture}
\end{tkzexample}


\subsection{示例 1：来自印度尼西亚}

\begin{tkzexample}[vbox,small]
\begin{tikzpicture}[scale=3]
   \tkzDefPoints{0/0/A,2/0/B}
   \tkzDefSquare(A,B) \tkzGetPoints{C}{D}
   \tkzDefPointBy[rotation=center D angle 45](C)\tkzGetPoint{G}
   \tkzDefSquare(G,D)\tkzGetPoints{E}{F}
   \tkzInterLL(B,C)(E,F)\tkzGetPoint{H}
   \tkzFillPolygon[gray!10](D,E,H,C,D)
   \tkzDrawPolygon(A,...,D)\tkzDrawPolygon(D,...,G)
   \tkzDrawSegment(B,E)
   \tkzMarkSegments[mark=|,size=3pt,color=gray](A,B B,C C,D D,A E,F F,G G,D D,E)
   \tkzMarkSegments[mark=||,size=3pt,color=gray](B,E E,H)
   \tkzLabelPoints[left](A,D)
   \tkzLabelPoints[right](B,C,F,H)
   \tkzLabelPoints[above](G)\tkzLabelPoints[below](E)
   \tkzMarkRightAngles(D,A,B D,G,F)
\end{tikzpicture}
\end{tkzexample}

\subsection{示例 2：来自印度尼西亚}
\begin{tkzexample}[vbox,overhang,small]
  \begin{tikzpicture}[pol/.style={fill=brown!40,opacity=.2},
      seg/.style={tkzdotted,color=gray}, hidden pt/.style={fill=gray!40},
       mra/.style={color=gray!70,tkzdotted,/tkzrightangle/size=.2},scale=2]
  \tkzDefPoints{0/0/A,2.5/0/B,1.33/0.75/D,0/2.5/E,2.5/2.5/F}
  \tkzDefLine[parallel=through D](A,B)  \tkzGetPoint{I1}
  \tkzDefLine[parallel=through B](A,D)  \tkzGetPoint{I2}
  \tkzInterLL(D,I1)(B,I2)               \tkzGetPoint{C}
  \tkzDefLine[parallel=through E](A,D)  \tkzGetPoint{I3}
  \tkzDefLine[parallel=through D](A,E)  \tkzGetPoint{I4}
  \tkzInterLL(E,I3)(D,I4)               \tkzGetPoint{H}
  \tkzDefLine[parallel=through F](E,H)  \tkzGetPoint{I5}
  \tkzDefLine[parallel=through H](E,F)  \tkzGetPoint{I6}
  \tkzInterLL(F,I5)(H,I6)               \tkzGetPoint{G}
  \tkzDefMidPoint(G,H) \tkzGetPoint{P}  \tkzDefMidPoint(G,C) \tkzGetPoint{Q}
  \tkzDefMidPoint(B,C) \tkzGetPoint{R}  \tkzDefMidPoint(A,B) \tkzGetPoint{S}
  \tkzDefMidPoint(A,E) \tkzGetPoint{T}  \tkzDefMidPoint(E,H) \tkzGetPoint{U}
  \tkzDefMidPoint(A,D) \tkzGetPoint{M}  \tkzDefMidPoint(D,C) \tkzGetPoint{N}
  \tkzInterLL(B,D)(S,R)\tkzGetPoint{L} \tkzInterLL(H,F)(U,P) \tkzGetPoint{K}
  \tkzDefLine[parallel=through K](D,H)  \tkzGetPoint{I7}
  \tkzInterLL(K,I7)(B,D)                \tkzGetPoint{O}
  \tkzFillPolygon[pol](P,Q,R,S,T,U)
  \tkzDrawSegments[seg](K,O K,L P,Q R,S T,U C,D H,D A,D M,N B,D)
  \tkzDrawSegments(E,H B,C G,F G,H G,C Q,R S,T U,P H,F)
  \tkzDrawPolygon(A,B,F,E)
  \tkzDrawPoints(A,B,C,E,F,G,H,P,Q,R,S,T,U,K) \tkzDrawPoints[hidden pt](M,N,O,D)
  \tkzMarkRightAngle[mra](L,O,K)
  \tkzMarkSegments[mark=|,size=1pt,thick,color=gray](A,S B,S B,R C,R
                    Q,C Q,G G,P H,P E,U H,U E,T A,T)
  \tkzLabelAngle[pos=.3](K,L,O){$\alpha$}
  \tkzLabelPoints[below](O,A,S,B)    \tkzLabelPoints[above](H,P,G)
  \tkzLabelPoints[left](T,E)         \tkzLabelPoints[right](C,Q)
  \tkzLabelPoints[above left](U,D,M) \tkzLabelPoints[above right](L,N)
  \tkzLabelPoints[below right](F,R)  \tkzLabelPoints[below left](K)
\end{tikzpicture}
\end{tkzexample}
\newpage

\subsection{Nicolas François 绘制的 Morley 定理示意图}
\begin{tkzexample}[vbox,small]
  \begin{tikzpicture}
    \tkzInit[ymin=-3,ymax=5,xmin=-5,xmax=7]
    \tkzClip
    \tkzDefPoints{-2.5/-2/A,2/4/B,5/-1/C}
    \tkzFindAngle(C,A,B) \tkzGetAngle{anglea}
    \tkzDefPointBy[rotation=center A angle 1*\anglea/3](C) \tkzGetPoint{TA1}
    \tkzDefPointBy[rotation=center A angle 2*\anglea/3](C) \tkzGetPoint{TA2}
    \tkzFindAngle(A,B,C) \tkzGetAngle{angleb}
    \tkzDefPointBy[rotation=center B angle 1*\angleb/3](A) \tkzGetPoint{TB1}
    \tkzDefPointBy[rotation=center B angle 2*\angleb/3](A) \tkzGetPoint{TB2}
    \tkzFindAngle(B,C,A) \tkzGetAngle{anglec}
    \tkzDefPointBy[rotation=center C angle 1*\anglec/3](B) \tkzGetPoint{TC1}
    \tkzDefPointBy[rotation=center C angle 2*\anglec/3](B) \tkzGetPoint{TC2}
    \tkzInterLL(A,TA1)(B,TB2) \tkzGetPoint{U1}
    \tkzInterLL(A,TA2)(B,TB1) \tkzGetPoint{V1}
    \tkzInterLL(B,TB1)(C,TC2) \tkzGetPoint{U2}
    \tkzInterLL(B,TB2)(C,TC1) \tkzGetPoint{V2}
    \tkzInterLL(C,TC1)(A,TA2) \tkzGetPoint{U3}
    \tkzInterLL(C,TC2)(A,TA1) \tkzGetPoint{V3}
    \tkzDrawPolygons(A,B,C U1,U2,U3 V1,V2,V3)
    \tkzDrawLines[add=2 and 2,very thin,dashed](A,TA1 B,TB1 C,TC1 A,TA2 B,TB2 C,TC2)
    \tkzDrawPoints(U1,U2,U3,V1,V2,V3)
    \tkzLabelPoint[left](V1){$s_a$} \tkzLabelPoint[right](V2){$s_b$}
    \tkzLabelPoint[below](V3){$s_c$} \tkzLabelPoint[above left](A){$A$}
    \tkzLabelPoints[above right](B,C) \tkzLabelPoint(U1){$t_a$}
    \tkzLabelPoint[below left](U2){$t_b$} \tkzLabelPoint[above](U3){$t_c$}
  \end{tikzpicture}
  \end{tkzexample}

\newpage
\subsection{赵爽的勾股定理（毕达哥拉斯定理）}
\begin{tikzpicture}
\node [mybox,title={赵爽的勾股定理 / 毕达哥拉斯定理}] (box){%
\begin{minipage}{0.90\textwidth}
  {\emph{毕达哥拉斯并非世界上最早发现这一定理的人。古代中国远比他更早发现了这一定理。因此，在中国，毕达哥拉斯定理还有另一个名字——勾股定理。
赵爽是中国古代数学家。他重新发现了"勾股定理"，这实际上是"毕达哥拉斯定理"的中国版本。赵爽使用了一种名为"出入相补原理"的方法，创造了一幅"勾股圆方图"。
下图用于说明"勾股定理"的证明。（代码来自 Nan Geng）
}} 
\end{minipage}
};
\end{tikzpicture}
  
\begin{tkzexample}[latex=7cm,small]
\begin{tikzpicture}[scale=.8]
  \tkzDefPoint(0,0){A} \tkzDefPoint(4,0){A'}
  \tkzInterCC[R](A, 5)(A', 3)
  \tkzGetSecondPoint{B}
  \tkzDefSquare(A,B)   \tkzGetPoints{C}{D}
  \tkzCalcLength(A,A') \tkzGetLength{lA}
  \tkzCalcLength(A',B) \tkzGetLength{lB}
  \pgfmathparse{\lA-\lB}
  \tkzInterLC[R](A,A')(A',\pgfmathresult)
  \tkzGetFirstPoint{D'}
  \tkzDefSquare(D',A')\tkzGetPoints{B'}{C'}
  \tkzDefLine[orthogonal=through D](D,D') 
   \tkzGetPoint{d}
  \tkzDefLine[orthogonal=through A](A,A')
   \tkzGetPoint{a}
  \tkzDefLine[orthogonal=through C](C,C') 
   \tkzGetPoint{c}  
  \tkzInterLL(D,d)(C,c) \tkzGetPoint{E}
  \tkzInterLL(D,d)(A,a) \tkzGetPoint{F}
  \tkzDefSquare(E,F)\tkzGetPoints{G}{H}
  \tkzDrawPolygons[fill=teal!10](A,B,A' B,C,B'
     C,D,C' A,D',D)  
  \tkzDrawPolygons(A,B,C,D E,F,G,H)
  \tkzDrawPolygon[fill=green!10](A',B',C',D')
  \tkzDrawSegment[dim={$a$,-10pt,}](D,C')
  \tkzDrawSegment[dim={$b$,-10pt,}](C,C')
  \tkzDrawSegment[dim={$c$,-10pt,}](C,D)
  \tkzDrawPoints[size=2](A,B,C,D,A',B',C',D')
  \tkzLabelPoints[left](A)
  \tkzLabelPoints[below](B)
  \tkzLabelPoints[right](C)
  \tkzLabelPoints[above](D)
  \tkzLabelPoints[right](A')
  \tkzLabelPoints[below right](B')
  \tkzLabelPoints[below left](C') 
  \tkzLabelPoints[below](D')
 \end{tikzpicture}
\end{tkzexample}

\newpage
\subsection{Reuleaux 三角形}
\begin{tikzpicture}
\node [mybox,title={Stefan Kottwitz 的 Reuleaux 三角形}] (box){%
\begin{minipage}{0.90\textwidth}
  {\emph{数学中一个众所周知的经典领域是几何。
你可能在学校学过欧几里得几何，使用圆规和直尺进行
作图。数学老师可能对
绘制几何作图和讲解非常感兴趣。基础
作图可以帮助我们处理一般绘图中需要
直线与圆的交点和切线的情形，即使
它看起来不像是几何。
所以，在这里，我们将回忆学校里的几何绘图。
我们将使用 tkz-euclide 宏包，它基于 TikZ 工作。
我们将构造一个等边三角形。
然后将其扩展得到一个 Reuleaux 三角形，并添加标注。
该代码在 LaTeX Cookbook 第 10 章
高等数学、绘制几何图形中有完整解释。
 Stefan Kottwitz
}} 
\end{minipage}
};
\end{tikzpicture}

\begin{tkzexample}[vbox,small]
  \begin{tikzpicture}
    \tkzDefPoint(0,0){A} \tkzDefPoint(4,1){B}
    \tkzInterCC(A,B)(B,A) \tkzGetPoints{C}{D}
    \tkzInterLC(A,B)(B,A) \tkzGetPoints{F}{E}
    \tkzDrawCircles[dashed](A,B B,A)
    \tkzDrawPolygons(A,B,C A,E,D)
    \tkzCompasss[color=red, very thick](A,C B,C A,D B,D)
    \begin{scope}
      \tkzSetUpArc[thick,delta=0]
      \tkzDrawArc[fill=blue!10](A,B)(C)
      \tkzDrawArc[fill=blue!10](B,C)(A)
      \tkzDrawArc[fill=blue!10](C,A)(B)
    \end{scope}
    \tkzMarkAngles(D,A,E A,E,D)
    \tkzFillAngles[fill=yellow,opacity=0.5](D,A,E A,E,D) 
    \tkzMarkRightAngle[size=0.65,fill=red!20,opacity=0.2](A,D,E) 
    \tkzLabelAngle[pos=0.7](D,A,E){$\alpha$}
    \tkzLabelAngle[pos=0.8](A,E,D){$\beta$}
    \tkzLabelAngle[pos=0.5,xshift=-1.4mm](A,D,D){$90^\circ$}
    \begin{scope}[font=\small]
      \tkzLabelSegment[below=0.6cm,align=center](A,B){Reuleaux\\triangle}
      \tkzLabelSegment[above right,sloped](A,E){hypotenuse}
      \tkzLabelSegment[below,sloped](D,E){opposite}
      \tkzLabelSegment[below,sloped](A,D){adjacent}
      \tkzLabelSegment[below right=4cm](A,E){Thales circle}
    \end{scope}
    \tkzLabelPoints[below left](A)
    \tkzLabelPoints(B,D)
    \tkzLabelPoint[above](C){$C$}
    \tkzLabelPoints(E)
    \tkzDrawPoints(A,...,E)

  \end{tikzpicture}
\end{tkzexample}




\endinput

\section{一些有趣的示例}

\subsection{整数的平方根}
\begin{tikzpicture}
\node [mybox,title={整数的平方根}] (box){%
\begin{minipage}{0.90\textwidth}
  {\emph{如何用直尺和圆规作出 $1$、$\sqrt{2}$、$\sqrt{3}$。
}} 
\end{minipage}
};
\end{tikzpicture}

\begin{tkzexample}[latex=7cm,small]
\begin{tikzpicture}
  \tkzDefPoint(0,0){O}
  \tkzDefPoint(1,0){a0}
   \tkzDrawSegment(O,a0)
  \foreach \i [count=\j] in {0,...,16}{%
    \tkzDefPointWith[orthogonal normed](a\i,O)
    \tkzGetPoint{a\j}
       \pgfmathsetmacro{\c}{5*\i} 
    \tkzDrawPolySeg[fill=teal!\c](a\i,a\j,O)}
 \end{tikzpicture}
\end{tkzexample}

\subsection{关于直角三角形}
\begin{tikzpicture}
\node [mybox,title={关于直角三角形}] (box){%
\begin{minipage}{0.90\textwidth}
  {\emph{给定线段 $[AB]$，试确定一点 $C$，使 $AC=8$~cm 且 $ABC$ 在 $B$ 处为直角三角形。
}} 
\end{minipage}
};
\end{tikzpicture}

\begin{tkzexample}[latex=7cm,small]
\begin{tikzpicture}[scale=.5]
  \tkzDefPoint["$A$" left](2,1){A}
  \tkzDefPoint["$B$" right](6,4){B}
  \tkzDefPointWith[orthogonal,K=-1](B,A)
  \tkzDrawLine[add = .5 and .5](B,tkzPointResult)
  \tkzInterLC[R](B,tkzPointResult)(A,8)
  \tkzGetPoints{J}{C}
  \tkzDrawSegment(A,B)
  \tkzDrawPoints(A,B,C)
  \tkzCompass(A,C)
  \tkzMarkRightAngle(A,B,C)
  \tkzDrawLine[color=gray,style=dashed](A,C)
  \tkzLabelPoint[above](C){$C$}
\end{tikzpicture}
\end{tkzexample}

\subsection{阿基米德}
\begin{tikzpicture}
\node [mybox,title={阿基米德}] (box){%
\begin{minipage}{0.90\textwidth}
  {\emph{这是一个古老的问题，由伟大的希腊数学家阿基米德证明。
下图给出一个以 $AB$ 为直径的半圆。过 $B$ 作半圆的一条切线，再过半圆上一点 $C$ 作另一条切线。将点 $C$ 投影到线段 $[AB]$ 上得到点 $D$。两条切线相交于点 $T$。证明直线 $(AT)$ 平分 $(CD)$。
}} 
\end{minipage}
};
\end{tikzpicture}

\begin{tkzexample}[]
\begin{tikzpicture}[scale=1]
  \tkzDefPoint(0,0){A}\tkzDefPoint(6,0){D}
  \tkzDefPoint(8,0){B}\tkzDefPoint(4,0){I}
  \tkzDefLine[orthogonal=through D](A,D)
  \tkzInterLC[R](D,tkzPointResult)(I,4) \tkzGetSecondPoint{C}
  \tkzDefLine[orthogonal=through C](I,C)    \tkzGetPoint{c}
  \tkzDefLine[orthogonal=through B](A,B)    \tkzGetPoint{b}
  \tkzInterLL(C,c)(B,b) \tkzGetPoint{T}
  \tkzInterLL(A,T)(C,D) \tkzGetPoint{P}
  \tkzDrawArc(I,B)(A)
  \tkzDrawSegments(A,B A,T C,D I,C) \tkzDrawSegment[new](I,C)
  \tkzDrawLine[add = 1 and 0](C,T)   \tkzDrawLine[add = 0 and 1](B,T)
  \tkzMarkRightAngle(I,C,T)
  \tkzDrawPoints(A,B,I,D,C,T)
  \tkzLabelPoints(A,B,I,D)  \tkzLabelPoints[above right](C,T)
  \tkzMarkSegment[pos=.25,mark=s|](C,D) \tkzMarkSegment[pos=.75,mark=s|](C,D)
\end{tikzpicture}
\end{tkzexample}

\newpage
\subsubsection{等面积的正方形与矩形；黄金分割}

\begin{tikzpicture}
\node [mybox,title={第二卷，命题 XI \_Euclid's Elements\_}] (box){%
\begin{minipage}{0.90\textwidth}
{\emph{作一个正方形和一个矩形，使二者面积相等。}
} 
\end{minipage}
};
\end{tikzpicture}% 

\begin{tkzexample}[vbox,small]
 \begin{tikzpicture}[scale=.75]
  \tkzDefPoint(0,0){D} \tkzDefPoint(8,0){A}
  \tkzDefSquare(D,A) \tkzGetPoints{B}{C}
  \tkzDefMidPoint(D,A) \tkzGetPoint{E}
  \tkzInterLC(D,A)(E,B)\tkzGetSecondPoint{F}
  \tkzInterLC[near](B,A)(A,F)\tkzGetFirstPoint{G}
  \tkzDefSquare(A,F)\tkzGetFirstPoint{H}
  \tkzInterLL(C,D)(H,G)\tkzGetPoint{I}
  \tkzFillPolygon[teal!10](I,G,B,C)
  \tkzFillPolygon[teal!10](A,F,H,G)
  \tkzDrawArc[angles](E,B)(0,120)
  \tkzDrawSemiCircle(A,F)
  \tkzDrawSegments(A,F E,B H,I F,H)
  \tkzDrawPolygons(A,B,C,D)
  \tkzDrawPoints(A,...,I)
  \tkzLabelPoints[below right](A,E,D,F,I)
  \tkzLabelPoints[above right](C,B,G,H)
 \end{tikzpicture}
\end{tkzexample}

\newpage

\subsubsection{施泰纳线与西姆松线}

\begin{tikzpicture}
\node [mybox,title={施泰纳线与西姆松线}] (box){%
\begin{minipage}{0.90\textwidth}
{\emph{考虑三角形 $ABC$ 及其外接圆上一点 $M$。$M$ 到三角形各边的垂足共线（施泰纳线）。在直线 $AB$、$AC$、$BC$ 上距 $M$ 最近的三个点共线，这就是西姆松线。
}} 
\end{minipage}
};
\end{tikzpicture}%

\begin{tkzexample}[latex=7cm,small]
\begin{tikzpicture}[scale=.75,rotate=-20]
  \tkzDefPoint(0,0){B} 
  \tkzDefPoint(2,4){A} \tkzDefPoint(7,0){C}
  \tkzDefCircle[circum](A,B,C)  
  \tkzGetPoint{O}
  \tkzDrawCircle(O,A) 
  \tkzCalcLength(O,A)  
  \tkzGetLength{rOA} 
  \tkzDefShiftPoint[O](40:\rOA){M}
  \tkzDefShiftPoint[O](60:\rOA){N}  
  \tkzDefTriangleCenter[orthic](A,B,C)
  \tkzGetPoint{H}
  \tkzDefSpcTriangle[orthic,name=H](A,B,C){a,b,c}
  \tkzDefPointsBy[reflection=over A--B](M,N){P,P'}
  \tkzDefPointsBy[reflection=over A--C](M,N){Q,Q'}
  \tkzDefPointsBy[reflection=over C--B](M,N){R,R'}
  \tkzDefMidPoint(M,P)\tkzGetPoint{I}
  \tkzDefMidPoint(M,Q)\tkzGetPoint{J}
  \tkzDefMidPoint(M,R)\tkzGetPoint{K} 
  \tkzDrawSegments[new](P,R M,P M,Q M,R N,P'
   N,Q' N,R' P',R' I,K)
  \tkzDrawPolygons(A,B,C)
  \tkzDrawPoints(A,B,C,H,M,N,P,Q,R,P',Q',R',I,J,K)
  \tkzLabelPoints(A,B,C,H,M,N,P,Q,R,P',Q',R',I,J,K)
\end{tikzpicture}
\end{tkzexample}

\newpage
\subsection{希波克拉底月牙}

\begin{tikzpicture}
\node [mybox,title={希波克拉底月牙}] (box){%
\begin{minipage}{0.90\textwidth}
  { \emph{来自维基百科：在几何学中，希波克拉底月牙以希俄斯的希波克拉底命名，是由两段圆弧围成的月牙形，其中较小的圆以大圆上张成直角的一条弦为直径。在第一个图中，月牙的面积等于三角形 $ABC$ 的面积。希俄斯的希波克拉底（古希腊数学家）
}} 
\end{minipage}
};
\end{tikzpicture}% 

\begin{tkzexample}[latex=7cm,small]
\begin{tikzpicture}
 \tkzInit[xmin=-2,xmax=5,ymin=-1,ymax=6]
 \tkzClip % allows you to define a bounding box 
   % large enough
  \tkzDefPoint(0,0){A}\tkzDefPoint(4,0){B}
  \tkzDefSquare(A,B) 
  \tkzGetFirstPoint{C} 
  \tkzDrawPolygon[fill=green!5](A,B,C)
   \begin{scope}
     \tkzClipCircle[out](B,A)
     \tkzDefMidPoint(C,A) \tkzGetPoint{M}
     \tkzDrawSemiCircle[fill=teal!5](M,C)
   \end{scope}
   \tkzDrawArc[delta=0](B,C)(A)
\end{tikzpicture}
\end{tkzexample}

\subsection{海什木月牙}

\begin{tikzpicture}
\node [mybox,title={海什木月牙}] (box){%
\begin{minipage}{0.90\textwidth}
  { \emph{来自维基百科：阿拉伯数学家海什木（Hasan Ibn al-Haytham，拉丁化名字为 Alhazen）证明，在直角三角形两条直角边上分别作两个月牙形，其外边界为半圆，内边界由三角形的外接圆构成，则这两个月牙的面积之和等于该三角形的面积。以这种方式由直角三角形作出的月牙形被称为海什木月牙。
}} 
\end{minipage}};
\end{tikzpicture}% 

\begin{tkzexample}[latex=7cm,small]
\begin{tikzpicture}[scale=.5,rotate=180]
  \tkzInit[xmin=-1,xmax=11,ymin=-4,ymax=7]
  \tkzClip
  \tkzDefPoints{0/0/A,8/0/B}
  \tkzDefTriangle[pythagore,swap](A,B) 
  \tkzGetPoint{C}
  \tkzDrawPolygon[fill=green!5](A,B,C)
  \tkzDefMidPoint(C,A) \tkzGetPoint{I}
  \begin{scope}
    \tkzClipCircle[out](I,A)
    \tkzDefMidPoint(B,A) \tkzGetPoint{x}
    \tkzDrawSemiCircle[fill=teal!5](x,A)
    \tkzDefMidPoint(B,C) \tkzGetPoint{y}
    \tkzDrawSemiCircle[fill=teal!5](y,B)
  \end{scope}
  \tkzSetUpCompass[/tkzcompass/delta=0]
      \tkzDefMidPoint(C,A) \tkzGetPoint{z}
  \tkzDrawSemiCircle(z,A)
\end{tikzpicture}
\end{tkzexample}

\newpage
\subsection{关于圆的裁剪}\label{About clipping circles}
\begin{tikzpicture}
\node [mybox,title={关于圆的裁剪}] (box){%
\begin{minipage}{0.90\textwidth}
  { \emph{问题在于 bounding box 的管理。首先需要定义一个矩形，图形将被放置在其中。这由前两行完成。
}} 
\end{minipage}
};
\end{tikzpicture}% 

\begin{tkzexample}[latex=7cm,small]
\begin{tikzpicture}
  \tkzInit[xmin=0,xmax=6,ymin=0,ymax=6]
  \tkzClip
  \tkzDefPoints{0/0/A, 6/0/B}
  \tkzDefSquare(A,B)      \tkzGetPoints{C}{D}
  \tkzDefMidPoint(A,B)        \tkzGetPoint{M}
  \tkzDefMidPoint(A,D)        \tkzGetPoint{N}
  \tkzDefMidPoint(B,C)        \tkzGetPoint{O}
  \tkzDefMidPoint(C,D)        \tkzGetPoint{P}
 \begin{scope}
  \tkzClipCircle[out](M,B) \tkzClipCircle[out](P,D)
  \tkzFillPolygon[teal!20](M,N,P,O)
 \end{scope}
 \begin{scope}
   \tkzClipCircle[out](N,A) \tkzClipCircle[out](O,C)
   \tkzFillPolygon[teal!20](M,N,P,O)
 \end{scope}
 \begin{scope}
   \tkzClipCircle(P,C) \tkzClipCircle(N,A)      
   \tkzFillPolygon[teal!20](N,P,D)
 \end{scope}
 \begin{scope}
     \tkzClipCircle(O,C) \tkzClipCircle(P,C) 
     \tkzFillPolygon[teal!20](P,C,O)
 \end{scope}
 \begin{scope}
     \tkzClipCircle(M,B)  \tkzClipCircle(O,B)
     \tkzFillPolygon[teal!20](O,B,M)
 \end{scope}
 \begin{scope}
     \tkzClipCircle(N,A) \tkzClipCircle(M,A)  
     \tkzFillPolygon[teal!20](A,M,N)
 \end{scope}
 \tkzDrawSemiCircles(M,B N,A O,C P,D)
 \tkzDrawPolygons(A,B,C,D M,N,P,O)
 \end{tikzpicture}
 \end{tkzexample}

\newpage
\subsection{相似等腰三角形}

\begin{tikzpicture}
\node [mybox,title={相似等腰三角形}] (box){%
\begin{minipage}{0.90\textwidth}
  { \emph{以下内容出自优秀网站 \textbf{Descartes et les Mathématiques}。我未修改原文，仅是图形编程的作者。
\url{http://debart.pagesperso-orange.fr/seconde/triangle.html}
}} 
\end{minipage}
};
\end{tikzpicture}% 

以下内容出自优秀网站 \textbf{Descartes et les Mathématiques}。我未修改原文；仅编写了图形程序。

\url{http://debart.pagesperso-orange.fr/seconde/triangle.html}

参考文献：

\begin{itemize}
\item   Géométrie au Bac - Tangente, 特刊第 8 期 - 习题 11，第 11 页

\item   Elisabeth Busser 与 Gilles Cohen：200 nouveaux problèmes du "Monde" - POLE 2007（《世界报》200 道新题）

\item   Affaire de logique n° 364 - Le Monde，2004 年 2 月 17 日
\end{itemize}


共给出两则题目，一则来自 \textit{Tangente} 杂志，另一则来自 \textit{Le Monde}。

\vspace*{2cm}
\emph{《Tangente》杂志编辑}：\textcolor{orange}{作两个相似的等腰三角形 $AXB$ 和 $BYC$，主顶点分别为 $X$ 和 $Y$，使 $A$、$B$、$C$ 三点共线，且这两个三角形为 ``indirect'' 型。设顶角 $\widehat{AXB}$ = $\widehat{BYC}$ 为 $\alpha$。再作第三个等腰三角形 $XZY$，与前两个相似，主顶点为 $Z$，同样为 ``indirect'' 型。证明点 $Z$ 在直线 $(AC)$ 上。}

\vspace*{2cm}
\emph{《Le Monde》编辑}：\textcolor{orange}{作两个相似的等腰三角形 $AXB$ 和 $BYC$，主顶点分别为 $X$ 和 $Y$，使 $A$、$B$、$C$ 三点共线，且这两个三角形为 ``indirect'' 型。设顶角 $\widehat{AXB}$ = $\widehat{BYC}$ 为 $\alpha$。线段 $[AC]$ 上的点 $Z$ 到两顶点 $X$ 和 $Y$ 的距离相等。\\
问：$Z$ 以多大的角度观察这两个顶点？}

\vspace*{2cm} 作法及其代码见接下来两页，但你可以先自行探索。两种情形中，程序都遵循我的思路（我想是这样……）。

\subsection{《Tangente》修订版}
\begin{tkzexample}[]
\begin{tikzpicture}[scale=.8,rotate=60]
  \tkzDefPoint(6,0){X}   \tkzDefPoint(3,3){Y}
  \tkzDefShiftPoint[X](-110:6){A}    \tkzDefShiftPoint[X](-70:6){B}
  \tkzDefShiftPoint[Y](-110:4.2){A'} \tkzDefShiftPoint[Y](-70:4.2){B'}
  \tkzDefPointBy[translation= from A' to B ](Y) \tkzGetPoint{Y}
  \tkzDefPointBy[translation= from A' to B ](B') \tkzGetPoint{C}
  \tkzInterLL(A,B)(X,Y) \tkzGetPoint{O}
  \tkzDefMidPoint(X,Y) \tkzGetPoint{I}
  \tkzDefPointWith[orthogonal](I,Y)
  \tkzInterLL(I,tkzPointResult)(A,B) \tkzGetPoint{Z}
  \tkzDefCircle[circum](X,Y,B) \tkzGetPoint{O}
  \tkzDrawCircle(O,X)
  \tkzDrawLines[add = 0 and 1.5](A,C) \tkzDrawLines[add = 0 and 3](X,Y)
  \tkzDrawSegments(A,X B,X B,Y C,Y)   \tkzDrawSegments[color=red](X,Z Y,Z)
  \tkzDrawPoints(A,B,C,X,Y,O,Z)
  \tkzLabelPoints(A,B,C,Z)   \tkzLabelPoints[above right](X,Y,O)
\end{tikzpicture}
\end{tkzexample}

\subsection{《Le Monde》版}

\begin{tkzexample}[]
\begin{tikzpicture}[scale=1.25]
  \tkzDefPoint(0,0){A}
  \tkzDefPoint(3,0){B}
  \tkzDefPoint(9,0){C}
  \tkzDefPoint(1.5,2){X}
  \tkzDefPoint(6,4){Y}
  \tkzDefCircle[circum](X,Y,B) \tkzGetPoint{O}
  \tkzDefMidPoint(X,Y)               \tkzGetPoint{I}
  \tkzDefPointWith[orthogonal](I,Y)  \tkzGetPoint{i}
  \tkzDrawLines[add = 2 and 1,color=orange](I,i)
  \tkzInterLL(I,i)(A,B)              \tkzGetPoint{Z}
  \tkzInterLC(I,i)(O,B)              \tkzGetFirstPoint{M}
  \tkzDefPointWith[orthogonal](B,Z)  \tkzGetPoint{b}
  \tkzDrawCircle(O,B)
  \tkzDrawLines[add = 0 and 2,color=orange](B,b)
  \tkzDrawSegments(A,X B,X B,Y C,Y A,C X,Y)
  \tkzDrawSegments[color=red](X,Z Y,Z)
  \tkzDrawPoints(A,B,C,X,Y,Z,M,I)
  \tkzLabelPoints(A,B,C,Z)
  \tkzLabelPoints[above right](X,Y,M,I)
\end{tikzpicture}
\end{tkzexample}

\subsection{三角形的高}

\begin{tikzpicture}
\node [mybox,title={三角形的高}] (box){%
\begin{minipage}{0.90\textwidth}
  { \emph{来自维基百科：以下内容同样出自优秀网站 \textbf{Descartes et les Mathématiques}。
\url{http://debart.pagesperso-orange.fr/geoplan/geometrie_triangle.html}。
三角形的三条高相交于同一点 $H$。
}} 
\end{minipage}
};
\end{tikzpicture}% 

\begin{tkzexample}[vbox,small]
\begin{tikzpicture}
   \tkzDefPoint(0,0){C} \tkzDefPoint(7,0){B}
   \tkzDefPoint(5,6){A}
   \tkzDefMidPoint(C,B) \tkzGetPoint{I}
   \tkzInterLC(A,C)(I,B)
   \tkzGetFirstPoint{B'}
   \tkzInterLC(A,B)(I,B)
   \tkzGetSecondPoint{C'}
   \tkzInterLL(B,B')(C,C') \tkzGetPoint{H}
   \tkzInterLL(A,H)(C,B) \tkzGetPoint{A'}
   \tkzDefCircle[circum](A,B',C') \tkzGetPoint{O}
   \tkzDrawArc(I,B)(C)
   \tkzDrawPolygon(A,B,C)
   \tkzDrawCircle[color=red](O,A)
   \tkzDrawSegments[color=orange](B,B' C,C' A,A')
   \tkzMarkRightAngles(C,B',B B,C',C C,A',A)
   \tkzDrawPoints(A,B,C,A',B',C',H)
   \tkzLabelPoints[above right](A,B',C',H)
   \tkzLabelPoints[below right](B,C,A')
\end{tikzpicture}
\end{tkzexample}

\subsection{高——另一种作法}

\begin{tkzexample}[vbox,small]
\begin{tikzpicture}
\tkzDefPoint(0,0){A} \tkzDefPoint(8,0){B} 
\tkzDefPoint(5,6){C} 
\tkzDefMidPoint(A,B)\tkzGetPoint{O} 
\tkzDefPointBy[projection=onto A--B](C) \tkzGetPoint{P}
\tkzInterLC[common=A](C,A)(O,A)
\tkzGetFirstPoint{M}
\tkzInterLC(C,B)(O,A)
\tkzGetSecondPoint{N}
\tkzInterLL(B,M)(A,N)\tkzGetPoint{I}
\tkzDefCircle[diameter](A,B)\tkzGetPoint{x}
\tkzDefCircle[diameter](I,C)\tkzGetPoint{y}
\tkzDrawCircles(x,A y,C)
\tkzDrawSegments(C,A C,B A,B B,M A,N)
\tkzMarkRightAngles[fill=brown!20](A,M,B A,N,B A,P,C)
\tkzDrawSegment[style=dashed,color=orange](C,P)
\tkzLabelPoints(O,A,B,P)
\tkzLabelPoint[left](M){$M$} 
\tkzLabelPoint[right](N){$N$} 
\tkzLabelPoint[above](C){$C$} 
\tkzLabelPoint[above right](I){$I$} 
\tkzDrawPoints[color=red](M,N,P,I) 
\tkzDrawPoints[color=brown](O,A,B,C)
\end{tikzpicture}
\end{tkzexample}

\newpage
\subsection{等边三角形中的三个圆}
\begin{tikzpicture}
\node [mybox,title={等边三角形中的三个圆}] (box){%
\begin{minipage}{0.90\textwidth}
  { \emph{来自维基百科：在几何学中，Malfatti 圆是给定三角形内的三个圆，使得每个圆都与另外两个圆以及三角形的两条边相切。它们以 Gian Francesco Malfatti 的名字命名，他最早研究了构造这三个圆的问题，并错误地认为这样能使三角形内三个互不相交的圆的总面积最大。下面研究一个特例：等边三角形与三个等圆。
}} 
\end{minipage}
};
\end{tikzpicture}% 
\begin{tkzexample}[latex=7cm,small]
\begin{tikzpicture}[scale=.8]
  \tkzDefPoints{0/0/A,8/0/B,0/4/a,8/4/b,8/8/c}
  \tkzDefTriangle[equilateral](A,B) \tkzGetPoint{C}
  \tkzDefMidPoint(A,B) \tkzGetPoint{M}
  \tkzDefMidPoint(B,C) \tkzGetPoint{N}
  \tkzDefMidPoint(A,C) \tkzGetPoint{P}
  \tkzInterLL(A,N)(M,a) \tkzGetPoint{Ia}
  \tkzDefPointBy[projection = onto A--B](Ia)
  \tkzGetPoint{ha}
  \tkzInterLL(B,P)(M,b) \tkzGetPoint{Ib}
  \tkzDefPointBy[projection = onto A--B](Ib)
  \tkzGetPoint{hb}
  \tkzInterLL(A,c)(M,C) \tkzGetPoint{Ic}
  \tkzDefPointBy[projection = onto A--C](Ic)
  \tkzGetPoint{hc}
  \tkzInterLL(A,Ia)(B,Ib) \tkzGetPoint{G}
  \tkzDefSquare(A,B) \tkzGetPoints{D}{E}
  \tkzDrawPolygon(A,B,C)
  \tkzClipBB
  \tkzDrawSemiCircles[gray,dashed](M,B A,M 
  A,B B,A G,Ia)
  \tkzDrawCircles[gray](Ia,ha Ib,hb Ic,hc)
  \tkzDrawPolySeg(A,E,D,B)
  \tkzDrawPoints(A,B,C,G,Ia,Ib,Ic)
  \tkzDrawSegments[gray,dashed](C,M A,N B,P
   M,a M,b A,a a,b b,B A,D Ia,ha)
\end{tikzpicture}
\end{tkzexample}

\newpage
\subsection{正弦定理}
\begin{tikzpicture}
\node [mybox,title={正弦定理}] (box){%
\begin{minipage}{0.90\textwidth}
  {来自维基百科：\emph{在三角学中，正弦定理（law of sines、sine law、sine formula 或 sine rule）是一个将三角形（任意形状）三边长度与其对角的正弦值联系起来的等式。
}} 
\end{minipage}
};
\end{tikzpicture}% 

\begin{tkzexample}[latex=7cm,small]
  \begin{tikzpicture}
  \tkzDefPoints{0/0/A,5/1/B,2/6/C}
  \tkzDefTriangleCenter[circum](A,B,C)
   \tkzGetPoint{O} 
  \tkzDefPointBy[symmetry= center O](B) 
   \tkzGetPoint{D} 
  \tkzDrawPolygon[color=brown](A,B,C)
  \tkzDrawCircle(O,A)
  \tkzDrawPoints(A,B,C,D,O)
  \tkzDrawSegments[dashed](B,D A,D)
  \tkzLabelPoint[left](D){$D$}
  \tkzLabelPoint[below](A){$A$}
  \tkzLabelPoint[above](C){$C$}
  \tkzLabelPoint[right](B){$B$}
  \tkzLabelPoint[below](O){$O$}
  \tkzLabelSegment(B,C){$a$}
  \tkzLabelSegment[left](A,C){$b$}
  \tkzLabelSegment(A,B){$c$}
  \end{tikzpicture}
\end{tkzexample}

在三角形 $ABC$ 中：

\begin{equation}
\frac{a}{\sin A} = \frac{b}{\sin B} =\frac{c}{\sin C}
\end{equation}

\[\widehat{C} = \widehat{D}\] 
\begin{equation}
\frac{c}{2R} = \sin D = \sin C 
\end{equation}

因此 \[ \frac{c}{\sin C} = 2R\]

\newpage
\subsection{生命之花}
\begin{tikzpicture}
\node [mybox,title={第四卷，命题 XI \_Euclid's Elements\_}] (box){%
\begin{minipage}{0.90\textwidth}
  {\emph{神圣几何可以被描述为一种信仰体系，它将宗教或文化价值赋予空间与时间的诸多基本形式。按照这一信仰体系，存在的基本模式被视为神圣的，因为沉思它们即是沉思万物的本源。通过研究这些形式的本质及其相互关系，人们可以寻求领悟宇宙的科学、哲学、心理学、美学与神秘法则。
生命之花被认为是神圣几何的象征，据说蕴含着描绘空间与时间基本形式的古老宗教价值。从这个意义上说，它是生命贯穿全人类之联系的视觉表达，有人相信它包含了一切生物基本信息的阿卡西记录。
}} 
\end{minipage}
};
\end{tikzpicture}% 

在埃及阿拜多斯的欧西里斯神庙中发现的精美圆形排列之一（Rawles 1997）。\\
Weisstein, Eric W. ``Flower of Life.'' 摘自 MathWorld--A Wolfram Web Resource。\\ \url{http://mathworld.wolfram.com/FlowerofLife.html}
 
\begin{tkzexample}[vbox,small]
\begin{tikzpicture}[scale=.75]
  \tkzSetUpLine[line width=2pt,color=teal!80!black]
  \tkzSetUpCompass[line width=2pt,color=teal!80!black]
   \tkzDefPoint(0,0){O}  \tkzDefPoint(2.25,0){A}
   \tkzDrawCircle(O,A)
\foreach \i in {0,...,5}{
   \tkzDefPointBy[rotation= center O angle 30+60*\i](A)\tkzGetPoint{a\i}
   \tkzDefPointBy[rotation= center {a\i} angle  120](O)\tkzGetPoint{b\i}
   \tkzDefPointBy[rotation= center {a\i} angle  180](O)\tkzGetPoint{c\i}
   \tkzDefPointBy[rotation= center {c\i} angle  120](a\i)\tkzGetPoint{d\i}
   \tkzDefPointBy[rotation= center {c\i} angle   60](d\i)\tkzGetPoint{f\i}
   \tkzDefPointBy[rotation= center {d\i} angle   60](b\i)\tkzGetPoint{e\i} 
   \tkzDefPointBy[rotation= center {f\i} angle   60](d\i)\tkzGetPoint{g\i} 
   \tkzDefPointBy[rotation= center {d\i} angle   60](e\i)\tkzGetPoint{h\i}
   \tkzDefPointBy[rotation= center {e\i} angle  180](b\i)\tkzGetPoint{k\i}   
   \tkzDrawCircle(a\i,O)
   \tkzDrawCircle(b\i,a\i)
   \tkzDrawCircle(c\i,a\i)
   \tkzDrawArc[rotate](f\i,d\i)(-120)
   \tkzDrawArc[rotate](e\i,d\i)(180)
   \tkzDrawArc[rotate](d\i,f\i)(180)
   \tkzDrawArc[rotate](g\i,f\i)(60)
   \tkzDrawArc[rotate](h\i,d\i)(60)
   \tkzDrawArc[rotate](k\i,e\i)(60) 
}
   \tkzClipCircle(O,f0)
\end{tikzpicture}
\end{tkzexample}


\newpage
\subsection{圆内接五边形}
\begin{tikzpicture}
\node [mybox,title={第四卷，命题 XI \_Euclid's Elements\_}] (box){%
\begin{minipage}{0.90\textwidth}
  {\emph{在给定圆内作一个等边且等角的五边形。
}} 
\end{minipage}
};
\end{tikzpicture}% 

\begin{tkzexample}[vbox,small]
\begin{tikzpicture}[scale=.75]
   \tkzDefPoint(0,0){O} 
   \tkzDefPoint(5,0){A}
   \tkzDefPoint(0,5){B}
   \tkzDefPoint(-5,0){C} 
   \tkzDefPoint(0,-5){D}
   \tkzDefMidPoint(A,O)             \tkzGetPoint{I}
   \tkzInterLC(I,B)(I,A)            \tkzGetPoints{F}{E}
   \tkzInterCC(O,C)(B,E)            \tkzGetPoints{D3}{D2}
   \tkzInterCC(O,C)(B,F)            \tkzGetPoints{D4}{D1}
   \tkzDrawArc[angles](B,E)(180,360)
   \tkzDrawArc[angles](B,F)(220,340)
   \tkzDrawLine[add=.5 and .5](B,I)
   \tkzDrawCircle(O,A)
   \tkzDefCircle[diameter](O,A)     \tkzGetPoint{x}
   \tkzDrawCircle(x,A)
   \tkzDrawSegments(B,D C,A) 
   \tkzDrawPolygon[new](D,D1,D2,D3,D4)
   \tkzDrawPoints(A,...,D,O)
   \tkzDrawPoints[new](E,F,I,D1,D2,D4,D3)
   \tkzLabelPoints[below left](A,...,D,O)
   \tkzLabelPoints[new,below right](I,E,F,D1,D2,D4,D3)  
\end{tikzpicture}
\end{tkzexample}

 \newpage
 \subsection{正方形内作五边形}
 \begin{tikzpicture}
 \node [mybox,title={正方形内作五边形}] (box){%
 \begin{minipage}{0.90\textwidth}
   { \emph{在给定正方形内作一个五边形。
 }} 
 \end{minipage}
 };
 \end{tikzpicture}%
    
\begin{tkzexample}[vbox,small]
\begin{tikzpicture}[scale=.75]
  \tkzDefPoints{0/0/O,-5/-5/A,5/-5/B}
  \tkzDefSquare(A,B)   \tkzGetPoints{C}{D}
  \tkzDefMidPoint(A,B) \tkzGetPoint{F}
  \tkzDefMidPoint(C,D) \tkzGetPoint{E}
  \tkzDefMidPoint(B,C) \tkzGetPoint{G}
  \tkzDefMidPoint(A,D) \tkzGetPoint{K}
  \tkzInterLC[near](C,D)(E,B)              \tkzGetFirstPoint{T}
  \tkzDefMidPoint(D,T)                     \tkzGetPoint{I}
  \tkzInterCC[with nodes](O,D,I)(E,D,I)    \tkzGetSecondPoint{H}
  \tkzInterLC(O,H)(O,E)                    \tkzGetSecondPoint{M}
  \tkzInterCC(O,E)(E,M)                    \tkzGetFirstPoint{Q}
  \tkzInterCC[with nodes](O,O,E)(Q,E,M)    \tkzGetFirstPoint{P}
  \tkzInterCC[with nodes](O,O,E)(P,E,M)    \tkzGetFirstPoint{N}
  \tkzCompasss(O,H E,H)
  \tkzDrawArc(E,B)(T)
  \tkzDrawPolygons[purple](A,B,C,D M,E,Q,P,N) 
  \tkzDrawCircle(O,E)
  \tkzDrawSegments(T,I O,H E,H E,F G,K)
  \tkzDrawPoints(T,M,Q,P,N,I)
  \tkzLabelPoints(A,B,O,N,P,Q,M,H)
  \tkzLabelPoints[above right](C,D,E,I,T)
\end{tikzpicture} 
\end{tkzexample}

\newpage
\subsection{内接六边形}
 \begin{tikzpicture}
 \node [mybox,title={内接六边形}] (box){%
 \begin{minipage}{0.90\textwidth}
   {\emph{在给定的等边三角形内作一个完美内切的正六边形（不留边距）。
 }} 
 \end{minipage}
 };
 \end{tikzpicture}%
 
\subsubsection{内接六边形：版本 1} % (fold)
\label{ssub:hexagon_inscribed_version_1}
\begin{tkzexample}[latex=7cm,small]
  \begin{tikzpicture}[scale=.5]
   \pgfmathsetmacro{\c}{6} 
   \tkzDefPoints{0/0/A,\c/0/B}
   \tkzDefTriangle[equilateral](A,B)\tkzGetPoint{C}
   \tkzDefTriangleCenter[centroid](A,B,C) 
   \tkzGetPoint{I}
   \tkzDefPointBy[homothety=center A ratio 1./3](B) 
   \tkzGetPoint{c1}
   \tkzInterLC(B,C)(I,c1) \tkzGetPoints{a1}{a2}
   \tkzInterLC(A,C)(I,c1) \tkzGetPoints{b1}{b2}
   \tkzInterLC(A,B)(I,c1) \tkzGetPoints{c1}{c2}
   \tkzDrawPolygon(A,B,C)
   \tkzDrawCircle[thin,orange](I,c1)
   \tkzDrawPolygon[red,thick](a2,a1,b2,b1,c2,c1)
 \end{tikzpicture} 
\end{tkzexample}
% subsubsection hexagon_inscribed_version_1 (end)

\subsubsection{内接六边形：版本 2} % (fold)
\label{ssub:hexagon_inscribed_version_2}
\begin{tkzexample}[latex=7cm,small]
\begin{tikzpicture}[scale=.5]
 \pgfmathsetmacro{\c}{6} 
 \tkzDefPoints{0/0/A,\c/0/B}
 \tkzDefTriangle[equilateral](A,B)\tkzGetPoint{C}
 \tkzDefTriangleCenter[centroid](A,B,C) 
 \tkzGetPoint{I}
 \tkzDefPointsBy[rotation= center I
                 angle 60](A,B,C){a,b,c}
 \tkzDrawPolygon[fill=teal!20,opacity=.5](A,B,C)
 \tkzDrawPolygon[fill=purple!20,opacity=.5](a,b,c)
\end{tikzpicture} 
\end{tkzexample}
% subsubsection hexagon_inscribed_version_2 (end)

\newpage
\subsection{点关于圆的幂}

\begin{tikzpicture}
\node [mybox,title={点关于圆的幂}] (box){%
\begin{minipage}{0.90\textwidth}
  {\emph{$\overline{MA} \times \overline{MB}={MT}^2={MO}^2-{OT}^2$} } 
\end{minipage}
};
\end{tikzpicture}% 

\begin{tkzexample}[vbox,small]
\begin{tikzpicture}
 \pgfmathsetmacro{\r}{2}%
 \pgfmathsetmacro{\xO}{6}%
 \pgfmathsetmacro{\xE}{\xO-\r}%
 \tkzDefPoints{0/0/M,\xO/0/O,\xE/0/E}
 \tkzDefCircle[diameter](M,O)
 \tkzGetPoint{I}
 \tkzInterCC(I,O)(O,E) \tkzGetPoints{T}{T'}
 \tkzDefShiftPoint[O](45:2){B}
 \tkzInterLC(M,B)(O,E) \tkzGetPoints{A}{B}
 \tkzDrawCircle(O,E)
 \tkzDrawSemiCircle[dashed](I,O)
 \tkzDrawLine(M,O)
 \tkzDrawLines(M,T O,T M,B)
 \tkzDrawPoints(A,B,T)
 \tkzLabelPoints[above](A,B,O,M,T)
\end{tikzpicture}
\end{tkzexample}

\newpage
\subsection{两个不同心圆的根轴}
\begin{tikzpicture}
\node [mybox,title={两个不同心圆的根轴}] (box){%
\begin{minipage}{0.90\textwidth}
  {来自维基百科：\emph{在几何学中，两个不同心圆的根轴是指关于两圆幂相等的点的集合。因此根轴也被称为两圆的幂线或幂平分线。根轴（radical axis）这一术语由法国数学家 M. Chasles 以 axe radical 的形式使用。
}} 
\end{minipage}
};
\end{tikzpicture}% 

\begin{tkzexample}[vbox,small]
\begin{tikzpicture}
\tkzDefPoints{0/0/A,4/2/B,2/3/K}
\tkzDefCircle[R](A,1)\tkzGetPoint{a}
\tkzDefCircle[R](B,2)\tkzGetPoint{b}
\tkzDefCircle[R](K,3)\tkzGetPoint{k}
\tkzDrawCircles(A,a B,b)
\tkzDrawCircle[dashed,new](K,k)
\tkzInterCC(A,a)(K,k) \tkzGetPoints{a}{a'}
\tkzInterCC(B,b)(K,k) \tkzGetPoints{b}{b'}
\tkzDrawLines[new,add=2 and 2](a,a')
\tkzDrawLines[new,add=1 and 1](b,b')
\tkzInterLL(a,a')(b,b') \tkzGetPoint{X}
\tkzDefPointBy[projection= onto A--B](X) \tkzGetPoint{H}
\tkzDrawPoints(A,B,H,X,a,b,a',b')
\tkzDrawLine(A,B)
\tkzDrawLine[add= 1 and 2,new](X,H)
\tkzLabelPoints(A,B,H,X,a,b,a',b')
\end{tikzpicture}
\end{tkzexample}

\newpage
\subsection{外位似中心}
\begin{tikzpicture}
\node [mybox,title={外位似中心}] (box){%
\begin{minipage}{0.90\textwidth}
  {来自维基百科：\emph{给定两个不同心的圆，作方向相同的平行半径。则连接半径外端的直线经过连心线上一个定点，该点按半径比外分连心线。这个点称为外位似中心，或称外相似中心（Johnson 1929，第 19-20 页和第 41 页）。
}} 
\end{minipage}
};
\end{tikzpicture}% 

\begin{tkzexample}[vbox,small]
\begin{tikzpicture}
\tkzDefPoints{0/0/A,4/2/B,2/3/K}
\tkzDefCircle[R](A,1)\tkzGetPoint{a}
\tkzDefCircle[R](B,2)\tkzGetPoint{b}
\tkzDrawCircles(A,a B,b)
\tkzDrawLine(A,B)
\tkzDefShiftPoint[A](60:1){M}
\tkzDefShiftPoint[B](60:2){M'}
\tkzInterLL(A,B)(M,M') \tkzGetPoint{O}
\tkzDefLine[tangent from = O](B,M') \tkzGetPoints{X}{T'}
\tkzDefLine[tangent from = O](A,M) \tkzGetPoints{X}{T}
\tkzDrawPoints(A,B,O,T,T',M,M')
\tkzDrawLines[new](O,B O,T' O,M')
\tkzDrawSegments[new](A,M B,M')
\tkzLabelPoints(A,B,O,T,T',M,M')
\end{tikzpicture}
\end{tkzexample}

\newpage
\subsection{两圆的公切线}

\begin{tikzpicture}
\node [mybox,title={两圆的公切线}] (box){%
\begin{minipage}{0.90\textwidth}
  {\emph{对于两个圆，当两圆外离时，一般有四条不同的公切线。其中两条外公切线，两圆位于直线的同一侧；外公切线相交于外位似中心。}}
\end{minipage}
};
\end{tikzpicture}%

\begin{tkzexample}[vbox,small]
\begin{tikzpicture}
 \pgfmathsetmacro{\r}{1}%
 \pgfmathsetmacro{\R}{2}%
 \pgfmathsetmacro{\rt}{\R-\r}%
 \tkzDefPoints{0/0/A,4/2/B,2/3/K}
 \tkzDefMidPoint(A,B) \tkzGetPoint{I}
 \tkzInterLC[R](A,B)(B,\rt) \tkzGetPoints{E}{F}
 \tkzInterCC(I,B)(B,F) \tkzGetPoints{a}{a'}
 \tkzInterLC[R](B,a)(B,\R) \tkzGetPoints{X'}{T'}
 \tkzDefLine[tangent at=T'](B) \tkzGetPoint{h}
 \tkzInterLL(T',h)(A,B) \tkzGetPoint{O}
 \tkzInterLC[R](O,T')(A,\r) \tkzGetPoints{T}{T}
 \tkzDefCircle[R](A,\r)  \tkzGetPoint{a}         
 \tkzDefCircle[R](B,\R)  \tkzGetPoint{b}
 \tkzDefCircle[R](B,\rt)  \tkzGetPoint{c}
 \tkzDrawCircles(A,a)  
 \tkzDrawCircles[orange](B,b B,c)           
 \tkzDrawCircle[orange,dashed](I,B)
 \tkzDrawPoints(O,A,B,a,a',E,F,T',T)
 \tkzDrawLines(O,B A,a B,T' A,T)
 \tkzDrawLines[add= 1 and 8](T',h)
 \tkzLabelPoints(O,A,B,a,a',E,F,T,T')
\end{tikzpicture}
\end{tkzexample}

\newpage
\subsection{用根轴作两圆的公切线}

\begin{tikzpicture}
\node [mybox,title={用根轴作两圆的公切线}] (box){%
\begin{minipage}{0.90\textwidth}
  {\emph{只要两个圆不同心，就可以构造它们的根轴，即关于两圆幂相等的点的集合。我们知道根轴是垂直于连心线的一条直线。注意：若 $P$ 和 $Q$ 为某条外公切线与两圆的切点，$D$ 和 $E$ 为两圆在 $[AB]$ 外侧的点，则 $(DP)$ 与 $(EQ)$ 相交于两圆的根轴上。我们将证明这一性质始终成立，并且它使我们能够构造公切线，即使两圆半径相等。}}
\end{minipage}
};
\end{tikzpicture}%


\begin{tkzexample}[vbox,small]
\begin{tikzpicture}
\tkzDefPoints{0/0/A,4/2/B,2/3/K}
\tkzDefCircle[R](A,1) \tkzGetPoint{a}
\tkzDefCircle[R](B,3) \tkzGetPoint{b}
\tkzInterCC[R](A,1)(K,3) \tkzGetPoints{a}{a'}
\tkzInterCC[R](B,3)(K,3) \tkzGetPoints{b}{b'}
\tkzInterLL(a,a')(b,b')  \tkzGetPoint{X}
\tkzDefPointBy[projection= onto A--B](X) \tkzGetPoint{H}
\tkzGetPoint{C}
\tkzInterLC[R](A,B)(B,3) \tkzGetPoints{b1}{E}
\tkzInterLC[R](A,B)(A,1) \tkzGetPoints{D}{a2}
\tkzDefMidPoint(D,E) \tkzGetPoint{I}
\tkzDrawCircle[orange](I,D)
\tkzInterLC(X,H)(I,D) \tkzGetPoints{M}{M'}
\tkzInterLC(M,D)(A,D) \tkzGetPoints{P}{P'}
\tkzInterLC(M,E)(B,E) \tkzGetPoints{Q'}{Q}
\tkzInterLL(P,Q)(A,B) \tkzGetPoint{O}
\tkzDrawCircles(A,a B,b)
\tkzDrawSegments[orange](A,P I,M B,Q)
\tkzDrawPoints(A,B,D,E,M,I,O,P,Q,X,H)
\tkzDrawLines(O,E M,D M,E O,Q)
\tkzDrawLine[add= 3 and 4,orange](X,H)
\tkzLabelPoints(A,B,D,E,M,I,O,P,Q,X,H)
\end{tikzpicture}
\end{tkzexample}

\newpage
\subsection{用圆规求线段中点}

\begin{tikzpicture}
\node [mybox,title={用圆规求线段中点}] (box){%
\begin{minipage}{0.90\textwidth}
  {\emph{本例仅使用圆规来确定一条线段的中点。}}
\end{minipage}
};
\end{tikzpicture}%

\begin{tkzexample}[vbox,small]
\begin{tikzpicture}
\tkzDefPoint(0,0){A}
\tkzDefRandPointOn[circle= center A radius 4]    \tkzGetPoint{B}
\tkzDefPointBy[rotation= center A angle 180](B)  \tkzGetPoint{C}
\tkzInterCC(A,B)(B,A)                            \tkzGetPoints{I}{I'}
\tkzInterCC(A,I)(I,A)                            \tkzGetPoints{J}{B}
\tkzInterCC(B,A)(C,B)                            \tkzGetPoints{D}{E}
\tkzInterCC(D,B)(E,B)                            \tkzGetPoints{M}{M'}
\tkzSetUpArc[color=orange,style=solid,delta=10]
\tkzDrawArc(C,D)(E)
\tkzDrawArc(B,E)(D)
\tkzDrawCircle[color=teal,line width=.2pt](A,B)
\tkzDrawArc(D,B)(M) 
\tkzDrawArc(E,M)(B)
\tkzCompasss[color=orange,style=solid](B,I I,J J,C)
\tkzDrawPoints(A,B,C,D,E,M)
\tkzLabelPoints(A,B,M)
\end{tikzpicture}
\end{tkzexample}
 
\newpage

\subsection{圆的定义，\_Apollonius\_}

\begin{tikzpicture}
\node [mybox,title={圆的定义\_Apollonius\_}] (box){%
\begin{minipage}{0.90\textwidth}
  {来自维基百科：\emph{Apollonius 证明，圆可以定义为平面上到两个定点（称为焦点）距离之比为定值的点的集合。这个阿波罗尼奥斯圆是阿波罗尼奥斯追及问题的基础。……该问题的解有时被称为阿波罗尼奥斯圆。}} 
\end{minipage}
};
\end{tikzpicture}% 

说明

圆是平面上到定点 $O$ 等距的点的集合。到圆心的距离 $r$ 称为半径，点 $O$ 称为圆心。这是最简单的定义，但并非唯一的定义。佩尔加的阿波罗尼奥斯给出了另一种定义：
到两个定点距离之比为常数的所有点的集合是一个圆。

使用 \pkg{tkz-euclide} 可以很容易地说明后一种定义。

\begin{tkzexample}[vbox, small]
\begin{tikzpicture}[scale=1.5]
    % Firstly we defined two fixed point. 
    % The figure depends of these points and the ratio K
\tkzDefPoint(0,0){A}
\tkzDefPoint(4,0){B}
    % tkz-euclide.sty knows about the apollonius's circle
    % with K=2 we search some points like  I such as IA=2 x IB
\tkzDefCircle[apollonius,K=2](A,B) \tkzGetPoints{K1}{k}
\tkzDefPointOnCircle[through=  center K1 angle 30 point k]
\tkzGetPoint{I}
\tkzDefPointOnCircle[through= center K1 angle 280  point k]
\tkzGetPoint{J}
\tkzDrawSegments[new](A,I I,B A,J J,B)  
\tkzDrawCircle[color = teal,fill=teal!20,opacity=.4](K1,k)
\tkzDrawPoints(A,B,K1,I,J)
\tkzDrawSegment(A,B)
\tkzLabelPoints[below,font=\scriptsize](A,B,K1,I,J)
\end{tikzpicture}
\end{tkzexample}

\subsection{反演的应用：\tkzname{Pappus 链}}\label{pappus}
\begin{tikzpicture}
\node [mybox,title={Pappus 链}] (box){%
\begin{minipage}{0.90\textwidth}
来自维基百科 {\emph{在几何学中，Pappus 链是夹在两个相切圆之间的一串圆，由亚历山大的帕普斯在公元 3 世纪研究。}}
\end{minipage}
};
\end{tikzpicture}%

\begin{tkzexample}[vbox,small]
\begin{tikzpicture}[ultra thin]
  \pgfmathsetmacro{\xB}{6}%
  \pgfmathsetmacro{\xC}{9}%
  \pgfmathsetmacro{\xD}{(\xC*\xC)/\xB}%
  \pgfmathsetmacro{\xJ}{(\xC+\xD)/2}%
  \pgfmathsetmacro{\r}{\xD-\xJ}%
  \pgfmathsetmacro{\nc}{16}%
  \tkzDefPoints{0/0/A,\xB/0/B,\xC/0/C,\xD/0/D}
  \tkzDefCircle[diameter](A,C) \tkzGetPoint{x}
  \tkzDrawCircle[fill=teal!30](x,C)
  \tkzDefCircle[diameter](A,B) \tkzGetPoint{y}
  \tkzDrawCircle[fill=teal!30](y,B)
  \foreach \i in {-\nc,...,0,...,\nc}
  {\tkzDefPoint(\xJ,2*\r*\i){J}
   \tkzDefPoint(\xJ,2*\r*\i-\r){H}
   \tkzDefCircleBy[inversion = center A through C](J,H)
   \tkzDrawCircle[fill=teal](tkzFirstPointResult,tkzSecondPointResult)}
\end{tikzpicture}
\end{tkzexample}

\subsection{阿基米德《引理集》命题 1}
\begin{tikzpicture}
\node [mybox,title={《引理集》命题 1 阿基米德}] (box){%
\begin{minipage}{0.90\textwidth}
  {\emph{若两圆相切于 $A$，且 $[CD]$、$[EF]$ 分别为两圆内平行的直径，则 $A$、$C$、$E$ 三点共线。}}
\end{minipage}
};
\end{tikzpicture}%

\begin{tkzexample}[latex=7cm,small]
  \begin{tikzpicture}[scale=.75]
    \tkzDefPoints{0/0/O_1,0/1/O_2,0/3/A}
    \tkzDefPoint(15:3){F}
     \tkzDefPointBy[symmetry=center O_1](F) \tkzGetPoint{E}
     \tkzDefLine[parallel=through O_2](E,F) \tkzGetPoint{x}
     \tkzInterLC(x,O_2)(O_2,A) \tkzGetPoints{D}{C}
     \tkzDrawCircles(O_1,A O_2,A)
     \tkzDrawSegments[orange](O_1,A E,F C,D)
     \tkzDrawSegments[purple](A,E A,F)
     \tkzDrawPoints(A,O_1,O_2,E,F,x,C,D)
     \tkzLabelPoints(A,O_1,O_2,E,F,x,C,D)
  \end{tikzpicture}
\end{tkzexample}

$(CD) \parallel (EF)$，而 $(AO_1)$ 与这两条直线相交，故
$\widehat{AO_2C} = \widehat{AO_1E}$。

由于三角形 $AO_2C$ 与 $AO_1E$ 均为等腰三角形，它们的底角相等：$\widehat{ACO_2} = \widehat{AEO_1} = \widehat{CAO_2} = \widehat{EAO_1}$。因此 $A$、$C$、$E$ 三点共线。

\subsection{阿基米德《引理集》命题 6}
\begin{tikzpicture}
\node [mybox,title={《引理集》命题 6 阿基米德}] (box){%
\begin{minipage}{0.90\textwidth}
  {\emph{设半圆的直径 $AC$ 在 $B$ 处被分割，使 $AC/AB =\phi$ 或取任意比值。在大半圆内、以 $AB$、$BC$ 为直径分别作半圆，再作一个与这三个半圆都相切的圆。若 $GH$ 为该圆的直径，求 $GH$ 与 $AC$ 之间的关系。}}
\end{minipage}
};
\end{tikzpicture}%


\begin{tkzexample}[vbox,overhang,small]
\begin{tikzpicture}
\tkzDefPoints{0/0/A,12/0/C}
\tkzDefGoldenRatio(A,C)                  \tkzGetPoint{B}
\tkzDefMidPoint(A,C)                     \tkzGetPoint{O}
\tkzDefMidPoint(A,B)                     \tkzGetPoint{O_1}
\tkzDefMidPoint(B,C)                     \tkzGetPoint{O_2}
\tkzDefExtSimilitudeCenter(O_1,A)(O_2,B) \tkzGetPoint{M_0}
\tkzDefIntSimilitudeCenter(O,A)(O_1,A)   \tkzGetPoint{M_1}
\tkzDefIntSimilitudeCenter(O,C)(O_2,C)   \tkzGetPoint{M_2}
\tkzInterCC(O_1,A)(M_2,C)                \tkzGetFirstPoint{E}
\tkzInterCC(O_2,C)(M_1,A)                \tkzGetSecondPoint{F}
\tkzInterCC(O,A)(M_0,B)                  \tkzGetFirstPoint{D}
\tkzInterLL(O_1,E)(O_2,F)                \tkzGetPoint{O_3}
\tkzDefCircle[circum](E,F,B)             \tkzGetPoint{0_4}
\tkzInterLC(A,D)(O_1,A)                  \tkzGetFirstPoint{I}
\tkzInterLC(C,D)(O_2,B)                  \tkzGetSecondPoint{K}
\tkzInterLC[common=D](A,D)(O_3,D)        \tkzGetFirstPoint{G}
\tkzInterLC[common=D](C,D)(O_3,D)        \tkzGetFirstPoint{H}
\tkzInterLL(C,G)(B,K)                    \tkzGetPoint{M}
\tkzInterLL(A,H)(B,I)                    \tkzGetPoint{L}
\tkzInterLL(L,G)(A,C)                    \tkzGetPoint{N}
\tkzInterLL(M,H)(A,C)                    \tkzGetPoint{P}  
\tkzDrawCircles[red,thin](O_3,F)
\tkzDrawCircles[new,thin](0_4,B)
\tkzDrawSemiCircles[teal](O,C O_1,B O_2,C)
\tkzDrawSemiCircles[green](M_2,C)
\tkzDrawSemiCircles[green,swap](M_1,A)
\tkzDrawSegment(A,C)
\tkzDrawSegments[new](O_1,O_3 O_2,O_3)
\tkzDrawSegments[new,very thin](B,H C,G A,H G,N H,P)
\tkzDrawSegments[new,very thin](B,D A,D C,D G,H I,B K,B B,G)
\tkzDrawPoints(A,B,C,M_1,M_2,E,O_3,F,D,0_4,O_1,O_2,I,K,G,H,L,P,N,M)  
\tkzLabelPoints[font=\scriptsize](A,B,C,M_1,M_2,F,O_1,O_2,I,K,G,H,L,M,N)
\tkzLabelPoints[font=\scriptsize,right](E,O_3,D,0_4,P)
\end{tikzpicture}
\end{tkzexample}

设 $GH$ 为与 $AC$ 平行的圆的直径，且该圆分别在 $D$、$E$、$F$ 处与以 $AC$、$AB$、$BC$ 为直径的半圆相切。

则由命题 1，$A$、$G$、$D$ 共线，$D$、$H$、$C$ 共线。\\
同理，$A$、$E$、$H$ 共线，$C$、$F$、$G$ 共线，$B$、$E$、$G$ 与 $B$、$F$、$H$ 也分别共线。
 
设 $(AD)$ 交以 $[AC]$ 为直径的半圆于 $I$，$(BD)$ 交以 $[BC]$ 为直径的半圆于 $K$。连接 $CI$、$CK$，分别交 $AE$、$BF$ 于 $L$、$M$，延长 $GL$、$HM$，分别交 $AB$ 于 $N$、$P$。

在三角形 $AGB$ 中，从 $A$、$C$ 向对边所作的垂线交于 $L$。因此由三角形的性质，$(GN)$ 垂直于 $(AC)$。
同理 $(HP)$ 垂直于 $(BC)$。\\
又因为 $I$、$K$、$D$ 处的角为直角，$(CK)$ 平行于 $(AD)$，$(CI)$ 平行于 $(BD)$。

因此\\
\[\frac{AB}{BC} = \frac{AL}{LH}    =  \frac{AN}{NP}  \quad\text{且} \quad \frac{BC}{AB} = \frac{CM}{MG}    =  \frac{PC}{NP} \]

从而

\[ \frac{AN}{NP}    =  \frac{NP}{PC} \quad\text{故} \quad {NP}^2 = AN \times PC  \]

现设 $B$ 按黄金分割比分割 $[AC]$，即：
\[\phi = \frac{AB}{BC} =  \frac{AC}{AB} \quad\text{则} \quad AN = \phi NP \text{且}\quad  NP = \phi PC \]

我们有
\[ AC = AN + NP + PC\quad \text{即} \quad AB + BC = AN + NP + PC \quad \text{或} \quad (\phi + 1) BC = AN + NP + PC \]

于是

\[ (\phi + 1) BC = \phi NP + NP + PC =(\phi + 1)NP + PC = \phi(\phi + 1)PC + PC = ({\phi}^2 + \phi + 1)PC \]

由于
\[ {\phi}^2 = \phi + 1 \quad \text{故} \quad (\phi + 1) BC = 2(\phi + 1) PC \quad\text{即}\quad BC = 2 PC \]

也就是说，
$P$ 是线段 $BC$ 的中点。

部分证明摘自 \url{https://www.cut-the-knot.org}


\subsection{阿波罗尼奥斯圆}

\begin{tikzpicture}
\node [mybox,title={三角形的阿波罗尼奥斯圆\_Apollonius\_}] (box){%
\begin{minipage}{0.90\textwidth}
  {\emph{与三角形的三个旁切圆都相切且将它们包含在内的圆，通常被称为阿波罗尼奥斯圆（Kimberling 1998，第 102 页）。}}
\end{minipage}
};
\end{tikzpicture}%

说明

前面几个示例的目的是展示这些命题可以多么简单地被重现。使用 TikZ 需要计算和三角学，而使用 \pkg{tkz-euclide} 只需构建简单的对象。

但请不要忘记，TikZ 位于 \pkg{tkz-euclide} 之下，且远在其之上：我只是在 TikZ 与我的宏包用户之间创建了一个接口。

最后一个示例非常复杂；它展示了使用 \pkg{tkz-euclide} 所能做的一切。


\begin{tkzexample}[vbox,small]
\begin{tikzpicture}[scale=.6]
\tkzDefPoints{0/0/A,6/0/B,0.8/4/C}
\tkzDefTriangleCenter[euler](A,B,C)        \tkzGetPoint{N} 
\tkzDefTriangleCenter[circum](A,B,C)       \tkzGetPoint{O} 
\tkzDefTriangleCenter[lemoine](A,B,C)      \tkzGetPoint{K}
\tkzDefTriangleCenter[ortho](A,B,C)        \tkzGetPoint{H}
\tkzDefSpcTriangle[excentral,name=J](A,B,C){a,b,c}
\tkzDefSpcTriangle[centroid,name=M](A,B,C){a,b,c}
\tkzDefCircle[in](Ma,Mb,Mc)                \tkzGetPoint{Sp}  % Sp Spieker center
\tkzDefProjExcenter[name=J](A,B,C)(a,b,c){Y,Z,X}
\tkzDefLine[parallel=through Za](A,B)      \tkzGetPoint{Xc}
\tkzInterLL(Za,Xc)(C,B)                    \tkzGetPoint{C'}
\tkzDefLine[parallel=through Zc](B,C)      \tkzGetPoint{Ya}
\tkzInterLL(Zc,Ya)(A,B)                    \tkzGetPoint{A'}
\tkzDefPointBy[reflection= over Ja--Jc](C')\tkzGetPoint{Ab}
\tkzDefPointBy[reflection= over Ja--Jc](A')\tkzGetPoint{Cb}
\tkzInterLL(K,O)(N,Sp)                     \tkzGetPoint{Q}
\tkzInterLC(A,B)(Q,Cb)                     \tkzGetFirstPoint{Ba}
\tkzInterLC(A,C)(Q,Cb)                     \tkzGetPoints{Ac}{Ca}
\tkzInterLC(B,C')(Q,Cb)                    \tkzGetFirstPoint{Bc}
\tkzInterLC[next to=Ja](Ja,Q)(Q,Cb)        \tkzGetFirstPoint{F'a}
\tkzInterLC[next to=Jc](Jc,Q)(Q,Cb)        \tkzGetFirstPoint{F'c}
\tkzInterLC[next to=Jb](Jb,Q)(Q,Cb)        \tkzGetFirstPoint{F'b}
\tkzInterLC[common=F'a](Sp,F'a)(Ja,F'a)    \tkzGetFirstPoint{Fa}
\tkzInterLC[common=F'b](Sp,F'b)(Jb,F'b)    \tkzGetFirstPoint{Fb}
\tkzInterLC[common=F'c](Sp,F'c)(Jc,F'c)    \tkzGetFirstPoint{Fc}
\tkzInterLC(Mc,Sp)(Q,Cb)                   \tkzGetFirstPoint{A''}
\tkzDefCircle[euler](A,B,C)                \tkzGetPoints{E}{e}
\tkzDefCircle[ex](C,A,B)                   \tkzGetPoints{Ea}{a}
\tkzDefCircle[ex](A,B,C)                   \tkzGetPoints{Eb}{b}
\tkzDefCircle[ex](B,C,A)                   \tkzGetPoints{Ec}{c}
% Calculations are done, now you can draw, mark and label
\tkzDrawCircles(Q,Cb E,e)%
\tkzDrawCircles(Eb,b Ea,a Ec,c)
\tkzDrawPolygon(A,B,C)
\tkzDrawSegments[dashed](A,A' C,C' A',Zc Za,C' B,Cb B,Ab A,Ca)
\tkzDrawSegments[dashed](C,Ac Ja,Xa Jb,Yb Jc,Zc)
\begin{scope}
   \tkzClipCircle(Q,Cb) % We limit the drawing of the lines
   \tkzDrawLine[add=5 and 12,orange](K,O)
   \tkzDrawLine[add=12 and 28,red!50!black](N,Sp)
\end{scope}
\tkzDrawPoints(A,B,C,K,Ja,Jb,Jc,Q,N,O,Sp,Mc,Xa,Xb,Yb,Yc,Za,Zc)
\tkzDrawPoints(A',C',A'',Ab,Cb,Bc,Ca,Ac,Ba,Fa,Fb,Fc,F'a,F'b,F'c)
\tkzLabelPoints(Ja,Jb,Jc,Q,Xa,Xb,Za,Zc,Ab,Cb,Bc,Ca,Ac,Ba,F'b)
\tkzLabelPoints[above](O,K,F'a,Fa,A'')
\tkzLabelPoints[below](B,F'c,Yc,N,Sp,Fc,Mc)
\tkzLabelPoints[left](A',C',Fb)
\tkzLabelPoints[right](C)
\tkzLabelPoints[below right](A)
\tkzLabelPoints[above right](Yb)
\tkzDrawSegments(Fc,F'c Fb,F'b Fa,F'a)
\tkzDrawSegments[color=green!50!black](Mc,N Mc,A'' A'',Q)
\tkzDrawSegments[color=red,dashed](Ac,Ab Ca,Cb Ba,Bc Ja,Jc A',Cb C',Ab)
\tkzDrawSegments[color=red](Cb,Ab Bc,Ac Ba,Ca A',C')
\tkzMarkSegments[color=red,mark=|](Cb,Ab Bc,Ac Ba,Ca)
\tkzMarkRightAngles(Jc,Zc,A Ja,Xa,B Jb,Yb,C)
\tkzDrawSegments[green,dashed](A,F'a B,F'b C,F'c)
\end{tikzpicture}
\end{tkzexample}

\endinput


\part{常见问题}
\section{常见问题}

\subsection{最常见的错误}
这些是我自己最常犯的错误，因为语法已经修改过好几次。本节将持续扩充。

\begin{itemize}\setlength{\itemsep}{10pt}
  \item 忘记在绘制多个对象的宏后加 ``s''，例如 \tkzcname{tkzDrawSegments} 或 \tkzcname{tkzDrawCircles}：\tkzcname{tkzDrawPoint(A,B)} 是错误的，应当写成 \tkzcname{tkzDrawPoints(A,B)}。

  \item 自第 4 版起，单位始终为 cm，因此不要再写单位：\tkzcname{tkzDrawCircle[R](O,3cm)} 应写成 \tkzcname{tkzDrawCircle[R](O,3)}。常用的 \tkzname{TikZ} 选项仍保留其单位，例如 \tkzname{below right = 12pt}；另一方面，在 \tkzcname{tkzMarkAngle} 中设置圆弧半径时应写成 \tkzname{size=1.2}。

  \item 圆括号与花括号混用。创建的点其名称放在花括号中，而作为选项或参数使用的点其名称放在圆括号中：应写 \tkzcname{tkzGetPoint\{A\}}，而不是 \tkzcname{tkzGetPoint(A)}。

  \item 用 \tkzcname{tkzGetPoint\{A\}} 代替了 \tkzcname{tkzGetFirstPoint\{A\}}。当一个宏给出两个点时，要么用 \tkzcname{tkzGetPoints\{A\}\{B\}} 同时获取两个点，要么用 \tkzcname{tkzGetFirstPoint\{A\}} 或 \tkzcname{tkzGetSecondPoint\{A\}} 只获取其中一个。这些点也可以通过引用 \tkzname{tkzFirstPointResult} 和 \tkzname{tkzSecondPointResult} 来使用。有些宏还会给出第三个点 \tkzname{tkzPointResult}。

  \item 立即获取结果：\tkzcname{tkzGetPoint}、\tkzcname{tkzGetPoints} 或 \tkzcname{tkzGetLength} 必须紧跟在计算结果的宏之后，因为下一次计算会覆盖它。

  \item 直线与圆、或两圆的交点顺序在不同版本间有时会发生变化；这会导致图形错误或报错。选项 \tkzname{common} 和 \tkzname{near} 允许你选择所需的点。

  \item 选项与参数混用：所有使用圆的宏都需要知道其半径。当半径由一个测量值给出时，选项中要包含 \tkzname{R}。

  \item 角度以度为单位，少数情况下以弧度为单位。

  \item 点列表不能在行尾被 \code{\%} 截断：|\tkzDrawSegments(A,B%| 紧接着 |C,D)| 会得到点 |BC,D|，因为 \code{\%} 也会删掉两个点之间的空格。直接换行，不要加 \code{\%}。

  \item 第 3 版中用于特殊圆的语法，例如 |\tkzDrawCircle[euler](A,B,C)|，已不再存在：先用 \tkzcname{tkzDefCircle[euler](A,B,C)} 定义圆，再用 \tkzcname{tkzGetPoints\{N\}\{P\}} 获取圆心和圆上一点，最后用 \tkzcname{tkzDrawCircle(N,P)} 绘制。自第 5.14 版起，旧语法会给出错误信息。

  \item 如果传参时计算出错，最好先完成计算再调用宏。

  \item 不要混用 \tkzNamePack{pgfmath} 和 \tkzNamePack{xfp} 的语法。我通常选择 \tkzNamePack{xfp}；如果你更喜欢 \tkzNamePack{pgfmath}，请在传参之前完成计算。

  \item ``dimension too large''（维度太大）错误在某些情况下会出现，例如使用 \tkzcname{tkzDrawArc} 或 \tkzcname{tkzMarkAngle} 且角度非常小时：它来自 \NameLib{decorations} 库，在弧上放置标记时触发，即使没有标记也会发生。一种避免方法是使用选项 \tkzname{veclen}（参见第 \ref{opt-veclen} 节）：在使用它的作用域内，函数 \code{veclen} 由 \tkzNamePack{xfp} 计算。在此作用域内不要使用求交宏。

  \item 自第 5.14 版起，\tkzcname{tkzSetUpCompass}\code{[delta=...]} 不再扩展用 \tkzcname{tkzDrawArc} 绘制的圆弧：对这些圆弧请使用 \tkzcname{tkzSetUpArc}\code{[delta=...]}。
\end{itemize}
\endinput


\part{新特性、兼容性与选项}
\section{新特性与兼容性}\label{news}

\subsection{第 5.14 版}
\begingroup\sloppy

本版本仅修正缺陷；语法没有变化，也没有新增宏。

\begin{itemize}
  \item \tkzcname{tkzDefCircle} 在使用选项 \tkzname{orthogonal from} 或 \tkzname{orthogonal through} 时，不再影响后续的 \tkzcname{tkzDefCircle} 调用。
  \item 下列宏的选项现在仅作用于当前调用：\tkzcname{tkzShowLine}、\tkzcname{tkzDefSimilitudeCenter}、\tkzcname{tkzCalcLength}、\tkzcname{tkzGrid}、\tkzcname{tkzClip}、\tkzcname{tkzDefEquiPoints}、\tkzcname{tkzLegend}。
  \item \tkzcname{tkzIsLinear} 曾给出错误答案（始终为真）；\tkzcname{tkzIsLinear} 和 \tkzcname{tkzIsOrtho} 现在使用随长度变化的容差。
  \item \tkzcname{tkzPowerCircle} 的结果不再在其他计算之后发生改变。
  \item 圆的反演（\tkzcname{tkzDefCircleBy}）仅当原圆经过反演中心时才得到一条直线。
  \item \tkzcname{tkzDefCircle}\code{[spieker]} 现在以 cm 为单位给出半径，与其他选项一致。
  \item 用 \tkzcname{tkzDefTriangleCenter} 和 \tkzcname{tkzDefSpcTriangle} 的 \tkzname{feuerbach} 选项得到的 Feuerbach 点和 Feuerbach 三角形现在是精确的。
  \item \tkzcname{tkzDrawPolySeg}：路径起点处的箭头方向正确。
  \item \tkzcname{tkzCompass} 接受 \tkzname{delta} 的小数值；\tkzcname{tkzProtractor}\code{[return]} 正常工作；\tkzcname{tkzLabelRegPolygon} 使用其 \tkzname{sep} 选项；\tkzcname{tkzDrawX} 使用其 \tkzname{step} 选项。
  \item 第 3 版的旧语法，例如 |\tkzDrawCircle[euler](A,B,C)|，现在会给出错误信息，而不是陷入无休止的编译。
  \item 别名 \tkzcname{tkzDefNinePointsCircle}、\tkzcname{tkzFeuerBachCircle} 和 \tkzcname{tkzIncentralTriangle} 恢复正常。
\end{itemize}

有两处变化可能会影响你的图形：
\begin{itemize}
  \item \tkzcname{tkzSetUpCompass}\code{[delta=...]} 现在只作用于圆规标记。要扩展由 \tkzcname{tkzDrawArc}、\tkzcname{tkzMarkArc} 和 \tkzcname{tkzLabelArc} 绘制的圆弧，请使用 \tkzcname{tkzSetUpArc}\code{[delta=...]}。在此版本之前，两者使用同一个值。
  \item \TIKZ\ 样式 \tkzname{help lines} 不再被本宏包修改。
\end{itemize}

\endgroup

\subsection{第 5.10 版}
\begin{itemize}
  \item 新增法语文档。
  \item 新增 \code{mini} 选项。你可以将此选项与 \tkzNamePack{tkz-elements} 宏包配合使用。只会加载绘图所需的模块。目前此选项仅在使用 \tkzNamePack{tkz-elements} 时有用。
\end{itemize}

\subsection{从第 3 版到第 5 版}

\tkzname{\tkznameofpack} 4 是其前身的简化版本，独立于 \tkzname{tkz-base}。\tkzname{tkz-euclide} 3 的大多数宏都被保留，但特殊圆现在先用 \tkzcname{tkzDefCircle} 定义，再用 \tkzcname{tkzDrawCircle} 绘制。单位现在是 \tkzname{cm}。如果你需要使用 \tkzname{tkz-base} 中的某些宏，可能需要使用 \tkzcname{tkzInit}。

第 5 版对用户没有任何改变。使用选项 \code{lua}（|\usepackage[lua]{tkz-euclide}|）时，计算由 Lua 完成：之后必须用 LuaLaTeX 编译。参见第 \ref{calc_with_lua} 节。

% \subsection{With 5.06 version}
%      \begin{itemize}
%       \item 修正了宏 |\tkzLabelAngle| 与选项 \code{angle} 的缺陷。
%        \item 新增 |\tkzSetUpCircle|。
%        \item 修正了若干笔误。
%        \item 移除了一些法文文本。
%      \end{itemize}
%
%
% \subsection{With 5.05 version}
%
%  修正了 ``Complete but minimal example'' 中的文档。
%
% \subsection{With 5.03 version}
%
% \begin{itemize}
%
%     \item 修正了文件 tkz-obj-lua-points-spc.tex 中宏 |\tkzDefBarycentricPointTwo| 的缺陷。
%
%   \item 新增宏 |\tkzDrawEllipse|。
%
%   \item 删除宏 |\tkzDrawSectorAngles| 和 |\tkzDrawSectorRwithNodesAngles|。
% \end{itemize}
%
%
% \subsection{With 5.0 version} % (fold)
%
% \begin{itemize}
%
%     \item 最后，我为宏包 \tkzname{\tkznameofpack} 新增了 \code{lua} 选项。它允许用 lua 完成主要函数的计算；（参见 \ref{calc_with_lua}）。语法不变。对用户而言一切如常。
%
%   \item 选项 \code{xfp} 已改为 \code{veclen}，参见 \ref{opt-veclen}。
%
% \end{itemize}
%
% \subsection{With 4.2 version} % (fold)
% \label{sub:with_4_2_version}
%
% 为使语法更统一，特别是为了将坐标的定义与求解和其余部分（即绘图、标记与标注）区分开，做了一些修改。
% 现在定义宏被单独隔离，今后引入使用 \tkzimp{Lua} 的坐标计算阶段会更容易。
%
% 以下是部分改动。
% \vspace{1cm}
%  \begin{itemize}\setlength{\itemsep}{10pt}
%
%
% \item 我最近发现在使用 \code{scale} 选项时存在一个问题。当用某些工具绘制某些图形时，大量使用 |pgfmathreciprocal| 会引入微小的计算误差，这些误差会累积起来，使图形无法使用。以下是避免这些问题的做法：
% \begin{enumerate}
%
%   \item 我这边引入了 Muzimuzhi 提出的补丁，修改了 |pgfmathreciprocal|；
%
% \item  Muzimuzhi 提出的另一个思路是，在加载 |fpu| 库之后，把 |/pgf/fpu/install only={reciprocal}| 作为 |tikzpicture| 环境的选项传入；
%
% \item 在我所选择的宏定义方法中，尽量避免使用 |pgfmathreciprocal|；
%
% \item 还有一种万无一失的方法，即避免使用 |scale = ...|。如果你和我一样，只在代码开头固定若干定点，那么这很容易做到。图形的大小只取决于这些定点，因此你只需调整这些点的坐标即可。
% \end{enumerate}
%
% \item 现在 |\tkzDefCircle| 给出两个点作为结果：圆心和圆上一点。当圆上一点已知时，用 |\tkzGetPoint| 或 |\tkzGetFirstPoint|
% 即可得到圆心；否则 |\tkzGetPoints| 会给出圆心和圆上一点。你始终可以用 |\tkzGetLength| 获取半径长度。我希望优先使用节点，并杜绝在代码中出现数字。
%
% \item 为了隔离定义，我删除或修改了以下宏：|\tkzDrawLine|、|\tkzDrawTriangle|、|\tkzDrawCircle|、|\tkzDrawSemiCircle| 和 |\tkzDrawRectangle|；
%
% 因此 |\tkzDrawSquare(A,B)| 变为 |\tkzDefSquare(A,B)||\tkzGetPoints{C}{D}|，然后
%
%  |\tkzDrawPolygon(A,B,C,D)|；
%
% 如果你想画圆，不能直接写 |\tkzDrawCircle[R](A,1)|。必须先定义圆经过的点，因此要执行
% |\tkzDefCircle[R](A,1)| |\tkzGetPoint{a}|，最后 |\tkzDrawCircle(A,a)|。另一种方法是定义圆上一点 |\tkzDefShiftPoint[A](1,O){a}|；
%
%
% \item 下列宏 |tkzDefCircleBy[orthogonal through]| 和 |\tkzDefCircleBy[orthogonal from]| 变为 |tkzDefCircle[orthogonal through]| 和 |\tkzDefCircle[orthogonal from]|；
%
%
% \item |\tkzDefLine[euler](A,B,C)| 是一个宏，用于在可能时得到 \tkzname{Euler} 直线。|\tkzDefLine[altitude](A,B,C)| 再次可用，|\tkzDefLine[tangent at=A](O)| 和 |\tkzDefLine[tangent from=P](O,A)| 同样可用，此前它们无法工作；
%
%
% \item | \tkzDefTangent| 被替换为 |\tkzDefLine[tangent from = ...]| 或 |\tkzDefLine[tangent at = ...]|；
%
%
% \item 我为更喜欢直接使用 \TIKZ 的人新增了宏 |\tkzPicAngle[tikz options](A,B,C)|；
%
%
% \item 宏 |\tkzMarkAngle| 已修正；
%
% \item 与 |\tkzDefCircle| 命令的 \tkzname{apollonius} 选项相关的宏已重写；
%
% \item (4.23) 宏 |\tkzDrawSemiCircle| 已修正；
%
% \item
% 宏 \tkzcname{tkzDefPointOnCircle} 的参数顺序已改变：现在是圆心、角度、点或半径。
% 我新增了两个用于弧度制的选项：\tkzname{through in rad} 和 \tkzname{R in rad}。
%
%
% \item 我为圆弧路径新增了选项 \tkzname{reverse}。它可以反转路径，并在必要时反转已有的箭头。
%
%
% \item 我统一了标签的样式。现在只剩 \tkzname{label style} 一个，对点、线段、直线、圆和角都有效。我删除了 \tkzname{label seg style}、\tkzname{label line style} 和 \tkzname{label angle style}。
%
% \item 我新增了宏 |tkzFillAngles|，用于同时处理多个角。
%
% \item 修正 \tkzname{return} 选项与 \tkzcname{tkzProtractor} 的配合。
%
% 提醒一下，此前还做过以下改动：
%
%  \item  |\tkzDrawMedian|、|\tkzDrawBisector|、|\tkzDrawAltitude|、|\tkzDrawMedians|、|\tkzDrawBisectors| 和 |\tkzDrawAltitudes| 已不复存在。由于不遵守创建与绘制分离的原则，建议先用 |\tkzDefSpcTriangle[median]| 创建这些点的坐标，再用 |\tkzDrawSegments| 或 |\tkzDrawLines| 选择要绘制的那些；
%
% \item |\tkzDrawTriangle| 已删除。|\tkzDrawTriangle[equilateral]| 虽然方便，但更好的做法是用 |\tkzDefTriangle[equilateral]| 得到第三个点，再用 |\tkzDrawPolygon| 绘制；|\tkzDrawSquare| 和 |\tkzDrawGoldRectangle| 同理；
%
%
% \item 圆的反演此前定义有误，因此我重写了该宏。输入参数始终是圆心和圆上一点，输出参数是像圆的圆心和像圆上一点；若原圆经过反演极点，则输出为像直线上的两个点。如果圆经过反演中心，像就是一条直线；该过程是否有效取决于原圆上所取点的选择；
%
% \item 修正黄金三角形与欧几里得三角形的正确归属；
%
%
% \item 我为求交 LC 和 CC 新增了选项 \code{next to}；
%
%
% \item 修正等腰直角选项；
%
% \item (4.22 和 4.23) 修正宏 |\tkzMarkAngle|；
%
%
% \item |\tkzDefMidArc(O,A,B)| 给出以 $O$ 为圆心、从 $A$ 到 $B$ 的弧的中点；
%
% \item 好消息：新增了一些有用的工具。它们目前以实验性方式提供，无疑还需要改进；
%
%
% \item 选项 \code{orthogonal from} 和 \code{through} 现在归 \tkzcname{tkzDefCircleBy} 管辖。
%
% \begin{enumerate}
%
%   \item |\tkzDotProduct(A,B,C)| 在正交参考系中计算向量 $\overrightarrow{A,B}$ 与 $\overrightarrow{A,C}$ 的数量积。
%
%   |\tkzDotProduct(A,B,C)=aa'+bb' 当 vec{AB} =(a,b) 且 vec{AC} =(a',b')|
%
%
%   \item |\tkzPowerCircle(A)(B,C)| 点 $A$ 关于以 $B$ 为圆心且经过 $C$ 的圆的幂；
%
%
%   \item |\tkzDefRadicalAxis(A,B)(C,D)| 圆心分别为 $A$ 和 $C$ 的两圆的根轴；
%
%
% \item (4.23) 宏 |tkzDefRadicalAxis| 已完善
%
%   \item 一些测试：|\tkzIsOrtho(A,B,C)| 和 |\tkzIsLinear(A,B,C)|。前者判断直线 $(A,B)$ 与 $(A,C)$ 是否正交。后者判断点 $A$、$B$、$C$ 是否共线；
%
%  |\tkzIsLinear(A,B,C)| 若 $A$、$B$、$C$ 共线，则 |\tkzLineartrue|
%   你可以使用 |\iftkzLinear|（|\tkzIsOrtho| 同理）；
%
% \item 新增了一种向量样式，你当然可以修改它
%
% |tikzset{vector style/.style={>=Latex,->}}|；
%
%
% \item 现在可以用选项 |tkz arrow| 在直线或圆上加箭头。
% \end{enumerate}
% \end{itemize}
%
% % subsection with_4_2_version (end)
% \subsection{Changes with previous versions} % (fold)
% \label{sub:changes_with_previous_versions}
%
% \vspace{1cm}
%  \begin{itemize}\setlength{\itemsep}{10pt}
%
% \item 提醒一下，一个重要的新变化是最近用 \tkzNamePack{xfp} 替换了 \tkzNamePack{fp} 宏包。这是为了进一步改进计算并更易用；
%
%
% \item 首先，你不必再用 \TIKZ 处理 bounding box 的大小。早期版本的 \tkzname{\tkznameofpack} 不控制 bounding box 的大小；现在每个宏（希望如此）都会控制 bounding box，以避免先使用 \tkzcname{tkzInit} 再使用 \tkzcname{tkzClip}；
%
% \item 使用 \tkzimp{tkz-euclide} 会加载所有对象，因此不需要再写 \tkzcname{usetkzobj\{all\}}；
%
% \item 新增了关于 bounding box 的宏：\tkzcname{tkzSaveBB}、\tkzcname{tkzClipBB} 等；
%
% \item 按常理，大多数宏都接受 \TIKZ 选项。因此我在可能的情况下移除了重复的选项，例如 \code{label options} 选项被删除；
%
% \item 单位现在是 cm；
%
% \item |\tkzCalcLength| |\tkzGetLength| 以 cm 为单位给出结果；
%
% \item  |\tkzMarkArc| 和 |\tkzLabelArc| 是新宏；
%
% \item 现在 |\tkzClipCircle| 和 |\tkzClipPolygon| 有一个 \tkzimp{out} 选项。要使用此选项，你必须有一个包含待裁剪对象的 bounding box。可以使用一个包围整个图形的对象，或使用宏 \tkzcname{tkzInit} 来实现；
%
%
% \item 此前用于给直线标注的选项 \tkzname{end} 和 \tkzname{start} 已被移除。现在必须使用宏 \tkzcname{tkzLabelLine}；
%
% \item 引入库 \NameLib{quotes} 和 \NameLib{angles}；它允许给点加标签，尽管我并不赞成这种做法；
%
% \item 向量的概念消失了；要画向量，只需把 ``->'' 作为选项传给 \tkzcname{tkzDrawSegment}；
%
%
% \item |\tkzDefIntSimilitudeCenter| 和 |\tkzDefExtSimilitudeCenter| 已不复存在，现在需要使用 |\tkzDefSimilitudeCenter[int]| 或 |\tkzDefSimilitudeCenter[ext]|；
%
% \item |\tkzDefRandPointOn| 被 |\tkzGetRandPointOn| 替换；
%
%
% \item 宏 \tkzcname{tkzDefTriangle} 的一个选项已改变：上一版中该选项是带 ``e'' 的 \code{euclide}，现在是 \code{euclid}；
%
% \item 随机点现在位于 \tkzname{\tkznameofpack} 中，宏 \tkzcname{tkzGetRandPointOn} 被替换为
%
%  \tkzcname{tkzDefRandPointOn}。出于一致性考虑，点必须用 \tkzcname{tkzGetPoint} 获取；
%
% \item 新增了以下宏：\tkzcname{tkzDrawSemiCircles}、\tkzcname{tkzDrawPolygons}、\tkzcname{tkzDrawTriangles}；
%
%
% \item 选项 \code{isosceles right} 是宏 \tkzcname{tkzDefTriangle} 的一个新选项；
%
% \item 出现了宏 \tkzcname{usetkztool}，用于加载新的 \code{tools}；
%
% \item 可以借助以下宏修改样式：\tkzcname{tkzSetUpPoint}、\tkzcname{tkzSetUpLine}、\tkzcname{tkzSetUpArc}、\tkzcname{tkzSetUpCompass}、\tkzcname{tkzSetUpLabel} 和 \tkzcname{tkzSetUpStyle}。最后一个允许你创建新样式。
% \end{itemize}
% % subsection changes_with_previous_versions (end)

\endinput

\newpage
\section{使用 lua} \label{calc_with_lua}

如果你希望由 Lua 完成计算，现在推荐的方式是使用 \tkzNamePack{tkz-elements} 宏包：它用 Lua 计算点，而 \tkzname{\tkznameofpack} 负责绘制（参见下文的 \code{mini} 选项）。\tkzname{\tkznameofpack} 的 \code{lua} 选项仍然可用并维护，但新的 Lua 开发已转入 \tkzNamePack{tkz-elements}。

\subsection{选项 \code{lua}} % (fold)
\label{ssub:option_code_lua}

% subsubsection option_code_lua (end)
自第 5 版起，你可以使用 \tkzname{\tkznameofpack} 的选项 \ItkzPopt{tkz-euclide}{lua}。只需在导言区写入

 |\usepackage[lua]{tkz-euclide}|。
 之后必须用 LuaLaTeX 编译。语法保持不变。

若不使用此选项，\tkzname{\tkznameofpack} 将使用第 4 版的代码，即配合 \tkzNamePack{xfp}。

你正在阅读的本文档就是用此选项编译的。

关于方法和结果略作说明。我考虑过两种方案。第一种是尽可能地把 \code{xfp} 所做的计算替换为 \code{lua} 计算：我就是这样从 \code{fp} 过渡到 \code{xfp}，现在再到 \code{lua}。第二种更大胆的方案是给每个点关联一个复数，并用 \code{lua} 进行复数运算；这正是 \tkzNamePack{tkz-elements} 现在所做的事情。

结果令人满意：本文档（当时为 236 页）用 \code{LuaLaTeX} 加 \code{xfp} 编译需 47\,s，而用 \code{lua} 仅需 30\,s。

另一份 61 页的文档用 \code{pdfLaTeX} 加 \code{xfp} 编译需 16\,s，用 \code{LuaLaTeX} 加 \code{xfp} 需 13\,s。

本文档使用 |\usepackage{tkz-base}| 和 |\usepackage[lua]{tkz-euclide}| 编译，但我没有对所有交互做充分测试。

\subsection{选项 \code{mini}} % (fold)
\label{ssub:option_code_mini}

当你使用 \tkzNamePack{tkz-elements} 来计算图形中的点时，并不需要 \tkzname{\tkznameofpack} 的全部模块。使用选项 \ItkzPopt{tkz-euclide}{mini}，即 |\usepackage[mini]{tkz-euclide}|，只会加载绘图所需的模块。有些宏，例如 ``The Show'' 一节中的那些，仍然需要完整宏包。

% subsubsection option_code_mini (end)
\endinput


\clearpage\newpage
\small\printindex
\end{document}
